\documentclass[11pt]{article}

% arXiv/template formatting
\usepackage[noblocks]{authblk}
\usepackage[letterpaper, margin=1in]{geometry}
\usepackage[title]{appendix}
\usepackage{setspace}
\usepackage{url}

% Manuscript packages
\usepackage{amsmath}
\usepackage{amssymb}
\usepackage{mathtools}
\usepackage{graphicx}
\usepackage{flafter}
\usepackage{booktabs}
\usepackage{algorithm}
\usepackage{algpseudocode}
\usepackage[authoryear, round]{natbib}
\usepackage{xcolor}
\usepackage{comment}
\usepackage{bm}
\usepackage{tikz}
\usepackage{pgfplots}
\usepackage{mathrsfs}
\pgfplotsset{compat=1.18}
\usetikzlibrary{plotmarks, arrows.meta, positioning}
\usepgfplotslibrary{patchplots}
\usepackage{hyperref}
\usepackage[noabbrev]{cleveref}
\crefname{section}{\S}{\S\S}
\Crefname{section}{Section}{Sections}

% Print \label names in the margins while drafting, comment out for the final version.
\usepackage[notref,notcite]{showkeys}

\graphicspath{{./}{figures/}}

\definecolor{darkblue}{rgb}{0,0,0.80}
\hypersetup{
    colorlinks=true,
    linkcolor={darkblue},
    citecolor={darkblue},
    filecolor=black,
    urlcolor={darkblue},
    plainpages=false,
}

\makeatletter
\newcommand\affiliation[1]{\gdef\@affiliation{\let\aff\aff@inst#1}}
\gdef\@affiliation{}
\def\email#1{Email address for correspondence: #1}
\def\aff#1{\ignorespaces\textsuperscript{#1}}
\def\corresp#1{\unskip\thanks{#1}}
\makeatother
\setlength{\abovecaptionskip}{8pt plus 0pt minus 2pt}
\setlength{\bibsep}{2pt plus 0.3ex}
\numberwithin{equation}{section}
\renewcommand\Authsep{, }
\renewcommand\Authands{, }
\renewenvironment{abstract}
{\begin{quote}
\noindent \rule{\linewidth}{.5pt}\par{\bfseries \abstractname.}}
{\medskip\noindent \rule{\linewidth}{.5pt}
\end{quote}}

% Custom fonts from the arXiv template
\DeclareTextFontCommand\textsfi{\usefont{OT1}{cmss}{m}{sl}}
\DeclareMathAlphabet\mathsfi{OT1}{cmss}{m}{sl}
\DeclareTextFontCommand\textsfb{\usefont{OT1}{cmss}{bx}{n}}
\DeclareMathAlphabet\mathsfb{OT1}{cmss}{bx}{n}
\DeclareTextFontCommand\textsfbi{\usefont{OT1}{cmss}{m}{sl}}
\DeclareMathAlphabet\mathsfbi{OT1}{cmss}{m}{sl}
\IfFileExists{t1phv.fd}
{\DeclareTextFontCommand\textsfbi{\usefont{T1}{phv}{b}{it}}
\DeclareMathAlphabet\mathsfbi{T1}{phv}{b}{it}}
{}
\IfFileExists{ot1phv.fd}
{\DeclareTextFontCommand\textsfbi{\usefont{OT1}{phv}{b}{it}}
\DeclareMathAlphabet\mathsfbi{OT1}{phv}{b}{it}}
{}

\newcommand{\tc}[1]{\mathsfbi{#1}}
\newcommand{\td}[1]{{\bf #1}}
\newcommand{\vc}[1]{\boldsymbol{#1}}
\newcommand{\vd}[1]{{\bf #1}}
\newcommand{\ts}[1]{\mathsfi{#1}}
\newcommand{\ii}{\mathrm{i}}


\DeclareRobustCommand{\solidlegend}{%
    \tikz[baseline=-0.6ex]\draw[line width=0.5pt] (0,0) -- (0.25,0);%
}
\DeclareRobustCommand{\dashlegend}{%
    \tikz[baseline=-0.6ex]\draw[line width=0.5pt,dashed] (0,0) -- (0.25,0);%
}
\DeclareRobustCommand{\dottedlegend}{%
    \tikz[baseline=-0.6ex]\draw[line width=0.5pt,dotted] (0,0) -- (0.25,0);%
}
\DeclareRobustCommand{\circlegend}{%
    \tikz[baseline=-0.6ex]\draw[line width=0.5pt] (0.45,0) circle (0.11);%
}
\DeclareRobustCommand{\circlegendfilled}{%
    \tikz[baseline=-0.6ex]\filldraw[line width=0.5pt] (0.45,0) circle (0.11);%
}
\DeclareRobustCommand{\squarelegend}{%
    \tikz[baseline=-0.6ex]\draw[line width=0.5pt] (0.34,-0.11) rectangle (0.56,0.11);%
}

\newcommand{\diff}[2]{\frac{\partial #1}{\partial #2}}
\newcommand{\difftwo}[2]{\frac{\partial^2 #1}{\partial #2^2}}
\newcommand{\difft}[2]{\frac{\mathrm{d} #1}{\mathrm{d} #2}}
\newcommand{\diffp}[2]{\frac{\partial \left( #1 \right)}{\partial #2}}
\newcommand{\divv}[1]{\nabla\!\cdot\!\mathbf{#1}}
\newcommand{\mean}[1]{\left\langle #1 \right\rangle}
\newcommand{\Gm}{\mathscr{G}^{\mathrm{mean}}}
\newcommand{\kstar}{k_\ast(t)}
\newcommand{\sigmastar}{\overline{\sigma}_\ast(t)}
\newcommand{\Atw}{\mathrm{A_t}}
\newcommand{\Sch}{\mathrm{Sc}}
\newcommand{\figpath}{figures/}
\newcommand{\Frd}{\mathrm{Fr}}
\newcommand{\Pe}{\mathrm{Pe}}
\newcommand{\Amu}{\mathrm{A}_{\mu}}
\newcommand{\CO}{\mathrm{CO}_2}
\newcommand{\bnabla}{\boldsymbol{\nabla}}
\newcommand{\bcdot}{\boldsymbol{\cdot}}
\newcommand{\Rey}{\mathit{Re}}
\newcommand{\Real}{\operatorname{Re}}
\newcommand{\Imag}{\operatorname{Im}}
\DeclareMathOperator{\erf}{erf}
\DeclareMathOperator{\tr}{tr}
\DeclareMathOperator{\LDA}{LDA}


% Figure loading flag: comment out \loadfigtrue for fast caption-only compilation.
\newif\ifloadfig
\loadfigtrue
\newcommand{\figfile}[1]{\ifloadfig\IfFileExists{#1}{\input{#1}}{}\fi}

% Tikz figures are externalised: the first compile turns each picture into
% figures_pdf/<name>.pdf, every later compile just includes that PDF. The
% tracked figures_pdf/README.txt keeps this required directory present in
% GitHub and Overleaf imports; do not remove that marker file.
% Needs shell escape, so compile with
%     pdflatex -shell-escape main.tex
% A figure is rebuilt automatically when its source changes. To force the lot,
% delete figures_pdf/ or run  pdflatex -shell-escape -jobname ...
% Comment out \tikzexternalize to send everything back through pgfplots.
\usetikzlibrary{external}
\tikzexternalize[prefix=figures_pdf/]

% \tikzfig{name} draws figures/name.tex and pins its cached PDF to that name
\newcommand{\tikzfig}[1]{\tikzsetnextfilename{#1}\input{figures/#1.tex}}

\title{\bf On the diffusion of Rayleigh--Taylor Instability}
\author[1]{\bf Marildo Kola\corresp{\email{marildo@umich.edu}}}
\author[2,3]{\bf Daniel Israel}
\author[1]{\bf Aaron Towne}
\affil[1]{\normalsize Department of Mechanical Engineering, University of Michigan, Ann Arbor, MI, USA}
\affil[2]{\normalsize Los Alamos National Laboratory Michigan SPARC, Ann Arbor, MI, USA}
\affil[3]{\normalsize Department of Aerospace Engineering, University of Michigan, Ann Arbor, MI, USA \vspace{-1cm}}
\date{}
\setcounter{Maxaffil}{0}

\begin{document}
\maketitle


\begin{abstract}
    abstract
    \citet{dimonte-ramaprabhu-young-et-al-2005}
\end{abstract}

\section{Introduction}
Understanding transition has been a long-standing problem in fluid mechanics since the pipe-flow experiments of \citet{reynolds-1883} and has continued to vex generations of scientists seeking to unravel the mechanisms leading to turbulence. A cornerstone of this effort has been linear stability theory, in which the governing equations are linearized about a fixed base state and the system is classified as linearly stable or unstable through the eigenspectrum of the resulting operator. This approach has provided the foundation for the study of hydrodynamic instabilities, including the Orr--Sommerfeld problem for parallel shear flows~\citep{orr-1907,sommerfeld-1908} and Taylor--Couette flow between rotating cylinders~\citep{taylor-1923}. These and many other examples are collected in the classical texts of \citet{chandrasekhar-1961} and \citet{drazin-reid-2004}.

Despite its broad success, and its continued extension to complex two- and three-dimensional flows~\citep{theofilis-2011}, the predictive limitations of modal stability analysis became apparent early on, particularly for parallel shear flows. In fact, for plane Couette flow, the system is predicted to remain linearly stable for all Reynolds numbers~\citep{romanov-1973}, even though transition is observed in experiments.


To address this discrepancy, subsequent developments in hydrodynamic stability theory shifted the focus beyond classical eigenvalue analysis toward non-normal stability theory, where the non-orthogonality of the eigenvectors was recognized as a mechanism capable of producing substantial transient amplification, even when all eigenmodes are asymptotically stable~\citep{trefethen-1993}. Disturbances may therefore grow sufficiently to trigger nonlinear effects and transition well before the late-time dynamics predicted by the eigenspectrum become relevant. By identifying the optimal energy amplification attainable over all possible initial perturbations, non-modal stability theory has found considerable success in describing transition in boundary layers~\citep{hoepffner-brandt-henningson-2005,monokrousos-et-al-2010} and pipe flows~\citep{meseguer-trefethen-2003}, and many other set-ups. A comprehensive review of the theoretical framework and applications is provided by \citet{schmid-2007}.

Although many of these analyses considered perturbations evolving about steady base states, the initial-value formulation underlying non-modal stability theory is not restricted to autonomous systems. It extends naturally to flows whose base state evolves in time~\citep{schmid-2007}, enabling transient amplification to be investigated in a broad range of time-dependent configurations. Examples include temporally evolving mixing layers undergoing Kelvin--Helmholtz roll-up and pairing~\citep{arratia-caulfield-chomaz-2013,kaminski-caulfield-taylor-2014}, oscillatory and solitary-wave boundary layers~\citep{biau-2016,verschaeve-pedersen-tropea-2018}, and pulsatile internal flows~\citep{xu-song-avila-2021,pier-schmid-2021}. Nevertheless, optimal-growth analysis remains intrinsically tied to a prescribed initial time and optimization horizon. It therefore identifies the disturbance that maximizes amplification over a particular interval, rather than providing a continuously evolving description of the dominant stability mechanisms along the underlying base-state trajectory.

A framework that addresses this limitation was introduced by \citet{babaee-sapsis-2016} through optimally time-dependent (OTD) modes. The approach employs a minimization principle to construct a finite-dimensional, time-dependent orthonormal basis that evolves along the underlying trajectory and spans the dynamically most unstable directions of the tangent space.

\citet{kern-et-al-2021} applied this method to pulsating Poiseuille flow and identified subharmonic perturbation cycles that would remain hidden if the disturbances were restricted to the periodicity of the base flow, as in Floquet theory~\citep{davis-1976}. A related application was presented by \citet{beneitez-et-al-2023}, who investigated the finite-time stability of a fully three-dimensional, unsteady trajectory in the Blasius boundary layer initiated by the minimal seed~\citep{pringle-kerswell-2010,cherubini-et-al-2011}. By using OTD modes to construct a reduced linearized operator along this trajectory, they showed that meaningful instantaneous stability information such as finite-time Lyapunov exponents and coherent flow structures could be extracted even from a strongly aperiodic base flow. More recently, this strategy has been used to examine the onset of global instability in configurations of direct relevance to external aerodynamics such as the flow over a pitching airfoil~\citep{kern-et-al-2024}.

Despite these promising applications, the OTD framework has several acknowledged limitations~\citep{kern-et-al-2021}. The dimension \(r\) of the OTD subspace must be prescribed \emph{a priori}, and no general criterion is available for determining how many modes should be retained. Excluding relevant stable directions may cause the reduced system to underestimate both non-normality and energy growth, particularly in strongly non-normal flows. Increasing the subspace dimension can alleviate this problem, but also increases the computational cost. Because the method requires, in principle, the solution of a full three-dimensional problem for each mode, its implementation may involve \(r+1\) coupled direct numerical simulations and become prohibitively expensive as the number of retained modes increases~\citep{beneitez-et-al-2023}. A further limitation concerns the initialization of the OTD basis, for which no general procedure is currently available~\citep{kern-et-al-2024}. Consequently, the short-time dynamics may remain strongly dependent on the chosen initialization until the basis aligns with the dynamically dominant subspace.

In summary, classical non-modal growth theory does not provide a continuously evolving description along the base-state trajectory, but only the maximal gain attainable over a prescribed time horizon. OTD modes overcome this restriction, yet introduce computational compromises that may neglect relevant non-normal behaviour and remain strongly dependent on how the basis is initialized, a choice that is itself tied to the initial conditions imposed on the system.

Another fundamental shortcoming of non-normal optimal-growth theory emerges in the context of diffusive instabilities, particularly the Rayleigh--Taylor instability (RTI), which occurs when two superposed fluids of distinct densities are subjected to an acceleration directed against their density stratification~\citep{rayleigh-1883,taylor-1950}, and will be the main focus of this paper. Among the several configurations in which this instability arises, the diffusive RTI encountered during carbon dioxide (\(\mathrm{CO}_2\)) sequestration in deep geological formations~\citep{orr-2009} provides an especially clear example of the limitations of non-modal growth theory. In this physical configuration, dissolved \(\mathrm{CO}_2\) diffuses into the underlying brine, producing a dense diffusive boundary layer above a lighter fluid. Although the resulting stratification is gravitationally unstable, the earliest dynamics remain dominated by molecular diffusion, since the convective motion has not yet developed, and the base state continues to broaden as diffusion proceeds.

\citet{riaz-et-al-2006} analyzed this onset problem using the quasi-steady-state approximation (QSSA), which is commonly employed for this class of problems. In the QSSA, the diffusive base state is frozen at a prescribed time, and a classical eigenvalue analysis is performed with time treated as a parameter. The approximation assumes that the base state evolves much more slowly than the instability and may therefore be regarded as fixed during perturbation growth~\citep{tan-homsy-1986}. This separation of time scales is, however, violated during the early diffusion-dominated stage. Later studies therefore sought to overcome the QSSA, and in particular \citet{rapaka-et-al-2008} used non-modal stability theory to compute the maximum finite-time amplification over all admissible initial perturbations. Despite this attempt, \citet{don-nils-amir-2013} showed that convection is instead more likely to be triggered by suboptimal perturbations than by the mathematically optimal disturbance obtained from the unconstrained amplification problem.

This suggests that mathematical optimality does not necessarily correspond to a physically realizable initial condition.

Meaningful predictions must instead consider the complete initial-value problem initialized with physically relevant disturbances, requiring an assessment of the effect of initial conditions. An instructive example is provided by \citet{daniel-riaz-tchelepi-2015}, who examined how a permeability field varying periodically across the aquifer thickness affects the onset of natural convection. They used ensembles of random initial perturbations localized within the diffusive boundary layer to characterize the resulting distribution of instantaneous growth rates. The ensemble-mean growth agreed closely with that obtained from the boundary-layer-localized perturbation profile of \citet{riaz-et-al-2006}. Nevertheless, there is no guarantee that this agreement would persist under a different correlation structure of the initial perturbations or for fully three-dimensional initial conditions. Consequently, the statistical specification of the initial disturbances must itself be regarded as an essential part of the stability problem.

A related, but distinct, manifestation of this sensitivity was identified by \citet{prathama-pantano-2021} for a baroclinically generated, variable-density shear layer with an analogous diffusively evolving base state \citep{gat-et-al-2017}. Even within a fully non-modal treatment, the early disturbance evolution remained dominated by the diffusive spreading of the base state, with genuine convective growth emerging only at later times. Comparison with direct numerical simulations further showed that the wavenumber assumed to be representative of the disturbance did not drive the observed growth. Instead, the growth was triggered by a subharmonic wavenumber already present in the broadband initial condition.


The available methods for analyzing transient instabilities therefore offer different advantages, but each leaves an important part of the problem unresolved. For transient instabilities, classical modal theory takes the form of the QSSA, and cannot represent the simultaneous evolution of the base state and the perturbations. Non-modal optimal-growth theory retains this evolution and captures transient amplification, but remains tied to a prescribed initial time and optimization horizon, and the resulting optimal perturbation need not be physically realizable or representative of what is actually observed, as demonstrated for both porous-media RTI and a diffusively evolving shear layer. OTD modes overcome the fixed-horizon restriction by continuously tracking the dynamically dominant directions along the trajectory. They nevertheless require a prescribed subspace dimension, depend strongly on the initialization of the basis, and incur a computational cost that increases with the number of retained modes.

It is within this context that we position the present work. We seek a description that retains the non-autonomous evolution of disturbances along a time-dependent base state while explicitly accounting for the statistics of physically admissible initial conditions, rather than isolating a single optimally amplified disturbance. Such a description should provide not only an integrated measure of amplification, but also the spatially resolved statistics needed to explain how the instability mechanism develops as the base state evolves.


The starting point is the work of \citet{frame-towne-2024}, who showed that the large optimal gains commonly invoked to explain transient amplification in non-normal systems correspond to rare events when disturbances are instead drawn from a random ensemble. They further showed that the resulting mean amplification depends on the spatial correlation of the imposed initial conditions. This result motivates a shift from the deterministic question of which individual disturbance experiences the greatest gain to the statistical question of what amplification is expected when the system is perturbed by an ensemble of initial conditions. The literature discussed above suggests that this is a natural perspective from which to analyze transient instabilities. Nevertheless, the original formulation was developed for a fixed base state and left unresolved the physical specification of the initial correlations.

\section{Governing equations}%
\label{sec:governing_equation}
In this work we consider two physical models to describe the Rayleigh--Taylor instability that develops between two fluids of different densities. The first model is the Boussinesq approximation that provides the simplest model for examining the effect of buoyant stratification on the flow. Density variations relative to a reference state are assumed to be small and affect only the gravitational buoyancy term, while the velocity field remains incompressible~\citep{boffetta-mazzino-2017}. The other model that we employ is the incompressible variable-density (IVD) formulation of \citet{sandoval-1995}. It applies in the low-Mach-number regime, where acoustic waves are suppressed and the energy equation decouples from the problem~\citep{livescu-ristorcelli-2007}. Each fluid remains individually incompressible, although the mixture velocity is not divergence-free~\citep{livescu-2013}.

For both models we consider miscible fluids with reference densities \(\rho_1\) and \(\rho_2\)
and dynamic viscosities \(\mu_1\) and \(\mu_2\). We nondimensionalize the governing equations using the reference density
\(\rho_r\), length \(L_r\), velocity \(U_r\), dynamic viscosity \(\mu_r\), gravitational
acceleration \(g\), and mass diffusivity \(\mathcal{D}\). The resulting Reynolds, Froude, and Schmidt numbers
are
\begin{equation}
    \label{eq:Re_Fr_definitions}
    \Rey = \frac{\rho_r L_r U_r}{\mu_r}, \qquad \Frd = \frac{U_r}{\sqrt{g L_r}}, \qquad \Sch = \frac{\mu_r/\rho_r}{\mathcal{D}}.
\end{equation}
The present problem has no externally imposed length or velocity scale, so the reference length and velocity are set by the classical viscous-gravitational scales
of~\citet{chandrasekhar-1961}. With this choice, the Reynolds and Froude numbers reduce to unity, and the reference scales become
\begin{equation}
    \label{eq:chandrasekhar_scaling}
    U_r = (\nu_r g)^{1/3}, \qquad T_r = \left(\frac{\nu_r}{g^{2}}\right)^{1/3}, \qquad L_r = \left(\frac{\nu_r^{2}}{g}\right)^{1/3},
\end{equation}
where \(\nu_r=\mu_r/\rho_r\) is the reference kinematic viscosity. The reference density
and dynamic viscosity are taken as the arithmetic means of the two fluids,
\begin{equation}
    \label{eq:reference_quantities}
    \rho_r = \frac{\rho_1 + \rho_2}{2}, \qquad \mu_r = \frac{\mu_1 + \mu_2}{2},
\end{equation}
and the density and viscosity contrasts are measured by the density Atwood number \(\Atw\) and a corresponding viscosity Atwood number \(A_{\mu}\),
\begin{equation}
    \label{eq:atwood_numbers}
    \Atw = \frac{\rho_2 - \rho_1}{\rho_2 + \rho_1}, \qquad A_{\mu} = \frac{\mu_2 - \mu_1}{\mu_2 + \mu_1}.
\end{equation}
We consider two models of the Rayleigh--Taylor instability. The first is the Boussinesq formulation presented in \cref{sec:bq_equation}. It is valid for weak stratification and provides the simplest setting in which to develop the statistical approach, interpret its predictions, and validate it against three-dimensional direct numerical simulations. The second is the full variable-density formulation presented in \cref{sec:vd_equation}, which retains non-Boussinesq effects throughout the transient instability.

\subsection{Boussinesq approximation}
\label{sec:bq_equation}
In this model, the density is taken to vary linearly with the local mixture concentration, and the governing equations are
\begin{subequations}
    \label{eq:bq_nl_continuous}
    \begin{align}
        \label{eq:bq_nl_continuous_momentum}
        \diff{\boldsymbol{u}}{t} + (\boldsymbol{u}\bcdot\bnabla)\boldsymbol{u} &= -\bnabla p + \nabla^{2}\boldsymbol{u} - c\,\boldsymbol{e}_{z}, \\
        \label{eq:bq_nl_continuous_scalar}
        \diff{c}{t} + \boldsymbol{u}\bcdot\bnabla c &= \Sch^{-1}\nabla^{2}c, \\
        \label{eq:bq_nl_continuous_continuity}
        \bnabla\bcdot\boldsymbol{u} &= 0,
    \end{align}
\end{subequations}
where \(\boldsymbol{u}(\boldsymbol{x},t)\) is the velocity field, \(p(\boldsymbol{x},t)\) is the pressure field, and \(c(\boldsymbol{x},t)\) is the scalar concentration field. The viscosity is taken constant throughout the mixture by setting \(A_{\mu}=0\); the effect of viscosity stratification is deferred to the variable-density model of \cref{sec:vd_equation}, which recovers the Boussinesq approximation in the limit of weak stratification.

Introducing the state vector \(\boldsymbol{Q}(\boldsymbol{x},t) = (u, v, w, p, c)^\top\), the nonlinear system can be written compactly as
\begin{equation}
    \label{eq:compact_nl_continuous}
    \mathscr{M}\,\diff{\boldsymbol{Q}}{t} = \mathcal{N}\!\left( \boldsymbol{Q} \right),
\end{equation}
where \(\mathscr{M}\) is the identity matrix with the entry corresponding to pressure set to zero, since the pressure is not dynamically evolved but instead enforces the incompressibility constraint.
We introduce the perturbation expansion
\begin{equation}
    \label{eq:perturbation_expansion}
    \boldsymbol{Q}(\boldsymbol{x},t) = \boldsymbol{Q}_0(\boldsymbol{x},t) + \epsilon\,\boldsymbol{Q}'(\boldsymbol{x},t) + O(\epsilon^2),
    \qquad 0 < \epsilon \ll 1,
\end{equation}
where \(\boldsymbol{Q}_0\) is a solution of the nonlinear equations, hereafter referred to as the base state, and \(\boldsymbol{Q}'\) is the perturbation about this base state. Substituting this expansion into the compact form and collecting terms of equal order in \(\epsilon\), we obtain the governing equations
\begin{subequations}
    \label{eq:compact_lin_continuous}
    \begin{align}
        \label{eq:compact_lin_continuous_0th}
        \mathscr{M}\,\diff{\boldsymbol{Q}_0}{t} &= \mathcal{N}\!\left( \boldsymbol{Q}_0 \right), \\
        \label{eq:compact_lin_continuous_1th}
        \mathscr{M}\,\diff{\boldsymbol{Q}'}{t} &= \mathcal{L}(\boldsymbol{Q}_0)\,\boldsymbol{Q}',
    \end{align}
\end{subequations}
where \(\mathcal{L}\) is the Jacobian of the nonlinear Boussinesq operator. Since no intrinsic length or time scale is imposed on the system, the base-state evolution is assumed to depend only on the vertical coordinate and time~\citep{prathama-pantano-2021}, leading to \(\boldsymbol{u}_0=0\) and
\begin{equation}
    \label{eq:bq_base_state}
    c_0(z,t) = \frac{1}{2}\left[1-\erf\!\left(\frac{z}{\delta(t)}\right)\right], \qquad \delta(t) = \sqrt{4t/\Sch}.
\end{equation}
The Boussinesq linearized equations can be further simplified by eliminating the pressure term. This can be done easely by taking the difference between the vertical derivative of the poisson equation and the laplacian of the vertical momentum to get rid of the pressure term and exploting the zero divergence constraint. The resulting system for the reduced perturbation state \(\boldsymbol{q}=(w,c)^{\top}\) is
\begin{equation}
    \label{eq:bq_lin_continuous}
    \begin{pmatrix}
        \nabla^{2} & 0 \\
        0 & 1
    \end{pmatrix}
    \begin{pmatrix}
        \partial_t\, w \\[4pt]
        \partial_t\, c
    \end{pmatrix}
    =
    \begin{pmatrix}
        \nabla^{4} & \nabla_{xy}^{2} \\[4pt]
        -\partial_{z}c_{0} & \Sch^{-1}\nabla^{2}
    \end{pmatrix}
    \begin{pmatrix}
        w \\[4pt]
        c
    \end{pmatrix}.
\end{equation}
The boundary conditions enforce vanishing vertical velocity and its vertical gradient at the
domain boundaries, and require the concentration perturbation to vanish in the far field,
\begin{equation}
    \label{eq:bq_boundary_conditions}
    w(\pm L_z) = 0, \qquad \left.\frac{\partial w}{\partial z}\right|_{\pm L_z} = 0, \qquad
    c(\pm L_z) = 0.
\end{equation}
Assuming homogeneity in the horizontal directions \((x,y)\), the perturbation fields are decomposed into Fourier modes. Fourier-transforming the reduced system then yields, for the
state vector \(\hat{\boldsymbol{q}}=(\hat{w},\hat{c})^{\top}\),
\begin{equation}
    \label{eq:lti_dz_k_boussinesq}
    \begin{pmatrix}
        \partial_{zz}-k^{2} & 0 \\
        0 & 1
    \end{pmatrix}
    \begin{pmatrix}
        \partial_t\, \hat{w} \\[4pt]
        \partial_t\, \hat{c}
    \end{pmatrix}
    =
    \begin{pmatrix}
        \left(\partial_{zz}-k^{2}\right)^{2} & -k^{2} \\[4pt]
        -\partial_{z}c_{0} & \Sch^{-1}\left(\partial_{zz}-k^{2}\right)
    \end{pmatrix}
    \begin{pmatrix}
        \hat{w} \\[4pt]
        \hat{c}
    \end{pmatrix},
\end{equation}
where \(\hat{w}(k,z,t)\) and \(\hat{c}(k,z,t)\) are the Fourier-transformed velocity and
concentration perturbations. After discretizing the inhomogeneous vertical direction, the state associated with a given wavenumber is denoted by \(\hat{\mathbf{q}}_{k}(t)\in\mathbb{C}^{2N_z}\), where \(N_z\) is the number of degrees of freedom. The continuous differential operators then become the time-dependent matrices \(\mathsfbi{A}_{k}(t)\) and \(\mathsfbi{B}_{k}(t)\), and the linearized dynamics take the form
\begin{equation}
    \label{eq:semidiscrete_system}
    \mathsfbi{B}_{k}(t)\,\difft{\hat{\mathbf{q}}_{k}}{t}
    =
    \mathsfbi{A}_{k}(t)\,\hat{\mathbf{q}}_{k},
    \qquad
    \hat{\mathbf{q}}_{k}(0)=\hat{\mathbf{q}}_{0,k}.
\end{equation}
This semi-discrete initial-value problem defines the discrete propagator
\(\boldsymbol{\Phi}_{k}(t,0)\) through
\begin{equation}
    \label{eq:discrete_propagator}
    \hat{\mathbf{q}}_{k}(t)
    =
    \boldsymbol{\Phi}_{k}(t,0)\,\hat{\mathbf{q}}_{0,k},
    \qquad
    \boldsymbol{\Phi}_{k}(0,0)=\mathsfbi{I}.
\end{equation}
The two-argument form \(\boldsymbol{\Phi}_{k}(t,0)\) is cumbersome for sustained use. We therefore adopt the condensed notation
\begin{equation}
    \label{eq:propagator_shorthand}
    \boldsymbol{\Phi}_{t,k} \equiv \boldsymbol{\Phi}_{k}(t,0),
\end{equation}
so that the discrete propagator takes the concise form \(\hat{\mathbf{q}}_{k}(t)=\boldsymbol{\Phi}_{t,k}\,\hat{\mathbf{q}}_{0,k}\). This shorthand is used throughout the remainder of this work.

\subsection{Variable-density formulation}
\label{sec:vd_equation}
For the IVD, with the nondimensionalization introduced above, the governing equations are
\begin{subequations}
    \label{eq:vd_nl_continuous}
    \begin{align}
        \label{eq:vd_nl_continuous_continuity}
        \partial_t \rho + \bnabla\bcdot(\rho \boldsymbol{u}) &= 0, \\
        \label{eq:vd_nl_continuous_momentum}
        \rho\,\partial_t \boldsymbol{u} + \rho\,(\boldsymbol{u}\bcdot\bnabla)\boldsymbol{u} &= -\,\bnabla p + \bnabla\bcdot\boldsymbol{\tau} - \rho\,\boldsymbol{e}_z, \\
        \label{eq:vd_nl_continuous_divergence}
        \bnabla \bcdot \boldsymbol{u} &= -\Sch^{-1}\,\nabla^2 \log\rho,
    \end{align}
\end{subequations}
where \(\rho(\boldsymbol{x},t)\) is the mixture density, \(\boldsymbol{u}(\boldsymbol{x},t)\) is the
velocity field, \(p(\boldsymbol{x},t)\) is the pressure, and \(\boldsymbol{e}_z\) points
opposite to gravity. The viscous stress tensor \(\boldsymbol{\tau}\) is given by
\begin{equation}
    \label{eq:vd_stress_tensor}
    \boldsymbol{\tau} =  2\,\mu\,\mathsfbi{S} - \frac{2}{3}\mu\,(\bnabla\bcdot\boldsymbol{u})\mathsfbi{I}, \qquad 2\,\mathsfbi{S} = \bnabla\boldsymbol{u} + (\bnabla\boldsymbol{u})^{\!\top},
\end{equation}
where \(\mu\) is the dynamic viscosity, \(\mathsfbi{S}\) is the strain-rate tensor, and
\(\mathsfbi{I}\) is the identity tensor.

Introducing the state vector \(\boldsymbol{Q}(\boldsymbol{x},t) = (u, v, w, p, \rho)^\top\), the nonlinear system can be written compactly as
\begin{equation}
    \label{eq:vd_compact_nl_continuous}
    \mathscr{M}\,\diff{\boldsymbol{Q}}{t} = \mathcal{N}\!\left( \boldsymbol{Q} \right),
\end{equation}
where \(\mathscr{M}\) is again the identity matrix with the entry corresponding to \(p\) set to zero, since the pressure is not dynamically evolved but instead enforces the divergence constraint.
We introduce the perturbation expansion
\begin{equation}
    \label{eq:vd_perturbation_expansion}
    \boldsymbol{Q}(\boldsymbol{x},t) = \boldsymbol{Q}_0(\boldsymbol{x},t) + \epsilon\,\boldsymbol{Q}'(\boldsymbol{x},t) + O(\epsilon^2),
    \qquad 0 < \epsilon \ll 1,
\end{equation}
where \(\boldsymbol{Q}_0\) is a solution of the nonlinear equations, hereafter referred to as the base state, and \(\boldsymbol{Q}'\) is the perturbation about this base state. Substituting this expansion into the compact form and collecting terms of equal order in \(\epsilon\), we obtain
\begin{subequations}
    \label{eq:vd_compact_lin_continuous}
    \begin{align}
        \label{eq:vd_compact_lin_continuous_0th}
        \mathscr{M}\,\diff{\boldsymbol{Q}_0}{t} &= \mathcal{N}\!\left( \boldsymbol{Q}_0 \right), \\
        \label{eq:vd_compact_lin_continuous_1th}
        \mathscr{M}\,\diff{\boldsymbol{Q}'}{t} &= \mathcal{L}(\boldsymbol{Q}_0)\,\boldsymbol{Q}',
    \end{align}
\end{subequations}
where \(\mathcal{L}\) is the Jacobian of the nonlinear variable-density Navier--Stokes operator.

The base-state evolution and the linearized system were fully derived and characterized by \citet{kola-israel-towne-2026}; we report here only the results needed for the present analysis. The base state is a one-dimensional flow \(\boldsymbol{U}_0 = (0, 0, w_0)\) with density \(\rho_0\),
\begin{subequations}\label{eq:vd_base_state}
\begin{align}
    \rho_{0}(z,t) &= 1 + \Atw\,\erf\!\left(\frac{z}{\delta}\right), \label{eq:vd_base_state_rho0} \\[4pt]
    w_0(z,t) &= \frac{2\,\Atw}{\sqrt{\pi}\,\delta}\,
               \frac{\exp\!\left(-z^2/\delta^2\right)}{\rho_0(z,t)},
               \label{eq:vd_base_state_w0}
\end{align}
\end{subequations}
with diffusive layer thickness \(\delta(t)=\sqrt{4t/\Sch}\). Linearizing about this base state and reducing to the two-component perturbation state \(\boldsymbol{q} = (w, \rho)^\top\) gives
\begin{equation}
    \label{eq:vd_lin_continuous}
    \begin{pmatrix}
        \mathscr{B}_{ww} & \mathscr{B}_{w\rho} \\
        0 & 1
    \end{pmatrix}
    \begin{pmatrix}
        \partial_t\, w \\[4pt]
        \partial_t\,\rho
    \end{pmatrix}
    =
    \begin{pmatrix}
        \mathscr{A}_{w_0w}+\mathscr{A}_{ww} & \mathscr{A}_{w_0\rho}+\mathscr{A}_{w\rho} \\[4pt]
        \mathscr{A}_{\rho w} & \mathscr{A}_{\rho\rho}
    \end{pmatrix}
    \begin{pmatrix}
        w \\[4pt]
        \rho
    \end{pmatrix},
\end{equation}
where the operator blocks \(\mathscr{A}\) and \(\mathscr{B}\) are given in \Cref{app:operators}. The boundary conditions enforce vanishing vertical velocity and its vertical gradient at the domain boundaries, and require the density perturbation to vanish in the far field, so that the base-state stratification remains undisturbed away from the diffusive layer,
\begin{equation}
    \label{eq:vd_boundary_conditions}
    w(\pm L_z) = 0, \qquad \left.\frac{\partial w}{\partial z}\right|_{\pm L_z} = 0, \qquad
    \rho(\pm L_z) = 0.
\end{equation}
Assuming homogeneity in the horizontal directions \((x,y)\), the perturbation fields are decomposed into Fourier modes, and for the state vector \(\hat{\boldsymbol{q}}=(\hat{w},\hat{\rho})^{\top}\) the reduced system becomes
\begin{equation}
    \label{eq:vd_lti_dz_k}
    \begin{pmatrix}
        \widehat{\mathscr{B}}_{ww} & \widehat{\mathscr{B}}_{w\rho} \\
        0 & 1
    \end{pmatrix}
    \begin{pmatrix}
        \partial_t\, \hat{w} \\[4pt]
        \partial_t\,\hat{\rho}
    \end{pmatrix}
    =
    \begin{pmatrix}
        \widehat{\mathscr{A}}_{w_0w}+\widehat{\mathscr{A}}_{ww} & \widehat{\mathscr{A}}_{w_0\rho}+\widehat{\mathscr{A}}_{w\rho} \\[4pt]
        \widehat{\mathscr{A}}_{\rho w} & \widehat{\mathscr{A}}_{\rho\rho}
    \end{pmatrix}
    \begin{pmatrix}
        \hat{w} \\[4pt]
        \hat{\rho}
    \end{pmatrix},
\end{equation}
where \(\hat{w}(k,z,t)\) and \(\hat{\rho}(k,z,t)\) are the Fourier-transformed vertical velocity and density perturbations. Discretizing the inhomogeneous vertical direction with \(N_z\) points, the state at a given wavenumber is \(\hat{\mathbf{q}}_{k}(t)\in\mathbb{C}^{2N_z}\), the operators become the time-dependent matrices \(\mathsfbi{A}_{k}(t)\) and \(\mathsfbi{B}_{k}(t)\), and the linearized dynamics take the form
\begin{equation}
    \label{eq:vd_semidiscrete_system}
    \mathsfbi{B}_{k}(t)\,\difft{\hat{\mathbf{q}}_{k}}{t}
    =
    \mathsfbi{A}_{k}(t)\,\hat{\mathbf{q}}_{k},
    \qquad
    \hat{\mathbf{q}}_{k}(0)=\hat{\mathbf{q}}_{0,k},
\end{equation}
which defines the discrete propagator \(\boldsymbol{\Phi}_{t,k}\) through \(\hat{\mathbf{q}}_{k}(t)=\boldsymbol{\Phi}_{t,k}\,\hat{\mathbf{q}}_{0,k}\), with the shorthand adopted previously. Concerning the numerical solution, see \citet{kola-israel-towne-2026}, which employs spectral collocation methods with the nonlinear transformation used to cluster points within the diffusive layer.

\section{Statistics of the initial perturbations}
\label{sec:statistics_of_the_initial_perturbations}

Having introduced the governing equations in \cref{sec:governing_equation}, we now turn our attention to the statistical description of the initial conditions. Our objective is to relate their spatial covariance to a physically realizable perturbation that can be imposed in a controlled experiment or numerical simulation. In fact, one contribution made here is to overcome an apparent limitation in the formulation introduced by \citet{frame-towne-2024} for Poiseuille flow, where the initial covariance was treated primarily as a mathematical model and did not necessarily correspond to a directly measurable experimental quantity. Here, instead, we show that the methodology is considerably more general and that the initial correlation has a precise physical meaning.

In physical space, the covariance of the initial perturbation field is
\begin{equation}
    \mathsfbi{C}_{00}(\boldsymbol{x},\boldsymbol{r})
    =
    \mean{
        \boldsymbol{q}'(\boldsymbol{x}),
        \boldsymbol{q}'(\boldsymbol{x}+\boldsymbol{r})
    }
    =
    \begin{pmatrix}
        C_{00}^{ww}(\boldsymbol{x},\boldsymbol{r})
        &
        C_{00}^{w\rho}(\boldsymbol{x},\boldsymbol{r})
        \\[4pt]
        C_{00}^{\rho w}(\boldsymbol{x},\boldsymbol{r})
        &
        C_{00}^{\rho\rho}(\boldsymbol{x},\boldsymbol{r})
    \end{pmatrix}.
    \label{eq:physical_covariance_general}
\end{equation}
For the initial conditions that we wish to describe here, we assume that they can be represented through a perturbation of the interface separating the two fluids \citep{cook-dimotakis-2001,wei-livescu-2012}, and that they are homogeneous and isotropic. In principle, this is not the only possibility, since the disturbance may also be introduced through the vertical velocity, as discussed by \citet{mueschke-schilling-2009}, and need not necessarily be isotropic \citep{dalziel-linden-youngs-1999}. Even though these assumptions do not encompass all possible initial conditions that may arise in experiments or numerical initializations, they still provide a physically relevant setting that captures the fundamental nature of the problem, as reflected by the large number of studies that have adopted them \citep{kord-capecelatro-2019}; moreover, the framework that we develop can readily accommodate more general prescriptions. The hypotheses we adopt are intended solely to keep the exposition clear from additional cumbersome notations, without sacrificing the physical interpretation of the prescribed statistics of the initial conditions. Under these assumptions, the initial vertical-velocity perturbation is set to zero, so that all covariance components involving \(w\) vanish and only the density-density component remains nonzero, namely
\begin{equation}
    \mathsfbi{C}_{00}(\boldsymbol{x},\boldsymbol{r})
    =
    \mean{
        \boldsymbol{q}'(\boldsymbol{x}),
        \boldsymbol{q}'(\boldsymbol{x}+\boldsymbol{r})
    }
    =
    \begin{pmatrix}
        0 & 0
        \\[4pt]
        0 & C_{00}^{\rho\rho}(\boldsymbol{x},\boldsymbol{r})
    \end{pmatrix}.
    \label{eq:physical_covariance_density_only}
\end{equation}
Denoting the horizontal coordinate by \(\boldsymbol{x}_h=(x,y)\) and the local displacement of the interface by \(\eta(\boldsymbol{x}_h)\), the displaced base density profile admits the expansion
\begin{equation}
    \rho_0(z-\eta)
    =
    \rho_0(z)
    -
    \rho_{0,z}(z)\,\eta(\boldsymbol{x}_h)
    +
    O\!\left(\eta^2\right).
    \label{eq:initial_density_displacement}
\end{equation}
Consequently, at leading order, the initial density perturbation is determined directly by the local displacement of the base density profile according to
\begin{equation}
    \rho'(\boldsymbol{x},0)
    =
    -
    \rho_{0,z}(z)\,\eta(\boldsymbol{x}_h).
    \label{eq:initial_density_perturbation_eta}
\end{equation}
Substitution of \eqref{eq:initial_density_perturbation_eta} into the density covariance then gives
\begin{equation}
    C_{00}^{\rho\rho}(\boldsymbol{x}_h,\boldsymbol{r}_h,z,z')
    =
    \underbrace{
        \rho_{0,z}(z)\,\rho_{0,z}(z')
    }_{
        C_{00}^{z}(z,z')
    }
    \underbrace{
        \mean{
            \eta(\boldsymbol{x}_h)
            \eta(\boldsymbol{x}_h+\boldsymbol{r}_h)
        }
    }_{
        C_{00}^{xy}(\boldsymbol{x}_h,\boldsymbol{r}_h)
    }.
    \label{eq:initial_density_covariance_separable}
\end{equation}
Equation \eqref{eq:initial_density_covariance_separable} separates the initial density covariance into two physically distinct contributions. The vertical factor is determined entirely by the base state, whereas the horizontal factor contains the prescribed statistics of the interface displacement. In particular, \(C_{00}^{z}\) is fixed by \(\rho_{0,z}\) and retains the inhomogeneity of the diffusive layer, so that the initial perturbations are localized within this layer and their vertical correlation length is associated with its initial thickness. For the diffusive profile considered here, the vertical covariance factor is
\begin{equation}
    C_{00}^{z}(z,z')
    =
    \frac{4\,\Atw^{2}}{\pi\,\delta_0^{2}}
    \exp\!\left(
        -\frac{z^2+{z'}^2}{\delta_0^2}
    \right).
    \label{eq:vertical_covariance_factor_erf}
\end{equation}
The horizontal correlation \(C_{00}^{xy}\), on the other hand, carries the statistical information associated with the displacement of the interface and, although this correlation may be prescribed directly in physical space, perturbations imposed experimentally or numerically are often characterized more naturally through their spectral content; this is where we draw a direct connection between the statistical description of the perturbations and the classical spectra used to initialize direct numerical simulations \citep{rogallo-1981}.


Since the perturbation field is homogeneous in the horizontal plane, we write
\begin{equation}
    \eta(\boldsymbol{x}_h)
    =
    \int_{\mathbb{R}^2}
        \hat{\eta}(\boldsymbol{k}_h)
        \exp(i\boldsymbol{k}_h\bcdot\boldsymbol{x}_h)
    \,\mathrm{d}\boldsymbol{k}_h,
    \qquad
    \boldsymbol{k}_h=(k_x,k_y).
    \label{eq:eta_fourier_transform}
\end{equation}
The corresponding spectral density \(\Phi_{00}^{xy}\) is defined by
\begin{equation}
    \mean{
        \hat{\eta}(\boldsymbol{k}_h)
        \hat{\eta}^{\ast}(\boldsymbol{k}'_h)
    }
    =
    \Phi_{00}^{xy}(\boldsymbol{k}_h)
    \delta(\boldsymbol{k}_h-\boldsymbol{k}'_h),
    \label{eq:eta_spectral_density_definition}
\end{equation}
so that the physical-space correlation is obtained as the Fourier transform of the spectral density,
\begin{equation}
    C_{00}^{xy}(\boldsymbol{r}_h)
    =
    \int_{\mathbb{R}^2}
        \Phi_{00}^{xy}(\boldsymbol{k}_h)
        \exp(i\boldsymbol{k}_h\bcdot\boldsymbol{r}_h)
    \,\mathrm{d}\boldsymbol{k}_h.
    \label{eq:eta_correlation_spectral_density}
\end{equation}
Invoking the isotropy assumption introduced above, the correlation depends only on \(r=\lvert\boldsymbol{r}_h\rvert\), while the spectral density depends only on \(k=\lvert\boldsymbol{k}_h\rvert\), and therefore
\begin{equation}
    C_{00}^{xy}(r)
    =
    2\pi
    \int_0^\infty
        \Phi_{00}^{xy}(k)
        J_0(kr)\,
        k
    \,\mathrm{d}k,
    \label{eq:eta_correlation_hankel_phi}
\end{equation}
where
\begin{equation}
    J_0(kr)
    =
    \frac{1}{2\pi}
    \int_0^{2\pi}
        \exp(i k r\cos\theta)
    \,\mathrm{d}\theta
    \label{eq:bessel_j0}
\end{equation}
is the Bessel function of the first kind of order zero. The two-dimensional isotropic spectrum of the interface displacement is then defined as the circularly integrated spectral density
\begin{equation}
    E_{\eta}(k)
    =
    \pi k\,\Phi_{00}^{xy}(k),
    \label{eq:eta_2d_spectrum_definition}
\end{equation}
with the normalization chosen such that
\begin{equation}
    \mean{\eta_0^2}
    =
    2
    \int_0^\infty
        E_{\eta}(k)
    \,\mathrm{d}k.
    \label{eq:eta_variance_spectrum}
\end{equation}
It is immediate to show that the in-plane correlation and the isotropic spectrum form the Hankel transform pair
\begin{equation}
    C_{00}^{xy}(r)
    =
    2
    \int_0^\infty
        E_{\eta}(k)
        J_0(kr)
    \,\mathrm{d}k,
    \qquad
    E_{\eta}(k)
    =
    \frac{k}{2}
    \int_0^\infty
        C_{00}^{xy}(r)
        J_0(kr)\,
        r
    \,\mathrm{d}r.
    \label{eq:eta_hankel_transform_pair}
\end{equation}
This equivalence provides the desired connection between the two descriptions, since prescribing the correlation \(C_{00}^{xy}(r)\) in physical space is equivalent to prescribing the spectrum \(E_{\eta}(k)\) in Fourier space, with both representing the same random displacement of the interface. Consequently, modelling the spatial correlation is equivalent to imposing the corresponding spectrum when initializing the perturbations in a numerical simulation or controlled experiment.

\begin{figure}
    \centering
    \begin{minipage}{0.48\textwidth}
        \centering
        % This file was created by matlab2tikz.
%
\definecolor{mycolor1}{rgb}{0.48200,0.16500,0.51400}%
\definecolor{mycolor2}{rgb}{0.10600,0.47100,0.21600}%
\definecolor{mycolor3}{rgb}{0.68600,0.55300,0.76500}%
%
\begin{tikzpicture}[trim axis left,trim axis right]

\begin{axis}[%
width=11.252in,
height=7.131in,
at={(1.887in,0.962in)},
scale only axis,
clip=true,
separate axis lines,
every outer x axis line/.append style={black},
every x tick label/.append style={font=\color{black}},
every x tick/.append style={black},
xmin=0,
xmax=8,
xlabel={$r/\lambda_0$},
every outer y axis line/.append style={black},
every y tick label/.append style={font=\color{black}},
every y tick/.append style={black},
ymin=-0.4,
ymax=1,
ylabel={$C_{00,n}^{xy}(r)$},
axis background/.style={fill=white},
legend style={legend cell align=left, align=left},
clip=true,
clip mode=individual,
axis lines=box,
xtick pos=bottom,
ytick pos=left,
tick align=inside,
major tick length=2pt,
width=0.7\linewidth,
height=0.65\linewidth,
every axis/.append style={font=\fontsize{9}{9}\selectfont},xlabel style={font=\fontsize{9}{9}\selectfont},title style={font=\fontsize{9}{9}\selectfont},
legend style={font=\fontsize{8}{9}\selectfont,inner sep=1pt,row sep=1pt,column sep=2pt,nodes={scale=1.00},draw=none},legend image post style={xscale=0.35},
legend columns=1,
scaled x ticks=false,
tick label style={/pgf/number format/fixed,/pgf/number format/precision=2},
every x tick label/.append style={font=\fontsize{9}{9}\selectfont\color{black}},
every y tick label/.append style={font=\fontsize{9}{9}\selectfont\color{black}},
,,
colormap={sigmamap}{rgb(0.0000)=(0.0200,0.1880,0.3800); rgb(0.1000)=(0.0743,0.3456,0.5616); rgb(0.2000)=(0.1918,0.4980,0.7060); rgb(0.3000)=(0.3890,0.6708,0.8226); rgb(0.4000)=(0.7552,0.8682,0.9228); rgb(0.5000)=(0.9690,0.9689,0.9689); rgb(0.6000)=(0.9425,0.8044,0.7614); rgb(0.7000)=(0.8767,0.4910,0.4020); rgb(0.8000)=(0.7869,0.2866,0.2470); rgb(0.9000)=(0.6186,0.1239,0.1568); rgb(1.0000)=(0.4040,0.0000,0.1220)}
]
\addplot [color=mycolor1, line width=1.4pt]
  table[row sep=crcr]{%
1e-06	0.999999999999\\
0.0400210050025012	0.998399601164802\\
0.0800410100050025	0.993613914988808\\
0.120061015007504	0.985688746313652\\
0.160081020010005	0.974699624420412\\
0.200101025012506	0.960750604625268\\
0.240121030015008	0.943972627966267\\
0.280141035017509	0.924521475140297\\
0.32016104002001	0.902575359394744\\
0.360181045022511	0.878332210368438\\
0.400201050025012	0.852006706724764\\
0.440221055027514	0.823827119687516\\
0.484243060530265	0.79097308236634\\
0.524263065532766	0.759684729113178\\
0.564283070535268	0.727300616844434\\
0.604303075537769	0.6940701748391\\
0.64432308054027	0.660239763357353\\
0.684343085542771	0.626049737144811\\
0.724363090545273	0.59173174103289\\
0.764383095547774	0.557506276865767\\
0.804403100550275	0.523580573109958\\
0.844423105552776	0.490146780255911\\
0.884443110555278	0.457380506752258\\
0.924463115557779	0.425439701962636\\
0.96448312056028	0.39446388472475\\
1.00450312556278	0.364573708720333\\
1.04452313056528	0.335870849201984\\
1.08454313556778	0.308438189804935\\
1.12456314057028	0.282340283299189\\
1.16458314557279	0.257624056274803\\
1.20460315057529	0.234319724928809\\
1.24462315557779	0.212441887331191\\
1.28464316058029	0.191990756752799\\
1.32466316558279	0.172953500776418\\
1.36868517108554	0.153616030642925\\
1.40870517608804	0.137456156692945\\
1.44872518109055	0.122602894026061\\
1.48874518609305	0.109004923383705\\
1.52876519109555	0.096605170968143\\
1.56878519609805	0.0853421343566971\\
1.60880520110055	0.0751511263337177\\
1.64882520610305	0.0659654235207782\\
1.68884521110555	0.0577173101953688\\
1.72886521610805	0.0503390110203401\\
1.76888522111056	0.0437635094916467\\
1.80890522611306	0.0379252516972433\\
1.84892523111556	0.0327607374274328\\
1.88894523611806	0.0282090027635224\\
1.92896524112056	0.0242119999879278\\
1.96898524612306	0.0207148820077474\\
2.00900525112556	0.0176661994786927\\
2.04902525612806	0.0150180194791661\\
2.08904526113057	0.012725974944051\\
2.12906526613307	0.0107492541581323\\
2.16908527113557	0.00905053946678608\\
2.20910527613807	0.00759590402492638\\
2.25312728164082	0.00624120162348047\\
2.29314728664332	0.00520297616216849\\
2.33316729164582	0.00432358823816512\\
2.37318729664832	0.00358134108139998\\
2.41320730165083	0.00295703121831066\\
2.45322730665333	0.0024337446047919\\
2.49324731165583	0.00199665468359556\\
2.53326731665833	0.00163282555159215\\
2.57328732166083	0.00133102275996024\\
2.61330732666333	0.00108153366651469\\
2.65332733166583	0.000875998719344083\\
2.69334733666833	0.000707254577226614\\
2.73336734167084	0.000569189565086317\\
2.77338734667334	0.000456611620243819\\
2.81340735167584	0.000365128604090978\\
2.85342735667834	0.000291040629623266\\
2.89344736168084	0.000231243882767755\\
2.93346736668334	0.000183145288956025\\
2.97348737168584	0.000144587290025496\\
3.01350737668834	0.000113781944435656\\
3.05352738169085	8.92535403114787e-05\\
3.09354738669335	6.97889106716043e-05\\
3.1375693921961	5.3046498924736e-05\\
3.1775893971986	4.1199914486186e-05\\
3.2176094022011	3.18966278072457e-05\\
3.2576294072036	2.46151277540626e-05\\
3.2976494122061	1.89351297201011e-05\\
3.3376694172086	1.45192218940205e-05\\
3.37768942221111	1.1097554015231e-05\\
3.41770942721361	8.45512553019439e-06\\
3.45772943221611	6.42128164882296e-06\\
3.49774943721861	4.8610741519019e-06\\
3.53776944222111	3.66818845661831e-06\\
3.57778944722361	2.75917915890061e-06\\
3.61780945222611	2.06879295846269e-06\\
3.65782945722861	1.54619058659399e-06\\
3.69784946223112	1.15190824727759e-06\\
3.73786946723362	8.55424368192136e-07\\
3.77788947223612	6.33219404399189e-07\\
3.81790947723862	4.6723533522549e-07\\
3.85792948224112	3.43657646102901e-07\\
3.89794948724362	2.51956292962435e-07\\
3.93796949224612	1.84133698326139e-07\\
3.97798949724862	1.34137500465833e-07\\
4.02201150275138	9.43196276680939e-08\\
4.06203150775388	6.82492105497248e-08\\
4.10205151275638	4.92268508511979e-08\\
4.14207151775888	3.53928342989794e-08\\
4.18209152276138	2.53651537421737e-08\\
4.22211152776388	1.8120430961328e-08\\
4.26213153276638	1.29035263483008e-08\\
4.30215153776888	9.1591921633164e-09\\
4.34217154277139	6.48059372303923e-09\\
4.38219154777389	4.57068496023989e-09\\
4.42221155277639	3.21333975834547e-09\\
4.46223155777889	2.25185734354156e-09\\
4.50225156278139	1.57301899928504e-09\\
4.54227156778389	1.09530716668187e-09\\
4.58229157278639	7.60233070441163e-10\\
4.62231157778889	5.2597663614046e-10\\
4.6623315827914	3.6273963865497e-10\\
4.7023515877939	2.49363247995418e-10\\
4.7423715927964	1.70875084877425e-10\\
4.7823915977989	1.16716943962175e-10\\
4.8224116028014	7.9469035013004e-11\\
4.8624316078039	5.39350192834225e-11\\
4.90645361330665	3.50836095281114e-11\\
4.94647361830916	2.36513420422409e-11\\
4.98649362331166	1.58933766667526e-11\\
5.02651362831416	1.06459747216356e-11\\
5.06653363331666	7.1082640536508e-12\\
5.10655363831916	4.73097378204817e-12\\
5.14657364332166	3.13867535609458e-12\\
5.18659364832416	2.07563564786554e-12\\
5.22661365332666	1.36824765173743e-12\\
5.26663365832917	8.99056880984078e-13\\
5.30665366333167	5.88868731658669e-13\\
5.34667366833417	3.84466661568461e-13\\
5.38669367333667	2.50211776318582e-13\\
5.42671367833917	1.62317622470872e-13\\
5.46673368334167	1.04962089263502e-13\\
5.50675368834417	6.76562820137053e-14\\
5.54677369334667	4.34703018874741e-14\\
5.58679369834918	2.78410783066888e-14\\
5.62681370335168	1.77741284597066e-14\\
5.66683370835418	1.13109594300301e-14\\
5.70685371335668	7.17495930073217e-15\\
5.74687371835918	4.53678612841743e-15\\
5.79089572386193	2.73000409792804e-15\\
5.83091572886443	1.71463144696031e-15\\
5.87093573386693	1.07346314241637e-15\\
5.91095573886944	6.69903670730517e-16\\
5.95097574387194	4.1672201538937e-16\\
5.99099574887444	2.58398149592029e-16\\
6.03101575387694	1.59713350705079e-16\\
6.07103575887944	9.84015433414923e-17\\
6.11105576388194	6.04326259970597e-17\\
6.15107576888444	3.69955841911029e-17\\
6.19109577388694	2.25754897281505e-17\\
6.23111577888945	1.37319855590568e-17\\
6.27113578389195	8.32603750963023e-18\\
6.31115578889445	5.03213467220562e-18\\
6.35117579389695	3.03162175470186e-18\\
6.39119579889945	1.82056689621065e-18\\
6.43121580390195	1.08980083300828e-18\\
6.47123580890445	6.50274159901648e-19\\
6.51125581390695	3.86771731816392e-19\\
6.55127581890946	2.29309380300723e-19\\
6.59129582391196	1.35518249587828e-19\\
6.63131582891446	7.98330301173825e-20\\
6.67533783441721	4.44405237063416e-20\\
6.71535783941971	2.60041474190959e-20\\
6.75537784442221	1.51675326273018e-20\\
6.79539784942471	8.81852851361355e-21\\
6.83541785442721	5.11076811135279e-21\\
6.87543785942972	2.95246730096635e-21\\
6.91545786443222	1.70017209473838e-21\\
6.95547786943472	9.75909483367566e-22\\
6.99549787443722	5.5838669301376e-22\\
7.03551787943972	3.18470675439909e-22\\
7.07553788444222	1.81055909834295e-22\\
7.11555788944472	1.02604126765943e-22\\
7.15557789444722	5.79596586689341e-23\\
7.19559789944972	3.2635906260576e-23\\
7.23561790445223	1.83178459397958e-23\\
7.27563790945473	1.02485398559846e-23\\
7.31565791445723	5.71555560072707e-24\\
7.35567791945973	3.17734066173813e-24\\
7.39569792446223	1.76067004143708e-24\\
7.43571792946473	9.72525585628352e-25\\
7.47573793446723	5.35467311999438e-25\\
7.51575793946973	2.93882523083718e-25\\
7.55977994497249	1.51339354258516e-25\\
7.59979994997499	8.2503294939145e-26\\
7.63981995497749	4.48331834106133e-26\\
7.67983995997999	2.4284922874267e-26\\
7.71985996498249	1.31124168431141e-26\\
7.75987996998499	7.05728461628706e-27\\
7.79989997498749	3.78618111991713e-27\\
7.83991997999	2.02476211218377e-27\\
7.8799399849925	1.07933309282375e-27\\
7.919959989995	5.73516408144493e-28\\
7.9599799949975	3.0377013151799e-28\\
8	1.60381089054864e-28\\
};
\addlegendentry{$n=1$}

\addplot [color=mycolor2, line width=1.4pt]
  table[row sep=crcr]{%
1e-06	0.999999999998\\
0.0400210050025012	0.996800483651545\\
0.0800410100050025	0.98724826456394\\
0.120061015007504	0.971480390663682\\
0.160081020010005	0.949722037181622\\
0.200101025012506	0.922281746698046\\
0.240121030015008	0.889544951237612\\
0.280141035017509	0.851965954754393\\
0.32016104002001	0.810058594702914\\
0.360181045022511	0.764385834389475\\
0.400201050025012	0.715548562433984\\
0.440221055027514	0.664173891268017\\
0.484243060530265	0.605496743056042\\
0.524263065532766	0.550884042841713\\
0.564283070535268	0.495716911871969\\
0.604303075537769	0.440608096506016\\
0.64432308054027	0.386139751867212\\
0.684343085542771	0.332854706808546\\
0.724363090545273	0.281248979948098\\
0.764383095547774	0.231765663815864\\
0.804403100550275	0.184790250853524\\
0.844423105552776	0.140647431754698\\
0.884443110555278	0.0995993548719855\\
0.924463115557779	0.0618452964606401\\
0.96448312056028	0.0275226565096869\\
1.00450312556278	-0.00329083525115842\\
1.04452313056528	-0.0305738432035671\\
1.08454313556778	-0.0543572382355185\\
1.12456314057028	-0.0747191698508242\\
1.16458314557279	-0.0917795755422772\\
1.20460315057529	-0.105694305512064\\
1.24462315557779	-0.116649035973267\\
1.28464316058029	-0.124853134646875\\
1.32466316558279	-0.130533628414198\\
1.36868517108554	-0.134152740929596\\
1.40870517608804	-0.135318751000766\\
1.44872518109055	-0.134716630099241\\
1.48874518609305	-0.132589471591537\\
1.52876519109555	-0.129172996938708\\
1.56878519609805	-0.124692282335108\\
1.60880520110055	-0.119359090162107\\
1.64882520610305	-0.113369796992132\\
1.68884521110555	-0.106903894997967\\
1.72886521610805	-0.100123031195721\\
1.76888522111056	-0.0931705390878027\\
1.80890522611306	-0.0861714099803679\\
1.84892523111556	-0.0792326464422981\\
1.88894523611806	-0.0724439378864841\\
1.92896524112056	-0.0658785978650915\\
1.96898524612306	-0.0595947041098592\\
2.00900525112556	-0.0536363853194966\\
2.04902525612806	-0.0480352028889651\\
2.08904526113057	-0.0428115808804694\\
2.12906526613307	-0.0379762432562971\\
2.16908527113557	-0.0335316234539607\\
2.20910527613807	-0.0294732175382701\\
2.25312728164082	-0.0254427736122996\\
2.29314728664332	-0.0221570013462918\\
2.33316729164582	-0.0192125976635722\\
2.37318729664832	-0.0165888361561136\\
2.41320730165083	-0.0142634455382412\\
2.45322730665333	-0.0122133194909206\\
2.49324731165583	-0.0104151141998878\\
2.53326731665833	-0.00884573984070343\\
2.57328732166083	-0.00748275392071925\\
2.61330732666333	-0.00630466551580028\\
2.65332733166583	-0.00529116009842465\\
2.69334733666833	-0.00442325490878138\\
2.73336734167084	-0.00368339473892676\\
2.77338734667334	-0.00305549764820509\\
2.81340735167584	-0.00252495956900456\\
2.85342735667834	-0.00207862605354012\\
2.89344736168084	-0.00170473860658529\\
2.93346736668334	-0.0013928621908675\\
2.97348737168584	-0.00113379961894743\\
3.01350737668834	-0.000919497688415777\\
3.05352738169085	-0.00074294909992025\\
3.09354738669335	-0.000598093437336176\\
3.1375693921961	-0.000469161361993449\\
3.1775893971986	-0.000374798686414659\\
3.2176094022011	-0.000298329487304294\\
3.2576294072036	-0.000236604284356521\\
3.2976494122061	-0.000186974780233718\\
3.3376694172086	-0.000147225449228501\\
3.37768942221111	-0.000115512063013765\\
3.41770942721361	-9.03069782526349e-05\\
3.45772943221611	-7.03508735527202e-05\\
3.49774943721861	-5.46105277624496e-05\\
3.53776944222111	-4.22421709444225e-05\\
3.57778944722361	-3.25599070282465e-05\\
3.61780945222611	-2.50086972552949e-05\\
3.65782945722861	-1.91414008667813e-05\\
3.69784946223112	-1.45993895410912e-05\\
3.73786946723362	-1.10962810350996e-05\\
3.77788947223612	-8.40437216485464e-06\\
3.81790947723862	-6.34338911944219e-06\\
3.85792948224112	-4.77121213076601e-06\\
3.89794948724362	-3.57627019354867e-06\\
3.93796949224612	-2.67133872715635e-06\\
3.97798949724862	-1.98850802095575e-06\\
4.02201150275138	-1.43144904742223e-06\\
4.06203150775388	-1.05786958639331e-06\\
4.10205151275638	-7.79104833141027e-07\\
4.14207151775888	-5.71833404135077e-07\\
4.18209152276138	-4.18268582479506e-07\\
4.22211152776388	-3.04898462094721e-07\\
4.26213153276638	-2.21498903580513e-07\\
4.30215153776888	-1.60363787927053e-07\\
4.34217154277139	-1.1570746062093e-07\\
4.38219154777389	-8.3202933363625e-08\\
4.42221155277639	-5.96265880118317e-08\\
4.46223155777889	-4.25860237411274e-08\\
4.50225156278139	-3.03124994700107e-08\\
4.54227156778389	-2.15033233073229e-08\\
4.58229157278639	-1.52026818060235e-08\\
4.62231157778889	-1.07119162105998e-08\\
4.6623315827914	-7.52225369036748e-09\\
4.7023515877939	-5.26458443515363e-09\\
4.7423715927964	-3.67212066641458e-09\\
4.7823915977989	-2.55274772432255e-09\\
4.8224116028014	-1.76863532048689e-09\\
4.8624316078039	-1.22126364755585e-09\\
4.90645361330665	-8.09494193726573e-10\\
4.94647361830916	-5.55040264231443e-10\\
4.98649362331166	-3.7929731998587e-10\\
5.02651362831416	-2.58333511314688e-10\\
5.06653363331666	-1.75359189953915e-10\\
5.10655363831916	-1.1863810941404e-10\\
5.14657364332166	-7.99961101446666e-11\\
5.18659364832416	-5.37605276299237e-11\\
5.22661365332666	-3.60088442767934e-11\\
5.26663365832917	-2.40384705031853e-11\\
5.30665366333167	-1.5994012234022e-11\\
5.34667366833417	-1.06062507722495e-11\\
5.38669367333667	-7.0100504071987e-12\\
5.42671367833917	-4.61780997013962e-12\\
5.46673368334167	-3.03184854394097e-12\\
5.50675368834417	-1.98397555873564e-12\\
5.54677369334667	-1.29396736586755e-12\\
5.58679369834918	-8.41142003055296e-13\\
5.62681370335168	-5.44973129513101e-13\\
5.66683370835418	-3.51917929136384e-13\\
5.70685371335668	-2.26500401716275e-13\\
5.74687371835918	-1.45297641964696e-13\\
5.79089572386193	-8.8819245391004e-14\\
5.83091572886443	-5.65821145817908e-14\\
5.87093573386693	-3.59265424955271e-14\\
5.91095573886944	-2.27361271330136e-14\\
5.95097574387194	-1.43411182372224e-14\\
5.99099574887444	-9.0160360037968e-15\\
6.03101575387694	-5.64956447496186e-15\\
6.07103575887944	-3.5284308985916e-15\\
6.11105576388194	-2.1964239459225e-15\\
6.15107576888444	-1.3627594656806e-15\\
6.19109577388694	-8.42735511237248e-16\\
6.23111577888945	-5.19437124212444e-16\\
6.27113578389195	-3.19113238747604e-16\\
6.31115578889445	-1.95401248369331e-16\\
6.35117579389695	-1.19256220582389e-16\\
6.39119579889945	-7.25448277374096e-17\\
6.43121580390195	-4.39849465341486e-17\\
6.47123580890445	-2.65811871863263e-17\\
6.51125581390695	-1.60109775371226e-17\\
6.55127581890946	-9.61246918119099e-18\\
6.59129582391196	-5.7520965835442e-18\\
6.63131582891446	-3.43077254771283e-18\\
6.67533783441721	-1.93583522116711e-18\\
6.71535783941971	-1.1466796884185e-18\\
6.75537784442221	-6.7700478788251e-19\\
6.79539784942471	-3.98398471663878e-19\\
6.83541785442721	-2.33679329627201e-19\\
6.87543785942972	-1.3661552106499e-19\\
6.91545786443222	-7.9608105794927e-20\\
6.95547786943472	-4.62372956969093e-20\\
6.99549787443722	-2.6767377604612e-20\\
7.03551787943972	-1.54453538209407e-20\\
7.07553788444222	-8.8831888975206e-21\\
7.11555788944472	-5.09236225060441e-21\\
7.15557789444722	-2.9097078828001e-21\\
7.19559789944972	-1.65714130853316e-21\\
7.23561790445223	-9.40697709566797e-22\\
7.27563790945473	-5.3225696419869e-22\\
7.31565791445723	-3.00174411384222e-22\\
7.35567791945973	-1.68735845730853e-22\\
7.39569792446223	-9.45415508882542e-23\\
7.43571792946473	-5.27983178868163e-23\\
7.47573793446723	-2.93900110444211e-23\\
7.51575793946973	-1.63065471202749e-23\\
7.55977994497249	-8.49774602947474e-24\\
7.59979994997499	-4.68263615129716e-24\\
7.63981995497749	-2.57193846041647e-24\\
7.67983995997999	-1.40803841513233e-24\\
7.71985996498249	-7.68338296506925e-25\\
7.75987996998499	-4.17902310820732e-25\\
7.79989997498749	-2.26559170333825e-25\\
7.83991997999	-1.22425915486547e-25\\
7.8799399849925	-6.59401868374441e-26\\
7.919959989995	-3.54007397457217e-26\\
7.9599799949975	-1.89434946891918e-26\\
8	-1.01040086104564e-26\\
};
\addlegendentry{$n=2$}

\addplot [color=mycolor3, line width=1.4pt]
  table[row sep=crcr]{%
1e-06	0.999999999997\\
0.0400210050025012	0.99520264677623\\
0.0800410100050025	0.980903005110212\\
0.120061015007504	0.957374439231591\\
0.160081020010005	0.925064486930967\\
0.200101025012506	0.884583043389047\\
0.240121030015008	0.836686372969184\\
0.280141035017509	0.782257476692557\\
0.32016104002001	0.722283457211489\\
0.360181045022511	0.65783061203409\\
0.400201050025012	0.590018046409692\\
0.440221055027514	0.519990627805751\\
0.484243060530265	0.44176670157236\\
0.524263065532766	0.37077797482217\\
0.564283070535268	0.301003114037397\\
0.604303075537769	0.233425938764405\\
0.64432308054027	0.168936356173771\\
0.684343085542771	0.108315025012104\\
0.724363090545273	0.0522216641657382\\
0.764383095547774	0.0011871604973946\\
0.804403100550275	-0.0443905018834978\\
0.844423105552776	-0.0842465949090469\\
0.884443110555278	-0.118246501613079\\
0.924463115557779	-0.146379392601539\\
0.96448312056028	-0.168749126189556\\
1.00450312556278	-0.185562837200974\\
1.04452313056528	-0.197117721187076\\
1.08454313556778	-0.203786541461311\\
1.12456314057028	-0.206002386620344\\
1.16458314557279	-0.204243187680562\\
1.20460315057529	-0.199016468821881\\
1.24462315557779	-0.19084475673101\\
1.28464316058029	-0.180252013757371\\
1.32466316558279	-0.167751392801878\\
1.36868517108554	-0.152383022457257\\
1.40870517608804	-0.137439538655591\\
1.44872518109055	-0.122004447297636\\
1.48874518609305	-0.106453520672244\\
1.52876519109555	-0.0911154892160205\\
1.56878519609805	-0.0762702136835455\\
1.60880520110055	-0.0621483666866192\\
1.64882520610305	-0.0489324500001596\\
1.68884521110555	-0.0367589519805289\\
1.72886521610805	-0.0257214369501888\\
1.76888522111056	-0.0158743547841\\
1.80890522611306	-0.00723736322042795\\
1.84892523111556	0.00020003356054727\\
1.88894523611806	0.00647371008769024\\
1.92896524112056	0.0116401679354758\\
1.96898524612306	0.0157719434090352\\
2.00900525112556	0.0189532859681485\\
2.04902525612806	0.0212761936768565\\
2.08904526113057	0.0228368675313141\\
2.12906526613307	0.023732623581764\\
2.16908527113557	0.0240592809691401\\
2.20910527613807	0.0239090258024204\\
2.25312728164082	0.0232964090070341\\
2.29314728664332	0.02241957687141\\
2.33316729164582	0.0213128264084734\\
2.37318729664832	0.0200403866838539\\
2.41320730165083	0.0186583991468383\\
2.45322730665333	0.0172150066990815\\
2.49324731165583	0.0157506456407219\\
2.53326731665833	0.0142984943898845\\
2.57328732166083	0.0128850362784983\\
2.61330732666333	0.0115306980002704\\
2.65332733166583	0.010250530105904\\
2.69334733666833	0.0090549010232812\\
2.73336734167084	0.0079501811855913\\
2.77338734667334	0.00693939878212585\\
2.81340735167584	0.00602285325316882\\
2.85342735667834	0.00519867682311919\\
2.89344736168084	0.00446333803466766\\
2.93346736668334	0.00381208437573296\\
2.97348737168584	0.00323932367309791\\
3.01350737668834	0.0027389459787743\\
3.05352738169085	0.00230458923142039\\
3.09354738669335	0.0019298530826504\\
3.1375693921961	0.00157902708535051\\
3.1775893971986	0.00130938711973568\\
3.2176094022011	0.00108086157736653\\
3.2576294072036	0.000888226011362204\\
3.2976494122061	0.000726698107604974\\
3.3376694172086	0.000591950701284288\\
3.37768942221111	0.000480109322492955\\
3.41770942721361	0.000387738033883611\\
3.45772943221611	0.000311816801066955\\
3.49774943721861	0.000249713126653047\\
3.53776944222111	0.000199150197589604\\
3.57778944722361	0.00015817335374278\\
3.61780945222611	0.000125116290255197\\
3.65782945722861	9.85680603628926e-05\\
3.69784946223112	7.73416495403549e-05\\
3.73786946723362	6.04446444466514e-05\\
3.77788947223612	4.70523180307922e-05\\
3.81790947723862	3.64832911902355e-05\\
3.85792948224112	2.8177806865108e-05\\
3.89794948724362	2.16785594716875e-05\\
3.93796949224612	1.66139562539033e-05\\
3.97798949724862	1.26836428194304e-05\\
4.02201150275138	9.38363914600103e-06\\
4.06203150775388	7.10654798048342e-06\\
4.10205151275638	5.36166029528959e-06\\
4.14207151775888	4.02995670136077e-06\\
4.18209152276138	3.01765019484769e-06\\
4.22211152776388	2.25118649987996e-06\\
4.26213153276638	1.67314841899372e-06\\
4.30215153776888	1.23892193619445e-06\\
4.34217154277139	9.13998992114544e-07\\
4.38219154777389	6.71808297922792e-07\\
4.42221155277639	4.91980884606539e-07\\
4.46223155777889	3.5897106462663e-07\\
4.50225156278139	2.60966002494617e-07\\
4.54227156778389	1.89028128279455e-07\\
4.58229157278639	1.3642422626838e-07\\
4.62231157778889	9.8103275963237e-08\\
4.6623315827914	7.02921268197431e-08\\
4.7023515877939	5.01839779887669e-08\\
4.7423715927964	3.56995405202621e-08\\
4.7823915977989	2.53048103913761e-08\\
4.8224116028014	1.78726977398493e-08\\
4.8624316078039	1.25784524767123e-08\\
4.90645361330665	8.51180995394734e-09\\
4.94647361830916	5.94586586592604e-09\\
4.98649362331166	4.13874376523606e-09\\
5.02651362831416	2.87068323127658e-09\\
5.06653363331666	1.9841215110809e-09\\
5.10655363831916	1.36653381710714e-09\\
5.14657364332166	9.37873792022771e-10\\
5.18659364832416	6.41420746278332e-10\\
5.22661365332666	4.37138236542148e-10\\
5.26663365832917	2.96875463343613e-10\\
5.30665366333167	2.00914822642777e-10\\
5.34667366833417	1.35498408044807e-10\\
5.38669367333667	9.10632727204254e-11\\
5.42671367833917	6.09875802073137e-11\\
5.46673368334167	4.07034115753854e-11\\
5.50675368834417	2.70715794329505e-11\\
5.54677369334667	1.79428656461657e-11\\
5.58679369834918	1.18513395143147e-11\\
5.62681370335168	7.80085921105328e-12\\
5.66683370835418	5.1170386085357e-12\\
5.70685371335668	3.34501476572425e-12\\
5.74687371835918	2.17912560999528e-12\\
5.79089572386193	1.35465943572789e-12\\
5.83091572886443	8.76153528176437e-13\\
5.87093573386693	5.64729447267646e-13\\
5.91095573886944	3.62754152027037e-13\\
5.95097574387194	2.32218947546531e-13\\
5.99099574887444	1.48148664568334e-13\\
6.03101575387694	9.41920275446321e-14\\
6.07103575887944	5.96826800249013e-14\\
6.11105576388194	3.76878766979996e-14\\
6.15107576888444	2.37178647307976e-14\\
6.19109577388694	1.48754946957223e-14\\
6.23111577888945	9.29801998798465e-15\\
6.27113578389195	5.79207327258321e-15\\
6.31115578889445	3.59586507998287e-15\\
6.35117579389695	2.22484481957682e-15\\
6.39119579889945	1.37190568341215e-15\\
6.43121580390195	8.43098177828627e-16\\
6.47123580890445	5.1637184654508e-16\\
6.51125581390695	3.15194470335858e-16\\
6.55127581890946	1.9174645657742e-16\\
6.59129582391196	1.16254533560458e-16\\
6.63131582891446	7.02469204081058e-17\\
6.67533783441721	4.02045664999467e-17\\
6.71535783941971	2.41223297308923e-17\\
6.75537784442221	1.44245097440607e-17\\
6.79539784942471	8.59650218177638e-18\\
6.83541785442721	5.10601794827208e-18\\
6.87543785942972	3.02262074322651e-18\\
6.91545786443222	1.78330916672957e-18\\
6.95547786943472	1.04860559212389e-18\\
6.99549787443722	6.14527192170623e-19\\
7.03551787943972	3.58933748357849e-19\\
7.07553788444222	2.0894528137649e-19\\
7.11555788944472	1.21226268888292e-19\\
7.15557789444722	7.0098319041038e-20\\
7.19559789944972	4.03985655585005e-20\\
7.23561790445223	2.3204516749918e-20\\
7.27563790945473	1.32839767309582e-20\\
7.31565791445723	7.57937536174555e-21\\
7.35567791945973	4.31011819777227e-21\\
7.39569792446223	2.4428461126369e-21\\
7.43571792946473	1.37992314818243e-21\\
7.47573793446723	7.76901992779973e-22\\
7.51575793946973	4.35944082245404e-22\\
7.55977994497249	2.30001963233174e-22\\
7.59979994997499	1.28162206881505e-22\\
7.63981995497749	7.11776475244733e-23\\
7.67983995997999	3.93988118118765e-23\\
7.71985996498249	2.1735972291881e-23\\
7.75987996998499	1.19517657425422e-23\\
7.79989997498749	6.55002135627364e-24\\
7.83991997999	3.57776311683082e-24\\
7.8799399849925	1.9477770377773e-24\\
7.919959989995	1.05688139523921e-24\\
7.9599799949975	5.71574922894909e-25\\
8	3.08092072074393e-25\\
};
\addlegendentry{$n=3$}

\node[overlay,right, align=left, inner sep=0]
at (rel axis cs:-0.15,1.1) {$(a)$};
\end{axis}
\end{tikzpicture}%
    \end{minipage}
    \hfill
    \begin{minipage}{0.48\textwidth}
        \centering
        % This file was created by matlab2tikz.
%
\definecolor{mycolor1}{rgb}{0.48200,0.16500,0.51400}%
\definecolor{mycolor2}{rgb}{0.10600,0.47100,0.21600}%
\definecolor{mycolor3}{rgb}{0.68600,0.55300,0.76500}%
%
\begin{tikzpicture}[trim axis left,trim axis right]

\begin{axis}[%
width=11.252in,
height=7.131in,
at={(1.887in,0.962in)},
scale only axis,
clip=true,
separate axis lines,
every outer x axis line/.append style={black},
every x tick label/.append style={font=\color{black}},
every x tick/.append style={black},
xmin=0,
xmax=2.5,
xlabel={$k/k_0$},
every outer y axis line/.append style={black},
every y tick label/.append style={font=\color{black}},
every y tick/.append style={black},
ymin=0,
ymax=0.35,
ylabel={$E_{\eta,n}(k)$},
axis background/.style={fill=white},
legend style={legend cell align=left, align=left},
clip=true,
clip mode=individual,
axis lines=box,
xtick pos=bottom,
ytick pos=left,
tick align=inside,
major tick length=2pt,
width=0.7\linewidth,
height=0.65\linewidth,
every axis/.append style={font=\fontsize{9}{9}\selectfont},xlabel style={font=\fontsize{9}{9}\selectfont},title style={font=\fontsize{9}{9}\selectfont},
legend style={font=\fontsize{8}{9}\selectfont,inner sep=1pt,row sep=1pt,column sep=2pt,nodes={scale=1.00},draw=none},legend image post style={xscale=0.35},
legend columns=1,
scaled x ticks=false,
tick label style={/pgf/number format/fixed,/pgf/number format/precision=2},
every x tick label/.append style={font=\fontsize{9}{9}\selectfont\color{black}},
every y tick label/.append style={font=\fontsize{9}{9}\selectfont\color{black}},
,,
colormap={sigmamap}{rgb(0.0000)=(0.0200,0.1880,0.3800); rgb(0.1000)=(0.0743,0.3456,0.5616); rgb(0.2000)=(0.1918,0.4980,0.7060); rgb(0.3000)=(0.3890,0.6708,0.8226); rgb(0.4000)=(0.7552,0.8682,0.9228); rgb(0.5000)=(0.9690,0.9689,0.9689); rgb(0.6000)=(0.9425,0.8044,0.7614); rgb(0.7000)=(0.8767,0.4910,0.4020); rgb(0.8000)=(0.7869,0.2866,0.2470); rgb(0.9000)=(0.6186,0.1239,0.1568); rgb(1.0000)=(0.4040,0.0000,0.1220)}
]
\addplot [color=mycolor1, line width=1.4pt]
  table[row sep=crcr]{%
1e-06	9.99999999998e-07\\
0.012507248124062	0.0125033356870642\\
0.0250134962481241	0.0249822151857268\\
0.0375197443721861	0.0374142575081492\\
0.0500259924962481	0.049776227988151\\
0.0625322406203102	0.0620451106388611\\
0.0750384887443722	0.074198179555138\\
0.0875447368684342	0.0862130690507178\\
0.100050984992496	0.0980678422271457\\
0.112557233116558	0.109741057680322\\
0.12506348124062	0.121211834060891\\
0.137569729364682	0.132459912216648\\
0.151326602301151	0.144552237739127\\
0.163832850425213	0.15526980755966\\
0.176339098549275	0.165706455442367\\
0.188845346673337	0.175845062145538\\
0.201351594797399	0.185669401013813\\
0.213857842921461	0.195164180365556\\
0.226364091045523	0.204315081990821\\
0.238870339169585	0.213108795630047\\
0.251376587293647	0.221533049324865\\
0.263882835417709	0.229576635554121\\
0.276389083541771	0.237229433090159\\
0.288895331665833	0.244482424532426\\
0.301401579789895	0.251327709497517\\
0.313907827913957	0.257758513466586\\
0.326414076038019	0.26376919231265\\
0.338920324162081	0.269355232551407\\
0.351426572286143	0.274513247379763\\
0.363932820410205	0.279240968586193\\
0.376439068534267	0.283537234436133\\
0.388945316658329	0.287401973653864\\
0.401451564782391	0.290836185639548\\
0.413957812906453	0.293841917076269\\
0.427714685842921	0.296657016968592\\
0.440220933966983	0.298774103791276\\
0.452727182091046	0.300475402147823\\
0.465233430215107	0.301766937984341\\
0.47773967833917	0.302655615812989\\
0.490245926463232	0.303149177432573\\
0.502752174587294	0.303256158493771\\
0.515258422711356	0.302985843147067\\
0.527764670835418	0.302348217015337\\
0.54027091895948	0.30135391873543\\
0.552777167083542	0.300014190314105\\
0.565283415207604	0.298340826543286\\
0.577789663331666	0.296346123717858\\
0.590295911455728	0.2940428278962\\
0.60280215957979	0.291444082939338\\
0.615308407703852	0.288563378559127\\
0.627814655827914	0.285414498599276\\
0.640320903951976	0.282011469765338\\
0.652827152076038	0.278368511011214\\
0.6653334002001	0.274499983780127\\
0.677839648324162	0.270420343287768\\
0.690345896448224	0.266144091024213\\
0.704102769384692	0.261230433566072\\
0.716609017508754	0.256588445151956\\
0.729115265632816	0.251794604352604\\
0.741621513756878	0.24686313763554\\
0.75412776188094	0.241808100063135\\
0.766634010005002	0.23664333833741\\
0.779140258129065	0.231382455957741\\
0.791646506253126	0.226038780563521\\
0.804152754377189	0.220625333520528\\
0.816659002501251	0.215154801796455\\
0.829165250625313	0.209639512158135\\
0.841671498749375	0.204091407710348\\
0.854177746873437	0.19852202678381\\
0.866683994997499	0.192942484168157\\
0.879190243121561	0.187363454674395\\
0.891696491245623	0.181795159000548\\
0.904202739369685	0.176247351864029\\
0.916708987493747	0.170729312354729\\
0.929215235617809	0.165249836453935\\
0.941721483741871	0.159817231655987\\
0.954227731865933	0.154439313622079\\
0.966733979989995	0.149123404788854\\
0.980490852926463	0.143355662283408\\
0.992997101050525	0.138191618138487\\
1.00550334917459	0.13310916348621\\
1.01800959729865	0.128113622286538\\
1.03051584542271	0.123209845332416\\
1.04302209354677	0.11840221880887\\
1.05552834167084	0.113694674126555\\
1.0680345897949	0.109090698924165\\
1.08054083791896	0.104593349133712\\
1.09304708604302	0.100205262002926\\
1.10555333416708	0.0959286699697512\\
1.11805958229115	0.0917654152852739\\
1.13056583041521	0.0877169652831892\\
1.14307207853927	0.0837844281962175\\
1.15557832666333	0.0799685694225552\\
1.16808457478739	0.0762698281485362\\
1.18059082291146	0.0726883342371032\\
1.19309707103552	0.0692239252954183\\
1.20560331915958	0.0658761638389392\\
1.21810956728364	0.0626443544735069\\
1.2306158154077	0.0595275610214065\\
1.24312206353177	0.0565246235219153\\
1.25687893646823	0.0533512586932931\\
1.2693851845923	0.050582742554742\\
1.28189143271636	0.04792324446996\\
1.29439768084042	0.0453708597065667\\
1.30690392896448	0.0429235422940296\\
1.31941017708854	0.0405791200311605\\
1.33191642521261	0.038335309011414\\
1.34442267333667	0.0361897276348124\\
1.35692892146073	0.0341399100797784\\
1.36943516958479	0.0321833192124766\\
1.38194141770885	0.0303173589154311\\
1.39444766583292	0.028539385821176\\
1.40695391395698	0.0268467204405133\\
1.41946016208104	0.025236657678563\\
1.4319664102051	0.0237064767352062\\
1.44447265832916	0.022253450389725\\
1.45697890645323	0.020874853672428\\
1.46948515457729	0.0195679719288212\\
1.48199140270135	0.0183301082844277\\
1.49449765082541	0.0171585905206906\\
1.50700389894947	0.0160507773744975\\
1.51951014707354	0.015004064275758\\
1.53326702001	0.0139201932227356\\
1.54577326813407	0.0129935043648222\\
1.55827951625813	0.0121201276874942\\
1.57078576438219	0.0112976577428226\\
1.58329201250625	0.0105237469402583\\
1.59579826063032	0.0097961082518346\\
1.60830450875438	0.00911251749448843\\
1.62081075687844	0.00847081521196432\\
1.6333170050025	0.00786890817898354\\
1.64582325312656	0.00730477055043469\\
1.65832950125063	0.00677644467828962\\
1.67083574937469	0.00628204161877771\\
1.68334199749875	0.00581974135207076\\
1.69584824562281	0.00538779273635341\\
1.70835449374687	0.00498451321768748\\
1.72086074187094	0.00460828831653454\\
1.733366989995	0.00425757091118747\\
1.74587323811906	0.00393088033768892\\
1.75837948624312	0.00362680132509161\\
1.77088573436718	0.00334398278414903\\
1.78339198249125	0.00308113646672591\\
1.79589823061531	0.00283703551239094\\
1.80965510355178	0.0025887842603378\\
1.82216135167584	0.00238031932361065\\
1.8346675997999	0.00218716939922618\\
1.84717384792396	0.00200834229180101\\
1.85968009604802	0.00184289856882796\\
1.87218634417209	0.00168994972555733\\
1.88469259229615	0.00154865632932159\\
1.89719884042021	0.00141822615287879\\
1.90970508854427	0.00129791230555245\\
1.92221133666833	0.00118701137017154\\
1.9347175847924	0.00108486155306661\\
1.94722383291646	0.000990840853659224\\
1.95973008104052	0.000904365259492857\\
1.97223632916458	0.000824886971896154\\
1.98474257728864	0.000751892666844868\\
1.99724882541271	0.000684901794997244\\
2.00975507353677	0.00062346492432005\\
2.02226132166083	0.000567162128198694\\
2.03476756978489	0.000515601421434883\\
2.04727381790895	0.000468417246079032\\
2.05978006603302	0.000425269008621212\\
2.07228631415708	0.00038583966967374\\
2.08604318709355	0.000346411460523943\\
2.09854943521761	0.000313856995799808\\
2.11105568334167	0.000284173937835262\\
2.12356193146573	0.000257128211923276\\
2.13606817958979	0.000232502938143593\\
2.14857442771386	0.00021009734700092\\
2.16108067583792	0.000189725746983587\\
2.17358692396198	0.000171216543082732\\
2.18609317208604	0.000154411305167679\\
2.19859942021011	0.00013916388498924\\
2.21110566833417	0.000125339580478016\\
2.22361191645823	0.000112814345917912\\
2.23611816458229	0.000101474046504826\\
2.24862441270635	9.12137557453812e-05\\
2.26113066083042	8.19370941095129e-05\\
2.27363690895448	7.35556073223929e-05\\
2.28614315707854	6.59881826644289e-05\\
2.2986494052026	5.91605016418098e-05\\
2.31115565332666	5.30045273931476e-05\\
2.32366190145073	4.74580252092345e-05\\
2.33616814957479	4.2464114561786e-05\\
2.34867439769885	3.79708510623788e-05\\
2.36243127063532	3.35500659514266e-05\\
2.37493751875938	2.99588643141167e-05\\
2.38744376688344	2.67345922799202e-05\\
2.3999500150075	2.38417517501924e-05\\
2.41245626313157	2.12480587809647e-05\\
2.42496251125563	1.89241760545548e-05\\
2.43746875937969	1.68434651930235e-05\\
2.44997500750375	1.49817576851635e-05\\
2.46248125562781	1.33171432506067e-05\\
2.47498750375188	1.182977451687e-05\\
2.48749375187594	1.05016869373549e-05\\
2.5	9.31663293019668e-06\\
};
\addlegendentry{$n=1$}

\addplot [color=mycolor2, line width=1.4pt]
  table[row sep=crcr]{%
1e-06	1.999999999996e-18\\
0.012507248124062	3.91182500235304e-06\\
0.0250134962481241	3.12614947005987e-05\\
0.0375197443721861	0.000105338436566632\\
0.0500259924962481	0.000249139968883221\\
0.0625322406203102	0.000485227649066178\\
0.0750384887443722	0.000835586478256945\\
0.0875447368684342	0.00132148788088522\\
0.100050984992496	0.00196335734967166\\
0.112557233116558	0.00278064761171072\\
0.12506348124062	0.00379171813014021\\
0.137569729364682	0.00501372170881961\\
0.151326602301151	0.00662041748434639\\
0.163832850425213	0.00833525681121348\\
0.176339098549275	0.0103054427723424\\
0.188845346673337	0.0125421719033817\\
0.201351594797399	0.0150549902831408\\
0.213857842921461	0.0178517366579682\\
0.226364091045523	0.02093849634428\\
0.238870339169585	0.0243195661344885\\
0.251376587293647	0.0279974303533997\\
0.263882835417709	0.0319727481335621\\
0.276389083541771	0.036244351899698\\
0.288895331665833	0.0408092569748465\\
0.301401579789895	0.0456626821448074\\
0.313907827913957	0.0507980809434863\\
0.326414076038019	0.056207183350368\\
0.338920324162081	0.0618800475231208\\
0.351426572286143	0.0678051211237688\\
0.363932820410205	0.0739693117364142\\
0.376439068534267	0.0803580658185823\\
0.388945316658329	0.0869554555772692\\
0.401451564782391	0.0937442731150242\\
0.413957812906453	0.100706131151188\\
0.427714685842921	0.108540781846169\\
0.440220933966983	0.115801538607845\\
0.452727182091046	0.123172019498758\\
0.465233430215107	0.130630166447205\\
0.47773967833917	0.138153330169548\\
0.490245926463232	0.145718394385814\\
0.502752174587294	0.153301901038872\\
0.515258422711356	0.160880175717106\\
0.527764670835418	0.168429452494566\\
0.54027091895948	0.175925997422113\\
0.552777167083542	0.183346229927788\\
0.565283415207604	0.190666841414204\\
0.577789663331666	0.197864910374866\\
0.590295911455728	0.204918013389573\\
0.60280215957979	0.211804331401061\\
0.615308407703852	0.218502750720376\\
0.627814655827914	0.224992958256707\\
0.640320903951976	0.231255530518063\\
0.652827152076038	0.237272015981847\\
0.6653334002001	0.243025010488528\\
0.677839648324162	0.248498225366866\\
0.690345896448224	0.253676548054941\\
0.704102769384692	0.259015570360991\\
0.716609017508754	0.263530950488702\\
0.729115265632816	0.267712591193295\\
0.741621513756878	0.271550670738605\\
0.75412776188094	0.275036731399896\\
0.766634010005002	0.278163692429428\\
0.779140258129065	0.280925855405849\\
0.791646506253126	0.283318902153342\\
0.804152754377189	0.285339885459202\\
0.816659002501251	0.286987212858155\\
0.829165250625313	0.288260623788246\\
0.841671498749375	0.289161160456131\\
0.854177746873437	0.289691132779161\\
0.866683994997499	0.289854077797518\\
0.879190243121561	0.28965471397187\\
0.891696491245623	0.289098890800455\\
0.904202739369685	0.288193534204232\\
0.916708987493747	0.286946588139744\\
0.929215235617809	0.285366952906717\\
0.941721483741871	0.283464420621233\\
0.954227731865933	0.281249608325641\\
0.966733979989995	0.278733889203472\\
0.980490852926463	0.275633462054939\\
0.992997101050525	0.272525822525016\\
1.00550334917459	0.269156574707533\\
1.01800959729865	0.265539449734544\\
1.03051584542271	0.261688571204898\\
1.04302209354677	0.25761838441254\\
1.05552834167084	0.253343586669873\\
1.0680345897949	0.248879059071574\\
1.08054083791896	0.244239800020315\\
1.09304708604302	0.239440860810677\\
1.10555333416708	0.234497283541543\\
1.11805958229115	0.229424041600589\\
1.13056583041521	0.224235982937377\\
1.14307207853927	0.218947776314266\\
1.15557832666333	0.213573860697076\\
1.16808457478739	0.208128397920342\\
1.18059082291146	0.202625228735368\\
1.19309707103552	0.197077832323188\\
1.20560331915958	0.191499289329223\\
1.21810956728364	0.18590224845199\\
1.2306158154077	0.18029889659484\\
1.24312206353177	0.17470093256743\\
1.25687893646823	0.168562732150052\\
1.2693851845923	0.163011866002044\\
1.28189143271636	0.157499325565096\\
1.29439768084042	0.152034607235555\\
1.30690392896448	0.146626638441337\\
1.31941017708854	0.141283771586688\\
1.33191642521261	0.136013780932139\\
1.34442267333667	0.130823862268147\\
1.35692892146073	0.125720635232066\\
1.36943516958479	0.120710148110709\\
1.38194141770885	0.115797884964886\\
1.39444766583292	0.110988774907839\\
1.40695391395698	0.106287203366435\\
1.41946016208104	0.101697025152242\\
1.4319664102051	0.0972215791690861\\
1.44447265832916	0.092863704584415\\
1.45697890645323	0.0886257582935311\\
1.46948515457729	0.084509633508598\\
1.48199140270135	0.0805167793080312\\
1.49449765082541	0.0766482209864793\\
1.50700389894947	0.0729045810509312\\
1.51951014707354	0.0692861007145002\\
1.53326702001	0.0654501803871208\\
1.54577326813407	0.0620937483722139\\
1.55827951625813	0.0588610377416395\\
1.57078576438219	0.0557509565170094\\
1.58329201250625	0.052762143839635\\
1.59579826063032	0.0498929917026545\\
1.60830450875438	0.0471416663392376\\
1.62081075687844	0.0445061291813547\\
1.6333170050025	0.0419841573119727\\
1.64582325312656	0.0395733633418135\\
1.65832950125063	0.0372712146499506\\
1.67083574937469	0.0350750519354708\\
1.68334199749875	0.0329821070351399\\
1.69584824562281	0.0309895199694631\\
1.70835449374687	0.0290943551866689\\
1.72086074187094	0.0272936169809641\\
1.733366989995	0.0255842640678607\\
1.74587323811906	0.0239632233054625\\
1.75837948624312	0.0224274025562953\\
1.77088573436718	0.020973702689564\\
1.78339198249125	0.0195990287286103\\
1.79589823061531	0.0183003001528321\\
1.80965510355178	0.0169557685220001\\
1.82216135167584	0.0158066151622171\\
1.8346675997999	0.0147240471498316\\
1.84717384792396	0.0137051335517257\\
1.85968009604802	0.0127469978986512\\
1.87218634417209	0.0118468237386303\\
1.88469259229615	0.0110018595047996\\
1.89719884042021	0.0102094227289079\\
1.90970508854427	0.00946690363279301\\
1.92221133666833	0.0087717681309781\\
1.9347175847924	0.00812156027807352\\
1.94722383291646	0.0075139041949537\\
1.95973008104052	0.00694650550772695\\
1.97223632916458	0.00641715233334719\\
1.98474257728864	0.00592371584535152\\
1.99724882541271	0.00546415045266323\\
2.00975507353677	0.00503649362369966\\
2.02226132166083	0.00463886538718285\\
2.03476756978489	0.0042694675400896\\
2.04727381790895	0.003926582592112\\
2.05978006603302	0.00360857247484645\\
2.07228631415708	0.00331387704270614\\
2.08604318709355	0.00301487171909493\\
2.09854943521761	0.00276439575655084\\
2.11105568334167	0.00253287419120097\\
2.12356193146573	0.00231904719951338\\
2.13606817958979	0.00212172289180142\\
2.14857442771386	0.00193977504995146\\
2.16108067583792	0.00177214081014392\\
2.17358692396198	0.00161781830665559\\
2.18609317208604	0.00147586429152291\\
2.19859942021011	0.00134539174357198\\
2.21110566833417	0.0012255674790868\\
2.22361191645823	0.00111560977519774\\
2.23611816458229	0.00101478601593084\\
2.24862441270635	0.00092241036976801\\
2.26113066083042	0.000837841506531585\\
2.27363690895448	0.000760480360424915\\
2.28614315707854	0.0006897679451351\\
2.2986494052026	0.00062518322603573\\
2.31115565332666	0.000566241053715621\\
2.32366190145073	0.000512490162304697\\
2.33616814957479	0.000463511235369099\\
2.34867439769885	0.00041891504150344\\
2.36243127063532	0.000374491305379288\\
2.37493751875938	0.000337955655540948\\
2.38744376688344	0.000304768349542068\\
2.3999500150075	0.000274645539675444\\
2.41245626313157	0.000247325076337819\\
2.42496251125563	0.000222565084071919\\
2.43746875937969	0.00020014260831892\\
2.44997500750375	0.00017985233160018\\
2.46248125562781	0.000161505357616685\\
2.47498750375188	0.000144928061558302\\
2.48749375187594	0.000129961004750119\\
2.5	0.000116457911627458\\
};
\addlegendentry{$n=2$}

\addplot [color=mycolor3, line width=1.4pt]
  table[row sep=crcr]{%
1e-06	1.999999999996e-30\\
0.012507248124062	6.11931696949723e-10\\
0.0250134962481241	1.95595355265756e-08\\
0.0375197443721861	1.48288205584267e-07\\
0.0500259924962481	6.23497667500118e-07\\
0.0625322406203102	1.89737651358801e-06\\
0.0750384887443722	4.70499927917366e-06\\
0.0875447368684342	1.01279900979929e-05\\
0.100050984992496	1.96535989523522e-05\\
0.112557233116558	3.5228388098081e-05\\
0.12506348124062	5.930578680631e-05\\
0.137569729364682	9.48868414331085e-05\\
0.151326602301151	0.000151605842816972\\
0.163832850425213	0.000223728319113764\\
0.176339098549275	0.000320452665680735\\
0.188845346673337	0.000447286020246002\\
0.201351594797399	0.000610366412526424\\
0.213857842921461	0.000816452335435037\\
0.226364091045523	0.00107290324553506\\
0.238870339169585	0.0013876510709499\\
0.251376587293647	0.0017691629054498\\
0.263882835417709	0.00222639516592098\\
0.276389083541771	0.00276873958580403\\
0.288895331665833	0.00340596150832547\\
0.301401579789895	0.00414813102945641\\
0.313907827913957	0.00500554762059392\\
0.326414076038019	0.00598865893412824\\
0.338920324162081	0.00710797456057011\\
0.351426572286143	0.0083739755630968\\
0.363932820410205	0.00979702066366017\\
0.376439068534267	0.0113872499937184\\
0.388945316658329	0.0131544873518456\\
0.401451564782391	0.0151081419296906\\
0.413957812906453	0.0172571104768691\\
0.427714685842921	0.0198564346196217\\
0.440220933966983	0.022441697881073\\
0.452727182091046	0.0252455713162482\\
0.465233430215107	0.0282738733739445\\
0.47773967833917	0.0315314529777793\\
0.490245926463232	0.0350221145942289\\
0.502752174587294	0.0387485500358184\\
0.515258422711356	0.0427122776924657\\
0.527764670835418	0.0469135898131991\\
0.54027091895948	0.0513515083839636\\
0.552777167083542	0.0560237500658531\\
0.565283415207604	0.0609267005727671\\
0.577789663331666	0.066055398779111\\
0.590295911455728	0.0714035307576907\\
0.60280215957979	0.0769634338563467\\
0.615308407703852	0.0827261108300827\\
0.627814655827914	0.0886812539544072\\
0.640320903951976	0.0948172789562263\\
0.652827152076038	0.101121368511796\\
0.6653334002001	0.107579524977781\\
0.677839648324162	0.114176631942163\\
0.690345896448224	0.120896524107343\\
0.704102769384692	0.128409743025713\\
0.716609017508754	0.135330649484918\\
0.729115265632816	0.142318441786513\\
0.741621513756878	0.149353539545975\\
0.75412776188094	0.156415776806873\\
0.766634010005002	0.163484508648307\\
0.779140258129065	0.170538721072961\\
0.791646506253126	0.177557143330142\\
0.804152754377189	0.184518361818804\\
0.816659002501251	0.191400934713995\\
0.829165250625313	0.19818350646635\\
0.841671498749375	0.204844921337904\\
0.854177746873437	0.211364335158293\\
0.866683994997499	0.217721324513027\\
0.879190243121561	0.223895992609469\\
0.891696491245623	0.229869071106017\\
0.904202739369685	0.235622017235191\\
0.916708987493747	0.241137105601301\\
0.929215235617809	0.246397514087601\\
0.941721483741871	0.251387403365527\\
0.954227731865933	0.25609198955938\\
0.966733979989995	0.26049760968273\\
0.980490852926463	0.264983622531054\\
0.992997101050525	0.26872224576062\\
1.00550334917459	0.272127251839416\\
1.01800959729865	0.275190093398575\\
1.03051584542271	0.277903555980051\\
1.04302209354677	0.280261774884728\\
1.05552834167084	0.282260243937691\\
1.0680345897949	0.283895816303325\\
1.08054083791896	0.285166697538785\\
1.09304708604302	0.286072431126841\\
1.10555333416708	0.286613876777935\\
1.11805958229115	0.286793181836096\\
1.13056583041521	0.286613746163924\\
1.14307207853927	0.286080180917951\\
1.15557832666333	0.285198261657216\\
1.16808457478739	0.283974876254653\\
1.18059082291146	0.282417968102941\\
1.19309707103552	0.280536475123704\\
1.20560331915958	0.278340265101477\\
1.21810956728364	0.275840067871729\\
1.2306158154077	0.273047404895582\\
1.24312206353177	0.269974516752834\\
1.25687893646823	0.266286076147086\\
1.2693851845923	0.262667335887439\\
1.28189143271636	0.258810080868057\\
1.29439768084042	0.254728717361335\\
1.30690392896448	0.250437987542292\\
1.31941017708854	0.24595289521349\\
1.33191642521261	0.241288632862561\\
1.34442267333667	0.236460510444566\\
1.35692892146073	0.231483886252473\\
1.36943516958479	0.226374100208726\\
1.38194141770885	0.221146409879324\\
1.39444766583292	0.21581592947951\\
1.40695391395698	0.210397572107358\\
1.41946016208104	0.204905995408433\\
1.4319664102051	0.199355550841804\\
1.44447265832916	0.193760236684988\\
1.45697890645323	0.188133654883479\\
1.46948515457729	0.182488971819262\\
1.48199140270135	0.176838883042654\\
1.49449765082541	0.171195581982911\\
1.50700389894947	0.165570732625596\\
1.51951014707354	0.159975446118827\\
1.53326702001	0.153867336615347\\
1.54577326813407	0.148367733548091\\
1.55827951625813	0.142928434970103\\
1.57078576438219	0.137558121484787\\
1.58329201250625	0.132264859577002\\
1.59579826063032	0.127056100088265\\
1.60830450875438	0.121938679765726\\
1.62081075687844	0.116918825705802\\
1.6333170050025	0.112002162504847\\
1.64582325312656	0.107193721922586\\
1.65832950125063	0.102497954859222\\
1.67083574937469	0.0979187454440442\\
1.68334199749875	0.0934594270319558\\
1.69584824562281	0.0891227999044873\\
1.70835449374687	0.0849111504734621\\
1.72086074187094	0.0808262717884554\\
1.733366989995	0.0768694851533864\\
1.74587323811906	0.0730416626629076\\
1.75837948624312	0.0693432504755988\\
1.77088573436718	0.0657742926481854\\
1.78339198249125	0.0623344553630029\\
1.79589823061531	0.0590230513895659\\
1.80965510355178	0.0555276255685623\\
1.82216135167584	0.0524822616041784\\
1.8346675997999	0.0495612192972269\\
1.84717384792396	0.0467626177164741\\
1.85968009604802	0.0440843457628693\\
1.87218634417209	0.0415240851759331\\
1.88469259229615	0.0390793329261024\\
1.89719884042021	0.0367474229148731\\
1.90970508854427	0.0345255469145283\\
1.92221133666833	0.0324107746889247\\
1.9347175847924	0.0304000732461812\\
1.94722383291646	0.0284903251831199\\
1.95973008104052	0.0266783460899086\\
1.97223632916458	0.0249609009915164\\
1.98474257728864	0.0233347198102868\\
1.99724882541271	0.0217965118411319\\
2.00975507353677	0.0203429792375503\\
2.02226132166083	0.0189708295128481\\
2.03476756978489	0.0176767870665974\\
2.04727381790895	0.0164576037515019\\
2.05978006603302	0.0153100685004512\\
2.07228631415708	0.014231016037647\\
2.08604318709355	0.0131194439538038\\
2.09854943521761	0.0121741493755179\\
2.11105568334167	0.011287895922692\\
2.12356193146573	0.0104577787737574\\
2.13606817958979	0.00968097019663025\\
2.14857442771386	0.00895472336544476\\
2.16108067583792	0.00827637550756156\\
2.17358692396198	0.00764335042112507\\
2.18609317208604	0.00705316040372525\\
2.19859942021011	0.00650340763270439\\
2.21110566833417	0.0059917850373634\\
2.22361191645823	0.00551607670279081\\
2.23611816458229	0.00507415784429085\\
2.24862441270635	0.0046639943904439\\
2.26113066083042	0.00428364221172457\\
2.27363690895448	0.00393124603034815\\
2.28614315707854	0.00360503804564062\\
2.2986494052026	0.00330333630775216\\
2.31115565332666	0.00302454287097773\\
2.32366190145073	0.00276714175633238\\
2.33616814957479	0.0025296967513686\\
2.34867439769885	0.00231084907353713\\
2.36243127063532	0.00209006649953723\\
2.37493751875938	0.001906180820384\\
2.38744376688344	0.00173714537910419\\
2.3999500150075	0.00158189241407161\\
2.41245626313157	0.00143941839619503\\
2.42496251125563	0.00130878133095834\\
2.43746875937969	0.0011890980628289\\
2.44997500750375	0.00107954159524468\\
2.46248125562781	0.00097933843798312\\
2.47498750375188	0.000887765992373468\\
2.48749375187594	0.000804149983541334\\
2.5	0.000727861947671615\\
};
\addlegendentry{$n=3$}

\node[overlay,right, align=left, inner sep=0]
at (rel axis cs:-0.15,1.1) {$(b)$};
\end{axis}
\end{tikzpicture}%
    \end{minipage}
    \caption{Two-dimensional isotropic spectra and correlations for the first three members of the family defined in \eqref{eq:eta_spectrum_family} and
    \eqref{eq:eta_correlation_family}. (a) Physical-space correlation function
    \(C_{00,n}^{xy}(r)\).
    (b) Corresponding two-dimensional isotropic wavenumber spectrum \(E_{\eta,n}(k)\).}
    \label{fig:eta_2d_family}
\end{figure}
To make this connection explicit, consider the family of isotropic spectra
\begin{equation}
    E_{\eta,n}(k)
    =
    \frac{2^n\mean{\eta_0^2}}{(n-1)!\,k_0}
    \left(\frac{k}{k_0}\right)^{2n-1}
    \exp\!\left[
        -2\left(\frac{k}{k_0}\right)^2
    \right],
    \qquad
    n=1,2,\ldots,
    \label{eq:eta_spectrum_family}
\end{equation}
where \(k_0\) is the characteristic horizontal wavenumber of the perturbation field. Applying \eqref{eq:eta_hankel_transform_pair} gives the corresponding isotropic correlation function
\begin{equation}
    C_{00,n}^{xy}(r)
    =
    \mean{\eta_0^2}
    \exp\!\left(-\frac{k_0^2 r^2}{8}\right)
    L_{n-1}^{(0)}
    \!\left(
        \frac{k_0^2 r^2}{8}
    \right),
    \label{eq:eta_correlation_family}
\end{equation}
where \(L_n^{(0)}\) denotes the Laguerre polynomial of degree \(n\) and order zero. \Cref{fig:eta_2d_family} shows the spectrum \(E_{\eta,n}(k)\) together with its physical-space counterpart \(C_{00,n}^{xy}(r)\) for the first three members of this family. Defining the reference length scale
\begin{equation}
    \lambda_0
    =
    \frac{2\sqrt{2}}{k_0},
    \label{eq:lambda0_k0_relation}
\end{equation}
the first three correlations are
\begin{equation}
\begin{aligned}
    C_{00,1}^{xy}(r)
    &=
    \mean{\eta_0^2}
    \exp\!\left(-\frac{r^2}{\lambda_0^2}\right),
    \\
    C_{00,2}^{xy}(r)
    &=
    \mean{\eta_0^2}
    \left(
        1-\frac{r^2}{\lambda_0^2}
    \right)
    \exp\!\left(-\frac{r^2}{\lambda_0^2}\right),
    \\
    C_{00,3}^{xy}(r)
    &=
    \mean{\eta_0^2}
    \left(
        1
        -
        2\frac{r^2}{\lambda_0^2}
        +
        \frac{1}{2}\frac{r^4}{\lambda_0^4}
    \right)
    \exp\!\left(-\frac{r^2}{\lambda_0^2}\right).
\end{aligned}
\label{eq:eta_correlation_family_examples}
\end{equation}
The first member of this family recovers the Gaussian correlation used by \citet{frame-towne-2024} to describe three-dimensional perturbation growth, but here it is tied directly to a physical quantity through the corresponding spectrum. The higher-order members produce correlations consisting of a Gaussian envelope multiplied by a polynomial, with the corresponding spectra distributing the initial variance differently across the horizontal wavenumbers.

The subsequent analysis is restricted to the family of spectra introduced in \eqref{eq:eta_spectrum_family}, although this choice is not fundamental to the formulation. Whenever either a model for the spectrum or a physical-space correlation is available, the Hankel transform pair in \eqref{eq:eta_hankel_transform_pair} provides the corresponding representation in the other domain; casting the initial statistics in spectral space therefore does more than provide an equivalent description, as it gives the covariance model a direct physical interpretation in terms of an interface perturbation that can be imposed experimentally or numerically. We stress that relaxing the isotropy assumption the same connection can be drawn, however losing the immediate connection that practitioners have with homogeneous-isotropic fields.

\section{Perturbation energy}
\label{sec:perturbation_energy}

\subsection{Instantaneous perturbation energy}
\label{subsec:instantaneous_perturbation_energy}

The main quantity that we will analyze in this work is the energy of the perturbations. Having introduced the governing equations in \cref{sec:governing_equation} and the statistics of the initial perturbations in \cref{sec:statistics_of_the_initial_perturbations}, we now provide a formal definition of the instantaneous energy for both the Boussinesq and IVD formulations. We have shown that both formulations admit reduced linearized dynamics in terms of the vertical velocity and either the concentration or density fluctuations. However, in our definition of the energy, we wish to measure the kinetic energy associated with all three velocity components, namely
\begin{equation}
    \label{eq:bq_kinetic_energy_dz_k}
    2\mathcal{K}_k(t) = \int_{-L_z}^{L_z}\left(|\hat{u}_k|^{2}+|\hat{v}_k|^{2}+|\hat{w}_k|^{2}\right)\mathrm{d}z.
\end{equation}

It is therefore necessary, once \(\hat{w}_k\) is known, to recover the horizontal velocity components \(\hat{u}_k\) and \(\hat{v}_k\). We will first derive the energy used for the Boussinesq equations, while the same arguments carry over to the incompressible variable-density formulation.

Fourier transforming the incompressibility constraint in~\eqref{eq:bq_nl_continuous_continuity} in the horizontal directions gives the first relation,
\begin{equation}
    \label{eq:bq_divergence_dz_k}
    ik_{x}\hat{u}_k+ik_{y}\hat{v}_k+\partial_{z}\hat{w}_k = 0.
\end{equation}
An additional condition is required to recover the horizontal velocity components, for which we impose vanishing vertical vorticity, \(ik_{x}\hat{v}_k-ik_{y}\hat{u}_k = 0\). Together with the incompressibility constraint, it is trivial to show that this condition is equivalent to seeking the horizontal velocity components \(\hat{u}_k\) and \(\hat{v}_k\), with \(\hat{w}_k\) prescribed by~\eqref{eq:lti_dz_k_boussinesq}, that solve the constrained optimization problem
\begin{equation}
    \label{eq:bq_kinetic_minimization}
    \begin{gathered}
        \min_{\hat{u}_k,\hat{v}_k}\left(|\hat{u}_k|^{2}+|\hat{v}_k|^{2}\right)\\
        \text{subject to}\quad ik_{x}\hat{u}_k+ik_{y}\hat{v}_k+\partial_{z}\hat{w}_k = 0.
    \end{gathered}
\end{equation}
Thus, \(\mathcal{K}_k^{BQ}(t)\) measures the lowest kinetic energy among all in-plane velocity perturbations compatible with a prescribed \(\hat{w}_k\). The kinetic energy is
\begin{equation}
    \label{eq:bq_kinetic_energy_reduced_dz_k}
    2\mathcal{K}_k^{BQ}(t) = \int_{-L_z}^{L_z}\left(|\hat{w}_k|^{2} + \frac{1}{k^{2}}\left|\partial_{z}\hat{w}_k\right|^{2}\right)\mathrm{d}z.
\end{equation}
Having dealt with the kinetic energy, it is now necessary to consider the energy of the disturbance as a whole, meaning that we need to find a consistent way to define the contribution from the concentration fluctuation.

The classical disturbance-energy norm of \citet{chu-1965}, which despite being a classical reference remains widely used~\citep{schmidt-et-al-2018,towne-schmidt-colonious-2018}, imposes two main requirements on the definition of a general disturbance energy in a flow, namely that it be a positive-definite quantity and that it decay in the absence of external forcing. For the RTI configuration, using the Boussinesq equations reported in~\eqref{eq:bq_nl_continuous}, the construction derived following the heuristic of \citet{hanifi-schmidt-henningson-1996} can be shown to remain positive definite only under stable stratification. No positive-definite quadratic energy norm exists for an unstably stratified flow, as established through several independent derivations~\citep{holliday-mcintyre-1981,shepherd-1993,winters-et-al-1995,roullet-klein-2009}. We therefore choose to represent the concentration contribution by a prescribed positive-definite quadratic measure of its amplitude~\citep{don-nils-amir-2013}. For a given horizontal wavenumber, the disturbance energy in the Boussinesq formulation is defined as
\begin{equation}
    \label{eq:bq_energy_norm_dz_k}
    \mathcal{E}_k^{BQ}(t) = \int_{-L_z}^{L_z}\left(|\hat{w}_k|^{2}+\frac{1}{k^{2}}\left|\partial_{z}\hat{w}_k\right|^{2}+|\hat{c}_k|^{2}\right)\mathrm{d}z.
\end{equation}

The energy norm for the variable-density formulation follows from the same arguments. The only modification concerns the kinematic relation used to recover the horizontal velocity components from \(\hat{w}_k\). In place of the divergence-free condition~\eqref{eq:bq_divergence_dz_k}, the divergence constraint in~\eqref{eq:vd_nl_continuous_divergence}, once linearized and Fourier transformed, constrains these components according to
\begin{equation}
    \label{eq:vd_divergence_dz_k}
    ik_{x}\hat{u}_k+ik_{y}\hat{v}_k+\partial_{z}\hat{w}_k = \widehat{\mathscr{L}}_{\rho}(t)\hat{\rho}_k.
\end{equation}
The operator \(\widehat{\mathscr{L}}_{\rho}(t)\) depends on time through the evolving base-state density \(\rho_0(z,t)\), and therefore changes as the diffusive layer evolves.
Repeating the same constrained minimization with~\eqref{eq:vd_divergence_dz_k} in place of~\eqref{eq:bq_divergence_dz_k} as the divergence constraint yields
\begin{equation}
    \label{eq:vd_kinetic_energy_reduced_dz_k}
    2\mathcal{K}_k^{IVD}(t) = \int_{-L_z}^{L_z}\left(|\hat{w}_k|^{2}+\frac{1}{k^{2}}\left|\partial_{z}\hat{w}_k-\widehat{\mathscr{L}}_{\rho}(t)\hat{\rho}_k\right|^{2}\right)\mathrm{d}z.
\end{equation}
The disturbance energy for the variable-density formulation is then
\begin{equation}
    \label{eq:vd_energy_norm_dz_k}
    \mathcal{E}_k^{IVD}(t) = \int_{-L_z}^{L_z}\left(|\hat{w}_k|^{2}+\frac{1}{k^{2}}\left|\partial_{z}\hat{w}_k-\widehat{\mathscr{L}}_{\rho}(t)\hat{\rho}_k\right|^{2}+|\hat{\rho}_k|^{2}\right)\mathrm{d}z.
\end{equation}
After discretizing the vertical direction with \(N_z\) grid points, the continuous energy expressions in \eqref{eq:bq_energy_norm_dz_k} and \eqref{eq:vd_energy_norm_dz_k} are represented by Hermitian positive-definite weight matrices in \(\mathbb{C}^{2N_z\times 2N_z}\). For the Boussinesq formulation, this matrix is independent of time and is denoted by \(\mathsfbi{W}_{k}\). For the variable-density formulation, however, the time dependence of \(\widehat{\mathscr{L}}_{\rho}(t)\) makes the corresponding weight matrix a function of time, which we denote by \(\mathsfbi{W}_{t,k}\). Using \(\mathsfbi{W}_{t,k}\) below, with \(\mathsfbi{W}_{t,k}=\mathsfbi{W}_{k}\) for the Boussinesq formulation, the induced weighted inner product on the discrete state vector \(\hat{\mathbf{q}}_k(t)\in\mathbb{C}^{2N_z}\) is \(\left\langle\mathbf{a},\mathbf{b}\right\rangle_{W_{t,k}}=\mathbf{a}^{*}\mathsfbi{W}_{t,k}\mathbf{b}\), so that
\begin{equation}
    \mathcal{E}_k(t)
    =
    \left\langle\hat{\mathbf{q}}_k(t),\hat{\mathbf{q}}_k(t)\right\rangle_{W_{t,k}}
    =
    \hat{\mathbf{q}}_k^{*}(t)\mathsfbi{W}_{t,k}\hat{\mathbf{q}}_k(t).
    \label{eq:instantaneous_energy}
\end{equation}
Throughout this work, the subscripted bracket \(\left\langle\cdot,\cdot\right\rangle_{W_{t,k}}\) denotes the \(\mathsfbi{W}_{t,k}\)-weighted inner product, while the unsubscripted bracket \(\langle\cdot\rangle\) denotes the ensemble average over random initial conditions.

A final remark concerns the total perturbation energy injected into the system. Throughout the analysis, the variance of the interface displacement is held constant, which does not correspond to keeping the initial energy constant as the thickness of the diffusive layer is varied, since the total energy injected across the diffusive layer is
\begin{equation}
    E_0
    =
    \mean{\eta_0^2}
    \int_{-\infty}^{\infty}
        C_{00}^{z}(z,z)
    \,\mathrm{d}z
    =
    \sqrt{\frac{8}{\pi}}\,
    \frac{\Atw^2}{\delta_0}
    \mean{\eta_0^2}.
    \label{eq:integrated_initial_density_variance}
\end{equation}
Therefore, for a fixed interface variance, the initial energy scales as \(\delta_0^{-1}\), so that a thinner diffusive layer corresponds to a larger amount of energy injected into the system. As we will see, this choice does not influence the evolution during the linear stage because the energy tracked throughout the analysis is normalized by its initial value; it may, however, influence the weakly nonlinear and nonlinear stages of the evolution, which lie outside the scope of this work.

\subsection{The mean energy growth}
\label{subsec:mean_energy_growth}

We now use the physical description of the initial correlations developed in \cref{sec:statistics_of_the_initial_perturbations} to define the mean energy growth introduced by \citet{frame-towne-2024}. This quantity measures perturbation amplification averaged over an ensemble of random initial conditions. With the linearized operators from \cref{sec:governing_equation}, the instantaneous disturbance energy defined above, and the initial statistics now specified, the mean energy growth at each wavenumber \(k\) is the ratio of the expected energy at time \(t\) to its initial value,
\begin{equation}
    \label{eq:mean_energy_growth_def}
    g(k,t)
    =
    \frac{\mean{\mathcal{E}_k(t)}}{\mean{\mathcal{E}_k(0)}}.
\end{equation}
Because the initial statistics separate into a horizontal and a vertical factor, only the vertical covariance \(C_{00}^{z}\) enters at a given wavenumber, as the discrete matrix \(\mathsfbi{C}_{00}^{z}\in\mathbb{C}^{2N_z\times 2N_z}\). Using the cyclic invariance of the trace, the per-wavenumber mean energy growth is then
\begin{equation}
    g(k,t)
    =
    \frac{
      \tr
      \left\{
        \mathsfbi{T}_{t,k}\mathsfbi{C}_{00}^{z}
      \right\}
    }{
      \tr
      \left\{
        \mathsfbi{C}_{00}^{z}\mathsfbi{W}_{0,k}
      \right\}
    },
    \qquad
    \mathsfbi{T}_{t,k}
    =
    \boldsymbol{\Phi}_{t,k}^{*}\mathsfbi{W}_{t,k}\boldsymbol{\Phi}_{t,k}.
    \label{eq:T_operator_definition}
\end{equation}
We define the corresponding one-dimensional growth rate associated with the per-wavenumber mean energy growth as
\begin{equation}
    \label{eq:sigma_1d}
    \overline{\sigma}_{1D}(k,t)
    =
    \frac{1}{2}
    \diff{
        \log\!\left(g(k,t)\right)
    }{t}.
\end{equation}
The three-dimensional mean energy growth follows by weighting this per-wavenumber growth with the horizontal spectral density \(\Phi_{00}^{xy}\) and integrating over all wavenumbers,
\begin{equation}
    \Gm(t)
    =
    \frac{
      \displaystyle
      \iint
      g(k,t)\,
      \Phi_{00}^{xy}(k_x,k_y)
      \,\mathrm{d}k_x\,\mathrm{d}k_y
    }{
      \displaystyle
      \iint
      \Phi_{00}^{xy}(k_x,k_y)
      \,\mathrm{d}k_x\,\mathrm{d}k_y
    }.
    \label{eq:gmean_continuous_fourier}
\end{equation}
Passing to polar coordinates gives
\begin{equation}
    \Gm(t)
    =
    \frac{2}{\mean{\eta_0^2}}
    \int_0^\infty
      g(k,t)\,
      E_{\eta}(k)\,
      \mathrm{d}k.
    \label{eq:gmean_eta_spectrum}
\end{equation}
We define the corresponding mean growth rate of the full three-dimensional perturbation field as
\begin{equation}
    \label{eq:sigma_3d}
    \overline{\sigma}_{3D}(t)
    =
    \frac{1}{2}
    \difft{
        \log\!\left(\Gm(t)\right)
    }{t}.
\end{equation}
The per-wavenumber mean energy \(g(k,t)\) accounts for the vertical correlation of perturbations within the diffusive layer when a stochastic displacement is imposed at the interface. The spectrum \(E_{\eta}(k)\) determines how the initial interface variance is distributed across horizontal scales. The three-dimensional evolution is therefore obtained by weighting the mean energy at each wavenumber by its share of the initial variance.

\subsection{Randomized trace estimation}
\label{subsec:randomized_trace_estimation}
The final expression in \eqref{eq:T_operator_definition} reduces the calculation of the per-wavenumber mean energy growth to the trace of the covariance-weighted operator \(\mathsfbi{T}_{t,k}\mathsfbi{C}_{00}^{z}\), normalized by the known initial energy. The difficulty is that the propagator \(\boldsymbol{\Phi}_{t,k}\) of the non-autonomous semi-discrete system \eqref{eq:semidiscrete_system} is defined through time integration and is not available as an explicitly formed matrix. Constructing it would require evolving a complete basis of initial conditions, after which both \(\mathsfbi{T}_{t,k}\) and its product with the covariance would have to be assembled. This calculation is precisely what we wish to avoid.

Instead, we require only the action of \(\mathsfbi{T}_{t,k}\) on a vector. This action can be obtained through the loop-adjoint procedure used in matrix-free optimal-growth calculations~\citep{luchini-2000,luchini-bottaro-2014,barkeley-blackburn-sherwin-2008}. Given an initial perturbation, the direct equations are first integrated to the target time, the resulting state defines the terminal condition for the adjoint equations, and their backward integration returns the action of \(\mathsfbi{T}_{t,k}\) at the initial time. Here, this procedure is not used within an iteration that isolates the optimally amplified disturbance, but as the building block required to evaluate the covariance-weighted trace. The resulting action is summarized in \cref{alg:loop-adjoint}.
\begin{algorithm}[H]
  \caption{Matrix-free loop-direct-adjoint (LDA) action of \(\mathsfbi{T}_{T,k}\)}
  \label{alg:loop-adjoint}

  \begin{algorithmic}[1]
    \Require Initial perturbation \(\hat{\mathbf{q}}_{0,k}\), final time \(T\),
             operators \(\mathsfbi{A}(t)\) and \(\mathsfbi{B}(t)\)
    \Ensure Matrix-free action
      \(\LDA(\hat{\mathbf{q}}_{0,k})
      :=
      \mathsfbi{T}_{T,k}\hat{\mathbf{q}}_{0,k}
      =
      \mathsfbi{B}^{\dagger}(0)\boldsymbol{\psi}(0)\)

    \State \(\hat{\mathbf{q}}(0)\gets\hat{\mathbf{q}}_{0,k}\)
      \Comment{Initial condition}

    \State Solve
      \(\mathsfbi{B}(t)\,\partial_t\hat{\mathbf{q}}(t)
      =
      \mathsfbi{A}(t)\hat{\mathbf{q}}(t)\),
      \(0\leq t\leq T\)
      \Comment{Forward solve}

    \State Define \(\boldsymbol{\psi}(T)\) by
      \(\mathsfbi{B}^{\dagger}(T)\boldsymbol{\psi}(T)
      =
      \hat{\mathbf{q}}(T)\)
      \Comment{Terminal adjoint condition}

    \State Solve
      \(\partial_t
      \left[
        \mathsfbi{B}^{\dagger}(t)\boldsymbol{\psi}(t)
      \right]
      =
      -\mathsfbi{A}^{\dagger}(t)\boldsymbol{\psi}(t)\),
      \(T\geq t\geq 0\)
      \Comment{Backward adjoint solve}

    \State \(\LDA(\hat{\mathbf{q}}_{0,k})
      \gets
      \mathsfbi{B}^{\dagger}(0)\boldsymbol{\psi}(0)\)
      \Comment{Matrix-free action}

    \State \Return \(\LDA(\hat{\mathbf{q}}_{0,k})\)
  \end{algorithmic}
\end{algorithm}

Trace estimation has been studied extensively in randomized numerical linear algebra~\citep{haim-sivan-2011}. These methods replace explicit matrix construction with random probing, in which the operator is applied to a set of random vectors and the resulting quantities are combined to estimate global properties. Randomized numerical linear algebra now provides a standard framework for approximating large matrix quantities using only matrix-vector products, particularly when the operator is too large to form or factorize explicitly~\citep{martinsson-tropp-2020}. Related ideas have been used in fluid mechanics to reduce the cost of large-scale input-output calculations and modal decompositions. For example, \citet{ribeiro-yeh-taira-2020} introduced randomized resolvent analysis, in which random test matrices sketch the resolvent operator and approximate its dominant forcing-response subspace. More recently, \citet{farghadan-martini-towne-2025} combined randomized singular value decomposition with time stepping to compute resolvent modes in fully three-dimensional flows with reduced memory and computational cost.

The randomization used here is related in spirit but different in purpose. Rather than constructing a low-rank approximation to recover dominant resolvent or propagator modes, we use random probes to estimate the trace of the covariance-weighted amplification operator. The Hutchinson estimator obtains this trace from random quadratic forms without explicitly constructing the matrix~\citep{hutchinson-1990}. Subsequent work has refined the theory and provided bounds on the number of samples required for a prescribed accuracy~\citep{khorasani-szekely-ascher-2015}. More recent algorithms such as Hutch++ reduce the sampling cost further by separating the trace into a low-rank contribution, which captures the dominant range of the operator, and a residual contribution estimated stochastically~\citep{meyer-et-al-2021}. Extensions from matrices to integral operators have also been proposed by \citet{zvonek-horning-townsend-2025}. Building on these algorithms, and using the loop-adjoint procedure to sketch the required operator, we propose the stochastic stepping trace estimator (SSTE) summarized in \cref{alg:stochastic_stepping_trace_estimator}.
\begin{algorithm}[htbp]
  \caption{Stochastic stepping trace estimator for \(\tr\left\{\mathsfbi{T}_{t,k}\mathsfbi{C}_{00}^{z}\right\}\)}
  \label{alg:stochastic_stepping_trace_estimator}

  \begin{algorithmic}[1]
    \Require Number of samples \(m\), loop-direct-adjoint action \(\LDA(\cdot)\),
             factor \(\mathsfbi{L}\) s.t.
             \(\mathsfbi{C}_{00}^{z}=\mathsfbi{L}\mathsfbi{L}^{*}\)
    \Ensure Trace estimate \(\widehat{t}_{\mathrm{H++}}\)

    \State Define kernel
      \(\mathsfbi{K}(\mathbf{g})
      =
      \mathsfbi{L}^{*}\LDA(\mathsfbi{L}\mathbf{g})\)
      \Comment{Effective operator \(\mathsfbi{K}=\mathsfbi{L}^{*}\mathsfbi{T}_{t,k}\mathsfbi{L}\) acting on a vector \(\mathbf{g}\)}

    \State Set \(s \gets m/3\)
      \Comment{Number of probes per stage}

    \For{\(i=1,\ldots,s\)}
      \State Draw \(\mathbf{g}_i \sim \mathcal{N}(\mathbf{0},\mathsfbi{I})\)
        \Comment{Gaussian probe}
      \State \(\mathbf{y}_i \gets \mathsfbi{K}(\mathbf{g}_i)\)
        \Comment{Sketch action}
    \EndFor

    \State Form
      \(\mathsfbi{Y}
      \gets
      [\,\mathbf{y}_1\ \cdots\ \mathbf{y}_s\,]\)
      \Comment{Range sketch}

    \State Compute
      \(\mathsfbi{Y}=\mathsfbi{Q}\mathsfbi{R}\)
      \Comment{Thin QR factorization}

    \State \(t_d \gets 0\)
    \For{\(j=1,\ldots,s\)}
      \State \(t_d \gets t_d + \mathbf{q}_j^{*}\mathsfbi{K}(\mathbf{q}_j)\)
        \Comment{Low-rank trace contribution}
    \EndFor

    \State \(t_r \gets 0\)
    \For{\(i=1,\ldots,s\)}
      \State Draw \(\mathbf{g}_i \sim \mathcal{N}(\mathbf{0},\mathsfbi{I})\)
        \Comment{Residual probe}
      \State \(\mathbf{r}_i \gets (\mathsfbi{I}-\mathsfbi{Q}\mathsfbi{Q}^{*})\mathbf{g}_i\)
        \Comment{Orthogonal projection}
      \State \(t_r \gets t_r + s^{-1}\mathbf{r}_i^{*}\mathsfbi{K}(\mathbf{r}_i)\)
        \Comment{Stochastic residual trace}
    \EndFor

    \State \(\widehat{t}_{\mathrm{H++}} \gets t_d+t_r\)
      \Comment{Hutch++ estimate}

    \State \Return \(\widehat{t}_{\mathrm{H++}}\)
  \end{algorithmic}
\end{algorithm}
Extensive verification against analytical solutions and reference problems is reported in \Cref{sec:appendix_main}, \cref{sec:appendix_verification}. We also offer a detailed description of other methods that can be used and of where the stochastic trace estimation offers advantages over them.

\section{Results and discussion for the Boussinesq approximation}
\label{sec:boussinesq}

\subsection{Verification of the method against direct numerical simulations}
\label{subsec:dns_verification}
Having described the framework, we now present the results of our analysis, beginning with the
Boussinesq model. Given the novelty and stochastic character of
the framework, we first validate the formulation by comparing the predicted mean energy growth
with ensembles of three-dimensional direct numerical simulations.
The simulations are performed using the triply periodic pseudospectral solver
PsDNS\footnote{\url{https://github.com/lanl/PsDNS.git}}, developed at LANL to solve the
incompressible Navier-Stokes equations with buoyancy forcing. All simulations use
\(\Sch=\Rey=1\) and \(\delta_0=1\), with the velocity field initially set to
zero. The triply periodic domain has dimensions
\(L_x\times L_y\times L_z=80\pi\times80\pi\times20\pi\) and is discretized using
\(N_x\times N_y\times N_z=256\times256\times256\) Fourier modes and a time step
\(\Delta t=0.008\). Time advancement uses an explicit fourth-order Runge-Kutta scheme. The
computational grid is chosen to represent all modes within \(2k_{\min}\leq k\leq k_{\max}\),
where \(k_{\min}=0.025\) and \(k_{\max}=2\), and to capture their evolution during the linear
stage. The simulations are stopped well before the fully nonlinear regime is reached, given the limited
computational grid.

The simulations are initialized by prescribing an interface perturbation
\(z_i(x,y)=L_z/2+\eta(x,y)\), where
\begin{equation}
  \eta(x,y)
  =
  \sum_{(p,q)\in\mathcal{I}}
  A_{pq}
  \cos\!\left(\frac{2\pi p}{L_x}x+\varphi_{pq}\right)
  \cos\!\left(\frac{2\pi q}{L_y}y+\psi_{pq}\right).
  \label{eq:dns_eta_construction}
\end{equation}
Here, \(\mathcal{I}\) denotes the set of index pairs associated with the wavenumbers represented on
the horizontal grid,
\begin{equation}
  \mathcal{I}
  =
  \left\{
    (p,q)\in\mathbb{Z}^2:
    k_{\min}^2
    \leq
    \left(\frac{2\pi p}{L_x}\right)^2
    +
    \left(\frac{2\pi q}{L_y}\right)^2
    \leq k_{\max}^2
  \right\}.
  \label{eq:dns_retained_wavevectors}
\end{equation}
The phases \(\varphi_{pq}\) and \(\psi_{pq}\) are
independently sampled from \(\mathcal{U}[0,2\pi)\), while the modal amplitudes \(A_{pq}\) are
zero-mean Gaussian random variables satisfying
\begin{equation}
  \mean{A_{pq}A_{p'q'}}
  =
  \mean{A_{pq}^{2}}\,\delta_{pp'}\delta_{qq'}.
  \label{eq:dns_uncorrelated_amplitudes}
\end{equation}
Both the phases and the amplitudes are sampled independently for every realization, and therefore
the two-point correlation in space can be shown to be
\begin{equation}
  C_{\eta\eta}(r_x,r_y)
  =
  \frac{1}{4}
  \sum_{(p,q)\in\mathcal{I}}
  \mean{A_{pq}^{2}}
  \cos\!\left(\frac{2\pi p}{L_x}r_x\right)
  \cos\!\left(\frac{2\pi q}{L_y}r_y\right),
  \label{eq:dns_eta_correlation}
\end{equation}
which depends only on the separation \((r_x,r_y)\).
\begin{figure}[t]
  \centering
  \begin{minipage}{0.48\textwidth}
    \centering
    % This file was created by matlab2tikz.
%
\definecolor{mycolor1}{rgb}{0.91400,0.56100,0.47800}%
\definecolor{mycolor2}{rgb}{0.83900,0.37600,0.30200}%
\definecolor{mycolor3}{rgb}{0.64700,0.12900,0.14500}%
\definecolor{mycolor4}{rgb}{0.40400,0.00000,0.12200}%
%
\begin{tikzpicture}

\begin{axis}[%
width=11.252in,
height=7.131in,
at={(1.887in,0.962in)},
scale only axis,
clip=true,
separate axis lines,
every outer x axis line/.append style={black},
every x tick label/.append style={font=\color{black}},
every x tick/.append style={black},
xmin=0,
xmax=10,
xlabel={$t$},
every outer y axis line/.append style={black},
every y tick label/.append style={font=\color{black}},
every y tick/.append style={black},
ymode=log,
ymin=0.1,
ymax=100,
yminorticks=true,
ylabel={$\mathscr{G}^{\mathrm{mean}}(t)$},
axis background/.style={fill=white},
legend style={at={(0.03,0.97)}, anchor=north west, legend cell align=left, align=left},
clip=true,
clip mode=individual,
axis lines=box,
xtick pos=bottom,
ytick pos=left,
tick align=inside,
major tick length=2pt,
width=0.8\linewidth,
height=0.8\linewidth,
every axis/.append style={font=\fontsize{9}{9}\selectfont},xlabel style={font=\fontsize{9}{9}\selectfont},title style={font=\fontsize{9}{9}\selectfont},
legend style={font=\fontsize{8}{9}\selectfont,inner sep=1pt,row sep=1pt,column sep=2pt,nodes={scale=0.85},draw=none},legend image post style={xscale=0.35},
legend columns=1,
scaled x ticks=false,
tick label style={/pgf/number format/fixed,/pgf/number format/precision=2},
every x tick label/.append style={font=\fontsize{9}{9}\selectfont\color{black}},
every y tick label/.append style={font=\fontsize{9}{9}\selectfont\color{black}},
,,
colormap={sigmamap}{rgb(0.0000)=(0.0200,0.1880,0.3800); rgb(0.1000)=(0.0743,0.3456,0.5616); rgb(0.2000)=(0.1918,0.4980,0.7060); rgb(0.3000)=(0.3890,0.6708,0.8226); rgb(0.4000)=(0.7552,0.8682,0.9228); rgb(0.5000)=(0.9690,0.9689,0.9689); rgb(0.6000)=(0.9425,0.8044,0.7614); rgb(0.7000)=(0.8767,0.4910,0.4020); rgb(0.8000)=(0.7869,0.2866,0.2470); rgb(0.9000)=(0.6186,0.1239,0.1568); rgb(1.0000)=(0.4040,0.0000,0.1220)}
]
\addplot [color=mycolor1, line width=1.8pt]
  table[row sep=crcr]{%
0	1\\
0.08	0.863137158521505\\
0.168	0.765296195121739\\
0.248	0.70474397683936\\
0.336	0.658588067607769\\
0.416	0.629965624874714\\
0.504	0.609519584109834\\
0.584	0.598892908179321\\
0.672	0.594352800620287\\
0.752000000000001	0.595723893901938\\
0.840000000000001	0.602453019766551\\
0.920000000000001	0.612775592295249\\
1	0.626711973705201\\
1.088	0.645872140261718\\
1.168	0.666532028500899\\
1.256	0.692615375910713\\
1.336	0.719241218543471\\
1.424	0.75161809759949\\
1.504	0.783785847991089\\
1.592	0.822122085993051\\
1.672	0.859627742041189\\
1.752	0.899650303696986\\
1.84	0.94658590309429\\
1.92	0.991917281242138\\
2.008	1.04474302259802\\
2.088	1.09549773591915\\
2.176	1.15438823944275\\
2.256	1.21076558561448\\
2.344	1.27598138132431\\
2.424	1.33825346035854\\
2.512	1.41013060368432\\
2.592	1.47863494697464\\
2.672	1.55024351361084\\
2.76	1.63271408325125\\
2.84	1.71116430192978\\
2.928	1.80141989337375\\
3.008	1.88719744834199\\
3.096	1.98580457936394\\
3.176	2.07945400275672\\
3.264	2.18704447946353\\
3.344	2.28917040355475\\
3.432	2.40644321391367\\
3.512	2.51771252197615\\
3.592	2.6335889910158\\
3.68	2.76658138177793\\
3.76	2.89270580804419\\
3.848	3.03742035263772\\
3.928	3.1746277528457\\
4.016	3.33202369135629\\
4.096	3.48122430690785\\
4.184	3.65234635855449\\
4.264	3.81453076764412\\
4.344	3.98327082575806\\
4.432	4.17675989630082\\
4.512	4.36010508177377\\
4.6	4.57031481138781\\
4.68	4.76948043259263\\
4.768	4.99780315637551\\
4.848	5.21410727490269\\
4.936	5.46205258737654\\
5.016	5.69692387401654\\
5.104	5.96612764098106\\
5.184	6.22111391381719\\
5.264	6.48625964894577\\
5.352	6.79012553863693\\
5.432	7.07790869573738\\
5.52	7.40769197688594\\
5.6	7.71999688768499\\
5.688	8.07785384695352\\
5.768	8.41671957503659\\
5.856	8.80498325633162\\
5.936	9.17261573798088\\
6.024	9.59380981525452\\
6.104	9.99259495760747\\
6.184	10.407109026888\\
6.272	10.8819664995123\\
6.352	11.3315143903061\\
6.44	11.8464706574774\\
6.52	12.3339470364914\\
6.608	12.8923125156637\\
6.68800000000001	13.4208459890154\\
6.77600000000001	14.0261982991293\\
6.85600000000001	14.5991696840653\\
6.93600000000001	15.1945683338136\\
7.02400000000001	15.876436477224\\
7.10400000000001	16.5217667285783\\
7.19200000000001	17.2607682106632\\
7.27200000000001	17.9601231757078\\
7.36000000000001	18.7609370534059\\
7.44000000000001	19.5187369929625\\
7.52800000000001	20.3864158406654\\
7.60800000000001	21.207433434276\\
7.69600000000001	22.1474321520764\\
7.77600000000001	23.0368196368196\\
7.85600000000001	23.9606924910561\\
7.94400000000001	25.0183446871354\\
8.024	26.018949051088\\
8.11199999999999	27.1643640821777\\
8.19199999999998	28.2479219793519\\
8.27999999999997	29.4882088762784\\
8.35999999999997	30.6614326209336\\
8.44799999999996	32.0042598205101\\
8.52799999999995	33.2743891696673\\
8.61599999999994	34.7280268165682\\
8.69599999999993	36.1028687705243\\
8.77599999999992	37.5304930532827\\
8.86399999999991	39.1642054407962\\
8.9439999999999	40.7091946766412\\
9.03199999999989	42.4770841710685\\
9.11199999999988	44.1488397483288\\
9.19999999999987	46.0616415823677\\
9.27999999999986	47.870293520584\\
9.36799999999985	49.9395735308096\\
9.44799999999985	51.8960342800602\\
9.52799999999984	53.9269569894389\\
9.61599999999983	56.2502719535687\\
9.69599999999982	58.4466649750828\\
9.78399999999981	60.9590746898465\\
9.8639999999998	63.3340433754666\\
9.95199999999979	66.0505038538592\\
10.0319999999998	68.6181527059567\\
10.1199999999998	71.5547578800136\\
10.1999999999998	74.3302622980936\\
10.2879999999998	77.504325438466\\
10.3679999999997	80.5040080619984\\
10.4479999999997	83.6165323444928\\
10.5359999999997	87.17556024128\\
10.6159999999997	90.5386279632922\\
10.7039999999997	94.3838016223578\\
10.7839999999997	98.0169362234784\\
10.8719999999997	102.17051665725\\
10.9519999999997	106.09469307459\\
11.0399999999997	110.580591951725\\
11.1199999999997	114.818336213295\\
11.2079999999997	119.66222855866\\
11.2879999999996	124.237723147983\\
11.3679999999996	128.983059192458\\
11.4559999999996	134.406361683826\\
11.5359999999996	139.528410766155\\
11.6239999999996	145.381653212772\\
11.7039999999996	150.909191711203\\
11.7919999999996	157.225149158398\\
11.8719999999996	163.189023685422\\
11.9599999999996	170.002820882748\\
12.0399999999996	176.436081620031\\
12.1199999999996	183.105150249691\\
12.2079999999995	190.723383948328\\
12.2879999999995	197.914952576998\\
12.3759999999995	206.129096813346\\
12.4559999999995	213.882288496777\\
12.5439999999995	222.736852337883\\
12.6239999999995	231.093508923673\\
12.7119999999995	240.636080951151\\
12.7919999999995	249.640926243071\\
12.8799999999995	259.922357664499\\
12.9599999999995	269.623167158641\\
13.0399999999995	279.673109577726\\
13.1279999999994	291.14553911633\\
13.2079999999994	301.967923572421\\
13.2959999999994	314.320409887543\\
13.3759999999994	325.971339675488\\
13.4639999999994	339.267599744497\\
13.5439999999994	351.806868281739\\
13.6319999999994	366.114785496255\\
13.7119999999994	379.60607263263\\
13.7919999999994	393.574261837076\\
13.8799999999994	409.508988640465\\
13.9599999999993	424.530730198167\\
14.0479999999993	441.664553958808\\
14.1279999999993	457.814025783034\\
14.2159999999993	476.231044880986\\
14.2959999999993	493.587027276652\\
14.3839999999993	513.376519419052\\
14.4639999999993	532.02260770615\\
14.5519999999993	553.279264672947\\
14.6319999999993	573.304083115019\\
14.7119999999993	594.018101769329\\
14.7999999999993	617.62566226783\\
14.8799999999992	639.858881816661\\
14.9679999999992	665.192894843786\\
15.0479999999992	689.047311133611\\
15.1359999999992	716.223115297006\\
15.2159999999992	741.806489943681\\
15.3039999999992	770.94587927322\\
15.3839999999992	798.371943886715\\
15.4719999999992	829.603389177732\\
15.5519999999992	858.992043137176\\
15.6319999999992	889.359242283781\\
15.7199999999992	923.928405098054\\
15.7999999999991	956.446859195702\\
15.8879999999991	993.456276516502\\
15.9679999999991	1028.26195046884\\
16.0559999999991	1067.86489056044\\
16.1359999999991	1105.10052651647\\
16.2239999999991	1147.45779408721\\
16.3039999999991	1187.27305534548\\
16.3839999999991	1228.37135237476\\
16.4719999999991	1275.10473645266\\
16.5519999999991	1319.01650020257\\
16.6399999999991	1368.93573177344\\
16.719999999999	1415.82825788342\\
16.807999999999	1469.12134783689\\
16.887999999999	1519.16919198034\\
16.975999999999	1576.032162967\\
17.055999999999	1629.41718024179\\
17.143999999999	1690.05409305365\\
17.223999999999	1746.9654457506\\
17.303999999999	1805.62460539367\\
17.391999999999	1872.22251949208\\
17.471999999999	1934.70054200388\\
17.559999999999	2005.61212405685\\
17.6399999999989	2072.11606621728\\
17.7279999999989	2147.57297085521\\
17.8079999999989	2218.31692861305\\
17.8959999999989	2298.55846287473\\
17.9759999999989	2373.76339420429\\
18.0639999999989	2459.03629466115\\
18.1439999999989	2538.92979045099\\
18.2239999999989	2621.13778212532\\
18.3119999999989	2714.30393956954\\
18.3919999999989	2801.54796650997\\
18.4799999999989	2900.38646595288\\
18.5599999999988	2992.90918580977\\
18.6479999999988	3097.68996918807\\
18.7279999999988	3195.73940996572\\
18.8159999999988	3306.73805286356\\
18.8959999999988	3410.56709685279\\
18.9759999999988	3517.23847523236\\
19.0639999999988	3637.9299897827\\
19.1439999999988	3750.76186920952\\
19.2319999999988	3878.3737456284\\
19.3119999999988	3997.62830110731\\
19.3999999999988	4132.45038329951\\
19.4799999999987	4258.39221937926\\
19.5679999999987	4400.71654244535\\
19.6479999999987	4533.61179447231\\
19.7359999999987	4683.73154260114\\
19.8159999999987	4823.84687966257\\
19.8959999999987	4967.49396010987\\
19.9839999999987	5129.65675901201\\
20.0639999999987	5280.91619927742\\
20.1519999999987	5451.59767458955\\
20.2319999999987	5610.73250814424\\
20.3199999999987	5790.22012642598\\
20.3999999999986	5957.48963623982\\
20.4879999999986	6146.06594037093\\
20.5679999999986	6321.72418019434\\
20.6559999999986	6519.66510112123\\
20.7359999999986	6703.95930470078\\
20.8159999999986	6892.45036347861\\
20.9039999999986	7104.70122936248\\
20.9839999999986	7302.17711952195\\
21.0719999999986	7524.43548575096\\
21.1519999999986	7731.11899942858\\
21.2399999999985	7963.62295259907\\
21.3199999999985	8179.72368748792\\
21.4079999999985	8422.69568105594\\
21.4879999999985	8648.40798319373\\
21.5679999999985	8878.76056699721\\
21.6559999999985	9137.55334768161\\
21.7359999999985	9377.77195965913\\
21.8239999999985	9647.50128907282\\
21.9039999999985	9897.73313659324\\
21.9919999999985	10178.5486148553\\
22.0719999999985	10438.9179203066\\
22.1599999999984	10730.9425847312\\
22.2399999999984	11001.5482536085\\
22.3279999999984	11304.876003637\\
22.4079999999984	11585.7892803033\\
22.4879999999984	11871.6273688968\\
22.5759999999984	12191.7444388079\\
22.6559999999984	12487.9398226717\\
22.7439999999984	12819.4506258254\\
22.8239999999984	13125.9961152975\\
22.9119999999984	13468.8734906897\\
22.9919999999984	13785.7263701486\\
23.0799999999983	14139.902931539\\
23.1599999999983	14466.9829120455\\
23.2479999999983	14832.3488873319\\
23.3279999999983	15169.536238794\\
23.4079999999983	15511.4826799179\\
23.4959999999983	15893.0714475953\\
23.5759999999983	16244.873025402\\
23.6639999999983	16637.1852018458\\
23.7439999999983	16998.6180422679\\
23.8319999999983	17401.3833149078\\
23.9119999999983	17772.1796362877\\
23.9999999999982	18185.0790511946\\
24.0799999999982	18564.92664146\\
24.1599999999982	18948.9616998966\\
24.2479999999982	19376.134378848\\
24.3279999999982	19768.6785234199\\
24.4159999999982	20204.9851312068\\
24.4959999999982	20605.616638405\\
24.5839999999982	21050.5695218467\\
24.6639999999982	21458.823763223\\
24.7519999999982	21911.8887809043\\
24.8319999999982	22327.2596114919\\
24.9199999999981	22787.8578890932\\
24.9999999999981	23209.7994473818\\
};
\addlegendentry{$\lambda_0=7.07$}

\addplot [color=mycolor2, line width=1.8pt]
  table[row sep=crcr]{%
0	1\\
0.08	0.832819704066107\\
0.168	0.714683253153248\\
0.248	0.64222558102442\\
0.336	0.587026730695588\\
0.416	0.552268248773759\\
0.504	0.526239770653693\\
0.584	0.510957051120003\\
0.672	0.501266897007025\\
0.752000000000001	0.497634260728043\\
0.840000000000001	0.498268853598173\\
0.920000000000001	0.502365080522346\\
1	0.509317320266773\\
1.088	0.519829321027364\\
1.168	0.531685881906803\\
1.256	0.546992780402698\\
1.336	0.562786767973008\\
1.424	0.582069226497625\\
1.504	0.601228389733645\\
1.592	0.624005492144949\\
1.672	0.646200559436219\\
1.752	0.669773530817573\\
1.84	0.697265048201602\\
1.92	0.723660903198896\\
2.008	0.754233967753364\\
2.088	0.783428720554656\\
2.176	0.817096398143766\\
2.256	0.849133119079866\\
2.344	0.885972878230917\\
2.424	0.920946386854043\\
2.512	0.961086678958327\\
2.592	0.999133822099627\\
2.672	1.03870327442045\\
2.76	1.08403907359157\\
2.84	1.12694859775924\\
2.928	1.17607504326586\\
3.008	1.22254394253983\\
3.096	1.27571855160338\\
3.176	1.32599548507158\\
3.264	1.38350763017961\\
3.344	1.4378700240055\\
3.432	1.50004071072415\\
3.512	1.55879489320505\\
3.592	1.61975661699255\\
3.68	1.68945912810072\\
3.76	1.75531924536129\\
3.848	1.8306156273392\\
3.928	1.90175594582726\\
4.016	1.98308413940995\\
4.096	2.05991956005876\\
4.184	2.14775496336989\\
4.264	2.23073547206347\\
4.344	2.31681111790574\\
4.432	2.41520624420687\\
4.512	2.50816027146618\\
4.6	2.61441684500497\\
4.68	2.71479653181068\\
4.768	2.82954045681336\\
4.848	2.93793729752015\\
4.936	3.06184486519279\\
5.016	3.17889782063192\\
5.104	3.31269944234446\\
5.184	3.43909847985704\\
5.264	3.57020629587451\\
5.352	3.72007268133342\\
5.432	3.86164654232403\\
5.52	4.02347541800629\\
5.6	4.17634880784447\\
5.688	4.35109252457545\\
5.768	4.51616469792011\\
5.856	4.70485067029701\\
5.936	4.88309168395083\\
6.024	5.08682813697658\\
6.104	5.27928427320722\\
6.184	5.47890083099329\\
6.272	5.70706604279049\\
6.352	5.92259403335802\\
6.44	6.16894259700932\\
6.52	6.40164314932235\\
6.608	6.66761551252459\\
6.68800000000001	6.91884828099154\\
6.77600000000001	7.20599734068389\\
6.85600000000001	7.4772280421583\\
6.93600000000001	7.75852952676515\\
7.02400000000001	8.08003625929708\\
7.10400000000001	8.38371051461798\\
7.19200000000001	8.73077992469615\\
7.27200000000001	9.0585914721146\\
7.36000000000001	9.43323849037609\\
7.44000000000001	9.78708878800768\\
7.52800000000001	10.1914845975571\\
7.60800000000001	10.5734224450824\\
7.69600000000001	11.0099063428524\\
7.77600000000001	11.4221390556265\\
7.85600000000001	11.8496226965924\\
7.94400000000001	12.338136595384\\
8.024	12.7994888607299\\
8.11199999999999	13.3266909544546\\
8.19199999999998	13.8245651887723\\
8.27999999999997	14.3934844311703\\
8.35999999999997	14.9307381188624\\
8.44799999999996	15.544636302913\\
8.52799999999995	16.1243463966578\\
8.61599999999994	16.7867356161208\\
8.69599999999993	17.4122149745842\\
8.77599999999992	18.0607264519982\\
8.86399999999991	18.8016913561581\\
8.9439999999999	19.5013313013094\\
9.03199999999989	20.3006845496707\\
9.11199999999988	21.0554285088383\\
9.19999999999987	21.9177068050186\\
9.27999999999986	22.7318331052257\\
9.36799999999985	23.6619183172831\\
9.44799999999985	24.5400304522582\\
9.52799999999984	25.4503303706021\\
9.61599999999983	26.4902256829824\\
9.69599999999982	27.471952531396\\
9.78399999999981	28.5933966072787\\
9.8639999999998	29.6520657613237\\
9.95199999999979	30.8613503027007\\
10.0319999999998	32.0028930444898\\
10.1199999999998	33.3067836212315\\
10.1999999999998	34.5375773981777\\
10.2879999999998	35.94334785353\\
10.3679999999997	37.2702486049331\\
10.4479999999997	38.6454717214855\\
10.5359999999997	40.216094478701\\
10.6159999999997	41.6984934484664\\
10.7039999999997	43.3914369334648\\
10.7839999999997	44.9892060444985\\
10.8719999999997	46.8138131333496\\
10.9519999999997	48.5357554331821\\
11.0399999999997	50.5020620132785\\
11.1199999999997	52.3576324831401\\
11.2079999999997	54.4764164688637\\
11.2879999999996	56.4757680051309\\
11.3679999999996	58.5473699201293\\
11.4559999999996	60.9126338752825\\
11.5359999999996	63.1443844395034\\
11.6239999999996	65.6923506033964\\
11.7039999999996	68.0963481299812\\
11.7919999999996	70.8408016715023\\
11.8719999999996	73.4300262183668\\
11.9599999999996	76.3857549725686\\
12.0399999999996	79.174128762386\\
12.1199999999996	82.0625182054445\\
12.2079999999995	85.3594431792668\\
12.2879999999995	88.4693923365339\\
12.3759999999995	92.0189746995615\\
12.4559999999995	95.3670226702141\\
12.5439999999995	99.1880950766263\\
12.6239999999995	102.791962620024\\
12.7119999999995	106.904699876317\\
12.7919999999995	110.7833674195\\
12.8799999999995	115.209374662207\\
12.9599999999995	119.383164650541\\
13.0399999999995	123.705043043727\\
13.1279999999994	128.636234107579\\
13.2079999999994	133.285872554186\\
13.2959999999994	138.590597272469\\
13.3759999999994	143.592024813775\\
13.4639999999994	149.29761665036\\
13.5439999999994	154.676526269617\\
13.6319999999994	160.812205186686\\
13.7119999999994	166.596057129652\\
13.7919999999994	172.582930166104\\
13.8799999999994	179.411175985334\\
13.9599999999993	185.846988391208\\
14.0479999999993	193.186554890728\\
14.1279999999993	200.103619506876\\
14.2159999999993	207.991224962986\\
14.2959999999993	215.42401975176\\
14.3839999999993	223.898832507869\\
14.4639999999993	231.884126200704\\
14.5519999999993	240.987905604438\\
14.6319999999993	249.564888284411\\
14.7119999999993	258.438138380534\\
14.7999999999993	268.55254952509\\
14.8799999999992	278.080054027215\\
14.9679999999992	288.938939211143\\
15.0479999999992	299.166476641893\\
15.1359999999992	310.821773632205\\
15.2159999999992	321.798034377384\\
15.3039999999992	334.304966384045\\
15.3839999999992	346.081703471949\\
15.4719999999992	359.49894713095\\
15.5519999999992	372.131130785839\\
15.6319999999992	385.190963610323\\
15.7199999999992	400.066963646752\\
15.7999999999991	414.069604704096\\
15.8879999999991	430.017186045795\\
15.9679999999991	445.026257624301\\
16.0559999999991	462.117460068169\\
16.1359999999991	478.200361895172\\
16.2239999999991	496.511471275604\\
16.3039999999991	513.739545957733\\
16.3839999999991	531.539201830076\\
16.4719999999991	551.799975664408\\
16.5519999999991	570.857677572132\\
16.6399999999991	592.546702768952\\
16.719999999999	612.944270612799\\
16.807999999999	636.154011476638\\
16.887999999999	657.977782269284\\
16.975999999999	682.805760913825\\
17.055999999999	706.146749751861\\
17.143999999999	732.695721163306\\
17.223999999999	757.649778244723\\
17.303999999999	783.407093022594\\
17.391999999999	812.695862676171\\
17.471999999999	840.216895538401\\
17.559999999999	871.504696271981\\
17.6399999999989	900.897908690188\\
17.7279999999989	934.306962753992\\
17.8079999999989	965.686146880608\\
17.8959999999989	1001.34460881\\
17.9759999999989	1034.82900925262\\
18.0639999999989	1072.87109322147\\
18.1439999999989	1108.58551636145\\
18.2239999999989	1145.40735911637\\
18.3119999999989	1187.22651772528\\
18.3919999999989	1226.47298320323\\
18.4799999999989	1271.0348660341\\
18.5599999999988	1312.84490592299\\
18.6479999999988	1360.30555354604\\
18.7279999999988	1404.82393606636\\
18.8159999999988	1455.34580594717\\
18.8959999999988	1502.72314416194\\
18.9759999999988	1551.51652370401\\
19.0639999999988	1606.86788436829\\
19.1439999999988	1658.75323700946\\
19.2319999999988	1717.59576102619\\
19.3119999999988	1772.73813143835\\
19.3999999999988	1835.25653424755\\
19.4799999999987	1893.82670321931\\
19.5679999999987	1960.21195865531\\
19.6479999999987	2022.38634103213\\
19.7359999999987	2092.83554912661\\
19.8159999999987	2158.79605475464\\
19.8959999999987	2226.6249628602\\
19.9839999999987	2303.44612166693\\
20.0639999999987	2375.3393392842\\
20.1519999999987	2456.73764333926\\
20.2319999999987	2532.88982302757\\
20.3199999999987	2619.0820061984\\
20.3999999999986	2699.69247871382\\
20.4879999999986	2790.90023855902\\
20.5679999999986	2876.17266815102\\
20.6559999999986	2972.6222567314\\
20.7359999999986	3062.76424098217\\
20.8159999999986	3155.29905285585\\
20.9039999999986	3259.90860328515\\
20.9839999999986	3357.62562395908\\
21.0719999999986	3468.05365652886\\
21.1519999999986	3571.16811317298\\
21.2399999999985	3687.65259934521\\
21.3199999999985	3796.38177877745\\
21.4079999999985	3919.16258473917\\
21.4879999999985	4033.72514112161\\
21.5679999999985	4151.14255264835\\
21.6559999999985	4283.65841360933\\
21.7359999999985	4407.23283276937\\
21.8239999999985	4546.64177245601\\
21.9039999999985	4676.59177382029\\
21.9919999999985	4823.13345958748\\
22.0719999999985	4959.67598023849\\
22.1599999999984	5113.58771399828\\
22.2399999999984	5256.93700710458\\
22.3279999999984	5418.45252807397\\
22.4079999999984	5568.81903229563\\
22.4879999999984	5722.59961858051\\
22.5759999999984	5895.75635116913\\
22.6559999999984	6056.8552298482\\
22.7439999999984	6238.17069437374\\
22.8239999999984	6406.78345125953\\
22.9119999999984	6596.46837152687\\
22.9919999999984	6772.78219587675\\
23.0799999999983	6971.03724592516\\
23.1599999999983	7155.2294326121\\
23.2479999999983	7362.2435481526\\
23.3279999999983	7554.47992673137\\
23.4079999999983	7750.60142019731\\
23.4959999999983	7970.86183631874\\
23.5759999999983	8175.24787935207\\
23.6639999999983	8404.67318025266\\
23.7439999999983	8617.4540406147\\
23.8319999999983	8856.17819664035\\
23.9119999999983	9077.46674987961\\
23.9999999999982	9325.60361414408\\
24.0799999999982	9555.49354482151\\
24.1599999999982	9789.50627172294\\
24.2479999999982	10051.6988461908\\
24.3279999999982	10294.4123213728\\
24.4159999999982	10566.2001332452\\
24.4959999999982	10817.6526950032\\
24.5839999999982	11099.0641393223\\
24.6639999999982	11359.2686401108\\
24.7519999999982	11650.3031085644\\
24.8319999999982	11919.2451672592\\
24.9199999999981	12219.8711603627\\
24.9999999999981	12497.5074839751\\
};
\addlegendentry{$\lambda_0=3.54$}

\addplot [color=mycolor3, line width=1.8pt]
  table[row sep=crcr]{%
0	1\\
0.08	0.787130537175682\\
0.168	0.641812316602565\\
0.248	0.554861742170152\\
0.336	0.48930183522127\\
0.416	0.447683888717454\\
0.504	0.415395635661876\\
0.584	0.394846070670004\\
0.672	0.379349754121433\\
0.752000000000001	0.370187242086496\\
0.840000000000001	0.3643070861304\\
0.920000000000001	0.362011817534288\\
1	0.362085889871243\\
1.088	0.364446050548875\\
1.168	0.368348884565332\\
1.256	0.374304587826573\\
1.336	0.381048679812777\\
1.424	0.389771796648539\\
1.504	0.398780956387305\\
1.592	0.409784859888156\\
1.672	0.420720356101917\\
1.752	0.432498409841689\\
1.84	0.446389190288313\\
1.92	0.459842273724704\\
2.008	0.47552959569955\\
2.088	0.4905894389037\\
2.176	0.508029665397057\\
2.256	0.524680709516023\\
2.344	0.543879246027102\\
2.424	0.562144098958161\\
2.512	0.583142778536841\\
2.592	0.603073146682235\\
2.672	0.623821731052117\\
2.76	0.647613294866366\\
2.84	0.670145373196348\\
2.928	0.695953324811225\\
3.008	0.720372399405414\\
3.096	0.748320070637292\\
3.176	0.774746570521144\\
3.264	0.804975322031563\\
3.344	0.833545689476253\\
3.432	0.866214235989974\\
3.512	0.897080568287754\\
3.592	0.929098220596062\\
3.68	0.965694922365792\\
3.76	1.00026185913691\\
3.848	1.03976596355549\\
3.928	1.07707389590871\\
4.016	1.11970554217061\\
4.096	1.15996321067151\\
4.184	1.20596173994776\\
4.264	1.24939581223081\\
4.344	1.29442759105286\\
4.432	1.345876988159\\
4.512	1.39445483927423\\
4.6	1.44995378926163\\
4.68	1.50235369279386\\
4.768	1.56221783157325\\
4.848	1.61873806414432\\
4.936	1.68330840885843\\
5.016	1.74427113479602\\
5.104	1.81391591619913\\
5.184	1.87966892883598\\
5.264	1.94783317961545\\
5.352	2.02570420894914\\
5.432	2.09922305301262\\
5.52	2.18321070776634\\
5.6	2.2625039057872\\
5.688	2.35308771125838\\
5.768	2.4386079577897\\
5.856	2.53630501403917\\
5.936	2.62854043439975\\
6.024	2.73390826935368\\
6.104	2.83338511326685\\
6.184	2.93650700345734\\
6.272	3.05431030436585\\
6.352	3.16552634660867\\
6.44	3.29257529284389\\
6.52	3.41251911511942\\
6.608	3.54953727681315\\
6.68800000000001	3.67889172308447\\
6.77600000000001	3.82665883497539\\
6.85600000000001	3.96615971825062\\
6.93600000000001	4.11076723387684\\
7.02400000000001	4.27595610153293\\
7.10400000000001	4.43190169054578\\
7.19200000000001	4.61004027513143\\
7.27200000000001	4.77820893937936\\
7.36000000000001	4.97030766447877\\
7.44000000000001	5.15165276827184\\
7.52800000000001	5.35880006256407\\
7.60800000000001	5.55434852739859\\
7.69600000000001	5.77771674867596\\
7.77600000000001	5.98857463608754\\
7.85600000000001	6.20713593631519\\
7.94400000000001	6.4567849477594\\
8.024	6.69244588535137\\
8.11199999999999	6.96162203453913\\
8.19199999999998	7.21571140745643\\
8.27999999999997	7.50593145740044\\
8.35999999999997	7.77988004407015\\
8.44799999999996	8.09277709221404\\
8.52799999999995	8.38812534238097\\
8.61599999999994	8.72545758684236\\
8.69599999999993	9.04386387887826\\
8.77599999999992	9.37386949719721\\
8.86399999999991	9.75077353418169\\
8.9439999999999	10.106519744675\\
9.03199999999989	10.5128131073139\\
9.11199999999988	10.8962898368524\\
9.19999999999987	11.3342433633575\\
9.27999999999986	11.7475921707616\\
9.36799999999985	12.2196496145268\\
9.44799999999985	12.6651750130385\\
9.52799999999984	13.1268828866729\\
9.61599999999983	13.6541475410798\\
9.69599999999982	14.1517576672186\\
9.78399999999981	14.7200066935816\\
9.8639999999998	15.2562812235038\\
9.95199999999979	15.8686659541854\\
10.0319999999998	16.4465761043058\\
10.1199999999998	17.1064861995222\\
10.1999999999998	17.7292277065898\\
10.2879999999998	18.4403087508468\\
10.3679999999997	19.1113184277698\\
10.4479999999997	19.8065952725136\\
10.5359999999997	20.6004642421724\\
10.6159999999997	21.3495611141851\\
10.7039999999997	22.2048538618727\\
10.7839999999997	23.0118835005741\\
10.8719999999997	23.9332906510792\\
10.9519999999997	24.8026740627602\\
11.0399999999997	25.7952374725451\\
11.1199999999997	26.7317263705127\\
11.2079999999997	27.8008646540396\\
11.2879999999996	28.8095654534244\\
11.3679999999996	29.8545569656299\\
11.4559999999996	31.0475007376103\\
11.5359999999996	32.1729452229135\\
11.6239999999996	33.4576826105528\\
11.7039999999996	34.6696789182956\\
11.7919999999996	36.0531628368136\\
11.8719999999996	37.358261453675\\
11.9599999999996	38.8479590575211\\
12.0399999999996	40.2531943270266\\
12.1199999999996	41.7087350802563\\
12.2079999999995	43.3700481074514\\
12.2879999999995	44.9370672258285\\
12.3759999999995	46.7255381140911\\
12.4559999999995	48.4124208241357\\
12.5439999999995	50.3376050662995\\
12.6239999999995	52.153349976566\\
12.7119999999995	54.22550098167\\
12.7919999999995	56.1797629134904\\
12.8799999999995	58.409880607148\\
12.9599999999995	60.5130161627528\\
13.0399999999995	62.6908997217316\\
13.1279999999994	65.1760146864655\\
13.2079999999994	67.5194461662676\\
13.2959999999994	70.1933182891266\\
13.3759999999994	72.7146069888535\\
13.4639999999994	75.5912546052707\\
13.5439999999994	78.3035933473953\\
13.6319999999994	81.3980389785611\\
13.7119999999994	84.3155637073944\\
13.7919999999994	87.3360538916131\\
13.8799999999994	90.7817544677395\\
13.9599999999993	94.0301563374406\\
14.0479999999993	97.7356182540901\\
14.1279999999993	101.228682699574\\
14.2159999999993	105.212970188832\\
14.2959999999993	108.968626899406\\
14.3839999999993	113.252142439851\\
14.4639999999993	117.289576493515\\
14.5519999999993	121.894147340322\\
14.6319999999993	126.233879973747\\
14.7119999999993	130.725185951388\\
14.7999999999993	135.846823043272\\
14.8799999999992	140.673344189884\\
14.9679999999992	146.176814118928\\
15.0479999999992	151.362755806784\\
15.1359999999992	157.275579285267\\
15.2159999999992	162.846798041752\\
15.3039999999992	169.198366895376\\
15.3839999999992	175.182471270219\\
15.4719999999992	182.004161977523\\
15.5519999999992	188.430618123296\\
15.6319999999992	195.078641145107\\
15.7199999999992	202.656152568749\\
15.7999999999991	209.793654950021\\
15.8879999999991	217.928298994414\\
15.9679999999991	225.58983160784\\
16.0559999999991	234.320841275233\\
16.1359999999991	242.543216594743\\
16.2239999999991	251.912381736901\\
16.3039999999991	260.734800587749\\
16.3839999999991	269.85745838112\\
16.4719999999991	280.250792714061\\
16.5519999999991	290.036006212154\\
16.6399999999991	301.182903132889\\
16.719999999999	311.676361554343\\
16.807999999999	323.628634557474\\
16.887999999999	334.878890412734\\
16.975999999999	347.691582216295\\
17.055999999999	359.750196937428\\
17.143999999999	373.481741633763\\
17.223999999999	386.403434673111\\
17.303999999999	399.75630429805\\
17.391999999999	414.958624475379\\
17.471999999999	429.261443467586\\
17.559999999999	445.542978137445\\
17.6399999999989	460.858946046622\\
17.7279999999989	478.291226617529\\
17.8079999999989	494.6872399298\\
17.8959999999989	513.345951594153\\
17.9759999999989	530.892762372718\\
18.0639999999989	550.85791890401\\
18.1439999999989	569.63029376586\\
18.2239999999989	589.013753716761\\
18.3119999999989	611.063290822764\\
18.3919999999989	631.790357151732\\
18.4799999999989	655.364217271166\\
18.5599999999988	677.520276126239\\
18.6479999999988	702.714869933341\\
18.7279999999988	726.389864484235\\
18.8159999999988	753.306697871869\\
18.8959999999988	778.595277118178\\
18.9759999999988	804.686898180512\\
19.0639999999988	834.342774170131\\
19.1439999999988	862.196602423997\\
19.2319999999988	893.849011597788\\
19.3119999999988	923.571950566369\\
19.3999999999988	957.341315002275\\
19.4799999999987	989.045439543296\\
19.5679999999987	1025.05793337862\\
19.6479999999987	1058.86060243213\\
19.7359999999987	1097.24826920973\\
19.8159999999987	1133.27222595233\\
19.8959999999987	1170.39866377657\\
19.9839999999987	1212.54653857949\\
20.0639999999987	1252.08547426275\\
20.1519999999987	1296.96136508499\\
20.2319999999987	1339.0492073044\\
20.3199999999987	1386.80624306283\\
20.3999999999986	1431.58500678263\\
20.4879999999986	1482.38249675317\\
20.5679999999986	1529.99982194529\\
20.6559999999986	1584.00326321939\\
20.7359999999986	1634.61240908427\\
20.8159999999986	1686.70172072841\\
20.9039999999986	1745.7532778068\\
20.9839999999986	1801.07084298198\\
21.0719999999986	1863.76457633388\\
21.1519999999986	1922.47736347849\\
21.2399999999985	1988.99991001223\\
21.3199999999985	2051.28025865526\\
21.4079999999985	2121.82408698602\\
21.4879999999985	2187.84955922264\\
21.5679999999985	2255.72230078778\\
21.6559999999985	2332.56597676339\\
21.7359999999985	2404.45509256605\\
21.8239999999985	2485.82044154641\\
21.9039999999985	2561.91552246627\\
21.9919999999985	2648.01354478012\\
22.0719999999985	2728.50855664564\\
22.1599999999984	2819.55487838109\\
22.2399999999984	2904.64780951612\\
22.3279999999984	3000.86226522773\\
22.4079999999984	3090.75471935957\\
22.4879999999984	3183.00445227271\\
22.5759999999984	3287.25760397398\\
22.6559999999984	3384.61007277223\\
22.7439999999984	3494.59058932196\\
22.8239999999984	3597.2542176917\\
22.9119999999984	3713.1924042996\\
22.9919999999984	3821.37734367479\\
23.0799999999983	3943.50503500289\\
23.1599999999983	4057.42247252942\\
23.2479999999983	4185.97223225781\\
23.3279999999983	4305.83362684303\\
23.4079999999983	4428.59966895139\\
23.4959999999983	4567.05411069683\\
23.5759999999983	4696.07496554323\\
23.6639999999983	4841.52463090505\\
23.7439999999983	4977.00865993169\\
23.8319999999983	5129.68125520841\\
23.9119999999983	5271.83384520376\\
23.9999999999982	5431.95320507326\\
24.0799999999982	5580.97565867119\\
24.1599999999982	5733.33874610732\\
24.2479999999982	5904.84882936568\\
24.3279999999982	6064.36894026343\\
24.4159999999982	6243.8550611007\\
24.4959999999982	6410.71810512614\\
24.5839999999982	6598.38025486152\\
24.6639999999982	6772.76359511876\\
24.7519999999982	6968.79157658831\\
24.8319999999982	7150.86258137793\\
24.9199999999981	7355.43437895018\\
24.9999999999981	7545.3489078087\\
};
\addlegendentry{$\lambda_0=2.36$}

\addplot [color=mycolor4, line width=1.8pt]
  table[row sep=crcr]{%
0	1\\
0.08	0.743470359435956\\
0.168	0.576117066798391\\
0.248	0.479642934855529\\
0.336	0.408811762327586\\
0.416	0.364561344340258\\
0.504	0.330338524359988\\
0.584	0.308289419628975\\
0.672	0.291081445996212\\
0.752000000000001	0.280158121223481\\
0.840000000000001	0.272038866716628\\
0.920000000000001	0.267407177271807\\
1	0.26485295974835\\
1.088	0.263989893371575\\
1.168	0.264669293847231\\
1.256	0.266767822810731\\
1.336	0.269731638154394\\
1.424	0.274004455791523\\
1.504	0.27870886531608\\
1.592	0.284697975259161\\
1.672	0.29082346419647\\
1.752	0.297553805893507\\
1.84	0.305617622430504\\
1.92	0.313522703969406\\
2.008	0.32282861008677\\
2.088	0.331830332571826\\
2.176	0.342318934754722\\
2.256	0.352383296698101\\
2.344	0.364035550954585\\
2.424	0.375159422835842\\
2.512	0.387985247574226\\
2.592	0.400188232505783\\
2.672	0.412917225311241\\
2.76	0.427538766866174\\
2.84	0.441407057066274\\
2.928	0.457311916896953\\
3.008	0.472377265762158\\
3.096	0.489635607373427\\
3.176	0.50596752405943\\
3.264	0.524661846079814\\
3.344	0.542340652028643\\
3.432	0.562565030844738\\
3.512	0.581681379023047\\
3.592	0.601517084481202\\
3.68	0.624195851986129\\
3.76	0.645621536133612\\
3.848	0.670111691580458\\
3.928	0.693243416859662\\
4.016	0.71967846655509\\
4.096	0.74464301265043\\
4.184	0.773168501316628\\
4.264	0.800103849841504\\
4.344	0.828029571620335\\
4.432	0.859933898833739\\
4.512	0.890055888274284\\
4.6	0.924467007479446\\
4.68	0.956953804663458\\
4.768	0.994064546860266\\
4.848	1.02909843577446\\
4.936	1.06911726037746\\
5.016	1.10689521981277\\
5.104	1.15004731307429\\
5.184	1.19078206314599\\
5.264	1.23300442678649\\
5.352	1.28123168890008\\
5.432	1.32675607677033\\
5.52	1.37875419561348\\
5.6	1.42783744549213\\
5.688	1.4838998349743\\
5.768	1.53681894783516\\
5.856	1.59726201991524\\
5.936	1.65431567709954\\
6.024	1.71948057634835\\
6.104	1.78099077796708\\
6.184	1.84474352517692\\
6.272	1.91755908608255\\
6.352	1.98629010768177\\
6.44	2.06479105199856\\
6.52	2.13888799080671\\
6.608	2.22351696897157\\
6.68800000000001	2.30339754375227\\
6.77600000000001	2.39463152585653\\
6.85600000000001	2.48074582024385\\
6.93600000000001	2.56999675062489\\
7.02400000000001	2.67193163400676\\
7.10400000000001	2.76814510889968\\
7.19200000000001	2.87803101986843\\
7.27200000000001	2.98174826577602\\
7.36000000000001	3.1002031083347\\
7.44000000000001	3.21200711333618\\
7.52800000000001	3.33969644359591\\
7.60800000000001	3.46021513102405\\
7.69600000000001	3.59785577714718\\
7.77600000000001	3.72776541773987\\
7.85600000000001	3.862399575797\\
7.94400000000001	4.01615825141499\\
8.024	4.16127786641281\\
8.11199999999999	4.32700921436304\\
8.19199999999998	4.48342664675097\\
8.27999999999997	4.66205792544116\\
8.35999999999997	4.83064785685359\\
8.44799999999996	5.02317742296469\\
8.52799999999995	5.20488156652133\\
8.61599999999994	5.41238424541112\\
8.69599999999993	5.60821640577552\\
8.77599999999992	5.81115479306449\\
8.86399999999991	6.04290076860931\\
8.9439999999999	6.26160711127599\\
9.03199999999989	6.5113548455803\\
9.11199999999988	6.74704571033657\\
9.19999999999987	7.01618354461257\\
9.27999999999986	7.27016827369722\\
9.36799999999985	7.5601903634255\\
9.44799999999985	7.83387803446804\\
9.52799999999984	8.11747411585883\\
9.61599999999983	8.44129917769126\\
9.69599999999982	8.74687650817078\\
9.78399999999981	9.09579337029708\\
9.8639999999998	9.42504130990313\\
9.95199999999979	9.80097749820868\\
10.0319999999998	10.1557135453726\\
10.1199999999998	10.5607426952924\\
10.1999999999998	10.9429220282489\\
10.2879999999998	11.3792747543705\\
10.3679999999997	11.7910004640787\\
10.4479999999997	12.21758071469\\
10.5359999999997	12.7046100787846\\
10.6159999999997	13.1641350384626\\
10.7039999999997	13.6887638783407\\
10.7839999999997	14.1837518085127\\
10.8719999999997	14.7488526514364\\
10.9519999999997	15.2820113847139\\
11.0399999999997	15.8906728338604\\
11.1199999999997	16.4649137334066\\
11.2079999999997	17.1204563605971\\
11.2879999999996	17.7389092724918\\
11.3679999999996	18.3795829734243\\
11.4559999999996	19.1109319041717\\
11.5359999999996	19.8008713718888\\
11.6239999999996	20.5884343806924\\
11.7039999999996	21.3313817896044\\
11.7919999999996	22.1794261568199\\
11.8719999999996	22.9794028440583\\
11.9599999999996	23.8925138914551\\
12.0399999999996	24.7538406238273\\
12.1199999999996	25.6459925916612\\
12.2079999999995	26.6642634167039\\
12.2879999999995	27.6247372343141\\
12.3759999999995	28.7209495282391\\
12.4559999999995	29.7549038748687\\
12.5439999999995	30.9349383853387\\
12.6239999999995	32.0479133448295\\
12.7119999999995	33.3180851292841\\
12.7919999999995	34.5160294749689\\
12.8799999999995	35.8831189124136\\
12.9599999999995	37.1724192338377\\
13.0399999999995	38.5076019150473\\
13.1279999999994	40.03121690435\\
13.2079999999994	41.4680499007916\\
13.2959999999994	43.1075925372543\\
13.3759999999994	44.6536847209131\\
13.4639999999994	46.4178249300737\\
13.5439999999994	48.0813401754795\\
13.6319999999994	49.9793793896437\\
13.7119999999994	51.769075362175\\
13.7919999999994	53.6221180893584\\
13.8799999999994	55.7362599914587\\
13.9599999999993	57.7295847820809\\
14.0479999999993	60.0036651629875\\
14.1279999999993	62.1476833450337\\
14.2159999999993	64.5935598909517\\
14.2959999999993	66.8994317530473\\
14.3839999999993	69.5298136727664\\
14.4639999999993	72.0094995173743\\
14.5519999999993	74.8380054012298\\
14.6319999999993	77.5043197253441\\
14.7119999999993	80.2642709177594\\
14.7999999999993	83.4122078739708\\
14.8799999999992	86.3793878985061\\
14.9679999999992	89.7634878053493\\
15.0479999999992	92.953080862236\\
15.1359999999992	96.5906250515255\\
15.2159999999992	100.018884287189\\
15.3039999999992	103.928367894044\\
15.3839999999992	107.612685371142\\
15.4719999999992	111.813896438205\\
15.5519999999992	115.772877025389\\
15.6319999999992	119.869560078944\\
15.7199999999992	124.540519310351\\
15.7999999999991	128.941712526306\\
15.8879999999991	133.959509586061\\
15.9679999999991	138.687165322899\\
16.0559999999991	144.076761633285\\
16.1359999999991	149.154334258738\\
16.2239999999991	154.942393978955\\
16.3039999999991	160.394932844985\\
16.3839999999991	166.035318436063\\
16.4719999999991	172.464174595208\\
16.5519999999991	178.519626824057\\
16.6399999999991	185.420984961522\\
16.719999999999	191.920932161725\\
16.807999999999	199.328223040718\\
16.887999999999	206.304047486509\\
16.975999999999	214.252913117244\\
17.055999999999	221.738064745098\\
17.143999999999	230.266486054508\\
17.223999999999	238.296600860448\\
17.303999999999	246.599412148355\\
17.391999999999	256.058053546416\\
17.471999999999	264.962690597738\\
17.559999999999	275.105866710262\\
17.6399999999989	284.653917683064\\
17.7279999999989	295.528805199244\\
17.8079999999989	305.764491227266\\
17.8959999999989	317.421238554331\\
17.9759999999989	328.391551910194\\
18.0639999999989	340.883433331412\\
18.1439999999989	352.638279482831\\
18.2239999999989	364.785134915351\\
18.3119999999989	378.6141966917\\
18.3919999999989	391.624888502212\\
18.4799999999989	406.435480177239\\
18.5599999999988	420.367761002939\\
18.6479999999988	436.225277472635\\
18.7279999999988	451.140332356766\\
18.8159999999988	468.11402908221\\
18.8959999999988	484.076630294433\\
18.9759999999988	500.561908890158\\
19.0639999999988	519.318457713057\\
19.1439999999988	536.953786114335\\
19.2319999999988	557.015721903897\\
19.3119999999988	575.8754286365\\
19.3999999999988	597.326764022613\\
19.4799999999987	617.489309159909\\
19.5679999999987	640.418687000969\\
19.6479999999987	661.966817192021\\
19.7359999999987	686.467679899568\\
19.8159999999987	709.488580939494\\
19.8959999999987	733.242770465368\\
19.9839999999987	760.244782309724\\
20.0639999999987	785.608873132757\\
20.1519999999987	814.435462490905\\
20.2319999999987	841.508263258122\\
20.3199999999987	872.270815357369\\
20.3999999999986	901.156068174385\\
20.4879999999986	933.971465715675\\
20.5679999999986	964.777972391395\\
20.6559999999986	999.768729170313\\
20.7359999999986	1032.61046750711\\
20.8159999999986	1066.46295230327\\
20.9039999999986	1104.90107854125\\
20.9839999999986	1140.96672608162\\
21.0719999999986	1181.90856286604\\
21.1519999999986	1220.31458871326\\
21.2399999999985	1263.90307765876\\
21.3199999999985	1304.78217449312\\
21.4079999999985	1351.16631154944\\
21.4879999999985	1394.65669795485\\
21.5679999999985	1439.44074653325\\
21.6559999999985	1490.23702422021\\
21.7359999999985	1537.84655204133\\
21.8239999999985	1591.83369174166\\
21.9039999999985	1642.42067055613\\
21.9919999999985	1699.76888223454\\
22.0719999999985	1753.49080094752\\
22.1599999999984	1814.37633234822\\
22.2399999999984	1871.39613283831\\
22.3279999999984	1936.0011813837\\
22.4079999999984	1996.48716179906\\
22.4879999999984	2058.68437806004\\
22.5759999999984	2129.12543176339\\
22.6559999999984	2195.04676811549\\
22.7439999999984	2269.68319841707\\
22.8239999999984	2339.50956534725\\
22.9119999999984	2418.54300711908\\
22.9919999999984	2492.4600401511\\
23.0799999999983	2576.09717874263\\
23.1599999999983	2654.29495928193\\
23.2479999999983	2742.74719668691\\
23.3279999999983	2825.41992340193\\
23.4079999999983	2910.29293238635\\
23.4959999999983	3006.24859567556\\
23.5759999999983	3095.88992836381\\
23.6639999999983	3197.2018145234\\
23.7439999999983	3291.81410393419\\
23.8319999999983	3398.70658370331\\
23.9119999999983	3498.49500648992\\
23.9999999999982	3611.19493302945\\
24.0799999999982	3716.36661961533\\
24.1599999999982	3824.17496915108\\
24.2479999999982	3945.86602116219\\
24.3279999999982	4059.36556108038\\
24.4159999999982	4187.43189670478\\
24.4959999999982	4306.83162104684\\
24.5839999999982	4441.50290308431\\
24.6639999999982	4567.01117237693\\
24.7519999999982	4708.5158774894\\
24.8319999999982	4840.33951054149\\
24.9199999999981	4988.90388843906\\
24.9999999999981	5127.24717789354\\
};
\addlegendentry{$\lambda_0=1.77$}

\addplot [color=black, only marks, mark size=1.5pt, mark=*, mark options={solid, fill=black, black}, forget plot]
  table[row sep=crcr]{%
0	1\\
0.5	0.608627008982342\\
1	0.626477585411236\\
1.5	0.783892678569528\\
2	1.04447905829334\\
2.5	1.40967378375594\\
3	1.89547621600287\\
3.5	2.52893543798223\\
4	3.34761866232078\\
4.5	4.40093459633577\\
5	5.75256398132912\\
5.5	7.48385895715207\\
6	9.69828627863822\\
6.5	12.5271008456559\\
7	16.1365216551613\\
7.5	20.7367679689801\\
8	26.5934109107459\\
8.5	34.0416124982757\\
9	43.5039671132983\\
9.5	55.5128367866284\\
10	70.7382899349903\\
};
\addplot [color=black, only marks, mark size=1.5pt, mark=*, mark options={solid, fill=black, black}, forget plot]
  table[row sep=crcr]{%
0	1\\
0.5	0.52644321399243\\
1	0.509613682034376\\
1.5	0.601918763103908\\
2	0.754984979208155\\
2.5	0.962223316549923\\
3	1.22883951951508\\
3.5	1.56686215480419\\
4	1.99344736842432\\
4.5	2.53104944322894\\
5	3.20830911184869\\
5.5	4.06142148374349\\
6	5.13595469903285\\
6.5	6.48917185295908\\
7	8.19295466429414\\
7.5	10.3374648684479\\
8	13.0357182028888\\
8.5	16.4292908213575\\
9	20.6954322321168\\
9.5	26.0559253058712\\
10	32.7881157550453\\
};
\addplot [color=black, only marks, mark size=1.5pt, mark=*, mark options={solid, fill=black, black}, forget plot]
  table[row sep=crcr]{%
0	1\\
0.5	0.414851284659067\\
1	0.361279998455926\\
1.5	0.398279029621783\\
2	0.474942714811734\\
2.5	0.582583395242735\\
3	0.722206752564905\\
3.5	0.899379168771463\\
4	1.1225585622544\\
4.5	1.40297127228764\\
5	1.75498141353041\\
5.5	2.19673883484128\\
6	2.75106347550766\\
6.5	3.44658224342495\\
7	4.31916272700269\\
7.5	5.4137083077048\\
8	6.78639885822517\\
8.5	8.50748327489697\\
9	10.6647564311049\\
9.5	13.3678852525256\\
10	16.753788089588\\
};
\addplot [color=black, only marks, mark size=1.5pt, mark=*, mark options={solid, fill=black, black}]
  table[row sep=crcr]{%
0	1\\
0.5	0.317323690958043\\
1	0.253745778453521\\
1.5	0.267345078687253\\
2	0.309695513539663\\
2.5	0.372305587092947\\
3	0.454767860500341\\
3.5	0.559998478563258\\
4	0.692815678437204\\
4.5	0.859743714622315\\
5	1.06918478620962\\
5.5	1.33177913080025\\
6	1.66091260612204\\
6.5	2.07337550175239\\
7	2.5901960819434\\
7.5	3.23768536270842\\
8	4.04874151256093\\
8.5	5.06447528002537\\
9	6.3362332166739\\
9.5	7.92811411575158\\
10	9.92009696718767\\
};
\addlegendentry{SSTE}

\node[right, align=left, inner sep=0]
at (rel axis cs:-0.15,1.15) {$(a)$};
\end{axis}
\end{tikzpicture}%
  \end{minipage}
  \hfill
  \begin{minipage}{0.48\textwidth}
    \centering
    % This file was created by matlab2tikz.
%
\begin{tikzpicture}

\begin{axis}[%
width=11.252in,
height=7.131in,
at={(1.887in,0.962in)},
scale only axis,
clip=true,
separate axis lines,
every outer x axis line/.append style={black},
every x tick label/.append style={font=\color{black}},
every x tick/.append style={black},
xmin=0.005,
xmax=6,
xlabel={$t$},
every outer y axis line/.append style={black},
every y tick label/.append style={font=\color{black}},
every y tick/.append style={black},
ymode=log,
ymin=0.1,
ymax=10,
yminorticks=true,
ylabel={$\mathscr{G}^{\mathrm{mean}}(t)$},
axis background/.style={fill=white},
clip=true,
clip mode=individual,
axis lines=box,
xtick pos=bottom,
ytick pos=left,
tick align=inside,
major tick length=2pt,
width=0.8\linewidth,
height=0.8\linewidth,
every axis/.append style={font=\fontsize{9}{9}\selectfont},xlabel style={font=\fontsize{9}{9}\selectfont},title style={font=\fontsize{9}{9}\selectfont},
legend style={font=\fontsize{8}{9}\selectfont,inner sep=1pt,row sep=1pt,column sep=2pt,nodes={scale=1.00},draw=none},legend image post style={xscale=0.35},
legend columns=1,
scaled x ticks=false,
tick label style={/pgf/number format/fixed,/pgf/number format/precision=2},
every x tick label/.append style={font=\fontsize{9}{9}\selectfont\color{black}},
every y tick label/.append style={font=\fontsize{9}{9}\selectfont\color{black}},
,,
colormap={sigmamap}{rgb(0.0000)=(0.0200,0.1880,0.3800); rgb(0.1000)=(0.0743,0.3456,0.5616); rgb(0.2000)=(0.1918,0.4980,0.7060); rgb(0.3000)=(0.3890,0.6708,0.8226); rgb(0.4000)=(0.7552,0.8682,0.9228); rgb(0.5000)=(0.9690,0.9689,0.9689); rgb(0.6000)=(0.9425,0.8044,0.7614); rgb(0.7000)=(0.8767,0.4910,0.4020); rgb(0.8000)=(0.7869,0.2866,0.2470); rgb(0.9000)=(0.6186,0.1239,0.1568); rgb(1.0000)=(0.4040,0.0000,0.1220)}
]
\addplot [color=white!55!black, dashed, line width=0.7pt, forget plot]
  table[row sep=crcr]{%
0.005	1\\
6	1\\
};
\addplot [color=black, line width=1.1pt, forget plot]
  table[row sep=crcr]{%
0.005	0.988188115174213\\
0.025	0.943747136082144\\
0.045	0.903410917403379\\
0.065	0.866723912918671\\
0.085	0.833292776372284\\
0.105	0.802776392246593\\
0.125	0.774877723002092\\
0.145	0.749337106046553\\
0.165	0.725926714105225\\
0.185	0.704445954226831\\
0.205	0.684717628040936\\
0.225	0.666584712537862\\
0.245	0.649907649146324\\
0.265	0.634562051162651\\
0.285	0.620436757085638\\
0.305	0.607432171225992\\
0.325	0.595458843918176\\
0.345	0.584436252397151\\
0.365	0.574291750396514\\
0.385	0.56495966015109\\
0.405	0.556380485033046\\
0.425	0.548500224739174\\
0.445	0.541269777953409\\
0.465	0.534644419868454\\
0.485	0.528583343970329\\
0.505	0.523049259156323\\
0.525	0.51800803463603\\
0.545	0.513428386209938\\
0.565	0.509281598475837\\
0.585	0.505541278311404\\
0.605	0.502183135651566\\
0.625	0.499184788143715\\
0.645	0.496525586739496\\
0.665	0.494186459685812\\
0.685	0.492149772719929\\
0.705	0.490399203565327\\
0.725	0.488919629073765\\
0.745	0.487697023572944\\
0.765	0.486718367160997\\
0.785	0.485971562847544\\
0.805	0.485445361576106\\
0.825	0.48512929428027\\
0.845	0.48501361022756\\
0.865	0.485089220992893\\
0.885	0.485347649481108\\
0.905	0.485780983483483\\
0.925	0.4863818333129\\
0.945	0.487143293112784\\
0.965	0.488058905479723\\
0.985	0.489122629078949\\
1.01	0.490651984609727\\
1.03	0.492028827351546\\
1.05	0.493536678062384\\
1.07	0.495170960724324\\
1.09	0.496927384323349\\
1.11	0.498801923759885\\
1.13	0.50079080224481\\
1.15	0.502890475050405\\
1.17	0.505097614498565\\
1.19	0.50740909608092\\
1.21	0.509821985612325\\
1.23	0.512333527334065\\
1.25	0.514941132885424\\
1.27	0.517642371073119\\
1.29	0.520434958373467\\
1.31	0.523316750108201\\
1.33	0.526285732240153\\
1.35	0.52934001373993\\
1.37	0.532477819478814\\
1.39	0.535697483607159\\
1.41	0.538997443380843\\
1.43	0.542376233403068\\
1.45	0.545832480246841\\
1.47	0.549364897434242\\
1.49	0.552972280742855\\
1.51	0.556653503816618\\
1.53	0.560407514058935\\
1.55	0.564233328787608\\
1.57	0.568130031633107\\
1.59	0.572096769162815\\
1.61	0.576132747715604\\
1.63	0.580237230432013\\
1.65	0.584409534466682\\
1.67	0.588649028370233\\
1.69	0.592955129631021\\
1.71	0.597327302361346\\
1.73	0.60176505512481\\
1.75	0.606267938890815\\
1.77	0.610835545109643\\
1.79	0.615467503900194\\
1.81	0.62016348234295\\
1.83	0.624923182871561\\
1.85	0.629746341756682\\
1.87	0.634632727676449\\
1.89	0.639582140368017\\
1.91	0.644594409355255\\
1.93	0.649669392747418\\
1.95	0.654806976107052\\
1.97	0.660007071376782\\
1.99	0.665269615869485\\
2.01	0.670594571311726\\
2.03	0.675981922941234\\
2.05	0.681431678652972\\
2.07	0.686943868191119\\
2.09	0.692518542385566\\
2.11	0.698155772429756\\
2.13	0.703855649197193\\
2.15	0.709618282595011\\
2.17	0.715443800951903\\
2.19	0.721332350439061\\
2.21	0.72728409452195\\
2.23	0.733299213441389\\
2.25	0.739377903722243\\
2.27	0.745520377708351\\
2.29	0.751726863122208\\
2.31	0.757997602648094\\
2.33	0.764332853537408\\
2.35	0.770732887235066\\
2.37	0.777197989025778\\
2.39	0.783728457699514\\
2.41	0.790324605234502\\
2.43	0.796986756497661\\
2.45	0.803715248960964\\
2.47	0.810510432433433\\
2.49	0.817372668807264\\
2.51	0.824302331818569\\
2.53	0.83129980682052\\
2.55	0.838365490569661\\
2.57	0.845499791023948\\
2.59	0.852703127152392\\
2.61	0.859975928754993\\
2.63	0.867318636293675\\
2.65	0.874731700732742\\
2.67	0.882215583388759\\
2.69	0.889770755789664\\
2.71	0.897397699541959\\
2.73	0.905096906206571\\
2.75	0.912868877182228\\
2.77	0.920714123596457\\
2.79	0.928633166203721\\
2.81	0.936626535290328\\
2.83	0.944694770585904\\
2.85	0.952838421181113\\
2.87	0.961058045451491\\
2.89	0.9693542109869\\
2.91	0.977727494526547\\
2.93	0.986178481899354\\
2.95	0.994707767969335\\
2.97	1.00331595658579\\
2.99	1.0120036605383\\
3.015	1.02297605686972\\
3.035	1.03184495688023\\
3.055	1.04079542507624\\
3.075	1.0498281118426\\
3.095	1.05894367633804\\
3.115	1.06814278647532\\
3.135	1.07742611890579\\
3.155	1.08679435900628\\
3.175	1.09624820087027\\
3.195	1.10578834730154\\
3.215	1.11541550981121\\
3.235	1.12513040861807\\
3.255	1.13493377265103\\
3.275	1.14482633955513\\
3.295	1.15480885569977\\
3.315	1.16488207618989\\
3.335	1.17504676487938\\
3.355	1.18530369438714\\
3.375	1.19565364611571\\
3.395	1.20609741027176\\
3.415	1.21663578588917\\
3.435	1.22726958085461\\
3.455	1.23799961193406\\
3.475	1.24882670480326\\
3.495	1.25975169407837\\
3.515	1.27077542334998\\
3.535	1.28189874521808\\
3.555	1.29312252132971\\
3.575	1.30444762241757\\
3.595	1.31587492834163\\
3.615	1.32740532813124\\
3.635	1.33903972002956\\
3.655	1.35077901153998\\
3.675	1.36262411947352\\
3.695	1.37457596999847\\
3.715	1.38663549869117\\
3.735	1.39880365058855\\
3.755	1.41108138024228\\
3.775	1.423469651774\\
3.795	1.43596943893308\\
3.815	1.44858172515453\\
3.835	1.46130750361888\\
3.855	1.47414777731397\\
3.875	1.48710355909772\\
3.895	1.5001758717618\\
3.915	1.51336574809767\\
3.935	1.52667423096325\\
3.955	1.54010237335096\\
3.975	1.55365123845779\\
3.995	1.56732189975526\\
4.015	1.5811154410623\\
4.035	1.5950329566188\\
4.055	1.60907555115854\\
4.075	1.62324433998806\\
4.095	1.63754044906067\\
4.115	1.65196501505789\\
4.135	1.66651918546619\\
4.155	1.6812041186608\\
4.175	1.69602098398466\\
4.195	1.71097096183474\\
4.215	1.72605524374396\\
4.235	1.74127503246955\\
4.255	1.75663154207699\\
4.275	1.77212599803074\\
4.295	1.7877596372814\\
4.315	1.80353370835858\\
4.335	1.81944947145929\\
4.355	1.83550819854458\\
4.375	1.85171117342975\\
4.395	1.86805969188294\\
4.415	1.88455506171868\\
4.435	1.90119860289783\\
4.455	1.91799164762363\\
4.475	1.93493554044472\\
4.495	1.95203163835334\\
4.515	1.96928131088947\\
4.535	1.98668594024366\\
4.555	2.00424692136052\\
4.575	2.02196566204689\\
4.595	2.03984358307595\\
4.615	2.05788211829815\\
4.635	2.07608271474818\\
4.655	2.09444683275741\\
4.675	2.11297594606381\\
4.695	2.13167154192694\\
4.715	2.15053512123966\\
4.735	2.169568198646\\
4.755	2.18877230265516\\
4.775	2.20814897576187\\
4.795	2.22769977456212\\
4.815	2.24742626987779\\
4.835	2.26733004687314\\
4.855	2.28741270518258\\
4.875	2.30767585902896\\
4.895	2.32812113735434\\
4.915	2.34875018394191\\
4.935	2.36956465754693\\
4.955	2.3905662320228\\
4.975	2.41175659645475\\
4.995	2.43313745528734\\
5.02	2.46013403328734\\
5.04	2.48194981687441\\
5.06	2.50396174201487\\
5.08	2.52617157993667\\
5.1	2.54858111794298\\
5.12	2.57119215954785\\
5.14	2.59400652462026\\
5.16	2.61702604952616\\
5.18	2.64025258727022\\
5.2	2.66368800764339\\
5.22	2.68733419736737\\
5.24	2.71119306024464\\
5.26	2.7352665173049\\
5.28	2.75955650695933\\
5.3	2.78406498514828\\
5.32	2.80879392550015\\
5.34	2.83374531948029\\
5.36	2.85892117655346\\
5.38	2.88432352433623\\
5.4	2.90995440876055\\
5.42	2.93581589423112\\
5.44	2.96191006379126\\
5.46	2.98823901928215\\
5.48	3.01480488151381\\
5.5	3.0416097904257\\
5.52	3.06865590526035\\
5.54	3.09594540472926\\
5.56	3.1234804871879\\
5.58	3.15126337080487\\
5.6	3.17929629374047\\
5.62	3.20758151431888\\
5.64	3.23612131120919\\
5.66	3.26491798360387\\
5.68	3.29397385139774\\
5.7	3.32329125537434\\
5.72	3.35287255738589\\
5.74	3.38272014054458\\
5.76	3.41283640940398\\
5.78	3.44322379015507\\
5.8	3.47388473081051\\
5.82	3.5048217014039\\
5.84	3.53603719417864\\
5.86	3.56753372378954\\
5.88	3.59931382749528\\
5.9	3.63138006536437\\
5.92	3.66373502047025\\
5.94	3.69638129910088\\
5.96	3.72932153095789\\
5.98	3.76255836937086\\
6	3.79609449149854\\
};
\addplot [color=black, line width=1.4pt, only marks, mark size=1.8pt, mark=*, mark options={solid, fill=black, black}, forget plot]
  table[row sep=crcr]{%
0.845	0.48501361022756\\
};
\addplot [color=black, line width=1.4pt, only marks, mark size=1.8pt, mark=*, mark options={solid, fill=black, black}, forget plot]
  table[row sep=crcr]{%
2.96231615336066	1\\
};
\node[right, align=left, inner sep=0]
at (axis cs:1.045,0.388) {$\mathscr{G}^{\mathrm{mean}}_{\min}$};
\node[right, align=left, inner sep=0]
at (axis cs:3.122,0.76) {$t_c$};
\node[right, align=left, inner sep=0]
at (rel axis cs:-0.15,1.15) {$(b)$};
\end{axis}
\end{tikzpicture}%
  \end{minipage}
  \caption{Mean energy growth in the Boussinesq limit for different Gaussian reference lengths
  \(\lambda_0\). (a)~Comparison between the DNS ensemble mean (\solidlegend) and the SSTE
  prediction (\circlegendfilled), using ten independent realizations for each ensemble.
  (b)~Representative evolution at \(\delta_0=2\) showing the
  minimum \(\Gm_{\min}\) and the critical time \(t_c\) at which the mean energy returns to its
  initial value.}
  \label{fig:compare_dns}
\end{figure}
The variance of the interface displacement is related to the modal amplitudes through
\begin{equation}
  \mean{\eta_0^{2}}
  =
  \frac{1}{4}
  \sum_{(p,q)\in\mathcal{I}}
  \mean{A_{pq}^{2}}.
  \label{eq:dns_variance_sum_rule}
\end{equation}
Given the relation between the correlation in space and the spectrum discussed in
\cref{sec:statistics_of_the_initial_perturbations}, we can build a strategy in which the
initialization of the DNS is equivalent to the computation of the mean energy growth. Denoting the
radial wavenumber of a mode by \(k_{p,q}=[(2\pi p/L_x)^2+(2\pi q/L_y)^2]^{1/2}\), we distribute
the retained modes over \(N_b\) radial bins, so that
\begin{subequations}
\label{eq:dns_radial_bins}
\begin{align}
  \Delta k
  &= \frac{k_{\max}-k_{\min}}{N_b},
  \label{eq:dns_radial_bin_width}
  \\
  k_\ell
  &= k_{\min}+\left(\ell-\frac{1}{2}\right)\Delta k,
  \qquad \ell=1,\ldots,N_b,
  \label{eq:dns_radial_bin_centers}
  \\
  \mathcal{I}_\ell
  &=
  \left\{
    (p,q)\in\mathcal{I}:
    \left|k_{p,q}-k_\ell\right|\leq\frac{\Delta k}{2}
  \right\}.
  \label{eq:dns_radial_bin_index_set}
\end{align}
\end{subequations}
Assuming that every Fourier mode in the bin \(\mathcal{I}_\ell\) carries the same amplitude
\(A_\ell\), a midpoint approximation of the energy integral gives
\begin{equation}
  \int_0^\infty E_{\eta}(k)\,\mathrm{d}k
  \approx
  \sum_{\ell=1}^{N_b} E_{\eta}(k_\ell)\Delta k
  \approx
  \sum_{\ell=1}^{N_b}\frac{A_\ell^2}{4}\,\left|\mathcal{I}_\ell\right| ,
  \label{eq:dns_midpoint_energy_integral}
\end{equation}
so that
\begin{equation}
  A_\ell
  =
  2\sqrt{\frac{E_{\eta}(k_\ell)\Delta k}{\left|\mathcal{I}_\ell\right|}} .
  \label{eq:dns_continuous_shell_amplitude}
\end{equation}
Letting \(\Delta k_x\) and \(\Delta k_y\) denote the Cartesian Fourier-grid spacings, the
number of grid points in a thin annulus is approximated by
\begin{equation}
  \left|\mathcal{I}_\ell\right|
  \approx
  \frac{2\pi k_\ell\Delta k}{\Delta k_x\Delta k_y},
  \label{eq:dns_shell_mode_count}
\end{equation}
and therefore
\begin{equation}
  A_\ell
  =
  \frac{2\sqrt{\Delta k_x\Delta k_y}}{\Delta k}
  \sqrt{\frac{E_{\eta}(k_\ell)}{2\pi k_\ell}} .
  \label{eq:dns_discrete_shell_amplitude}
\end{equation}
Choosing \(\Delta k=\sqrt{\Delta k_x\Delta k_y}\) gives
\begin{equation}
  \frac{A_\ell^{2}}{4}
  =
  \frac{E_{\eta}(k_\ell)}{2\pi k_\ell},
  \label{eq:dns_final_mode_amplitude}
\end{equation}
which defines the discrete mode amplitudes in terms of the prescribed spectrum. Thereby, for each
realization the amplitudes within each bin are assigned as discussed above, while the phases of
each mode are generated using different seeds across different realizations.

For the simulations, we use the \(n=1\) member of the family defined in
\eqref{eq:eta_spectrum_family}, which is
\begin{equation}
  E_{\eta,1}(k)
  =
  \frac{2\mean{\eta_0^2}}{k_0}
  \left(\frac{k}{k_0}\right)
  \exp\!\left[-2\left(\frac{k}{k_0}\right)^2\right],
  \qquad
  k_0=\frac{2\sqrt{2}}{\lambda_0}.
  \label{eq:dns_initial_spectrum}
\end{equation}
The amplitudes are normalized so that the ensemble variance satisfies
\(\mean{\eta_0^2}=10^{-2}\), which is sufficiently small for the interface displacement to remain
within the linear regime over the observation interval.


\Cref{fig:compare_dns}(a) shows good agreement between the mean energy growth predicted by the
stochastic theory and the DNS ensemble mean for all four values of \(\lambda_0\). The agreement
extends through both the initial decay and the subsequent growth, indicating that the statistical
formulation captures not only the late-time amplification but also the transient energy loss
preceding it.

\subsection{Characterisation of the transient instability}
\label{subsec:transient_characterisation}
Having validated the method and its numerical implementation, we proceed to characterise the
transient instability according to this stochastic approach. We define
\begin{equation}
  k_\ast(t)=\underset{k}{\operatorname{arg\,max}}\;\overline{\sigma}_{1D}(k,t),
  \qquad
  \overline{\sigma}_\ast(t)=\underset{k}{\max}\;\overline{\sigma}_{1D}(k,t).
  \label{eq:maximum_growth_definitions}
\end{equation}
Accordingly, \(k_\ast(t)\) is the wavenumber that maximises the mean growth-rate field, and
\(\overline{\sigma}_\ast(t)\) is the corresponding maximum value.

We do not restrict ourselves to the per-wavenumber description, but also consider the
three-dimensional mean growth, and define
\begin{equation}
  \Gm_{\min}=\min_{t\in[0,T]}\Gm(t),
  \qquad
  t_c=\min\left\{t\in(0,T]\mid\Gm(t)=1\right\}.
  \label{eq:minimum_growth_critical_time_definitions}
\end{equation}
Thus, \(\Gm_{\min}\) is the lowest mean energy growth reached during the observation interval,
while \(t_c\) is the first nonzero time at which the mean energy returns to its initial value. A
representative evolution is shown in \cref{fig:compare_dns}(b), together with \(\Gm_{\min}\) and
\(t_c\).

The first comparison is against the QSSA, reported in \cref{fig:BQ_P3_sigma_map}.
\Cref{fig:BQ_P3_sigma_map}(a) reports the non-autonomous mean growth-rate field
\(\overline{\sigma}_{1D}(k,t)\), whereas \cref{fig:BQ_P3_sigma_map}(b) reports the QSSA
growth rate obtained by solving the generalized eigenvalue problem for a base-state profile frozen
at the same time. During the early evolution, the non-autonomous mean growth rate is negative for
all wavenumbers and increases monotonically until the neutral contour is crossed around
\(t\sim0.5\). It then becomes positive, continues to increase to a maximum around \(t\sim2\),
and subsequently decays towards a slowly varying value. The comparison with the QSSA reveals two principal differences. First, the QSSA predicts an
unstable low-wavenumber band from the earliest time considered and therefore misses the initial
stable transient. Second, after crossing neutrality, the non-autonomous growth rate rises above the
QSSA prediction. Thus, during this stage, the instability is stronger than predicted by a classical
frozen-base-state eigenvalue analysis. At later times, as the diffusive layer broadens and the
relative variation of the base state slows, the non-autonomous growth rates tend towards the QSSA
prediction, as expected~\citep{tan-homsy-1986}.

\begin{figure}[t]
  \centering
  \begin{minipage}[t]{0.48\textwidth}
    \centering
    % This file was created by matlab2tikz.
%
\definecolor{mycolor1}{rgb}{0.20000,0.50000,0.30000}%
%
\begin{tikzpicture}

\begin{axis}[%
name=ax_main,width=11.252in,
height=7.131in,
at={(1.887in,0.962in)},
scale only axis,
point meta min=-0.5,
point meta max=0.5,
axis on top,
clip=true,
separate axis lines,
every outer x axis line/.append style={black},
every x tick label/.append style={font=\color{black}},
every x tick/.append style={black},
xmin=0.005,
xmax=15,
xlabel={$t$},
every outer y axis line/.append style={black},
every y tick label/.append style={font=\color{black}},
every y tick/.append style={black},
ymin=0,
ymax=2,
ytick={  0, 0.5,   1, 1.5,   2},
ylabel={$k$},
axis background/.style={fill=white},
clip=true,
clip mode=individual,
axis lines=box,
xtick pos=bottom,
ytick pos=left,
tick align=inside,
major tick length=2pt,
width=0.8\linewidth,
height=0.8\linewidth,
every axis/.append style={font=\fontsize{9}{9}\selectfont},xlabel style={font=\fontsize{9}{9}\selectfont},title style={font=\fontsize{9}{9}\selectfont},
legend style={font=\fontsize{8}{9}\selectfont,inner sep=1pt,row sep=1pt,column sep=2pt,nodes={scale=1.00},draw=none},legend image post style={xscale=0.35},
legend columns=1,
scaled x ticks=false,
tick label style={/pgf/number format/fixed,/pgf/number format/precision=2},
every x tick label/.append style={font=\fontsize{9}{9}\selectfont\color{black}},
every y tick label/.append style={font=\fontsize{9}{9}\selectfont\color{black}},
,,
colormap={sigmamap}{rgb(0.0000)=(0.7760,0.8510,0.9140); rgb(0.1000)=(0.8392,0.8918,0.9369); rgb(0.2000)=(0.8913,0.9259,0.9563); rgb(0.3000)=(0.9374,0.9569,0.9743); rgb(0.4000)=(0.9799,0.9861,0.9917); rgb(0.5000)=(1.0000,0.9999,0.9999); rgb(0.6000)=(0.9659,0.8254,0.7818); rgb(0.7000)=(0.8840,0.4951,0.4053); rgb(0.8000)=(0.7820,0.2849,0.2466); rgb(0.9000)=(0.6131,0.1208,0.1566); rgb(1.0000)=(0.4040,0.0000,0.1220)}
]
\addplot [forget plot] graphics [xmin=-0.00333704115684093, xmax=15.0033370411568, ymin=-0.00243892365706085, ymax=3.00250411920238] {BQ_P3_sigma_map-2.png};
\addplot [color=black, line width=0.8pt, forget plot]
  table[row sep=crcr]{%
14.995	0.563646335085671\\
14.92	0.563896840970456\\
14.845	0.564148632138408\\
14.77	0.56440172135392\\
14.695	0.564656121591333\\
14.62	0.564911846022406\\
14.545	0.565168908046938\\
14.47	0.565427321275749\\
14.39	0.565704466974204\\
14.315	0.565965716799437\\
14.245	0.566208733699811\\
14.17	0.566470140869557\\
14.095	0.566732955482606\\
14.02	0.566997192413921\\
13.945	0.567262866799346\\
13.87	0.567529994032781\\
13.795	0.567798589791099\\
13.72	0.5680686700258\\
13.645	0.568340250977025\\
13.57	0.568613349184761\\
13.495	0.568887981486633\\
13.42	0.569164165043781\\
13.34	0.569460490364186\\
13.265	0.569739935608849\\
13.19	0.570020986773734\\
13.115	0.570303662388231\\
13.04	0.570587981344453\\
12.965	0.570873962912375\\
12.895	0.57114112142818\\
12.82	0.571427700270563\\
12.745	0.571715982949784\\
12.67	0.572005990103955\\
12.595	0.572297742795926\\
12.52	0.572591262538562\\
12.445	0.572886571307742\\
12.37	0.573183691550926\\
12.29	0.573502642322558\\
12.215	0.573803579566892\\
12.14	0.574106400284032\\
12.065	0.574411129040206\\
11.99	0.574717790964893\\
11.915	0.575026411777074\\
11.84	0.575337017801407\\
11.765	0.575649635993926\\
11.69	0.575964293964492\\
11.62	0.576257615402508\\
11.545	0.576573428558759\\
11.47	0.576891345485095\\
11.395	0.577211396292183\\
11.32	0.577533611869785\\
11.24	0.577879730270883\\
11.165	0.578206521094079\\
11.09	0.578535576590663\\
11.015	0.578866931138757\\
10.94	0.579200620068978\\
10.865	0.579536679697902\\
10.79	0.579875147369384\\
10.715	0.58021606150411\\
10.64	0.580559461632833\\
10.565	0.580905388456385\\
10.495	0.581228453491394\\
10.42	0.581575840883396\\
10.345	0.581925851525882\\
10.27	0.582278530677524\\
10.195	0.582633925022898\\
10.115	0.583016059379349\\
10.04	0.583377219538565\\
9.965	0.583741247641285\\
9.89	0.584108196867011\\
9.815	0.58447812218252\\
9.74	0.5848510804382\\
9.665	0.585227130436299\\
9.59	0.585606333043588\\
9.515	0.585988751281857\\
9.445	0.586345430108745\\
9.37	0.586730350770998\\
9.295	0.587118647263667\\
9.22	0.587510391348723\\
9.145	0.587905657496404\\
9.065	0.588331243979841\\
8.99	0.588734037421103\\
8.915	0.58914059966246\\
8.84	0.589551017678788\\
8.765	0.589965381934569\\
8.69	0.590383786565297\\
8.615	0.590806329569331\\
8.545	0.59120291821945\\
8.47	0.591629607620087\\
8.395	0.592060700220832\\
8.32	0.592496310121986\\
8.245	0.592936556374917\\
8.17	0.593381563254477\\
8.095	0.593831460547832\\
8.015	0.594316894216766\\
7.94	0.594777334996464\\
7.865	0.595243101881472\\
7.79	0.59571435041222\\
7.72	0.596158224035005\\
7.645	0.596636092914729\\
7.57	0.597119883781652\\
7.495	0.597609781392549\\
7.42	0.598105979459426\\
7.345	0.598608681163365\\
7.27	0.599118099712261\\
7.195	0.599634458920538\\
7.12	0.60015799383559\\
7.045	0.600688951398602\\
6.97	0.601226061677109\\
6.895	0.601767843242043\\
6.82	0.602317790082459\\
6.745	0.602876201442447\\
6.67	0.603443392305362\\
6.595	0.604019694342415\\
6.52	0.604605456925402\\
6.445	0.605201048179853\\
6.37	0.605806856140978\\
6.3	0.60637902843677\\
6.225	0.607000225201181\\
6.15	0.607632797155262\\
6.075	0.608277218724596\\
6	0.608933990745294\\
5.925	0.609603642068285\\
5.845	0.610332762607761\\
5.77	0.611030836799365\\
5.7	0.611689442158706\\
5.625	0.612409335160704\\
5.55	0.613145070283563\\
5.475	0.613897376798839\\
5.4	0.614667025914437\\
5.325	0.615454833282812\\
5.255	0.616206030886656\\
5.18	0.617023529158985\\
5.105	0.617861653974311\\
5.03	0.618721413038019\\
4.955	0.619603872400332\\
4.88	0.620510159766851\\
4.805	0.621438216065709\\
4.73	0.622387023024689\\
4.655	0.623363186140366\\
4.58	0.624368102281762\\
4.505	0.625403247521404\\
4.435	0.6263953748757\\
4.36	0.627483833165658\\
4.285	0.628607029816474\\
4.21	0.629766773001691\\
4.135	0.630964967408388\\
4.065	0.632111236837287\\
3.99	0.633378844274568\\
3.915	0.634690693628427\\
3.84	0.636049073932631\\
3.765	0.637445086597728\\
3.69	0.638891556981338\\
3.615	0.640391582668245\\
3.545	0.64183551280659\\
3.47	0.643431624815294\\
3.395	0.645088757197963\\
3.325	0.646687849891222\\
3.25	0.64845486199518\\
3.175	0.650290684777905\\
3.105	0.652060579751013\\
3.03	0.654020175577614\\
2.955	0.656055948611998\\
2.885	0.658011130693741\\
2.81	0.660182420926922\\
2.735	0.662425012034154\\
2.66	0.664741856104613\\
2.59	0.666972205523608\\
2.515	0.66942848745481\\
2.445	0.671787276629842\\
2.37	0.674371215385532\\
2.3	0.676840685453892\\
2.225	0.679529920220235\\
2.155	0.682079292079041\\
2.08	0.684831946836596\\
2.01	0.687408154390079\\
1.935	0.690150608681882\\
1.865	0.692664812534902\\
1.79	0.695273931441511\\
1.72	0.697581452306701\\
1.64	0.699999935720575\\
1.57	0.701847259319141\\
1.495	0.703438990679207\\
1.42	0.704504762256797\\
1.345	0.704869554450413\\
1.27	0.704313126038761\\
1.195	0.702559882189189\\
1.125	0.699528356280903\\
1.055	0.694778503719184\\
0.985	0.687847317663844\\
0.92	0.678940793383244\\
0.855	0.667024900313387\\
0.795	0.652696794637406\\
0.74	0.636062679726888\\
0.68	0.613129488273384\\
0.63	0.589217924431571\\
0.581947132883608	0.560987896879085\\
0.54	0.530911906962008\\
0.5	0.495877198125237\\
0.463166488796296	0.455814890379004\\
0.430933003167371	0.410740744736113\\
0.405032374723676	0.360658360688455\\
0.388193562401176	0.305567738236032\\
0.386348986600311	0.240460638974078\\
0.406318851841981	0.185370016521654\\
0.441474533800984	0.145304109283528\\
0.487924743994969	0.115254678854934\\
0.54	0.0935858886679708\\
0.604959769085267	0.0751887716168081\\
0.665	0.0630874671528293\\
0.73	0.0532066653009179\\
0.799089494345231	0.0451393411882137\\
0.865	0.0391276399271436\\
0.935	0.0339871760140309\\
1.005	0.029791918751116\\
1.08	0.0261353579696706\\
1.15	0.0233012472509753\\
1.225	0.0207125274919451\\
1.295	0.0186994356405824\\
1.37	0.0168378774229096\\
1.445	0.0151827177949636\\
1.515	0.013923389396427\\
1.595	0.0126473808329802\\
1.67	0.0115741057375097\\
1.745	0.0105996180076262\\
1.815	0.00981133746552255\\
1.89	0.00911580955724423\\
1.965	0.00848449754697421\\
2.04	0.00790832065751883\\
2.115	0.00737965140478209\\
2.19	0.00689204558650669\\
2.265	0.00644003056606869\\
2.34	0.00601893506718137\\
2.415	0.00562475318198009\\
2.49	0.00525403608134915\\
2.56	0.0049661839453591\\
2.635	0.00472844459374244\\
2.715	0.004495937295368\\
2.79	0.00429567008613635\\
2.865	0.00411079842421719\\
2.94	0.00393982476373745\\
3.015	0.0037814272221527\\
3.09	0.00363443470794723\\
3.165	0.0034978068468364\\
3.24	0.0033706179464727\\
3.315	0.00325204102308165\\
3.39	0.00314133639564827\\
3.465	0.00303784100510815\\
3.54	0.00294095826992063\\
3.615	0.00285015161872368\\
3.69	0.00276493618832183\\
3.77	0.00267970963043034\\
3.845	0.00260470761170304\\
3.92	0.00253407876134529\\
3.995	0.00246749770282427\\
4.07	0.00240466830764199\\
4.145	0.00234532039406676\\
4.22	0.0022892065404239\\
4.295	0.00223610139293441\\
4.37	0.00218579782382131\\
4.445	0.00213810596311909\\
4.52	0.00209285154483201\\
4.595	0.00204987376303573\\
4.67	0.00200902482466622\\
4.745	0.00197016876515641\\
4.825	0.00193077710816715\\
4.9	0.00189565125406886\\
4.975	0.00186216158807596\\
5.05	0.00183020840209511\\
5.125	0.00179970018002497\\
5.2	0.00177055133733276\\
5.275	0.00174268249687697\\
5.35	0.00171601998246284\\
5.425	0.00169049456138717\\
5.5	0.0016660426611166\\
5.575	0.00164260436247496\\
5.65	0.00162012388492663\\
5.725	0.00159854927078198\\
5.8	0.00157783130311637\\
5.88	0.00155662555045656\\
5.955	0.00153753754575132\\
6.03	0.00151917557709821\\
6.105	0.00150150225333134\\
6.18	0.00148448263635859\\
6.255	0.00146808327242995\\
6.33	0.00145227365542823\\
6.405	0.00143702447065\\
6.48	0.00142230828366759\\
6.555	0.00140809941078447\\
6.63	0.00139437301529499\\
6.705	0.00138110661423629\\
6.78	0.00136827837691888\\
6.855	0.00135586784344404\\
6.93	0.00134385585054632\\
7.01	0.00133146142307437\\
7.085	0.00132021565045308\\
7.16	0.00130931547331513\\
7.235	0.00129874530989757\\
7.31	0.00128849071989418\\
7.385	0.00127853753760052\\
7.46	0.00126887264032147\\
7.535	0.00125948390633484\\
7.61	0.00125035935943349\\
7.685	0.00124148794707591\\
7.76	0.00123285868925301\\
7.835	0.00122446145747701\\
7.91	0.00121628695749535\\
7.985	0.00120832586555289\\
8.065	0.00120005964476853\\
8.14	0.00119251246787635\\
8.215	0.00118515328703815\\
8.29	0.00117797460349342\\
8.365	0.00117096941155973\\
8.44	0.00116413047870308\\
8.515	0.0011574517257042\\
8.59	0.00115092681248843\\
8.665	0.00114454982303479\\
8.74	0.00113831524878428\\
8.815	0.00113221727148652\\
8.89	0.00112625115928775\\
8.965	0.00112041184569202\\
9.04	0.00111469462122573\\
9.12	0.00110872574680078\\
9.195	0.00110324693190123\\
9.27	0.00109787671825805\\
9.345	0.00109261160146536\\
9.42	0.00108744767204951\\
9.495	0.00108238132238911\\
9.57	0.00107740924118627\\
9.645	0.00107252768259365\\
9.72	0.00106773390896712\\
9.795	0.00106302473191473\\
9.87	0.00105839723244736\\
9.945	0.00105384875938246\\
10.02	0.0010493761861185\\
10.095	0.00104497737662343\\
10.175	0.00104036366858159\\
10.25	0.00103610926771521\\
10.325	0.00103192120608695\\
10.4	0.00102779729386643\\
10.475	0.00102373543246813\\
10.55	0.00101973360418792\\
10.625	0.00101578986744176\\
10.7	0.00101190235931232\\
10.775	0.00100806928660703\\
10.85	0.00100428892236066\\
10.925	0.00100055960928056\\
11	0.000996879749592285\\
11.075	0.00099324780469729\\
11.15	0.000989662299122985\\
11.225	0.000986121808576344\\
11.305	0.000982393355169069\\
11.38	0.000978941603993089\\
11.455	0.000975530820254338\\
11.53	0.000972159781951955\\
11.605	0.000968827310315026\\
11.68	0.000965532272244083\\
11.755	0.000962273570182245\\
11.83	0.00095905015063787\\
11.905	0.000955860998326979\\
11.98	0.000952705130304319\\
12.055	0.000949581603121594\\
12.13	0.000946489503108448\\
12.205	0.000943427947244582\\
12.28	0.000940396091100855\\
12.36	0.000937193920022747\\
12.435	0.000934220870747487\\
12.51	0.000931275095284615\\
12.585	0.000928355856709106\\
12.66	0.000925462443107555\\
12.735	0.000922594171857362\\
12.81	0.000919750377073903\\
12.885	0.000916930422205428\\
12.96	0.000914133684590379\\
13.035	0.000911359572330856\\
13.11	0.000908607504757726\\
13.185	0.00090587692310052\\
13.26	0.000903167291590483\\
13.335	0.000900478085656384\\
13.415	0.000897631549772995\\
13.49	0.00089498314172409\\
13.565	0.000892353353184131\\
13.64	0.000889742200439732\\
13.715	0.000887149090321443\\
13.79	0.000884573594286209\\
13.865	0.000882015463098764\\
13.94	0.000879473816357012\\
14.015	0.000876948751167388\\
14.09	0.000874439734749317\\
14.165	0.000871946406438501\\
14.24	0.000869468576315629\\
14.315	0.000867005428213478\\
14.39	0.000864557111382507\\
14.47	0.000861961392180173\\
14.545	0.000859542412404054\\
14.62	0.000857137164642286\\
14.695	0.000854745527957702\\
14.77	0.000852366736373867\\
14.845	0.000850001008877983\\
14.92	0.000847647919342433\\
14.995	0.000845307217775319\\
};
\addplot [color=mycolor1, dashdotted, line width=0.9pt, forget plot]
  table[row sep=crcr]{%
0.01	0.0165664870989781\\
0.05	0.0717708079330276\\
0.085	0.10957376676504\\
0.125	0.144676514251908\\
0.16	0.17047853377217\\
0.2	0.196580576775226\\
0.235	0.213981938777264\\
0.275	0.232283371227682\\
0.31	0.245184380987814\\
0.35	0.257785367265151\\
0.385	0.267686142197344\\
0.425	0.274286658818807\\
0.46	0.282087269371444\\
0.5	0.286287598130557\\
0.535	0.288987809475701\\
0.575	0.292888114752019\\
0.61	0.295288302614369\\
0.65	0.296788420028338\\
0.685	0.297388466993926\\
0.725	0.297088443511132\\
0.76	0.29618837306275\\
0.8	0.294388232165988\\
0.835	0.292288067786432\\
0.875	0.289887879924082\\
0.91	0.287787715544526\\
0.95	0.286287598130557\\
0.985	0.285087504199382\\
1.025	0.282687316337032\\
1.06	0.279387058026301\\
1.1	0.275486752749982\\
1.135	0.272786541404838\\
1.175	0.271286423990869\\
1.21	0.269186259611313\\
1.25	0.265585977817788\\
1.285	0.26168567254147\\
1.325	0.258985461196326\\
1.36	0.257485343782357\\
1.4	0.255385179402801\\
1.435	0.25208492109207\\
1.475	0.248184615815751\\
1.51	0.246084451436195\\
1.55	0.244584334022226\\
1.585	0.242784193125463\\
1.625	0.239783958297526\\
1.66	0.236783723469589\\
1.7	0.234083512124445\\
1.74	0.232583394710476\\
1.775	0.231383300779301\\
1.815	0.229583159882539\\
1.85	0.226882948537395\\
1.89	0.224182737192251\\
1.925	0.222382596295489\\
1.965	0.22088247888152\\
2	0.219982408433139\\
2.04	0.218782314501964\\
2.075	0.217282197087995\\
2.115	0.214882009225645\\
2.15	0.213081868328883\\
2.19	0.21128172743212\\
2.225	0.210081633500945\\
2.265	0.209481586535358\\
2.3	0.20888153956977\\
2.34	0.207681445638595\\
2.375	0.206781375190214\\
2.415	0.205281257776245\\
2.45	0.203781140362276\\
2.49	0.202281022948308\\
2.525	0.201080929017133\\
2.565	0.200180858568751\\
2.6	0.199580811603164\\
2.64	0.198980764637576\\
2.675	0.198680741154783\\
2.715	0.198080694189195\\
2.75	0.197780670706401\\
2.79	0.197180623740814\\
2.825	0.196580576775226\\
2.865	0.195680506326845\\
2.9	0.195080459361258\\
2.94	0.194180388912876\\
2.975	0.193580341947289\\
3.015	0.192680271498908\\
3.05	0.19208022453332\\
3.09	0.191480177567733\\
3.125	0.190880130602145\\
3.165	0.190580107119351\\
3.2	0.190280083636558\\
3.24	0.18968003667097\\
3.275	0.18968003667097\\
3.315	0.189380013188177\\
3.355	0.189079989705383\\
3.39	0.189079989705383\\
3.43	0.188779966222589\\
3.465	0.188779966222589\\
3.505	0.188479942739795\\
3.54	0.188479942739795\\
3.58	0.188179919257002\\
3.615	0.188179919257002\\
3.655	0.188179919257002\\
3.69	0.187879895774208\\
3.73	0.187879895774208\\
3.765	0.187879895774208\\
3.805	0.187579872291414\\
3.84	0.187579872291414\\
3.88	0.187579872291414\\
3.915	0.187579872291414\\
3.955	0.18727984880862\\
3.99	0.18727984880862\\
4.03	0.18727984880862\\
4.065	0.18727984880862\\
4.105	0.186979825325827\\
4.14	0.186979825325827\\
4.18	0.186979825325827\\
4.215	0.186979825325827\\
4.255	0.186979825325827\\
4.29	0.186679801843033\\
4.33	0.186679801843033\\
4.365	0.186679801843033\\
4.405	0.186679801843033\\
4.44	0.186679801843033\\
4.48	0.186679801843033\\
4.515	0.186679801843033\\
4.555	0.186679801843033\\
4.59	0.186379778360239\\
4.63	0.186379778360239\\
4.665	0.186379778360239\\
4.705	0.186379778360239\\
4.74	0.186379778360239\\
4.78	0.186379778360239\\
4.815	0.186379778360239\\
4.855	0.186379778360239\\
4.89	0.186379778360239\\
4.93	0.186379778360239\\
4.965	0.186379778360239\\
5.005	0.186379778360239\\
5.045	0.186379778360239\\
5.08	0.186079754877445\\
5.12	0.186079754877445\\
5.155	0.186079754877445\\
5.195	0.186079754877445\\
5.23	0.186079754877445\\
5.27	0.186079754877445\\
5.305	0.186079754877445\\
5.345	0.186079754877445\\
5.38	0.186079754877445\\
5.42	0.186079754877445\\
5.455	0.186079754877445\\
5.495	0.186079754877445\\
5.53	0.186079754877445\\
5.57	0.185779731394652\\
5.605	0.185779731394652\\
5.645	0.185779731394652\\
5.68	0.185779731394652\\
5.72	0.185779731394652\\
5.755	0.185779731394652\\
5.795	0.185779731394652\\
5.83	0.185779731394652\\
5.87	0.185779731394652\\
5.905	0.185779731394652\\
5.945	0.185479707911858\\
5.98	0.185479707911858\\
6.02	0.185479707911858\\
6.055	0.185479707911858\\
6.095	0.185479707911858\\
6.13	0.185479707911858\\
6.17	0.185479707911858\\
6.205	0.185179684429064\\
6.245	0.185179684429064\\
6.28	0.185179684429064\\
6.32	0.185179684429064\\
6.355	0.185179684429064\\
6.395	0.185179684429064\\
6.43	0.18487966094627\\
6.47	0.18487966094627\\
6.505	0.18487966094627\\
6.545	0.18487966094627\\
6.58	0.18487966094627\\
6.62	0.184579637463477\\
6.655	0.184579637463477\\
6.695	0.184579637463477\\
6.735	0.184579637463477\\
6.77	0.184579637463477\\
6.81	0.184279613980683\\
6.845	0.184279613980683\\
6.885	0.184279613980683\\
6.92	0.184279613980683\\
6.96	0.183979590497889\\
6.995	0.183979590497889\\
7.035	0.183979590497889\\
7.07	0.183979590497889\\
7.11	0.183679567015095\\
7.145	0.183679567015095\\
7.185	0.183679567015095\\
7.22	0.183379543532302\\
7.26	0.183379543532302\\
7.295	0.183379543532302\\
7.335	0.183379543532302\\
7.37	0.183079520049508\\
7.41	0.183079520049508\\
7.445	0.183079520049508\\
7.485	0.182779496566714\\
7.52	0.182779496566714\\
7.56	0.182779496566714\\
7.595	0.18247947308392\\
7.635	0.18247947308392\\
7.67	0.18247947308392\\
7.71	0.18247947308392\\
7.745	0.182179449601127\\
7.785	0.182179449601127\\
7.82	0.182179449601127\\
7.86	0.181879426118333\\
7.895	0.181879426118333\\
7.935	0.181879426118333\\
7.97	0.181579402635539\\
8.01	0.181579402635539\\
8.045	0.181579402635539\\
8.085	0.181579402635539\\
8.12	0.181279379152745\\
8.16	0.181279379152745\\
8.195	0.181279379152745\\
8.235	0.180979355669952\\
8.27	0.180979355669952\\
8.31	0.180979355669952\\
8.35	0.180979355669952\\
8.385	0.180679332187158\\
8.425	0.180679332187158\\
8.46	0.180679332187158\\
8.5	0.180679332187158\\
8.535	0.180379308704364\\
8.575	0.180379308704364\\
8.61	0.180379308704364\\
8.65	0.180379308704364\\
8.685	0.18007928522157\\
8.725	0.18007928522157\\
8.76	0.18007928522157\\
8.8	0.18007928522157\\
8.835	0.18007928522157\\
8.875	0.179779261738777\\
8.91	0.179779261738777\\
8.95	0.179779261738777\\
8.985	0.179779261738777\\
9.025	0.179779261738777\\
9.06	0.179479238255983\\
9.1	0.179479238255983\\
9.135	0.179479238255983\\
9.175	0.179479238255983\\
9.21	0.179479238255983\\
9.25	0.179179214773189\\
9.285	0.179179214773189\\
9.325	0.179179214773189\\
9.36	0.179179214773189\\
9.4	0.179179214773189\\
9.435	0.179179214773189\\
9.475	0.179179214773189\\
9.51	0.178879191290395\\
9.55	0.178879191290395\\
9.585	0.178879191290395\\
9.625	0.178879191290395\\
9.66	0.178879191290395\\
9.7	0.178879191290395\\
9.735	0.178879191290395\\
9.775	0.178579167807602\\
9.81	0.178579167807602\\
9.85	0.178579167807602\\
9.885	0.178579167807602\\
9.925	0.178579167807602\\
9.96	0.178579167807602\\
10	0.178579167807602\\
10.04	0.178279144324808\\
10.075	0.178279144324808\\
10.115	0.178279144324808\\
10.15	0.178279144324808\\
10.19	0.178279144324808\\
10.225	0.178279144324808\\
10.265	0.177979120842014\\
10.3	0.177979120842014\\
10.34	0.177979120842014\\
10.375	0.177979120842014\\
10.415	0.177979120842014\\
10.45	0.17767909735922\\
10.49	0.17767909735922\\
10.525	0.17767909735922\\
10.565	0.17767909735922\\
10.6	0.17767909735922\\
10.64	0.177379073876427\\
10.675	0.177379073876427\\
10.715	0.177379073876427\\
10.75	0.177379073876427\\
10.79	0.177079050393633\\
10.825	0.177079050393633\\
10.865	0.177079050393633\\
10.9	0.177079050393633\\
10.94	0.176779026910839\\
10.975	0.176779026910839\\
11.015	0.176779026910839\\
11.05	0.176779026910839\\
11.09	0.176479003428045\\
11.125	0.176479003428045\\
11.165	0.176479003428045\\
11.2	0.176479003428045\\
11.24	0.176178979945252\\
11.275	0.176178979945252\\
11.315	0.176178979945252\\
11.35	0.175878956462458\\
11.39	0.175878956462458\\
11.425	0.175878956462458\\
11.465	0.175578932979664\\
11.5	0.175578932979664\\
11.54	0.175578932979664\\
11.575	0.17527890949687\\
11.615	0.17527890949687\\
11.65	0.174978886014077\\
11.69	0.174978886014077\\
11.73	0.174978886014077\\
11.765	0.174678862531283\\
11.805	0.174678862531283\\
11.84	0.174678862531283\\
11.88	0.174378839048489\\
11.915	0.174378839048489\\
11.955	0.174078815565695\\
11.99	0.174078815565695\\
12.03	0.174078815565695\\
12.065	0.173778792082902\\
12.105	0.173778792082902\\
12.14	0.173778792082902\\
12.18	0.173478768600108\\
12.215	0.173478768600108\\
12.255	0.173178745117314\\
12.29	0.173178745117314\\
12.33	0.173178745117314\\
12.365	0.17287872163452\\
12.405	0.17287872163452\\
12.44	0.17287872163452\\
12.48	0.172578698151727\\
12.515	0.172578698151727\\
12.555	0.172578698151727\\
12.59	0.172278674668933\\
12.63	0.172278674668933\\
12.665	0.172278674668933\\
12.705	0.171978651186139\\
12.74	0.171978651186139\\
12.78	0.171978651186139\\
12.815	0.171678627703345\\
12.855	0.171678627703345\\
12.89	0.171678627703345\\
12.93	0.171678627703345\\
12.965	0.171378604220552\\
13.005	0.171378604220552\\
13.04	0.171378604220552\\
13.08	0.171378604220552\\
13.115	0.171078580737758\\
13.155	0.171078580737758\\
13.19	0.171078580737758\\
13.23	0.171078580737758\\
13.265	0.170778557254964\\
13.305	0.170778557254964\\
13.345	0.170778557254964\\
13.38	0.170778557254964\\
13.42	0.17047853377217\\
13.455	0.17047853377217\\
13.495	0.17047853377217\\
13.53	0.17047853377217\\
13.57	0.17047853377217\\
13.605	0.17047853377217\\
13.645	0.170178510289377\\
13.68	0.170178510289377\\
13.72	0.170178510289377\\
13.755	0.170178510289377\\
13.795	0.170178510289377\\
13.83	0.170178510289377\\
13.87	0.169878486806583\\
13.905	0.169878486806583\\
13.945	0.169878486806583\\
13.98	0.169878486806583\\
14.02	0.169878486806583\\
14.055	0.169878486806583\\
14.095	0.169578463323789\\
14.13	0.169578463323789\\
14.17	0.169578463323789\\
14.205	0.169578463323789\\
14.245	0.169578463323789\\
14.28	0.169578463323789\\
14.32	0.169578463323789\\
14.355	0.169578463323789\\
14.395	0.169578463323789\\
14.43	0.169278439840995\\
14.47	0.169278439840995\\
14.505	0.169278439840995\\
14.545	0.169278439840995\\
14.58	0.169278439840995\\
14.62	0.169278439840995\\
14.655	0.169278439840995\\
14.695	0.169278439840995\\
14.73	0.168978416358202\\
14.77	0.168978416358202\\
14.805	0.168978416358202\\
14.845	0.168978416358202\\
14.88	0.168978416358202\\
14.92	0.168978416358202\\
14.955	0.168978416358202\\
14.995	0.168678392875408\\
};
\node[right, align=left, inner sep=0]
at (rel axis cs:-0.15,1.15) {$(a)$};
\coordinate (zoomlo) at (axis cs:2,0);
\coordinate (zoomhi) at (axis cs:2,2);
\end{axis}

\begin{axis}[%
name=ax_inset,
width=0.288\linewidth,
height=0.256\linewidth,
at={(ax_main.north east)}, anchor=north east, xshift=-10pt, yshift=-10pt,
scale only axis,
point meta min=-0.5,
point meta max=0.5,
axis on top,
separate axis lines,
every outer x axis line/.append style={black},
every x tick label/.append style={font=\color{black}},
every x tick/.append style={black},
xmin=0.005,
xmax=2,
every outer y axis line/.append style={black},
every y tick label/.append style={font=\color{black}},
every y tick/.append style={black},
ymin=0,
ymax=2,
axis background/.style={fill=white},
clip=true,
clip mode=individual,
axis lines=box,
xtick pos=bottom,
ytick pos=left,
tick align=inside,
major tick length=2pt,
every axis/.append style={font=\fontsize{6}{7}\selectfont},xlabel style={font=\fontsize{6}{7}\selectfont},title style={font=\fontsize{6}{7}\selectfont},
legend style={font=\fontsize{6}{7}\selectfont,inner sep=1pt,row sep=1pt,column sep=2pt,nodes={scale=1.00},draw=none},legend image post style={xscale=0.35},
legend columns=1,
scaled x ticks=false,
tick label style={/pgf/number format/fixed,/pgf/number format/precision=2},
every x tick label/.append style={font=\fontsize{6}{7}\selectfont\color{black}},
every y tick label/.append style={font=\fontsize{6}{7}\selectfont\color{black}},
,,
colormap={sigmamap}{rgb(0.0000)=(0.7760,0.8510,0.9140); rgb(0.1000)=(0.8392,0.8918,0.9369); rgb(0.2000)=(0.8913,0.9259,0.9563); rgb(0.3000)=(0.9374,0.9569,0.9743); rgb(0.4000)=(0.9799,0.9861,0.9917); rgb(0.5000)=(1.0000,0.9999,0.9999); rgb(0.6000)=(0.9659,0.8254,0.7818); rgb(0.7000)=(0.8840,0.4951,0.4053); rgb(0.8000)=(0.7820,0.2849,0.2466); rgb(0.9000)=(0.6131,0.1208,0.1566); rgb(1.0000)=(0.4040,0.0000,0.1220)}
]
\addplot [forget plot] graphics [xmin=0.0025, xmax=2.0025, ymin=-0.00243892365706085, ymax=3.00250411920238] {BQ_P3_sigma_map-1.png};
\addplot [color=black, line width=0.6pt, forget plot]
  table[row sep=crcr]{%
2	0.687774763111862\\
1.995	0.687958013565206\\
1.99	0.688141216352059\\
1.985	0.688324362478005\\
1.98	0.688507442784181\\
1.975	0.68869044794741\\
1.97	0.688873368476779\\
1.965	0.689056194710344\\
1.96	0.689238916813304\\
1.955	0.689421524774871\\
1.95	0.689604008405842\\
1.945	0.689786357335131\\
1.94	0.689968561007776\\
1.935	0.690150608681882\\
1.93	0.690332489425048\\
1.925	0.690514192111486\\
1.92	0.69069570542023\\
1.915	0.690877017830771\\
1.91	0.691058117620236\\
1.90601989263963	0.691202095402994\\
1.905	0.691238703281671\\
1.9	0.691417923685843\\
1.895	0.691596896232723\\
1.89	0.691775608447727\\
1.885	0.691954047638349\\
1.88	0.692132200892406\\
1.875	0.692310055074018\\
1.87	0.69248759681944\\
1.865	0.692664812534902\\
1.86	0.692841688392346\\
1.855	0.693018210325829\\
1.85	0.693194364027034\\
1.845	0.693370134942415\\
1.84	0.693545508269882\\
1.835	0.693720468953797\\
1.83	0.693895001680709\\
1.825	0.694069090876697\\
1.82	0.694242720702966\\
1.815	0.69441587505117\\
1.81	0.694588537539204\\
1.805	0.694760691507417\\
1.8	0.694932320014351\\
1.795	0.695103405832108\\
1.79	0.695273931441511\\
1.785	0.695443879028657\\
1.78	0.695613230478753\\
1.775	0.695781967372821\\
1.77	0.695950070983141\\
1.765	0.696117522267414\\
1.76221738511736	0.69621033380776\\
1.76	0.696283742622443\\
1.755	0.696448574548904\\
1.75	0.696612699667393\\
1.745	0.696776097776611\\
1.74	0.696938748339181\\
1.735	0.697100630476279\\
1.73	0.697261722961818\\
1.725	0.69742200421756\\
1.72	0.697581452306701\\
1.715	0.697740044930267\\
1.71	0.697897759420252\\
1.705	0.698054572733027\\
1.7	0.698210461445779\\
1.695	0.698365401749426\\
1.69	0.698519369442031\\
1.685	0.698672339924015\\
1.68	0.698824288191481\\
1.675	0.698975188829676\\
1.67	0.699125016007443\\
1.665	0.699273743470695\\
1.66	0.699421344535704\\
1.655	0.699567792081589\\
1.65	0.699713058545575\\
1.645	0.699857115914887\\
1.64	0.699999935720575\\
1.635	0.700141489029733\\
1.63	0.700281746438504\\
1.625	0.700420678066138\\
1.62	0.70055825354663\\
1.615	0.700694442021898\\
1.61	0.700829212133361\\
1.605	0.700962532015142\\
1.6	0.701094369286179\\
1.59523472051014	0.701218572212526\\
1.595	0.701224641747663\\
1.59	0.701352376602954\\
1.585	0.701478540300574\\
1.58	0.701603098603738\\
1.575	0.701726016725427\\
1.57	0.701847259319141\\
1.565	0.701966790472472\\
1.56	0.702084573697842\\
1.555	0.702200571924099\\
1.55	0.70231474748729\\
1.545	0.702427062121722\\
1.54	0.702537476951537\\
1.535	0.702645952481969\\
1.53	0.702752448588945\\
1.525	0.702856924510298\\
1.52	0.70295933883682\\
1.515	0.703059649502217\\
1.51	0.703157813772914\\
1.505	0.703253788239792\\
1.5	0.703347528806391\\
1.495	0.703438990679207\\
1.49	0.703528128358702\\
1.485	0.703614895627854\\
1.48	0.703699245542264\\
1.475	0.703781130418855\\
1.47	0.703860501825143\\
1.465	0.7039373105692\\
1.46	0.704011506688382\\
1.455	0.704083039437854\\
1.45	0.704151857279321\\
1.445	0.704217907869609\\
1.44	0.704281138049668\\
1.435	0.704341493832578\\
1.43	0.704398920390922\\
1.425	0.704453362046325\\
1.42	0.704504762256797\\
1.415	0.704553063603667\\
1.41	0.704598207780271\\
1.405	0.70464013557908\\
1.4	0.704678786879588\\
1.395	0.704714100634435\\
1.39	0.704746014857679\\
1.385	0.704774466611424\\
1.38	0.70479939199203\\
1.375	0.704820726117671\\
1.37	0.704838403115114\\
1.365	0.704852356104901\\
1.36	0.704862517188913\\
1.355	0.704868817436574\\
1.35	0.704871186869962\\
1.345	0.704869554450413\\
1.34	0.70486384806473\\
1.335	0.70485399451012\\
1.33	0.704839919480437\\
1.325	0.704821547551716\\
1.32	0.704798802166341\\
1.315	0.704771605620136\\
1.31	0.704739879046697\\
1.305	0.704703542401892\\
1.3	0.704662514449971\\
1.295	0.704616712747959\\
1.29	0.704566053630152\\
1.285	0.704510452192643\\
1.28	0.704449822279018\\
1.275	0.704384076464227\\
1.27	0.704313126038761\\
1.265	0.70423688099378\\
1.26	0.704155250004905\\
1.255	0.704068140416822\\
1.25	0.703975458227229\\
1.245	0.70387710807154\\
1.24	0.7037729932066\\
1.235	0.703663015495184\\
1.23	0.703547075389571\\
1.225	0.703425071916678\\
1.22	0.703296902661476\\
1.215	0.703162463750607\\
1.21	0.703021649837899\\
1.205	0.702874354087712\\
1.2	0.702720468158615\\
1.195	0.702559882189189\\
1.19	0.702392484781177\\
1.185	0.702218162983764\\
1.18	0.702036802278924\\
1.175	0.70184828656541\\
1.17	0.701652498143984\\
1.165	0.701449317701693\\
1.16	0.70123862429691\\
1.15954063187426	0.701218572212526\\
1.155	0.701018584023458\\
1.15	0.70079054055921\\
1.145	0.700554539595265\\
1.14	0.700310452203979\\
1.135	0.700058147735832\\
1.13	0.69979749380619\\
1.125	0.699528356280903\\
1.12	0.699250599262635\\
1.115	0.69896408507899\\
1.11	0.69866867426849\\
1.105	0.698364225568855\\
1.1	0.698050595904349\\
1.095	0.697727640374733\\
1.09	0.697395212244401\\
1.085	0.69705316293029\\
1.08	0.696701341993317\\
1.075	0.696339597128076\\
1.07326129481045	0.69621033380776\\
1.07	0.695965770305888\\
1.065	0.695580552330495\\
1.06	0.695184850616484\\
1.055	0.694778503719184\\
1.05	0.694361348266861\\
1.045	0.693933218954676\\
1.04	0.693493948539888\\
1.035	0.693043367837409\\
1.03	0.692581305716262\\
1.025	0.692107589097275\\
1.02	0.691622042951303\\
1.01577957564445	0.691202095402994\\
1.015	0.691123804512814\\
1.01	0.690609556691329\\
1.005	0.690082823924548\\
1	0.689543421398854\\
0.995	0.688991162332334\\
0.99	0.688425857984607\\
0.985	0.687847317663844\\
0.98	0.687255348731196\\
0.975	0.686649756611863\\
0.971319441945256	0.686193856998228\\
0.97	0.686028801713738\\
0.965	0.685389385643751\\
0.96	0.684735603336806\\
0.955	0.684067250211274\\
0.95	0.683384119794297\\
0.945	0.682686003741461\\
0.94	0.681972691863338\\
0.935	0.681243972145164\\
0.934607841125496	0.681185618593463\\
0.93	0.680492725167778\\
0.925	0.679724876700109\\
0.92	0.678940793383244\\
0.915	0.678140253404779\\
0.91	0.677323033259091\\
0.905	0.676488907783952\\
0.903169601562031	0.676177380188697\\
0.9	0.675631870323626\\
0.895	0.674753931430178\\
0.89	0.673858190014436\\
0.885	0.672944411601508\\
0.88	0.672012360275671\\
0.875564367539302	0.671169141783931\\
0.875	0.671060632501945\\
0.87	0.670080776472802\\
0.865	0.669081698841895\\
0.86	0.668063155277149\\
0.855	0.667024900313387\\
0.850917259494107	0.666160903379165\\
0.85	0.665964884330437\\
0.845	0.664876442656341\\
0.84	0.663767337762016\\
0.835	0.662637317717918\\
0.83	0.661486129873059\\
0.828577559587508	0.6611526649744\\
0.825	0.660305378108226\\
0.82	0.659099485873853\\
0.815	0.65787142874989\\
0.81	0.656620949358185\\
0.808127949555365	0.656144426569634\\
0.805	0.655339493716673\\
0.8	0.65402987652208\\
0.795	0.652696794637406\\
0.79	0.651339987667739\\
0.789261680463336	0.651136188164868\\
0.785	0.649946064478933\\
0.78	0.648525321837456\\
0.775	0.647079768933531\\
0.771763278089482	0.646127949760102\\
0.77	0.645602961025361\\
0.765	0.644089148971748\\
0.76	0.642549402608666\\
0.755434828471756	0.641119711355337\\
0.755	0.640981756093651\\
0.75	0.639369410002197\\
0.745	0.637729975279633\\
0.740144753166263	0.636111472950571\\
0.74	0.636062679726888\\
0.735	0.634350250137219\\
0.73	0.632609637017231\\
0.725742008691944	0.631103234545805\\
0.725	0.630837794631115\\
0.72	0.629020927846586\\
0.715	0.627174789801754\\
0.712120931818722	0.626094996141039\\
0.71	0.625289980075634\\
0.705	0.623362624564109\\
0.7	0.621404894332598\\
0.699199662784621	0.621086757736274\\
0.695	0.619395829472426\\
0.69	0.617351560171644\\
0.686932976180022	0.616078519331508\\
0.685	0.615265095287894\\
0.68	0.613129488273384\\
0.675250782351536	0.611070280926742\\
0.675	0.610959962577219\\
0.67	0.608728383498253\\
0.665	0.6064635019122\\
0.664126327913265	0.606062042521976\\
0.66	0.604142694573874\\
0.655	0.60178303049218\\
0.653476497701378	0.601053804117211\\
0.65	0.599369383378677\\
0.645	0.59691221790384\\
0.643260839184417	0.596045565712445\\
0.64	0.594399022537029\\
0.635	0.591838961066242\\
0.633455564977818	0.591037327307679\\
0.63	0.589217924431571\\
0.625	0.586549170743691\\
0.624038531456931	0.586029088902913\\
0.62	0.583810662183325\\
0.615	0.581026985598386\\
0.614989122928843	0.581020850498148\\
0.61	0.578162337061782\\
0.60629684631307	0.576012612093382\\
0.605	0.575250015652768\\
0.6	0.572271862642155\\
0.597898966452283	0.571004373688616\\
0.595	0.569231537115269\\
0.59	0.5661349918631\\
0.589778546903473	0.565996135283851\\
0.585	0.562954924544296\\
0.581947132883608	0.560987896879085\\
0.58	0.559712786320926\\
0.575	0.55639852242329\\
0.574375601995066	0.555979658474319\\
0.57	0.552992303296889\\
0.567074780471308	0.550971420069553\\
0.565	0.549515668383307\\
0.56	0.545966405590344\\
0.559995510225266	0.545963181664788\\
0.555	0.54232516831019\\
0.553139666803607	0.540954943260022\\
0.55	0.538605905243493\\
0.546485124810883	0.535946704855256\\
0.545	0.534803519470357\\
0.540033486246631	0.53093846645049\\
0.54	0.530911906962008\\
0.535	0.526903168366561\\
0.533799278515232	0.525930228045725\\
0.53	0.522794493761964\\
0.527754741057482	0.520921989640959\\
0.525	0.518589778950663\\
0.521871269339927	0.515913751236193\\
0.52	0.514286181397441\\
0.516151870274292	0.510905512831427\\
0.515	0.509874397903604\\
0.510600477049956	0.505897274426662\\
0.51	0.505343009147982\\
0.505220747681592	0.500889036021896\\
0.505	0.500679019229438\\
0.500003687166267	0.49588079761713\\
0.5	0.495877198125237\\
0.495	0.490951014420118\\
0.494921086257563	0.490872559212365\\
0.49	0.485889552440947\\
0.489975301132221	0.485864320807599\\
0.485174010691906	0.480856082402833\\
0.485	0.48067024220257\\
0.480522868436148	0.475847843998067\\
0.48	0.47527217233031\\
0.476007258560418	0.470839605593302\\
0.475	0.469701118842689\\
0.47160387161759	0.465831367188536\\
0.47	0.463965113045489\\
0.467320909456825	0.46082312878377\\
0.465	0.458034644874653\\
0.463166488796296	0.455814890379004\\
0.46	0.451876213134685\\
0.459146443140323	0.450806651974239\\
0.455232989641683	0.445798413569473\\
0.455	0.445493775115055\\
0.451427981803941	0.440790175164707\\
0.45	0.438861931979502\\
0.447734581669024	0.435781936759941\\
0.445	0.43195324857903\\
0.444163027295229	0.430773698355176\\
0.440707610475888	0.42576545995041\\
0.44	0.424716890168403\\
0.437344443268445	0.420757221545644\\
0.435	0.417167106500555\\
0.434079437190003	0.415748983140878\\
0.430933003167371	0.410740744736113\\
0.43	0.409203543042351\\
0.427904898466581	0.405732506331347\\
0.425	0.400803902293759\\
0.424953312776898	0.400724267926581\\
0.422106490834897	0.395716029521815\\
0.42	0.391871559245161\\
0.419365458018946	0.39070779111705\\
0.416744514367595	0.385699552712284\\
0.415	0.382273473477338\\
0.414198009634934	0.380691314307518\\
0.411748875249003	0.375683075902752\\
0.41	0.37194054931023\\
0.409410925038853	0.370674837497987\\
0.40718750942468	0.365666599093221\\
0.405032374723676	0.360658360688455\\
0.405	0.360579444613776\\
0.402984823961789	0.35565012228369\\
0.401057029440355	0.350641883878924\\
0.4	0.34776415123737\\
0.399219761543166	0.345633645474158\\
0.397466839946059	0.340625407069392\\
0.395825831778466	0.335617168664627\\
0.395	0.332889867142577\\
0.394310866070299	0.330608930259861\\
0.392868590888509	0.325600691855095\\
0.391529723836577	0.320592453450329\\
0.390331829588941	0.315584215045564\\
0.39	0.314098985202264\\
0.389214017858526	0.310575976640798\\
0.388193562401176	0.305567738236032\\
0.387327789339824	0.300559499831266\\
0.386548115234347	0.295551261426501\\
0.385858637548051	0.290543023021735\\
0.385348398614572	0.285534784616969\\
0.385	0.281230607678444\\
0.384942968026177	0.280526546212203\\
0.384621868295303	0.275518307807438\\
0.384464939298456	0.270510069402672\\
0.384451256040305	0.265501830997906\\
0.384541479960457	0.260493592593141\\
0.384744264591969	0.255485354188375\\
0.385	0.252323538508628\\
0.385151159456046	0.250477115783609\\
0.385681573713401	0.245468877378843\\
0.386348986600311	0.240460638974078\\
0.387223973564059	0.235452400569312\\
0.3882178759623	0.230444162164546\\
0.389421683713506	0.22543592375978\\
0.39	0.223304663435113\\
0.390791045966178	0.220427685355015\\
0.392355336893807	0.215419446950249\\
0.394126003519453	0.210411208545483\\
0.395	0.208165243865996\\
0.396090020891353	0.205402970140717\\
0.39829555859458	0.200394731735952\\
0.4	0.196845740725491\\
0.400711061413199	0.195386493331186\\
0.403394889351832	0.19037825492642\\
0.405	0.187610223537637\\
0.406318851841981	0.185370016521654\\
0.40953536980321	0.180361778116889\\
0.41	0.179687288285296\\
0.413033066832939	0.175353539712123\\
0.415	0.172755644169935\\
0.416854806663562	0.170345301307357\\
0.42	0.166531135576652\\
0.421001163213033	0.165337062902591\\
0.425	0.160886264579836\\
0.425509458262747	0.160328824497826\\
0.43	0.155729474410007\\
0.430406280027691	0.15532058609306\\
0.435	0.150973327625784\\
0.43571116402169	0.150312347688294\\
0.44	0.146567804974011\\
0.441474533800984	0.145304109283528\\
0.445	0.142451353517156\\
0.447715300242929	0.140295870878763\\
0.45	0.138582868869172\\
0.45448284418669	0.135287632473997\\
0.455	0.134928305900074\\
0.46	0.13151058329955\\
0.461836160282661	0.130279394069231\\
0.465	0.128267219446112\\
0.469811602287189	0.125271155664466\\
0.47	0.125159830668194\\
0.475	0.122248353243029\\
0.478483761043341	0.1202629172597\\
0.48	0.119442176258863\\
0.485	0.116781071814131\\
0.487924743994969	0.115254678854934\\
0.49	0.114224476279059\\
0.495	0.111785366174651\\
0.498216846177535	0.110246440450168\\
0.5	0.109433845515503\\
0.505	0.107192335894719\\
0.509452352994967	0.105238202045403\\
0.51	0.105009102533857\\
0.515	0.102945756483369\\
0.52	0.100926441913326\\
0.521751508390982	0.100229963640637\\
0.525	0.0989978114046441\\
0.53	0.0971349451152762\\
0.535	0.0953102162991138\\
0.53524549710534	0.0952217252358711\\
0.54	0.0935858886679708\\
0.545	0.0918992905160593\\
0.55	0.0902459350386614\\
0.550099285102202	0.0902134868311053\\
0.555	0.0886838511548485\\
0.56	0.0871531167594354\\
0.565	0.0856515288052999\\
0.566506292524277	0.0852052484263396\\
0.57	0.0842162765083621\\
0.575	0.0828240982487567\\
0.58	0.0814576136644576\\
0.584695983543427	0.0801970100215739\\
0.585	0.0801190161447827\\
0.59	0.0788506006288838\\
0.595	0.0776049489089468\\
0.6	0.07638136201233\\
0.604959769085267	0.0751887716168081\\
0.605	0.0751795257874264\\
0.61	0.0740423829421546\\
0.615	0.0729249001263936\\
0.62	0.0718264966547993\\
0.625	0.0707466124195741\\
0.627657002000467	0.0701805332120424\\
0.63	0.069703477481376\\
0.635	0.0686987839977837\\
0.64	0.0677106938533743\\
0.645	0.066738737538956\\
0.65	0.0657824614210998\\
0.653234690671652	0.0651722948072767\\
0.655	0.0648541488528058\\
0.66	0.063963738774519\\
0.665	0.0630874671528293\\
0.67	0.0622249506651828\\
0.675	0.061375818388415\\
0.68	0.0605397113088349\\
0.682273549088261	0.0601640564025109\\
0.685	0.0597338757859855\\
0.69	0.0589548696330712\\
0.695	0.0581876799169745\\
0.7	0.0574320007075131\\
0.705	0.0566875355136623\\
0.71	0.0559539969014057\\
0.715	0.0552311061522768\\
0.715525659999347	0.0551558179977452\\
0.72	0.0545443841759337\\
0.725	0.0538706639366692\\
0.73	0.0532066653009179\\
0.735	0.0525521487592025\\
0.74	0.0519068817982052\\
0.745	0.0512706386501369\\
0.75	0.050643200056245\\
0.754000113731909	0.0501475795929794\\
0.755	0.0500295480456699\\
0.76	0.0494449711292397\\
0.765	0.0488684799018159\\
0.77	0.0482998845732108\\
0.775	0.0477390006289757\\
0.78	0.0471856486513351\\
0.785	0.0466396541466163\\
0.79	0.0461008473815363\\
0.795	0.0455690632258098\\
0.799089494345231	0.0451393411882137\\
0.8	0.0450483302768625\\
0.805	0.0445530537069971\\
0.81	0.0440642703502822\\
0.815	0.0435818360716807\\
0.82	0.0431056105129254\\
0.825	0.0426354569740327\\
0.83	0.0421712422657495\\
0.835	0.0417128367215808\\
0.84	0.0412601139351788\\
0.845	0.0408129506590004\\
0.85	0.0403712268559695\\
0.852745512440666	0.0401311027834479\\
0.855	0.0399440630512307\\
0.86	0.0395333176846072\\
0.865	0.0391276399271436\\
0.87	0.0387269241028817\\
0.875	0.0383310671282323\\
0.88	0.0379399684385362\\
0.885	0.0375535299107098\\
0.89	0.0371716557905072\\
0.895	0.0367942526253063\\
0.9	0.0364212291973398\\
0.905	0.0360524964555719\\
0.91	0.0356879674552394\\
0.915	0.0353275572998439\\
0.917866759859521	0.0351228643786822\\
0.92	0.0349788756458123\\
0.925	0.0346444656616887\\
0.93	0.0343139242385796\\
0.935	0.0339871760140309\\
0.94	0.0336641473488575\\
0.945	0.0333447662795086\\
0.95	0.0330289624716783\\
0.955	0.0327166671845233\\
0.96	0.0324078131929685\\
0.965	0.0321023347927941\\
0.97	0.0318001677478176\\
0.975	0.0315012492223099\\
0.98	0.0312055177696774\\
0.985	0.0309129132877396\\
0.99	0.0306233769846062\\
0.995	0.0303368513429363\\
0.998915123876247	0.0301146259739165\\
1	0.0300567251653793\\
1.005	0.029791918751116\\
1.01	0.0295299665555397\\
1.015	0.0292708170839378\\
1.02	0.0290144199425307\\
1.025	0.0287607257523631\\
1.03	0.028509686238109\\
1.035	0.0282612540845209\\
1.04	0.0280153829692804\\
1.045	0.0277720275276145\\
1.05	0.0275311433285704\\
1.055	0.027292686850894\\
1.06	0.0270566154612146\\
1.065	0.0268228873917689\\
1.07	0.0265914617192527\\
1.075	0.0263622983436196\\
1.08	0.0261353579696706\\
1.085	0.0259106020868777\\
1.09	0.0256879929491328\\
1.095	0.0254674935578837\\
1.1	0.0252490676448323\\
1.10329310725095	0.0251063875691507\\
1.105	0.0250373981909201\\
1.11	0.0248368359052896\\
1.115	0.0246382620773991\\
1.12	0.024441643938175\\
1.125	0.0242469493498433\\
1.13	0.0240541467925347\\
1.135	0.0238632053484556\\
1.14	0.023674094686681\\
1.145	0.0234867850520171\\
1.15	0.0233012472509753\\
1.155	0.0231174526306951\\
1.16	0.0229353731058774\\
1.165	0.0227549810704799\\
1.17	0.0225762494602706\\
1.175	0.0223991517044669\\
1.18	0.0222236617236676\\
1.185	0.0220497539156133\\
1.19	0.0218774031482515\\
1.195	0.0217065847468518\\
1.2	0.0215372744849809\\
1.205	0.0213694485743143\\
1.21	0.0212030836552274\\
1.215	0.0210381567878198\\
1.22	0.0208746454426964\\
1.225	0.0207125274919451\\
1.23	0.0205517811994753\\
1.235	0.0203923852165064\\
1.24	0.0202343185688486\\
1.24434166389937	0.020098149164385\\
1.245	0.0200791275482917\\
1.25	0.0199355594793499\\
1.255	0.0197932851206102\\
1.26	0.0196522853396494\\
1.265	0.0195125413342609\\
1.27	0.0193740346267131\\
1.275	0.0192367470561317\\
1.28	0.0191006607633212\\
1.285	0.0189657582244066\\
1.29	0.018832022166908\\
1.295	0.0186994356405824\\
1.3	0.0185679819742475\\
1.305	0.0184376447749169\\
1.31	0.018308407924811\\
1.315	0.0181802555758314\\
1.32	0.0180531721411698\\
1.325	0.0179271422942917\\
1.33	0.0178021509601341\\
1.335	0.0176781833122617\\
1.34	0.0175552247667133\\
1.345	0.0174332609791436\\
1.35	0.017312277837294\\
1.355	0.0171922614592296\\
1.36	0.017073198186006\\
1.365	0.0169550745813984\\
1.37	0.0168378774229096\\
1.375	0.0167215937008874\\
1.38	0.0166062106128313\\
1.385	0.0164917155612849\\
1.39	0.0163780961464552\\
1.395	0.0162653401677295\\
1.4	0.0161534356138239\\
1.405	0.0160423706312533\\
1.41	0.0159321336804705\\
1.415	0.0158227132119978\\
1.42	0.0157140979817448\\
1.425	0.0156062768859386\\
1.43	0.015499238998754\\
1.435	0.0153929735579443\\
1.44	0.0152874699671368\\
1.445	0.0151827177949636\\
1.44946027699911	0.0150899107596192\\
1.45	0.0150797909043401\\
1.455	0.0149865580173009\\
1.46	0.0148940722681521\\
1.465	0.0148023240375087\\
1.47	0.0147113038517623\\
1.475	0.0146210023789976\\
1.48	0.0145314104274856\\
1.485	0.0144425189413005\\
1.49	0.0143543190005283\\
1.495	0.0142668018173424\\
1.5	0.0141799587342806\\
1.505	0.014093781218714\\
1.51	0.0140082608671095\\
1.515	0.013923389396427\\
1.52	0.0138391586459535\\
1.525	0.0137555605719259\\
1.53	0.0136725872510729\\
1.535	0.0135902308705113\\
1.54	0.0135084836924882\\
1.545	0.0134273382518797\\
1.55	0.0133467869475918\\
1.555	0.0132668224501924\\
1.56	0.0131874374924065\\
1.565	0.0131086249122157\\
1.57	0.0130303776476415\\
1.575	0.0129526887388584\\
1.58	0.0128755513216361\\
1.585	0.0127989586306803\\
1.59	0.0127229039934432\\
1.595	0.0126473808329802\\
1.6	0.012572382662058\\
1.605	0.0124979030853012\\
1.61	0.0124239357946955\\
1.615	0.0123504745706584\\
1.62	0.0122775132784752\\
1.625	0.0122050458694593\\
1.63	0.0121330663761199\\
1.635	0.0120615689142092\\
1.64	0.0119905476769971\\
1.645	0.0119199969412034\\
1.65	0.0118499110568641\\
1.655	0.0117802844535257\\
1.66	0.0117111116353212\\
1.665	0.0116423871343582\\
1.67	0.0115741057375097\\
1.675	0.0115062620297198\\
1.68	0.0114388508507697\\
1.685	0.0113718670623616\\
1.69	0.0113053055960184\\
1.695	0.0112391614472616\\
1.7	0.0111734296823937\\
1.705	0.0111081054286514\\
1.71	0.0110431838797502\\
1.715	0.0109786602918188\\
1.72	0.0109145299848061\\
1.725	0.0108507883357956\\
1.73	0.0107874307864085\\
1.735	0.0107244528333432\\
1.74	0.0106618500359484\\
1.745	0.0105996180076262\\
1.75	0.0105377524197928\\
1.755	0.010476249000264\\
1.76	0.0104151035291636\\
1.765	0.0103543118424253\\
1.77	0.0102938698306932\\
1.775	0.0102337734324819\\
1.78	0.0101740186422914\\
1.785	0.0101146015038561\\
1.78778453872093	0.0100816723548535\\
1.79	0.0100592707169268\\
1.795	0.0100089966879509\\
1.8	0.00995906997475516\\
1.805	0.00990948694573381\\
1.81	0.00986024396993109\\
1.815	0.00981133746552255\\
1.82	0.00976276388963685\\
1.825	0.00971451974837969\\
1.83	0.00966660158865409\\
1.835	0.00961900599885422\\
1.84	0.00957172961118433\\
1.845	0.00952476909937238\\
1.85	0.00947812117539493\\
1.855	0.00943178259374463\\
1.86	0.0093857501466726\\
1.865	0.00934002066749923\\
1.87	0.00929459102604614\\
1.875	0.00924945813100146\\
1.88	0.00920461892661855\\
1.885	0.00916007039617993\\
1.89	0.00911580955724423\\
1.895	0.00907183346468683\\
1.9	0.00902813920754529\\
1.905	0.00898472391089898\\
1.91	0.00894158473080544\\
1.915	0.00889871882343577\\
1.92	0.00885612352457127\\
1.925	0.00881379598108254\\
1.93	0.0087717335198349\\
1.935	0.00872993346336288\\
1.94	0.0086883931651144\\
1.945	0.0086471100086141\\
1.95	0.00860608140905841\\
1.955	0.00856530481076822\\
1.96	0.0085247776890572\\
1.965	0.00848449754697421\\
1.97	0.0084444619195335\\
1.975	0.00840466836519448\\
1.98	0.0083651144740761\\
1.985	0.00832579786292902\\
1.99	0.00828671617610448\\
1.995	0.00824786708554161\\
2	0.00820924825688923\\
};
\addplot [color=mycolor1, dashdotted, line width=0.7pt, forget plot]
  table[row sep=crcr]{%
0.01	0.0165664870989781\\
0.015	0.0246671211344093\\
0.02	0.0318676847214592\\
0.025	0.0396682952740966\\
0.03	0.0465688353783528\\
0.035	0.0531693519998153\\
0.04	0.0591698216556902\\
0.045	0.0660703617599464\\
0.05	0.0717708079330276\\
0.055	0.0774712541061088\\
0.06	0.0834717237619838\\
0.065	0.0891721699350649\\
0.07	0.0942725691425586\\
0.075	0.0993729683500524\\
0.08	0.104173344074752\\
0.085	0.10957376676504\\
0.09	0.114974189455327\\
0.095	0.119474541697233\\
0.1	0.123374846973552\\
0.105	0.12877526966384\\
0.11	0.132675574940158\\
0.115	0.136875903699271\\
0.12	0.141376255941177\\
0.125	0.144676514251908\\
0.13	0.149476889976608\\
0.135	0.152777148287339\\
0.14	0.156977477046452\\
0.145	0.160577758839977\\
0.15	0.163577993667914\\
0.155	0.16807834590982\\
0.16	0.17047853377217\\
0.165	0.174378839048489\\
0.17	0.178279144324808\\
0.175	0.18007928522157\\
0.18	0.184279613980683\\
0.185	0.187579872291414\\
0.19	0.189380013188177\\
0.195	0.192680271498908\\
0.2	0.196580576775226\\
0.205	0.198680741154783\\
0.21	0.200780905534339\\
0.215	0.20408116384507\\
0.22	0.207681445638595\\
0.225	0.209481586535358\\
0.23	0.210981703949326\\
0.235	0.213981938777264\\
0.24	0.217282197087995\\
0.245	0.219682384950345\\
0.25	0.221182502364314\\
0.255	0.222682619778282\\
0.26	0.225382831123426\\
0.265	0.228683089434157\\
0.27	0.230783253813714\\
0.275	0.232283371227682\\
0.28	0.233483465158857\\
0.285	0.23528360605562\\
0.29	0.23768379391797\\
0.295	0.240384005263114\\
0.3	0.24248416964267\\
0.305	0.243984287056638\\
0.31	0.245184380987814\\
0.315	0.246084451436195\\
0.32	0.247584568850163\\
0.325	0.24968473322972\\
0.33	0.251784897609276\\
0.335	0.254185085471626\\
0.34	0.255685202885595\\
0.345	0.25688529681677\\
0.35	0.257785367265151\\
0.355	0.258685437713532\\
0.36	0.259585508161913\\
0.365	0.260785602093088\\
0.37	0.262285719507057\\
0.375	0.26408586040382\\
0.38	0.265886001300582\\
0.385	0.267686142197344\\
0.39	0.268886236128519\\
0.395	0.270086330059695\\
0.4	0.270986400508076\\
0.405	0.271586447473663\\
0.41	0.272186494439251\\
0.415	0.272786541404838\\
0.42	0.273386588370426\\
0.425	0.274286658818807\\
0.43	0.275186729267188\\
0.435	0.276386823198363\\
0.44	0.277586917129538\\
0.445	0.278787011060713\\
0.45	0.279987104991888\\
0.455	0.281187198923063\\
0.46	0.282087269371444\\
0.465	0.282987339819826\\
0.47	0.283587386785413\\
0.475	0.284487457233794\\
0.48	0.284787480716588\\
0.485	0.285387527682176\\
0.49	0.285687551164969\\
0.495	0.285987574647763\\
0.5	0.286287598130557\\
0.505	0.286587621613351\\
0.51	0.286887645096144\\
0.515	0.287487692061732\\
0.52	0.287787715544526\\
0.525	0.288087739027319\\
0.53	0.288687785992907\\
0.535	0.288987809475701\\
0.54	0.289587856441288\\
0.545	0.289887879924082\\
0.55	0.290487926889669\\
0.555	0.290787950372463\\
0.56	0.291387997338051\\
0.565	0.291988044303638\\
0.57	0.292288067786432\\
0.575	0.292888114752019\\
0.58	0.293188138234813\\
0.585	0.293488161717607\\
0.59	0.294088208683194\\
0.595	0.294388232165988\\
0.6	0.294688255648782\\
0.605	0.294988279131575\\
0.61	0.295288302614369\\
0.615	0.295588326097163\\
0.62	0.295888349579957\\
0.625	0.29618837306275\\
0.63	0.29618837306275\\
0.635	0.296488396545544\\
0.64	0.296488396545544\\
0.645	0.296788420028338\\
0.65	0.296788420028338\\
0.655	0.297088443511132\\
0.66	0.297088443511132\\
0.665	0.297088443511132\\
0.67	0.297088443511132\\
0.675	0.297088443511132\\
0.68	0.297388466993926\\
0.685	0.297388466993926\\
0.69	0.297388466993926\\
0.695	0.297388466993926\\
0.7	0.297388466993926\\
0.705	0.297088443511132\\
0.71	0.297088443511132\\
0.715	0.297088443511132\\
0.72	0.297088443511132\\
0.725	0.297088443511132\\
0.73	0.296788420028338\\
0.735	0.296788420028338\\
0.74	0.296788420028338\\
0.745	0.296488396545544\\
0.75	0.296488396545544\\
0.755	0.29618837306275\\
0.76	0.29618837306275\\
0.765	0.295888349579957\\
0.77	0.295888349579957\\
0.775	0.295588326097163\\
0.78	0.295288302614369\\
0.785	0.295288302614369\\
0.79	0.294988279131575\\
0.795	0.294688255648782\\
0.8	0.294388232165988\\
0.805	0.294088208683194\\
0.81	0.293788185200401\\
0.815	0.293488161717607\\
0.82	0.293188138234813\\
0.825	0.292888114752019\\
0.83	0.292588091269226\\
0.835	0.292288067786432\\
0.84	0.291988044303638\\
0.845	0.291688020820844\\
0.85	0.291387997338051\\
0.855	0.291087973855257\\
0.86	0.290787950372463\\
0.865	0.290487926889669\\
0.87	0.290187903406876\\
0.875	0.289887879924082\\
0.88	0.289587856441288\\
0.885	0.289287832958494\\
0.89	0.288987809475701\\
0.895	0.288687785992907\\
0.9	0.288387762510113\\
0.905	0.288087739027319\\
0.91	0.287787715544526\\
0.915	0.287787715544526\\
0.92	0.287487692061732\\
0.925	0.287187668578938\\
0.93	0.287187668578938\\
0.935	0.286887645096144\\
0.94	0.286587621613351\\
0.945	0.286587621613351\\
0.95	0.286287598130557\\
0.955	0.285987574647763\\
0.96	0.285987574647763\\
0.965	0.285687551164969\\
0.97	0.285687551164969\\
0.975	0.285387527682176\\
0.98	0.285387527682176\\
0.985	0.285087504199382\\
0.99	0.284787480716588\\
0.995	0.284487457233794\\
1	0.284187433751001\\
1.005	0.283887410268207\\
1.01	0.283587386785413\\
1.015	0.283287363302619\\
1.02	0.282987339819826\\
1.025	0.282687316337032\\
1.03	0.282387292854238\\
1.035	0.281787245888651\\
1.04	0.281487222405857\\
1.045	0.280887175440269\\
1.05	0.280287128474682\\
1.055	0.279987104991888\\
1.06	0.279387058026301\\
1.065	0.278787011060713\\
1.07	0.278486987577919\\
1.075	0.277886940612332\\
1.08	0.277286893646744\\
1.085	0.276686846681157\\
1.09	0.276386823198363\\
1.095	0.275786776232776\\
1.1	0.275486752749982\\
1.105	0.274886705784394\\
1.11	0.274586682301601\\
1.115	0.274286658818807\\
1.12	0.273686611853219\\
1.125	0.273386588370426\\
1.13	0.273086564887632\\
1.135	0.272786541404838\\
1.14	0.272786541404838\\
1.145	0.272486517922044\\
1.15	0.272186494439251\\
1.155	0.271886470956457\\
1.16	0.271886470956457\\
1.165	0.271586447473663\\
1.17	0.271286423990869\\
1.175	0.271286423990869\\
1.18	0.270986400508076\\
1.185	0.270686377025282\\
1.19	0.270386353542488\\
1.195	0.270086330059695\\
1.2	0.269786306576901\\
1.205	0.269486283094107\\
1.21	0.269186259611313\\
1.215	0.268886236128519\\
1.22	0.268586212645726\\
1.225	0.267986165680138\\
1.23	0.267686142197344\\
1.235	0.267086095231757\\
1.24	0.26648604826617\\
1.245	0.266186024783376\\
1.25	0.265585977817788\\
1.255	0.264985930852201\\
1.26	0.264385883886613\\
1.265	0.263785836921026\\
1.27	0.263185789955438\\
1.275	0.262885766472645\\
1.28	0.262285719507057\\
1.285	0.26168567254147\\
1.29	0.261385649058676\\
1.295	0.261085625575882\\
1.3	0.260485578610295\\
1.305	0.260185555127501\\
1.31	0.259885531644707\\
1.315	0.259585508161913\\
1.32	0.25928548467912\\
1.325	0.258985461196326\\
1.33	0.258685437713532\\
1.335	0.258385414230738\\
1.34	0.258385414230738\\
1.345	0.258085390747945\\
1.35	0.257785367265151\\
1.355	0.257785367265151\\
1.36	0.257485343782357\\
1.365	0.257185320299563\\
1.37	0.257185320299563\\
1.375	0.25688529681677\\
1.38	0.256585273333976\\
1.385	0.256285249851182\\
1.39	0.255985226368388\\
1.395	0.255685202885595\\
1.4	0.255385179402801\\
1.405	0.255085155920007\\
1.41	0.25448510895442\\
1.415	0.254185085471626\\
1.42	0.253585038506038\\
1.425	0.253285015023245\\
1.43	0.252684968057657\\
1.435	0.25208492109207\\
1.44	0.251784897609276\\
1.445	0.251184850643688\\
1.45	0.250584803678101\\
1.455	0.250284780195307\\
1.46	0.24968473322972\\
1.465	0.249084686264132\\
1.47	0.248784662781338\\
1.475	0.248184615815751\\
1.48	0.247884592332957\\
1.485	0.247584568850163\\
1.49	0.24728454536737\\
1.495	0.246984521884576\\
1.5	0.246684498401782\\
1.505	0.246384474918988\\
1.51	0.246084451436195\\
1.515	0.245784427953401\\
1.52	0.245484404470607\\
1.525	0.245484404470607\\
1.53	0.245184380987814\\
1.535	0.24488435750502\\
1.54	0.24488435750502\\
1.545	0.244584334022226\\
1.55	0.244584334022226\\
1.555	0.244284310539432\\
1.56	0.243984287056638\\
1.565	0.243984287056638\\
1.57	0.243684263573845\\
1.575	0.243384240091051\\
1.58	0.243084216608257\\
1.585	0.242784193125463\\
1.59	0.24248416964267\\
1.595	0.242184146159876\\
1.6	0.241884122677082\\
1.605	0.241584099194288\\
1.61	0.241284075711495\\
1.615	0.240684028745907\\
1.62	0.240384005263114\\
1.625	0.239783958297526\\
1.63	0.239483934814732\\
1.635	0.238883887849145\\
1.64	0.238583864366351\\
1.645	0.237983817400764\\
1.65	0.23768379391797\\
1.655	0.237083746952382\\
1.66	0.236783723469589\\
1.665	0.236483699986795\\
1.67	0.235883653021207\\
1.675	0.235583629538414\\
1.68	0.23528360605562\\
1.685	0.234983582572826\\
1.69	0.234683559090032\\
1.695	0.234383535607239\\
1.7	0.234083512124445\\
1.705	0.233783488641651\\
1.71	0.233783488641651\\
1.715	0.233483465158857\\
1.72	0.233183441676064\\
1.725	0.233183441676064\\
1.73	0.23288341819327\\
1.735	0.232583394710476\\
1.74	0.232583394710476\\
1.745	0.232583394710476\\
1.75	0.232283371227682\\
1.755	0.232283371227682\\
1.76	0.231983347744889\\
1.765	0.231683324262095\\
1.77	0.231683324262095\\
1.775	0.231383300779301\\
1.78	0.231383300779301\\
1.785	0.231083277296507\\
1.79	0.230783253813714\\
1.795	0.23048323033092\\
1.8	0.230183206848126\\
1.805	0.230183206848126\\
1.81	0.229883183365332\\
1.815	0.229583159882539\\
1.82	0.229283136399745\\
1.825	0.228983112916951\\
1.83	0.228383065951364\\
1.835	0.22808304246857\\
1.84	0.227783018985776\\
1.845	0.227482995502982\\
1.85	0.226882948537395\\
1.855	0.226582925054601\\
1.86	0.226282901571807\\
1.865	0.225982878089014\\
1.87	0.225382831123426\\
1.875	0.225082807640632\\
1.88	0.224782784157839\\
1.885	0.224482760675045\\
1.89	0.224182737192251\\
1.895	0.223882713709457\\
1.9	0.223582690226664\\
1.905	0.22328266674387\\
1.91	0.222982643261076\\
1.915	0.222682619778282\\
1.92	0.222682619778282\\
1.925	0.222382596295489\\
1.93	0.222082572812695\\
1.935	0.221782549329901\\
1.94	0.221782549329901\\
1.945	0.221482525847107\\
1.95	0.221482525847107\\
1.955	0.221182502364314\\
1.96	0.221182502364314\\
1.965	0.22088247888152\\
1.97	0.22088247888152\\
1.975	0.220582455398726\\
1.98	0.220582455398726\\
1.985	0.220582455398726\\
1.99	0.220282431915932\\
1.995	0.220282431915932\\
2	0.219982408433139\\
};
\end{axis}
% magnifier connectors between the two named axes
\draw[gray,line width=0.3pt] (zoomlo) -- (ax_inset.south west);
\draw[gray,line width=0.3pt] (zoomhi) -- (ax_inset.north west);
\end{tikzpicture}%
  \end{minipage}
  \hfill
  \begin{minipage}[t]{0.48\textwidth}
    \centering
    % This file was created by matlab2tikz.
%
\definecolor{mycolor1}{rgb}{0.20000,0.50000,0.30000}%
%
\begin{tikzpicture}

\begin{axis}[%
width=11.252in,
height=7.131in,
at={(1.887in,0.962in)},
scale only axis,
point meta min=-0.5,
point meta max=0.5,
axis on top,
clip=true,
separate axis lines,
every outer x axis line/.append style={black},
every x tick label/.append style={font=\color{black}},
every x tick/.append style={black},
xmin=0.005,
xmax=15,
xlabel={$t$},
every outer y axis line/.append style={black},
every y tick label/.append style={font=\color{black}},
every y tick/.append style={black},
ymin=0,
ymax=2,
ytick={  0, 0.5,   1, 1.5,   2},
axis background/.style={fill=white},
clip=true,
clip mode=individual,
axis lines=box,
xtick pos=bottom,
ytick pos=left,
tick align=inside,
major tick length=2pt,
width=0.8\linewidth,
height=0.8\linewidth,
every axis/.append style={font=\fontsize{9}{9}\selectfont},xlabel style={font=\fontsize{9}{9}\selectfont},title style={font=\fontsize{9}{9}\selectfont},
legend style={font=\fontsize{8}{9}\selectfont,inner sep=1pt,row sep=1pt,column sep=2pt,nodes={scale=1.00},draw=none},legend image post style={xscale=0.35},
legend columns=1,
scaled x ticks=false,
tick label style={/pgf/number format/fixed,/pgf/number format/precision=2},
every x tick label/.append style={font=\fontsize{9}{9}\selectfont\color{black}},
every y tick label/.append style={font=\fontsize{9}{9}\selectfont\color{black}},
,,
colormap={sigmamap}{rgb(0.0000)=(0.7760,0.8510,0.9140); rgb(0.1000)=(0.8392,0.8918,0.9369); rgb(0.2000)=(0.8913,0.9259,0.9563); rgb(0.3000)=(0.9374,0.9569,0.9743); rgb(0.4000)=(0.9799,0.9861,0.9917); rgb(0.5000)=(1.0000,0.9999,0.9999); rgb(0.6000)=(0.9659,0.8254,0.7818); rgb(0.7000)=(0.8840,0.4951,0.4053); rgb(0.8000)=(0.7820,0.2849,0.2466); rgb(0.9000)=(0.6131,0.1208,0.1566); rgb(1.0000)=(0.4040,0.0000,0.1220)}
]
\addplot [forget plot] graphics [xmin=-0.00878446115288221, xmax=15.0187844611529, ymin=-0.00166944908180301, ymax=2.0016694490818] {BQ_P3_sigma_map_qssa-1.png};
\addplot [color=black, line width=0.8pt, forget plot]
  table[row sep=crcr]{%
0.01	0.546161165845511\\
0.0475689223057644	0.545484598740567\\
0.0851378446115288	0.544821000095272\\
0.118632649453542	0.54424040066778\\
0.122706766917293	0.54416930198413\\
0.160275689223058	0.543525294932562\\
0.197844611528822	0.542892904481303\\
0.235413533834586	0.542271714320357\\
0.310551378446115	0.541061178871774\\
0.32071185261835	0.540901502504174\\
0.34812030075188	0.540467656983913\\
0.385689223057644	0.539882790040356\\
0.423258145363409	0.539307345006525\\
0.460827067669173	0.538740820904269\\
0.498395989974937	0.538183475455224\\
0.535964912280702	0.537634251983742\\
0.540942632758624	0.537562604340568\\
0.573533834586466	0.537089907956987\\
0.611102756892231	0.536552651318848\\
0.648671679197995	0.536023415368614\\
0.686240601503759	0.535502052539666\\
0.723809523809524	0.534987530077304\\
0.761378446115288	0.534480785347708\\
0.780693662054715	0.534223706176962\\
0.836516290726817	0.533481005324573\\
0.874085213032581	0.532990335589935\\
0.911654135338346	0.532506496429839\\
0.94922305764411	0.532028232372601\\
0.986791979949875	0.531556104631557\\
1.02436090225564	0.531090702227141\\
1.04118461618918	0.530884808013356\\
1.0619298245614	0.530629239622541\\
1.09949874686717	0.53017138667989\\
1.13706766917293	0.52971907332712\\
1.1746365914787	0.529272762479861\\
1.21220551378446	0.528832307364121\\
1.24977443609023	0.528396301146548\\
1.28734335839599	0.527964981019615\\
1.32427683689341	0.52754590984975\\
1.32491228070175	0.527538654511904\\
1.40005012531328	0.526696618717139\\
1.43761904761905	0.526282255997085\\
1.47518796992481	0.525871964142777\\
1.51275689223058	0.525466015598158\\
1.55032581453634	0.525064681837858\\
1.58789473684211	0.524668233405281\\
1.62546365914787	0.524276419918909\\
1.63217510048647	0.524207011686144\\
1.66303258145363	0.523885754513156\\
1.7006015037594	0.523498149085558\\
1.73817042606516	0.523114301686316\\
1.77573934837093	0.522734418538351\\
1.81330827067669	0.522358705295641\\
1.85087719298246	0.521987367066331\\
1.88844611528822	0.521619924408591\\
1.92601503759399	0.521255393813175\\
1.96628828344396	0.520868113522538\\
2.00115288220551	0.520533156623955\\
2.03872180451128	0.520175595746973\\
2.07629072681704	0.519821576122525\\
2.11385964912281	0.519471252326915\\
2.15142857142857	0.51912477857399\\
2.18899749373434	0.518781633774958\\
2.2265664160401	0.518440979395729\\
2.26413533834586	0.518102923705672\\
2.30170426065163	0.517767582249679\\
2.32863388235687	0.517529215358932\\
2.33927318295739	0.517434333795342\\
2.37684210526316	0.517102190483403\\
2.41441102756892	0.516773133903986\\
2.45197994987469	0.516447279286597\\
2.48954887218045	0.516124741598385\\
2.56468671679198	0.515488041074519\\
2.60225563909774	0.515172943398648\\
2.63982456140351	0.514860102326118\\
2.67739348370927	0.514549603253674\\
2.71496240601504	0.514241531380095\\
2.72125784257892	0.514190317195326\\
2.7525313283208	0.513933888109457\\
2.79010025062657	0.513628446671214\\
2.82766917293233	0.513325712103004\\
2.86523809523809	0.513025769411683\\
2.90280701754386	0.512728703414742\\
2.94037593984962	0.51243402798078\\
2.97794486215539	0.512141125161325\\
3.01551378446115	0.511850056415731\\
3.05308270676692	0.511560884533809\\
3.09065162907268	0.511273672158963\\
3.14640687305472	0.510851419031719\\
3.16578947368421	0.510704124842197\\
3.20335839598997	0.510420761198053\\
3.24092731829574	0.510139628177208\\
3.2784962406015	0.509860788160893\\
3.31606516290727	0.50958430338862\\
3.35363408521303	0.509310235962734\\
3.3912030075188	0.509038508878037\\
3.42877192982456	0.508768365860406\\
3.46634085213033	0.508499721604743\\
3.50390977443609	0.508232621845585\\
3.54147869674185	0.507967112213196\\
3.57904761904762	0.507703238235655\\
3.60637337489859	0.507512520868113\\
3.61661654135338	0.507440447511823\\
3.65418546365915	0.507177804034646\\
3.72932330827068	0.506657929223634\\
3.76689223057644	0.506400788704097\\
3.80446115288221	0.506145573157691\\
3.84203007518797	0.505892327754978\\
3.87959899749373	0.505641097576756\\
3.9171679197995	0.50539192602963\\
3.95473684210526	0.505144344522298\\
3.99230576441103	0.504897979929063\\
4.02987468671679	0.504652865251023\\
4.06744360902256	0.504409033421217\\
4.10391160193612	0.504173622704507\\
4.10501253132832	0.504166469320236\\
4.14258145363408	0.503923673671515\\
4.18015037593985	0.50368226984505\\
4.21771929824561	0.503442290701683\\
4.25528822055138	0.503203769044283\\
4.33042606516291	0.502731229116777\\
4.36799498746867	0.502497276174318\\
4.40556390977444	0.502264911375845\\
4.4431328320802	0.502034167254316\\
4.48070175438596	0.501805076292893\\
4.51827067669173	0.50157767092641\\
4.55583959899749	0.501351951945181\\
4.59340852130326	0.501127402513511\\
4.63097744360902	0.500903848517712\\
4.64264572984922	0.500834724540901\\
4.66854636591479	0.500680242905305\\
4.70611528822055	0.500457205647351\\
4.74368421052632	0.500235242923041\\
4.78125313283208	0.500014378301708\\
4.81882205513784	0.499794635316748\\
4.85639097744361	0.499576037466287\\
4.93152882205514	0.499142370989054\\
4.9690977443609	0.498927349188302\\
5.00666666666667	0.498713566175505\\
5.04423558897243	0.498501045282802\\
5.0818045112782	0.498289809811307\\
5.11937343358396	0.498079883031855\\
5.15694235588972	0.497871288185776\\
5.19451127819549	0.497664048485667\\
5.22521030366265	0.497495826377296\\
5.23208020050125	0.497457909481658\\
5.26964912280702	0.497251942117938\\
5.30721804511278	0.497047195230255\\
5.34478696741855	0.496843261202439\\
5.38235588972431	0.496640145289027\\
5.41992481203008	0.496437864274179\\
5.45749373433584	0.496236434918854\\
5.53263157894737	0.495836198116681\\
5.57020050125313	0.495637424078875\\
5.6077694235589	0.495439568519396\\
5.64533834586466	0.495242648088484\\
5.68290726817043	0.495046679415334\\
5.72047619047619	0.494851679108484\\
5.75804511278196	0.494657663756196\\
5.79561403508772	0.494464649926856\\
5.83318295739348	0.494272654169377\\
5.85594915194473	0.49415692821369\\
5.87075187969925	0.494081108404014\\
5.90832080200501	0.493889723410103\\
5.94588972431078	0.493699415348911\\
5.98345864661654	0.493510200801723\\
6.02102756892231	0.493322096331884\\
6.05859649122807	0.493135118485245\\
6.1337343358396	0.492764597510381\\
6.17130325814536	0.492580716719104\\
6.20887218045113	0.49239748234437\\
6.24644110275689	0.492214906211383\\
6.28401002506266	0.492033000130734\\
6.32157894736842	0.491851775898587\\
6.35914786967419	0.491671245296864\\
6.39671679197995	0.491491420093443\\
6.43428571428571	0.491312312042348\\
6.47185463659148	0.491133932883946\\
6.50942355889724	0.490956294345155\\
6.5387887487295	0.490818030050083\\
6.54699248120301	0.490779092776668\\
6.58456140350877	0.49060153309758\\
6.62213032581454	0.490424756445916\\
6.6596992481203	0.49024877457692\\
6.73483709273183	0.489899242143463\\
6.77240601503759	0.489725715025728\\
6.80997493734336	0.489553029584585\\
6.84754385964912	0.489381197512917\\
6.88511278195489	0.489210230491852\\
6.92268170426065	0.489040140191\\
6.96025062656642	0.48887093826871\\
6.99781954887218	0.488702636372315\\
7.03538847117794	0.48853524613839\\
7.07295739348371	0.48836877919301\\
7.11052631578947	0.488203247152014\\
7.14809523809524	0.48803845796462\\
7.185664160401	0.487874180470783\\
7.22323308270677	0.487710422524178\\
7.26080200501253	0.487547192378805\\
7.27651733511448	0.487479131886477\\
7.33593984962406	0.487220151322345\\
7.37350877192982	0.487057172468274\\
7.41107769423559	0.486894760014395\\
7.44864661654135	0.486732922231025\\
7.48621553884712	0.486571667379757\\
7.52378446115288	0.486411003713573\\
7.56135338345865	0.486250939476954\\
7.59892230576441	0.486091482906003\\
7.63649122807018	0.485932642228554\\
7.67406015037594	0.485774425664297\\
7.7116290726817	0.485616841424894\\
7.74919799498747	0.485459897714104\\
7.78676691729323	0.485303602727906\\
7.824335839599	0.485147964654622\\
7.86190476190476	0.484992991675046\\
7.89947368421053	0.484838691962571\\
7.97461152882205	0.484532144996279\\
8.01218045112782	0.484379914053421\\
8.04974937343358	0.484228388999853\\
8.07170902870639	0.484140233722871\\
8.08731829573935	0.484077017231792\\
8.12488721804511	0.483925585711796\\
8.16245614035088	0.483774891816382\\
8.20002506265664	0.483624943722441\\
8.23759398496241	0.483475749600429\\
8.27516290726817	0.483327274672815\\
8.31273182957393	0.483179252884325\\
8.3503007518797	0.483031641839791\\
8.38786967418546	0.482884447309027\\
8.42543859649123	0.482737675056372\\
8.46300751879699	0.482591330840749\\
8.50057644110276	0.482445420415717\\
8.57571428571429	0.482154923925171\\
8.61328320802005	0.482010349340454\\
8.65085213032581	0.481866231508035\\
8.68842105263158	0.481722576155494\\
8.72598997493734	0.481579389005389\\
8.76355889724311	0.481436675775317\\
8.80112781954887	0.48129444217797\\
8.83869674185464	0.481152693921201\\
8.8762656641604	0.481011436708081\\
8.91383458646617	0.480870676236965\\
8.93240701087293	0.480801335559265\\
8.95140350877193	0.480729869238883\\
8.98897243107769	0.480589037923267\\
9.02654135338346	0.480448724896267\\
9.06411027568922	0.480308935875334\\
9.10167919799499	0.480169676573476\\
9.17681704260652	0.479892769957205\\
9.21438596491228	0.4797551340472\\
9.25195488721804	0.479618050665228\\
9.28952380952381	0.479481525503108\\
9.32709273182957	0.479345564248633\\
9.36466165413534	0.479210172585642\\
9.4022305764411	0.479075356194094\\
9.43979949874687	0.478941120750142\\
9.47736842105263	0.478807471926202\\
9.5149373433584	0.478674415391037\\
9.55250626566416	0.478541956809825\\
9.59007518796992	0.478410100102843\\
9.62764411027569	0.478278675253918\\
9.66521303258145	0.478147569867176\\
9.70278195488722	0.478016787913763\\
9.74035087719298	0.477886333361711\\
9.81548872180451	0.477626422318399\\
9.85305764411028	0.477496973747872\\
9.8631069163434	0.477462437395659\\
9.89062656641604	0.477367215409512\\
9.9281954887218	0.477237568463892\\
9.96576441102757	0.477108275406211\\
10.0033333333333	0.476979340207744\\
10.0409022556391	0.476850766836786\\
10.0784711779449	0.476722559258682\\
10.1160401002506	0.47659472143586\\
10.1536090225564	0.476467257327859\\
10.1911779448622	0.476340170891363\\
10.2287468671679	0.476213466080229\\
10.2663157894737	0.476087146845523\\
10.3038847117794	0.475961217135549\\
10.3414536340852	0.475835680895883\\
10.4165914786967	0.475585804596336\\
10.4541604010025	0.475461472414264\\
10.4917293233083	0.475337549458184\\
10.529298245614	0.475214039660531\\
10.5668671679198	0.475090946951213\\
10.6044360902256	0.474968275257646\\
10.6420050125313	0.474846028504792\\
10.6795739348371	0.474724210615189\\
10.7171428571429	0.474602825508994\\
10.7547117794486	0.474481877104011\\
10.7922807017544	0.474361369315735\\
10.8298496240602	0.474241306057384\\
10.8668380817063	0.474123539232053\\
10.8674185463659	0.474121677698849\\
10.9049874686717	0.474001642231496\\
10.9425563909774	0.473882066672766\\
11.017694235589	0.473644310992537\\
11.0552631578947	0.47352613872137\\
11.0928320802005	0.473408442059567\\
11.1304010025063	0.473291224755736\\
11.167969924812	0.473174364116877\\
11.2055388471178	0.473057754076458\\
11.2431077694236	0.472941397308185\\
11.2806766917293	0.472825296483956\\
11.3182456140351	0.472709454273877\\
11.3558145363409	0.472593873346274\\
11.3933834586466	0.472478556367704\\
11.4309523809524	0.472363506002977\\
11.4685213032581	0.47224872491516\\
11.5060902255639	0.4721342157656\\
11.5436591478697	0.472019981213933\\
11.5812280701754	0.471906023918102\\
11.656365914787	0.471678951717331\\
11.6939348370927	0.471565842119938\\
11.7315037593985	0.471453020393503\\
11.7690726817043	0.471340489187719\\
11.80664160401	0.471228251150678\\
11.8442105263158	0.471116308928882\\
11.8817794486216	0.47100466516726\\
11.9193483709273	0.470893322509188\\
11.9561196131706	0.470784641068447\\
11.9569172932331	0.470782265329912\\
11.9944862155388	0.470670674960913\\
12.0320551378446	0.470559396313421\\
12.0696240601504	0.470448432042149\\
12.1071929824561	0.47033778480028\\
12.1447619047619	0.470227457239484\\
12.1823308270677	0.470117452009939\\
12.2574686716792	0.469898419137932\\
12.295037593985	0.469789396788508\\
12.3326065162907	0.469680707356444\\
12.3701754385965	0.469572353484707\\
12.4077443609023	0.46946433781488\\
12.445313283208	0.469356662987176\\
12.4828822055138	0.469249331640457\\
12.5204511278195	0.469142346412256\\
12.5580200501253	0.469035709938791\\
12.5955889724311	0.468929424854989\\
12.6331578947368	0.468823493794501\\
12.6707268170426	0.468717919389725\\
12.7082957393484	0.468612704271822\\
12.7458646616541	0.468507851070737\\
12.7834335839599	0.468403362415221\\
12.8210025062657	0.468299240932848\\
12.8961403508772	0.468092109992065\\
12.933709273183	0.46798909233496\\
12.9712781954887	0.467886284791264\\
13.0088471177945	0.467783646828094\\
13.0464160401002	0.467681181755341\\
13.083984962406	0.467578892880511\\
13.1215538847118	0.467476783508744\\
13.1329946058905	0.467445742904841\\
13.1591228070175	0.467374277911693\\
13.1966917293233	0.46727170662493\\
13.2342606516291	0.467169326569583\\
13.2718295739348	0.467067141066598\\
13.3093984962406	0.466965153434547\\
13.3469674185464	0.46686336698964\\
13.3845363408521	0.466761785045743\\
13.4221052631579	0.466660410914388\\
13.4972431077694	0.466458299323867\\
13.5348120300752	0.466357568476242\\
13.572380952381	0.466257058664274\\
13.6099498746867	0.466156773188062\\
13.6475187969925	0.466056715345468\\
13.6850877192982	0.46595688843213\\
13.722656641604	0.46585729574148\\
13.7602255639098	0.46575794056476\\
13.7977944862155	0.46565882619104\\
13.8353634085213	0.465559955907234\\
13.8729323308271	0.46546133299812\\
13.9105012531328	0.465362960746356\\
13.9480701754386	0.465264842432498\\
13.9856390977444	0.46516698133502\\
14.0232080200501	0.465069380730332\\
14.0607769423559	0.464972043892797\\
14.1359147869674	0.464778174606539\\
14.1734837092732	0.464681648696493\\
14.2110526315789	0.464585399630997\\
14.2486215538847	0.464489430674484\\
14.2861904761905	0.464393745089461\\
14.3237593984962	0.464298346136531\\
14.361328320802	0.464203237074411\\
14.3988972431078	0.464108421159959\\
14.3995237930668	0.464106844741235\\
14.4364661654135	0.464013103250442\\
14.4740350877193	0.463918074290464\\
14.5116040100251	0.463823351100327\\
14.5491729323308	0.463728936954582\\
14.5867418546366	0.463634835125987\\
14.6243107769424	0.463541048885526\\
14.6618796992481	0.463447581502436\\
14.7370175438596	0.463261616376711\\
14.7745864661654	0.463169125164023\\
14.8121553884712	0.463076965868647\\
14.8497243107769	0.462985141751443\\
14.8872932330827	0.462893656071668\\
14.9248621553885	0.462802512087006\\
14.9624310776942	0.462711713053592\\
15	0.462621262226038\\
};
\addplot [color=black, line width=0.8pt, forget plot]
  table[row sep=crcr]{%
0.01	1.48804531780577e-05\\
0.0475689223057644	1.48813443937254e-05\\
0.0851378446115288	1.4882501653725e-05\\
0.122706766917293	1.48838250698126e-05\\
0.160275689223058	1.48848962451698e-05\\
0.197844611528822	1.48860377575492e-05\\
0.235413533834586	1.48870905950305e-05\\
0.272982456140351	1.48883644479456e-05\\
0.310551378446115	1.4889441161904e-05\\
0.34812030075188	1.48905419502061e-05\\
0.385689223057644	1.4891725696214e-05\\
0.423258145363409	1.48929539916182e-05\\
0.460827067669173	1.48941196933127e-05\\
0.498395989974937	1.48951914172189e-05\\
0.535964912280702	1.48962601846194e-05\\
0.573533834586466	1.48973348489726e-05\\
0.611102756892231	1.48985122409777e-05\\
0.648671679197995	1.48996838873613e-05\\
0.686240601503759	1.49008236014531e-05\\
0.723809523809524	1.49020225250992e-05\\
0.761378446115288	1.49031439978286e-05\\
0.798947368421053	1.49042210094447e-05\\
0.836516290726817	1.49053126024807e-05\\
0.874085213032581	1.49063823478308e-05\\
0.911654135338346	1.49074689409171e-05\\
0.94922305764411	1.49086422961422e-05\\
0.986791979949875	1.4909817187421e-05\\
1.02436090225564	1.49109059133534e-05\\
1.0619298245614	1.49119651005436e-05\\
1.09949874686717	1.49130587474536e-05\\
1.13706766917293	1.49141559537326e-05\\
1.1746365914787	1.49152257964114e-05\\
1.21220551378446	1.49162788591955e-05\\
1.24977443609023	1.49174022172779e-05\\
1.28734335839599	1.49185601988565e-05\\
1.32491228070175	1.49197085508651e-05\\
1.36248120300752	1.4920802996859e-05\\
1.40005012531328	1.49218328248687e-05\\
1.43761904761905	1.49228977311591e-05\\
1.47518796992481	1.49239895621987e-05\\
1.51275689223058	1.49250867677005e-05\\
1.55032581453634	1.49261677838295e-05\\
1.58789473684211	1.49272110333766e-05\\
1.62546365914787	1.49282265578536e-05\\
1.66303258145363	1.49292839917671e-05\\
1.7006015037594	1.49303708762804e-05\\
1.73817042606516	1.49314693823109e-05\\
1.77573934837093	1.49325616704298e-05\\
1.81330827067669	1.49336298907949e-05\\
1.85087719298246	1.49346561835132e-05\\
1.88844611528822	1.49356635656414e-05\\
1.92601503759399	1.49367068309997e-05\\
1.96358395989975	1.49377750440815e-05\\
2.00115288220551	1.49388561437526e-05\\
2.03872180451128	1.49399380619834e-05\\
2.07629072681704	1.49410087237393e-05\\
2.11385964912281	1.49420560470641e-05\\
2.15142857142857	1.49430679433609e-05\\
2.18899749373434	1.49440681190267e-05\\
2.2265664160401	1.49450981518926e-05\\
2.26413533834586	1.49461500680891e-05\\
2.30170426065163	1.49472154907537e-05\\
2.33927318295739	1.4948286038376e-05\\
2.37684210526316	1.49493533246974e-05\\
2.41441102756892	1.49504089586991e-05\\
2.45197994987469	1.49514445446919e-05\\
2.48954887218045	1.49524516824921e-05\\
2.52711779448622	1.49534387787749e-05\\
2.56468671679198	1.49544519979128e-05\\
2.60225563909774	1.49554884706846e-05\\
2.63982456140351	1.49565413021728e-05\\
2.67739348370927	1.49576035939135e-05\\
2.71496240601504	1.4958668443788e-05\\
2.7525313283208	1.49597289459775e-05\\
2.79010025062657	1.49607781909649e-05\\
2.82766917293233	1.49618092656011e-05\\
2.86523809523809	1.49628152532229e-05\\
2.90280701754386	1.49637892338352e-05\\
2.94037593984962	1.49647544977526e-05\\
2.97794486215539	1.49657421130616e-05\\
3.01551378446115	1.49667476398462e-05\\
3.05308270676692	1.49677665580365e-05\\
3.09065162907268	1.49687943451584e-05\\
3.12822055137845	1.4969826476266e-05\\
3.16578947368421	1.49708584238975e-05\\
3.20335839598997	1.49718856580633e-05\\
3.24092731829574	1.49729036462579e-05\\
3.2784962406015	1.49739078535061e-05\\
3.31606516290727	1.49748937424274e-05\\
3.35363408521303	1.49758567733393e-05\\
3.3912030075188	1.49767990567586e-05\\
3.42877192982456	1.49777554452929e-05\\
3.46634085213033	1.4978729028307e-05\\
3.50390977443609	1.49797166412325e-05\\
3.54147869674185	1.49807151178993e-05\\
3.57904761904762	1.49817212904877e-05\\
3.61661654135338	1.49827319894805e-05\\
3.65418546365915	1.49837440436428e-05\\
3.69175438596491	1.49847542800069e-05\\
3.72932330827068	1.49857595238725e-05\\
3.76689223057644	1.49867565988246e-05\\
3.80446115288221	1.49877423267571e-05\\
3.84203007518797	1.49887135279213e-05\\
3.87959899749373	1.49896670209753e-05\\
3.9171679197995	1.49905996949381e-05\\
3.95473684210526	1.49915317002899e-05\\
3.99230576441103	1.49924791478522e-05\\
4.02987468671679	1.49934398167685e-05\\
4.06744360902256	1.49944114851315e-05\\
4.10501253132832	1.49953919299436e-05\\
4.14258145363408	1.49963789270922e-05\\
4.18015037593985	1.49973702513241e-05\\
4.21771929824561	1.49983636762328e-05\\
4.25528822055138	1.4999356974247e-05\\
4.29285714285714	1.50003479166246e-05\\
4.33042606516291	1.5001334273465e-05\\
4.36799498746867	1.50023138137086e-05\\
4.40556390977444	1.50032843051627e-05\\
4.4431328320802	1.50042435145162e-05\\
4.48070175438596	1.50051892073757e-05\\
4.51827067669173	1.50061191482981e-05\\
4.55583959899749	1.50070324120596e-05\\
4.59340852130326	1.50079496901139e-05\\
4.63097744360902	1.50088777728713e-05\\
4.66854636591479	1.50098152891633e-05\\
4.70611528822055	1.50107608671681e-05\\
4.74368421052632	1.50117131343932e-05\\
4.78125313283208	1.5012670717662e-05\\
4.81882205513784	1.50136322431024e-05\\
4.85639097744361	1.50145963361391e-05\\
4.89395989974937	1.5015561621486e-05\\
4.93152882205514	1.5016526723143e-05\\
4.9690977443609	1.50174902643959e-05\\
5.00666666666667	1.50184508678152e-05\\
5.04423558897243	1.50194071552633e-05\\
5.0818045112782	1.50203577479e-05\\
5.11937343358396	1.50213012661904e-05\\
5.15694235588972	1.50222363299204e-05\\
5.19451127819549	1.5023161558205e-05\\
5.23208020050125	1.50240755695122e-05\\
5.26964912280702	1.50249769816781e-05\\
5.30721804511278	1.5025872792489e-05\\
5.34478696741855	1.50267783761595e-05\\
5.38235588972431	1.50276930836004e-05\\
5.41992481203008	1.50286158166653e-05\\
5.45749373433584	1.50295454767195e-05\\
5.4950626566416	1.50304809646241e-05\\
5.53263157894737	1.50314211807313e-05\\
5.57020050125313	1.5032365024877e-05\\
5.6077694235589	1.50333113963725e-05\\
5.64533834586466	1.50342591940028e-05\\
5.68290726817043	1.50352073160233e-05\\
5.72047619047619	1.50361546601583e-05\\
5.75804511278196	1.50371001236033e-05\\
5.79561403508772	1.5038042603022e-05\\
5.83318295739348	1.50389809945553e-05\\
5.87075187969925	1.50399141938214e-05\\
5.90832080200501	1.50408410959247e-05\\
5.94588972431078	1.50417605954633e-05\\
5.98345864661654	1.50426715865366e-05\\
6.02102756892231	1.50435729627577e-05\\
6.05859649122807	1.50444636172658e-05\\
6.09616541353383	1.50453424427369e-05\\
6.1337343358396	1.50462088020324e-05\\
6.17130325814536	1.50470766640005e-05\\
6.20887218045113	1.50479524158678e-05\\
6.24644110275689	1.5048835312476e-05\\
6.28401002506266	1.50497246083265e-05\\
6.32157894736842	1.50506195575831e-05\\
6.35914786967419	1.50515194140537e-05\\
6.39671679197995	1.50524234311948e-05\\
6.43428571428571	1.50533308621009e-05\\
6.47185463659148	1.50542409595015e-05\\
6.50942355889724	1.50551529757609e-05\\
6.54699248120301	1.50560661628697e-05\\
6.58456140350877	1.50569797724492e-05\\
6.62213032581454	1.5057893055748e-05\\
6.6596992481203	1.50588052636388e-05\\
6.69726817042607	1.50597156466257e-05\\
6.73483709273183	1.50606234548389e-05\\
6.77240601503759	1.50615279380416e-05\\
6.80997493734336	1.50624283456308e-05\\
6.84754385964912	1.50633239266403e-05\\
6.88511278195489	1.50642139297493e-05\\
6.92268170426065	1.50650976032808e-05\\
6.96025062656642	1.50659741952153e-05\\
6.99781954887218	1.50668429531923e-05\\
7.03538847117794	1.50677031245201e-05\\
7.07295739348371	1.50685539561856e-05\\
7.11052631578947	1.50693946948588e-05\\
7.14809523809524	1.50702330382244e-05\\
7.185664160401	1.50710783592667e-05\\
7.22323308270677	1.50719301418227e-05\\
7.26080200501253	1.50727878525407e-05\\
7.2983709273183	1.50736509578401e-05\\
7.33593984962406	1.50745189239042e-05\\
7.37350877192982	1.50753912166795e-05\\
7.41107769423559	1.50762673018722e-05\\
7.44864661654135	1.50771466449415e-05\\
7.48621553884712	1.50780287111016e-05\\
7.52378446115288	1.50789129653183e-05\\
7.56135338345865	1.50797988723059e-05\\
7.59892230576441	1.5080685896526e-05\\
7.63649122807018	1.50815735021877e-05\\
7.67406015037594	1.50824611532471e-05\\
7.7116290726817	1.50833483134041e-05\\
7.74919799498747	1.50842344461073e-05\\
7.78676691729323	1.50851190145501e-05\\
7.824335839599	1.50860014816747e-05\\
7.86190476190476	1.50868813101725e-05\\
7.89947368421053	1.50877579624836e-05\\
7.93704260651629	1.50886309008011e-05\\
7.97461152882205	1.50894995870755e-05\\
8.01218045112782	1.50903634830142e-05\\
8.04974937343358	1.5091222050085e-05\\
8.08731829573935	1.50920747495195e-05\\
8.12488721804511	1.50929210423209e-05\\
8.16245614035088	1.50937603892658e-05\\
8.20002506265664	1.50945922509061e-05\\
8.23759398496241	1.50954160875786e-05\\
8.27516290726817	1.50962331052825e-05\\
8.31273182957393	1.50970539360665e-05\\
8.3503007518797	1.50978801558405e-05\\
8.38786967418546	1.50987113847222e-05\\
8.42543859649123	1.50995472426732e-05\\
8.46300751879699	1.51003873494994e-05\\
8.50057644110276	1.5101231324846e-05\\
8.53814536340852	1.51020787881985e-05\\
8.57571428571429	1.51029293588775e-05\\
8.61328320802005	1.51037826560401e-05\\
8.65085213032581	1.51046382986754e-05\\
8.68842105263158	1.51054959056083e-05\\
8.72598997493734	1.51063550954916e-05\\
8.76355889724311	1.51072154868106e-05\\
8.80112781954887	1.51080766978805e-05\\
8.83869674185464	1.51089383468453e-05\\
8.8762656641604	1.51098000516769e-05\\
8.91383458646617	1.51106614301789e-05\\
8.95140350877193	1.51115220999812e-05\\
8.98897243107769	1.5112381678545e-05\\
9.02654135338346	1.51132397831584e-05\\
9.06411027568922	1.5114096030942e-05\\
9.10167919799499	1.51149500388459e-05\\
9.13924812030075	1.51158014236549e-05\\
9.17681704260652	1.5116649801983e-05\\
9.21438596491228	1.51174947902806e-05\\
9.25195488721804	1.51183360048335e-05\\
9.28952380952381	1.51191730617685e-05\\
9.32709273182957	1.5120005577047e-05\\
9.36466165413534	1.51208331664755e-05\\
9.4022305764411	1.51216554457037e-05\\
9.43979949874687	1.51224720302272e-05\\
9.47736842105263	1.51232825353928e-05\\
9.5149373433584	1.51240865763971e-05\\
9.55250626566416	1.51248837682951e-05\\
9.59007518796992	1.51256737976864e-05\\
9.62764411027569	1.51264635252189e-05\\
9.66521303258145	1.51272574834111e-05\\
9.70278195488722	1.51280554227745e-05\\
9.74035087719298	1.51288570937182e-05\\
9.77791979949875	1.51296622465501e-05\\
9.81548872180451	1.51304706314755e-05\\
9.85305764411028	1.51312819985962e-05\\
9.89062656641604	1.5132096097908e-05\\
9.9281954887218	1.51329126793004e-05\\
9.96576441102757	1.51337314925567e-05\\
10.0033333333333	1.51345522873517e-05\\
10.0409022556391	1.51353748132519e-05\\
10.0784711779449	1.51361988197148e-05\\
10.1160401002506	1.51370240560876e-05\\
10.1536090225564	1.51378502716076e-05\\
10.1911779448622	1.51386772154014e-05\\
10.2287468671679	1.51395046364849e-05\\
10.2663157894737	1.51403322837607e-05\\
10.3038847117794	1.51411599060248e-05\\
10.3414536340852	1.51419872519566e-05\\
10.379022556391	1.51428140701265e-05\\
10.4165914786967	1.51436401089945e-05\\
10.4541604010025	1.51444651169093e-05\\
10.4917293233083	1.51452888421059e-05\\
10.529298245614	1.51461110327136e-05\\
10.5668671679198	1.51469314367482e-05\\
10.6044360902256	1.51477498021172e-05\\
10.6420050125313	1.51485658766207e-05\\
10.6795739348371	1.5149379407949e-05\\
10.7171428571429	1.51501901436862e-05\\
10.7547117794486	1.51509978313096e-05\\
10.7922807017544	1.51518022181913e-05\\
10.8298496240602	1.51526030516009e-05\\
10.8674185463659	1.5153400078701e-05\\
10.9049874686717	1.51541930465565e-05\\
10.9425563909774	1.515498170213e-05\\
10.9801253132832	1.5155765792284e-05\\
11.017694235589	1.51565450637841e-05\\
11.0552631578947	1.5157319263302e-05\\
11.0928320802005	1.51580881374118e-05\\
11.1304010025063	1.51588514392671e-05\\
11.167969924812	1.5159613850917e-05\\
11.2055388471178	1.51603794394705e-05\\
11.2431077694236	1.51611480478167e-05\\
11.2806766917293	1.5161919518785e-05\\
11.3182456140351	1.51626936951443e-05\\
11.3558145363409	1.51634704196016e-05\\
11.3933834586466	1.51642495348049e-05\\
11.4309523809524	1.51650308833373e-05\\
11.4685213032581	1.51658143077212e-05\\
11.5060902255639	1.51665996504157e-05\\
11.5436591478697	1.51673867538182e-05\\
11.5812280701754	1.51681754602582e-05\\
11.6187969924812	1.51689656120056e-05\\
11.656365914787	1.51697570512635e-05\\
11.6939348370927	1.51705496201711e-05\\
11.7315037593985	1.51713431608019e-05\\
11.7690726817043	1.51721375151667e-05\\
11.80664160401	1.51729325252072e-05\\
11.8442105263158	1.51737280328027e-05\\
11.8817794486216	1.51745238797657e-05\\
11.9193483709273	1.51753199078429e-05\\
11.9569172932331	1.51761159587157e-05\\
11.9944862155388	1.51769118739999e-05\\
12.0320551378446	1.51777074952437e-05\\
12.0696240601504	1.51785026639324e-05\\
12.1071929824561	1.51792972214828e-05\\
12.1447619047619	1.51800910092492e-05\\
12.1823308270677	1.51808838685179e-05\\
12.2198997493734	1.51816756405105e-05\\
12.2574686716792	1.51824661663844e-05\\
12.295037593985	1.51832552872312e-05\\
12.3326065162907	1.51840428440778e-05\\
12.3701754385965	1.51848286778877e-05\\
12.4077443609023	1.51856126295608e-05\\
12.445313283208	1.51863945399315e-05\\
12.4828822055138	1.51871742497737e-05\\
12.5204511278195	1.51879515997953e-05\\
12.5580200501253	1.5188726430645e-05\\
12.5955889724311	1.51894985829068e-05\\
12.6331578947368	1.51902678971064e-05\\
12.6707268170426	1.5191034213706e-05\\
12.7082957393484	1.5191797373111e-05\\
12.7458646616541	1.51925572156627e-05\\
12.7834335839599	1.51933135816482e-05\\
12.8210025062657	1.51940663112925e-05\\
12.8585714285714	1.51948152447672e-05\\
12.8961403508772	1.51955602221837e-05\\
12.933709273183	1.5196301468776e-05\\
12.9712781954887	1.51970432678747e-05\\
13.0088471177945	1.519778678456e-05\\
13.0464160401002	1.519853193469e-05\\
13.083984962406	1.51992786340876e-05\\
13.1215538847118	1.52000267985464e-05\\
13.1591228070175	1.5200776343826e-05\\
13.1966917293233	1.52015271856526e-05\\
13.2342606516291	1.52022792397218e-05\\
13.2718295739348	1.52030324216935e-05\\
13.3093984962406	1.52037866471941e-05\\
13.3469674185464	1.52045418318203e-05\\
13.3845363408521	1.52052978911305e-05\\
13.4221052631579	1.52060547406532e-05\\
13.4596741854637	1.52068122958805e-05\\
13.4972431077694	1.52075704722724e-05\\
13.5348120300752	1.52083291852551e-05\\
13.572380952381	1.52090883502203e-05\\
13.6099498746867	1.52098478825246e-05\\
13.6475187969925	1.5210607697492e-05\\
13.6850877192982	1.52113677104127e-05\\
13.722656641604	1.52121278365405e-05\\
13.7602255639098	1.52128879910974e-05\\
13.7977944862155	1.52136480892702e-05\\
13.8353634085213	1.52144080462113e-05\\
13.8729323308271	1.52151677770399e-05\\
13.9105012531328	1.52159271968396e-05\\
13.9480701754386	1.52166862206603e-05\\
13.9856390977444	1.52174447635195e-05\\
14.0232080200501	1.5218202740397e-05\\
14.0607769423559	1.52189600662422e-05\\
14.0983458646617	1.5219716655968e-05\\
14.1359147869674	1.52204724244566e-05\\
14.1734837092732	1.52212272865513e-05\\
14.2110526315789	1.52219811570659e-05\\
14.2486215538847	1.522273395078e-05\\
14.2861904761905	1.52234855824392e-05\\
14.3237593984962	1.52242359667534e-05\\
14.361328320802	1.52249850184035e-05\\
14.3988972431078	1.52257326520343e-05\\
14.4364661654135	1.52264787822579e-05\\
14.4740350877193	1.5227223323656e-05\\
14.5116040100251	1.52279661907737e-05\\
14.5491729323308	1.52287072981276e-05\\
14.5867418546366	1.52294465601988e-05\\
14.6243107769424	1.52301838914387e-05\\
14.6618796992481	1.52309192062655e-05\\
14.6994486215539	1.52316524190659e-05\\
14.7370175438596	1.5232383444197e-05\\
14.7745864661654	1.52331121959807e-05\\
14.8121553884712	1.52338385887123e-05\\
14.8497243107769	1.52345625366547e-05\\
14.8872932330827	1.52352839540393e-05\\
14.9248621553885	1.52360027550692e-05\\
14.9624310776942	1.52367188539161e-05\\
15	1.52374321647237e-05\\
};
\addplot [color=mycolor1, dashdotted, line width=0.9pt, forget plot]
  table[row sep=crcr]{%
0.01	0.186978297161937\\
0.0475689223057644	0.186978297161937\\
0.0851378446115288	0.186978297161937\\
0.122706766917293	0.186978297161937\\
0.160275689223058	0.186978297161937\\
0.197844611528822	0.186978297161937\\
0.235413533834586	0.183639398998331\\
0.272982456140351	0.183639398998331\\
0.310551378446115	0.183639398998331\\
0.34812030075188	0.183639398998331\\
0.385689223057644	0.183639398998331\\
0.423258145363409	0.183639398998331\\
0.460827067669173	0.183639398998331\\
0.498395989974937	0.180300500834725\\
0.535964912280702	0.180300500834725\\
0.573533834586466	0.180300500834725\\
0.611102756892231	0.180300500834725\\
0.648671679197995	0.180300500834725\\
0.686240601503759	0.180300500834725\\
0.723809523809524	0.180300500834725\\
0.761378446115288	0.180300500834725\\
0.798947368421053	0.180300500834725\\
0.836516290726817	0.180300500834725\\
0.874085213032581	0.180300500834725\\
0.911654135338346	0.180300500834725\\
0.94922305764411	0.180300500834725\\
0.986791979949875	0.176961602671119\\
1.02436090225564	0.176961602671119\\
1.0619298245614	0.176961602671119\\
1.09949874686717	0.176961602671119\\
1.13706766917293	0.176961602671119\\
1.1746365914787	0.176961602671119\\
1.21220551378446	0.176961602671119\\
1.24977443609023	0.176961602671119\\
1.28734335839599	0.176961602671119\\
1.32491228070175	0.176961602671119\\
1.36248120300752	0.176961602671119\\
1.40005012531328	0.176961602671119\\
1.43761904761905	0.176961602671119\\
1.47518796992481	0.176961602671119\\
1.51275689223058	0.176961602671119\\
1.55032581453634	0.176961602671119\\
1.58789473684211	0.176961602671119\\
1.62546365914787	0.176961602671119\\
1.66303258145363	0.176961602671119\\
1.7006015037594	0.176961602671119\\
1.73817042606516	0.176961602671119\\
1.77573934837093	0.176961602671119\\
1.81330827067669	0.176961602671119\\
1.85087719298246	0.176961602671119\\
1.88844611528822	0.176961602671119\\
1.92601503759399	0.173622704507513\\
1.96358395989975	0.173622704507513\\
2.00115288220551	0.173622704507513\\
2.03872180451128	0.173622704507513\\
2.07629072681704	0.173622704507513\\
2.11385964912281	0.173622704507513\\
2.15142857142857	0.173622704507513\\
2.18899749373434	0.173622704507513\\
2.2265664160401	0.173622704507513\\
2.26413533834586	0.173622704507513\\
2.30170426065163	0.173622704507513\\
2.33927318295739	0.173622704507513\\
2.37684210526316	0.173622704507513\\
2.41441102756892	0.173622704507513\\
2.45197994987469	0.173622704507513\\
2.48954887218045	0.173622704507513\\
2.52711779448622	0.173622704507513\\
2.56468671679198	0.173622704507513\\
2.60225563909774	0.173622704507513\\
2.63982456140351	0.173622704507513\\
2.67739348370927	0.173622704507513\\
2.71496240601504	0.173622704507513\\
2.7525313283208	0.173622704507513\\
2.79010025062657	0.173622704507513\\
2.82766917293233	0.170283806343907\\
2.86523809523809	0.170283806343907\\
2.90280701754386	0.170283806343907\\
2.94037593984962	0.170283806343907\\
2.97794486215539	0.170283806343907\\
3.01551378446115	0.170283806343907\\
3.05308270676692	0.170283806343907\\
3.09065162907268	0.170283806343907\\
3.12822055137845	0.170283806343907\\
3.16578947368421	0.170283806343907\\
3.20335839598997	0.170283806343907\\
3.24092731829574	0.170283806343907\\
3.2784962406015	0.170283806343907\\
3.31606516290727	0.170283806343907\\
3.35363408521303	0.170283806343907\\
3.3912030075188	0.170283806343907\\
3.42877192982456	0.170283806343907\\
3.46634085213033	0.170283806343907\\
3.50390977443609	0.170283806343907\\
3.54147869674185	0.170283806343907\\
3.57904761904762	0.1669449081803\\
3.61661654135338	0.1669449081803\\
3.65418546365915	0.1669449081803\\
3.69175438596491	0.1669449081803\\
3.72932330827068	0.1669449081803\\
3.76689223057644	0.1669449081803\\
3.80446115288221	0.1669449081803\\
3.84203007518797	0.1669449081803\\
3.87959899749373	0.1669449081803\\
3.9171679197995	0.1669449081803\\
3.95473684210526	0.1669449081803\\
3.99230576441103	0.1669449081803\\
4.02987468671679	0.1669449081803\\
4.06744360902256	0.1669449081803\\
4.10501253132832	0.1669449081803\\
4.14258145363408	0.1669449081803\\
4.18015037593985	0.1669449081803\\
4.21771929824561	0.163606010016695\\
4.25528822055138	0.163606010016695\\
4.29285714285714	0.163606010016695\\
4.33042606516291	0.163606010016695\\
4.36799498746867	0.163606010016695\\
4.40556390977444	0.163606010016695\\
4.4431328320802	0.163606010016695\\
4.48070175438596	0.163606010016695\\
4.51827067669173	0.163606010016695\\
4.55583959899749	0.163606010016695\\
4.59340852130326	0.163606010016695\\
4.63097744360902	0.163606010016695\\
4.66854636591479	0.163606010016695\\
4.70611528822055	0.163606010016695\\
4.74368421052632	0.163606010016695\\
4.78125313283208	0.163606010016695\\
4.81882205513784	0.163606010016695\\
4.85639097744361	0.160267111853088\\
4.89395989974937	0.160267111853088\\
4.93152882205514	0.160267111853088\\
4.9690977443609	0.160267111853088\\
5.00666666666667	0.160267111853088\\
5.04423558897243	0.160267111853088\\
5.0818045112782	0.160267111853088\\
5.11937343358396	0.160267111853088\\
5.15694235588972	0.160267111853088\\
5.19451127819549	0.160267111853088\\
5.23208020050125	0.160267111853088\\
5.26964912280702	0.160267111853088\\
5.30721804511278	0.160267111853088\\
5.34478696741855	0.160267111853088\\
5.38235588972431	0.160267111853088\\
5.41992481203008	0.160267111853088\\
5.45749373433584	0.160267111853088\\
5.4950626566416	0.160267111853088\\
5.53263157894737	0.160267111853088\\
5.57020050125313	0.156928213689482\\
5.6077694235589	0.156928213689482\\
5.64533834586466	0.156928213689482\\
5.68290726817043	0.156928213689482\\
5.72047619047619	0.156928213689482\\
5.75804511278196	0.156928213689482\\
5.79561403508772	0.156928213689482\\
5.83318295739348	0.156928213689482\\
5.87075187969925	0.156928213689482\\
5.90832080200501	0.156928213689482\\
5.94588972431078	0.156928213689482\\
5.98345864661654	0.156928213689482\\
6.02102756892231	0.156928213689482\\
6.05859649122807	0.156928213689482\\
6.09616541353383	0.156928213689482\\
6.1337343358396	0.156928213689482\\
6.17130325814536	0.156928213689482\\
6.20887218045113	0.156928213689482\\
6.24644110275689	0.156928213689482\\
6.28401002506266	0.156928213689482\\
6.32157894736842	0.156928213689482\\
6.35914786967419	0.156928213689482\\
6.39671679197995	0.156928213689482\\
6.43428571428571	0.156928213689482\\
6.47185463659148	0.156928213689482\\
6.50942355889724	0.156928213689482\\
6.54699248120301	0.153589315525876\\
6.58456140350877	0.153589315525876\\
6.62213032581454	0.153589315525876\\
6.6596992481203	0.153589315525876\\
6.69726817042607	0.153589315525876\\
6.73483709273183	0.153589315525876\\
6.77240601503759	0.153589315525876\\
6.80997493734336	0.153589315525876\\
6.84754385964912	0.153589315525876\\
6.88511278195489	0.153589315525876\\
6.92268170426065	0.153589315525876\\
6.96025062656642	0.153589315525876\\
6.99781954887218	0.153589315525876\\
7.03538847117794	0.153589315525876\\
7.07295739348371	0.153589315525876\\
7.11052631578947	0.153589315525876\\
7.14809523809524	0.153589315525876\\
7.185664160401	0.153589315525876\\
7.22323308270677	0.153589315525876\\
7.26080200501253	0.153589315525876\\
7.2983709273183	0.153589315525876\\
7.33593984962406	0.153589315525876\\
7.37350877192982	0.153589315525876\\
7.41107769423559	0.153589315525876\\
7.44864661654135	0.153589315525876\\
7.48621553884712	0.153589315525876\\
7.52378446115288	0.153589315525876\\
7.56135338345865	0.153589315525876\\
7.59892230576441	0.153589315525876\\
7.63649122807018	0.153589315525876\\
7.67406015037594	0.153589315525876\\
7.7116290726817	0.153589315525876\\
7.74919799498747	0.153589315525876\\
7.78676691729323	0.153589315525876\\
7.824335839599	0.153589315525876\\
7.86190476190476	0.153589315525876\\
7.89947368421053	0.153589315525876\\
7.93704260651629	0.153589315525876\\
7.97461152882205	0.153589315525876\\
8.01218045112782	0.153589315525876\\
8.04974937343358	0.153589315525876\\
8.08731829573935	0.153589315525876\\
8.12488721804511	0.153589315525876\\
8.16245614035088	0.153589315525876\\
8.20002506265664	0.153589315525876\\
8.23759398496241	0.153589315525876\\
8.27516290726817	0.153589315525876\\
8.31273182957393	0.15025041736227\\
8.3503007518797	0.15025041736227\\
8.38786967418546	0.15025041736227\\
8.42543859649123	0.15025041736227\\
8.46300751879699	0.15025041736227\\
8.50057644110276	0.15025041736227\\
8.53814536340852	0.15025041736227\\
8.57571428571429	0.15025041736227\\
8.61328320802005	0.15025041736227\\
8.65085213032581	0.15025041736227\\
8.68842105263158	0.15025041736227\\
8.72598997493734	0.15025041736227\\
8.76355889724311	0.15025041736227\\
8.80112781954887	0.15025041736227\\
8.83869674185464	0.15025041736227\\
8.8762656641604	0.15025041736227\\
8.91383458646617	0.15025041736227\\
8.95140350877193	0.15025041736227\\
8.98897243107769	0.15025041736227\\
9.02654135338346	0.15025041736227\\
9.06411027568922	0.15025041736227\\
9.10167919799499	0.15025041736227\\
9.13924812030075	0.15025041736227\\
9.17681704260652	0.15025041736227\\
9.21438596491228	0.15025041736227\\
9.25195488721804	0.15025041736227\\
9.28952380952381	0.15025041736227\\
9.32709273182957	0.15025041736227\\
9.36466165413534	0.15025041736227\\
9.4022305764411	0.15025041736227\\
9.43979949874687	0.15025041736227\\
9.47736842105263	0.15025041736227\\
9.5149373433584	0.15025041736227\\
9.55250626566416	0.15025041736227\\
9.59007518796992	0.15025041736227\\
9.62764411027569	0.15025041736227\\
9.66521303258145	0.15025041736227\\
9.70278195488722	0.15025041736227\\
9.74035087719298	0.15025041736227\\
9.77791979949875	0.15025041736227\\
9.81548872180451	0.15025041736227\\
9.85305764411028	0.15025041736227\\
9.89062656641604	0.15025041736227\\
9.9281954887218	0.15025041736227\\
9.96576441102757	0.15025041736227\\
10.0033333333333	0.15025041736227\\
10.0409022556391	0.15025041736227\\
10.0784711779449	0.15025041736227\\
10.1160401002506	0.15025041736227\\
10.1536090225564	0.15025041736227\\
10.1911779448622	0.15025041736227\\
10.2287468671679	0.15025041736227\\
10.2663157894737	0.15025041736227\\
10.3038847117794	0.15025041736227\\
10.3414536340852	0.15025041736227\\
10.379022556391	0.15025041736227\\
10.4165914786967	0.15025041736227\\
10.4541604010025	0.15025041736227\\
10.4917293233083	0.15025041736227\\
10.529298245614	0.15025041736227\\
10.5668671679198	0.15025041736227\\
10.6044360902256	0.15025041736227\\
10.6420050125313	0.15025041736227\\
10.6795739348371	0.15025041736227\\
10.7171428571429	0.15025041736227\\
10.7547117794486	0.15025041736227\\
10.7922807017544	0.15025041736227\\
10.8298496240602	0.15025041736227\\
10.8674185463659	0.15025041736227\\
10.9049874686717	0.15025041736227\\
10.9425563909774	0.15025041736227\\
10.9801253132832	0.15025041736227\\
11.017694235589	0.15025041736227\\
11.0552631578947	0.15025041736227\\
11.0928320802005	0.15025041736227\\
11.1304010025063	0.15025041736227\\
11.167969924812	0.15025041736227\\
11.2055388471178	0.15025041736227\\
11.2431077694236	0.15025041736227\\
11.2806766917293	0.15025041736227\\
11.3182456140351	0.15025041736227\\
11.3558145363409	0.15025041736227\\
11.3933834586466	0.15025041736227\\
11.4309523809524	0.146911519198664\\
11.4685213032581	0.146911519198664\\
11.5060902255639	0.146911519198664\\
11.5436591478697	0.146911519198664\\
11.5812280701754	0.146911519198664\\
11.6187969924812	0.146911519198664\\
11.656365914787	0.146911519198664\\
11.6939348370927	0.146911519198664\\
11.7315037593985	0.146911519198664\\
11.7690726817043	0.146911519198664\\
11.80664160401	0.146911519198664\\
11.8442105263158	0.146911519198664\\
11.8817794486216	0.146911519198664\\
11.9193483709273	0.146911519198664\\
11.9569172932331	0.146911519198664\\
11.9944862155388	0.146911519198664\\
12.0320551378446	0.146911519198664\\
12.0696240601504	0.146911519198664\\
12.1071929824561	0.146911519198664\\
12.1447619047619	0.146911519198664\\
12.1823308270677	0.146911519198664\\
12.2198997493734	0.146911519198664\\
12.2574686716792	0.146911519198664\\
12.295037593985	0.146911519198664\\
12.3326065162907	0.146911519198664\\
12.3701754385965	0.146911519198664\\
12.4077443609023	0.146911519198664\\
12.445313283208	0.146911519198664\\
12.4828822055138	0.146911519198664\\
12.5204511278195	0.146911519198664\\
12.5580200501253	0.146911519198664\\
12.5955889724311	0.146911519198664\\
12.6331578947368	0.146911519198664\\
12.6707268170426	0.146911519198664\\
12.7082957393484	0.146911519198664\\
12.7458646616541	0.146911519198664\\
12.7834335839599	0.146911519198664\\
12.8210025062657	0.146911519198664\\
12.8585714285714	0.146911519198664\\
12.8961403508772	0.146911519198664\\
12.933709273183	0.146911519198664\\
12.9712781954887	0.146911519198664\\
13.0088471177945	0.146911519198664\\
13.0464160401002	0.146911519198664\\
13.083984962406	0.146911519198664\\
13.1215538847118	0.146911519198664\\
13.1591228070175	0.146911519198664\\
13.1966917293233	0.146911519198664\\
13.2342606516291	0.146911519198664\\
13.2718295739348	0.146911519198664\\
13.3093984962406	0.146911519198664\\
13.3469674185464	0.146911519198664\\
13.3845363408521	0.146911519198664\\
13.4221052631579	0.146911519198664\\
13.4596741854637	0.146911519198664\\
13.4972431077694	0.146911519198664\\
13.5348120300752	0.146911519198664\\
13.572380952381	0.146911519198664\\
13.6099498746867	0.146911519198664\\
13.6475187969925	0.146911519198664\\
13.6850877192982	0.146911519198664\\
13.722656641604	0.146911519198664\\
13.7602255639098	0.146911519198664\\
13.7977944862155	0.146911519198664\\
13.8353634085213	0.146911519198664\\
13.8729323308271	0.146911519198664\\
13.9105012531328	0.146911519198664\\
13.9480701754386	0.146911519198664\\
13.9856390977444	0.146911519198664\\
14.0232080200501	0.143572621035058\\
14.0607769423559	0.143572621035058\\
14.0983458646617	0.143572621035058\\
14.1359147869674	0.143572621035058\\
14.1734837092732	0.143572621035058\\
14.2110526315789	0.143572621035058\\
14.2486215538847	0.143572621035058\\
14.2861904761905	0.143572621035058\\
14.3237593984962	0.143572621035058\\
14.361328320802	0.143572621035058\\
14.3988972431078	0.143572621035058\\
14.4364661654135	0.143572621035058\\
14.4740350877193	0.143572621035058\\
14.5116040100251	0.143572621035058\\
14.5491729323308	0.143572621035058\\
14.5867418546366	0.143572621035058\\
14.6243107769424	0.143572621035058\\
14.6618796992481	0.143572621035058\\
14.6994486215539	0.143572621035058\\
14.7370175438596	0.143572621035058\\
14.7745864661654	0.143572621035058\\
14.8121553884712	0.143572621035058\\
14.8497243107769	0.143572621035058\\
14.8872932330827	0.143572621035058\\
14.9248621553885	0.143572621035058\\
14.9624310776942	0.143572621035058\\
15	0.143572621035058\\
};
\node[right, align=left, inner sep=0]
at (rel axis cs:-0.15,1.15) {$(b)$};
\end{axis}
\end{tikzpicture}%
  \end{minipage}

  \medskip
  \begin{tikzpicture}
\begin{axis}[
hide axis, scale only axis,
width=0.75\textwidth, height=0pt,
colormap={sigmamap}{rgb(0.0000)=(0.7760,0.8510,0.9140); rgb(0.1000)=(0.8392,0.8918,0.9369); rgb(0.2000)=(0.8913,0.9259,0.9563); rgb(0.3000)=(0.9374,0.9569,0.9743); rgb(0.4000)=(0.9799,0.9861,0.9917); rgb(0.5000)=(1.0000,0.9999,0.9999); rgb(0.6000)=(0.9659,0.8254,0.7818); rgb(0.7000)=(0.8840,0.4951,0.4053); rgb(0.8000)=(0.7820,0.2849,0.2466); rgb(0.9000)=(0.6131,0.1208,0.1566); rgb(1.0000)=(0.4040,0.0000,0.1220)},
point meta min=-0.5, point meta max=0.5,
colorbar horizontal,
colorbar style={
  width=0.75\textwidth, height=7pt,
  at={(0,0)}, anchor=north west,
  xlabel={$\sigma$},
  xlabel style={font=\footnotesize},
  tick label style={font=\footnotesize},
  xtick pos=bottom,
},
]
\addplot[draw=none] coordinates {(0,0)};
\end{axis}
\end{tikzpicture}

  \caption{Per-wavenumber growth-rate fields in the Boussinesq limit for \(\Sch=1\) and
  \(\delta_0=2\). (a)~Non-autonomous mean growth rate \(\overline{\sigma}_{1D}(k,t)\), with the
  early-time window \(t\leq2\) shown in the inset. (b)~QSSA growth rate
  \(\sigma_{\mathrm{QSSA}}(k,t_0)\).}
  \label{fig:BQ_P3_sigma_map}
\end{figure}






\subsection{Three-dimensional perturbation growth}
From the family of spectra introduced in~\cref{sec:statistics_of_the_initial_perturbations}, we fix
\(n=1\), the simplest member, which corresponds to a Gaussian spatial correlation. The
higher-order spectra produce the same qualitative behaviour and are therefore omitted here so that
we can isolate the effect of the horizontal correlation length. Keeping the initial variance
\(\mean{\eta_0^2}\) fixed, we vary the integral correlation length \(\lambda\) of the \(n=1\)
spectrum defined in~\eqref{eq:eta_spectrum_family}. For this spectrum,
\(\lambda_0=2\lambda/\sqrt{\pi}\), where \(\lambda_0\) is the Gaussian reference scale defined in
\eqref{eq:lambda0_k0_relation}. The mean energy growth defined in~\eqref{eq:gmean_eta_spectrum}
therefore reduces, after switching to polar coordinates \((k,\theta)\) and integrating by parts, to
\begin{equation}
    \label{eq:bq_gmean_isotropic}
    \Gm(t)
    =
    g(0,t)
    +
    \int_{0}^{\infty}
    \frac{\partial g(k,t)}{\partial k}
    \exp\!\left(-\frac{\lambda^{2}k^{2}}{\pi}\right)\mathrm{d}k.
\end{equation}

\Cref{fig:compare_dns}(b) shows a representative evolution of \(\Gm(t)\) for a
fixed correlation length. The mean energy initially decays, reflecting the dominance of stable
contributions at early times. This decay continues until the minimum \(\Gm_{\min}(\lambda)\) is
reached. Beyond this point, the trajectory reverses and exhibits growth driven by the unstable
modes represented in the spectrum. A key feature of this behaviour is the delay before the mean
energy returns to its initial level. In terms of the quantities defined in
\eqref{eq:minimum_growth_critical_time_definitions}, \(\Gm_{\min}\) measures the maximum transient
energy loss, whereas \(t_c\) measures how long the perturbation takes to recover its initial mean
energy.

\begin{figure}
  \centering
  \begin{minipage}{0.48\textwidth}
    \centering
    % This file was created by matlab2tikz.
%
\definecolor{mycolor1}{rgb}{0.75000,0.08000,0.08000}%
\definecolor{mycolor2}{rgb}{0.83368,0.83368,0.83368}%
\definecolor{mycolor3}{rgb}{0.78737,0.78737,0.78737}%
\definecolor{mycolor4}{rgb}{0.74105,0.74105,0.74105}%
\definecolor{mycolor5}{rgb}{0.69474,0.69474,0.69474}%
\definecolor{mycolor6}{rgb}{0.64842,0.64842,0.64842}%
\definecolor{mycolor7}{rgb}{0.60211,0.60211,0.60211}%
\definecolor{mycolor8}{rgb}{0.50947,0.50947,0.50947}%
\definecolor{mycolor9}{rgb}{0.46316,0.46316,0.46316}%
%
\begin{tikzpicture}

\begin{axis}[%
width=11.252in,
height=7.131in,
at={(1.887in,0.962in)},
scale only axis,
clip=true,
separate axis lines,
every outer x axis line/.append style={black},
every x tick label/.append style={font=\color{black}},
every x tick/.append style={black},
xmin=0.005,
xmax=8,
xlabel={$t$},
every outer y axis line/.append style={black},
every y tick label/.append style={font=\color{black}},
every y tick/.append style={black},
ymode=log,
ymin=0.1,
ymax=100,
yminorticks=true,
ylabel={$\mathscr{G}^{\mathrm{mean}}(t)$},
axis background/.style={fill=white},
clip=true,
clip mode=individual,
axis lines=box,
xtick pos=bottom,
ytick pos=left,
tick align=inside,
major tick length=2pt,
width=0.8\linewidth,
height=0.8\linewidth,
every axis/.append style={font=\fontsize{9}{9}\selectfont},xlabel style={font=\fontsize{9}{9}\selectfont},title style={font=\fontsize{9}{9}\selectfont},
legend style={font=\fontsize{8}{9}\selectfont,inner sep=1pt,row sep=1pt,column sep=2pt,nodes={scale=1.00},draw=none},legend image post style={xscale=0.35},
legend columns=1,
scaled x ticks=false,
tick label style={/pgf/number format/fixed,/pgf/number format/precision=2},
every x tick label/.append style={font=\fontsize{9}{9}\selectfont\color{black}},
every y tick label/.append style={font=\fontsize{9}{9}\selectfont\color{black}},
,,
colormap={sigmamap}{rgb(0.0000)=(0.0200,0.1880,0.3800); rgb(0.1000)=(0.0743,0.3456,0.5616); rgb(0.2000)=(0.1918,0.4980,0.7060); rgb(0.3000)=(0.3890,0.6708,0.8226); rgb(0.4000)=(0.7552,0.8682,0.9228); rgb(0.5000)=(0.9690,0.9689,0.9689); rgb(0.6000)=(0.9425,0.8044,0.7614); rgb(0.7000)=(0.8767,0.4910,0.4020); rgb(0.8000)=(0.7869,0.2866,0.2470); rgb(0.9000)=(0.6186,0.1239,0.1568); rgb(1.0000)=(0.4040,0.0000,0.1220)}
]
\addplot [color=white!55!black, dashed, line width=0.7pt, forget plot]
  table[row sep=crcr]{%
0.005	1\\
8	1\\
};
\addplot [color=white!88!black, line width=0.9pt, forget plot]
  table[row sep=crcr]{%
0.005	0.959051936801755\\
0.03	0.786622229327457\\
0.06	0.634528217945405\\
0.085	0.540096388507333\\
0.11	0.46667404537218\\
0.14	0.398746788498769\\
0.165	0.354524829253815\\
0.19	0.318695676138571\\
0.22	0.284098969809706\\
0.245	0.260638865906569\\
0.27	0.240970544521686\\
0.3	0.221314853434199\\
0.325	0.207557168028668\\
0.355	0.193552959443092\\
0.38	0.183586450015501\\
0.405	0.174896039567011\\
0.435	0.165870786594321\\
0.46	0.15933198135657\\
0.485	0.153549055135064\\
0.515	0.147461587181857\\
0.54	0.142999076105584\\
0.565	0.139016539204789\\
0.595	0.134789359848328\\
0.62	0.131669437223373\\
0.645	0.12887173268871\\
0.675	0.125890933092623\\
0.7	0.123685799654635\\
0.725	0.121706957760053\\
0.755	0.119600220215998\\
0.78	0.118045525724334\\
0.805	0.116655807448684\\
0.835	0.11518561448178\\
0.86	0.114110143165207\\
0.885	0.113158833833264\\
0.915	0.112167314561309\\
0.94	0.111455851055121\\
0.97	0.11072829471931\\
0.995	0.110219111163394\\
1.02	0.109791658689233\\
1.05	0.109378876680707\\
1.075	0.109112568650033\\
1.1	0.10891213100725\\
1.13	0.108752924365368\\
1.155	0.108683770989646\\
1.18	0.108668862142593\\
1.21	0.108718398088749\\
1.235	0.108812701218992\\
1.26	0.108952586010286\\
1.29	0.109177473426788\\
1.315	0.109410009553778\\
1.34	0.109681585253684\\
1.37	0.110056617017258\\
1.395	0.110408268474274\\
1.42	0.110793966202364\\
1.45	0.111299908584993\\
1.475	0.111756043151122\\
1.5	0.112242382286623\\
1.53	0.112864441011134\\
1.555	0.113413779524707\\
1.585	0.114108863521928\\
1.61	0.11471708588891\\
1.635	0.11535089579696\\
1.665	0.116144320338243\\
1.69	0.116832177355282\\
1.715	0.117543688643448\\
1.745	0.1184280118542\\
1.77	0.119189826327699\\
1.795	0.119973802103961\\
1.825	0.120943277414146\\
1.85	0.121774675606844\\
1.875	0.122627092107432\\
1.905	0.123677320132539\\
1.93	0.124574968490897\\
1.955	0.125492773722522\\
1.985	0.126620438678425\\
2.01	0.127581843743809\\
2.035	0.128562773448434\\
2.065	0.129765441470972\\
2.09	0.130788796268215\\
2.115	0.131831232222962\\
2.145	0.133107195599691\\
2.17	0.134191260638556\\
2.2	0.135516942366521\\
2.225	0.136642276446065\\
2.25	0.137786286793492\\
2.28	0.139183703282516\\
2.305	0.140368693422417\\
2.33	0.141572285620825\\
2.36	0.143041149397919\\
2.385	0.144285673838851\\
2.41	0.145548828080824\\
2.44	0.147089243678282\\
2.465	0.148393491592169\\
2.49	0.149716485732398\\
2.52	0.151328898624726\\
2.545	0.152693330731551\\
2.57	0.154076703563642\\
2.6	0.155761859345211\\
2.625	0.157187176396131\\
2.65	0.158631698241253\\
2.68	0.160390609908167\\
2.705	0.161877727410953\\
2.73	0.163384376522708\\
2.76	0.165218298218467\\
2.785	0.16676832615477\\
2.815	0.168654659839379\\
2.84	0.170248675742009\\
2.865	0.17186290230651\\
2.895	0.17382685254802\\
2.92	0.175486049222594\\
2.945	0.17716592782631\\
2.975	0.17920930384315\\
3	0.180935242154863\\
3.025	0.1826823810971\\
3.055	0.184807174351432\\
3.08	0.186601564252003\\
3.105	0.188417718589389\\
3.135	0.190626093826845\\
3.16	0.192490787511916\\
3.185	0.194377852749989\\
3.215	0.196672141093591\\
3.24	0.198609127769553\\
3.265	0.200569135077012\\
3.295	0.202951828842209\\
3.32	0.204963230844535\\
3.345	0.206998343549422\\
3.375	0.209472092511534\\
3.4	0.211560162567233\\
3.43	0.214098129649751\\
3.455	0.216240282638421\\
3.48	0.218407402104518\\
3.51	0.221041246486127\\
3.535	0.223264160173831\\
3.56	0.225512838115331\\
3.59	0.228245623259352\\
3.615	0.230551894358107\\
3.64	0.232884767354722\\
3.67	0.235719709845447\\
3.695	0.238112062823826\\
3.72	0.240531895414647\\
3.75	0.243472365732355\\
3.775	0.245953653655037\\
3.8	0.248463339347015\\
3.83	0.251512863293342\\
3.855	0.254086069181481\\
3.88	0.256688631961402\\
3.91	0.259850892666645\\
3.935	0.26251913132633\\
3.96	0.265217727626221\\
3.99	0.268496568083427\\
4.015	0.271263088359085\\
4.045	0.27462439113607\\
4.07	0.277460435515063\\
4.095	0.280328622932599\\
4.125	0.283813361630922\\
4.15	0.286753481533873\\
4.175	0.289726861304626\\
4.205	0.293339324711937\\
4.23	0.296387143302755\\
4.255	0.299469383665188\\
4.285	0.303214033670304\\
4.31	0.306373319623556\\
4.335	0.30956823550411\\
4.365	0.313449711617198\\
4.39	0.316724382989394\\
4.415	0.320035939967634\\
4.445	0.324059064200426\\
4.47	0.3274531925976\\
4.495	0.330885511163788\\
4.525	0.335055293257219\\
4.55	0.338573108288825\\
4.575	0.342130468386131\\
4.605	0.346452111375905\\
4.63	0.350098005387953\\
4.66	0.354527174172372\\
4.685	0.358263757344955\\
4.71	0.362042290428044\\
4.74	0.366632555134397\\
4.765	0.370505012477596\\
4.79	0.374420917746101\\
4.82	0.379178031017581\\
4.845	0.383191217802176\\
4.87	0.387249406943351\\
4.9	0.392179338075172\\
4.925	0.396338292116706\\
4.95	0.400543861209951\\
4.98	0.405652803243757\\
5.005	0.409962750907482\\
5.03	0.414320986514539\\
5.06	0.419615363659417\\
5.085	0.424081726138596\\
5.11	0.42859811166738\\
5.14	0.434084587049644\\
5.165	0.43871298691908\\
5.19	0.4433932092641\\
5.22	0.449078693011665\\
5.245	0.453874961060766\\
5.275	0.459701403853549\\
5.3	0.464616575190053\\
5.325	0.469586757540694\\
5.355	0.475624453277258\\
5.38	0.480717824371351\\
5.405	0.485868189885686\\
5.435	0.492124755214765\\
5.46	0.497402752513177\\
5.485	0.502739800497738\\
5.515	0.509223131863334\\
5.54	0.514692417729638\\
5.565	0.520222885942843\\
5.595	0.526941169321147\\
5.62	0.532608650255913\\
5.645	0.538339523227016\\
5.675	0.54530124423923\\
5.7	0.551174079417599\\
5.725	0.557112597084636\\
5.755	0.564326551496869\\
5.78	0.570412161631587\\
5.805	0.576565828307862\\
5.835	0.58404113292925\\
5.86	0.59034720946807\\
5.89	0.598007653054724\\
5.915	0.604469907136169\\
5.94	0.611004421804513\\
5.97	0.618942359922614\\
5.995	0.62563870040593\\
6.02	0.632409915103689\\
6.05	0.640635384574641\\
6.075	0.647574278786116\\
6.1	0.654590755851724\\
6.13	0.663114158361287\\
6.155	0.670304381339735\\
6.18	0.677574994190676\\
6.21	0.686407109217126\\
6.235	0.693857754621567\\
6.26	0.701391698785071\\
6.29	0.71054369699996\\
6.315	0.718264188415316\\
6.34	0.726070992961443\\
6.37	0.735554450119027\\
6.395	0.743554552773223\\
6.42	0.751644092162982\\
6.45	0.76147100349314\\
6.475	0.769760836398893\\
6.505	0.779831055753588\\
6.53	0.788326137884266\\
6.555	0.79691618642002\\
6.585	0.807351094918034\\
6.61	0.816153820651702\\
6.635	0.825054949596217\\
6.665	0.835867745217041\\
6.69	0.844989247974341\\
6.715	0.854212715010149\\
6.745	0.865417072922753\\
6.77	0.874868888591254\\
6.795	0.884426358292284\\
6.825	0.896036447844801\\
6.85	0.905830529111434\\
6.875	0.915734087433268\\
6.905	0.927764589752483\\
6.93	0.937913320945799\\
6.955	0.948175490233948\\
6.985	0.960641616452918\\
7.01	0.971157828919422\\
7.035	0.98179158345199\\
7.065	0.994709093590183\\
7.09	1.00560608161706\\
7.12	1.01884335481844\\
7.145	1.03001008621964\\
7.17	1.04130162352544\\
7.2	1.0550181751501\\
7.225	1.06658921222004\\
7.25	1.07828956903364\\
7.28	1.09250273306911\\
7.305	1.10449269760837\\
7.33	1.11661665835403\\
7.36	1.13134439329173\\
7.385	1.14376843381891\\
7.41	1.15633131543033\\
7.44	1.1715922265079\\
7.465	1.18446603701819\\
7.49	1.19748370790365\\
7.52	1.21329707013792\\
7.545	1.22663690952747\\
7.57	1.24012580920638\\
7.6	1.25651159124044\\
7.625	1.27033430341967\\
7.65	1.28431146285627\\
7.68	1.30129035165412\\
7.705	1.31561338637362\\
7.735	1.3330124216263\\
7.76	1.34768987544438\\
7.785	1.36253130985754\\
7.815	1.38056006167157\\
7.84	1.39576871873975\\
7.865	1.41114727972953\\
7.895	1.42982849478567\\
7.92	1.4455875420354\\
7.945	1.46152262932997\\
7.975	1.48087987148436\\
8	1.49720918505694\\
};
\addplot [color=mycolor1, line width=1.4pt, only marks, mark size=0.4pt, mark=*, mark options={solid, fill=mycolor1, mycolor1}, forget plot]
  table[row sep=crcr]{%
1.175	0.108667651227639\\
};
\addplot [color=mycolor2, line width=0.9pt, forget plot]
  table[row sep=crcr]{%
0.005	0.961707842717039\\
0.03	0.799460446285425\\
0.06	0.654646276767796\\
0.085	0.563645139599043\\
0.11	0.492129910071565\\
0.14	0.425214171316473\\
0.165	0.381170028601366\\
0.19	0.345150388287714\\
0.22	0.310037933930232\\
0.245	0.286016465240029\\
0.27	0.265730269316037\\
0.3	0.245310500935774\\
0.325	0.230924352181741\\
0.355	0.216196431151661\\
0.38	0.205661158619046\\
0.405	0.196437371813497\\
0.435	0.186820800469046\\
0.46	0.179829796995251\\
0.485	0.173630541747616\\
0.515	0.167088719904601\\
0.54	0.162283205266665\\
0.565	0.157988037693065\\
0.595	0.153423113413825\\
0.62	0.1500507354806\\
0.645	0.147025047234302\\
0.675	0.143800626219563\\
0.7	0.141415619918082\\
0.725	0.139276412771715\\
0.755	0.137001128315837\\
0.78	0.135324464279381\\
0.805	0.133828395776939\\
0.835	0.132249772299823\\
0.86	0.131098845276331\\
0.885	0.130084723755595\\
0.915	0.129033432135786\\
0.94	0.128284331827419\\
0.97	0.127525257490332\\
0.995	0.12700049330635\\
1.02	0.126566601181478\\
1.05	0.12615746884713\\
1.075	0.125903157108386\\
1.1	0.125722417917874\\
1.13	0.125596487074754\\
1.155	0.12556268994566\\
1.18	0.125589732203869\\
1.21	0.125697905866218\\
1.235	0.125847674035879\\
1.26	0.12604876194903\\
1.29	0.126354350022129\\
1.315	0.126659940532139\\
1.34	0.12700964281338\\
1.37	0.127484874562402\\
1.395	0.127925209520082\\
1.42	0.128404143902913\\
1.45	0.129027787549954\\
1.475	0.129586706068757\\
1.5	0.130179977356563\\
1.53	0.130935679911884\\
1.555	0.131600710582804\\
1.585	0.132439677017438\\
1.61	0.133171914364017\\
1.635	0.133933397328866\\
1.665	0.134884760496843\\
1.69	0.135708100861191\\
1.715	0.13655854959434\\
1.745	0.137614082569128\\
1.77	0.138522252544808\\
1.795	0.139455883074311\\
1.825	0.140609243307358\\
1.85	0.141597419615375\\
1.875	0.142609798679597\\
1.905	0.143856144879799\\
1.93	0.144920664367385\\
1.955	0.146008443347643\\
1.985	0.147344144674196\\
2.01	0.148482283388075\\
2.035	0.149642994968492\\
2.065	0.15106540803458\\
2.09	0.15227521125669\\
2.115	0.153507112545343\\
2.145	0.155014408855734\\
2.17	0.156294559195304\\
2.2	0.157859510122772\\
2.225	0.159187540987469\\
2.25	0.16053725520872\\
2.28	0.162185490426978\\
2.305	0.163582813599598\\
2.33	0.165001761373978\\
2.36	0.166733052471063\\
2.385	0.16819961080713\\
2.41	0.169687850355133\\
2.44	0.171502412486637\\
2.465	0.173038500947475\\
2.49	0.174596427810397\\
2.52	0.176494862804218\\
2.545	0.178101084566285\\
2.57	0.179729390673243\\
2.6	0.181712641520741\\
2.625	0.183389872836418\\
2.65	0.185089513710209\\
2.68	0.187158828083252\\
2.705	0.188908190097728\\
2.73	0.19068035854514\\
2.76	0.192837259523404\\
2.785	0.194660095662778\\
2.815	0.196878227663703\\
2.84	0.198752463272158\\
2.865	0.200650321126919\\
2.895	0.20295917130455\\
2.92	0.204909603384974\\
2.945	0.206884220247171\\
2.975	0.209285948240668\\
3	0.21131443769238\\
3.025	0.213367729084113\\
3.055	0.215864704392751\\
3.08	0.217973283916095\\
3.105	0.220107334584163\\
3.135	0.222702126656282\\
3.16	0.224892993195159\\
3.185	0.227110050154355\\
3.215	0.229805420786414\\
3.24	0.232080929781646\\
3.265	0.234383397072013\\
3.295	0.237182294745955\\
3.32	0.239544955911461\\
3.345	0.241935390800632\\
3.375	0.244840946659351\\
3.4	0.247293421049389\\
3.43	0.250274214341842\\
3.455	0.252790056681783\\
3.48	0.255335154861275\\
3.51	0.258428297593623\\
3.535	0.261038782222559\\
3.56	0.263679463387694\\
3.59	0.266888567845268\\
3.615	0.26959675606709\\
3.64	0.272336127992298\\
3.67	0.27566498494029\\
3.695	0.27847408702624\\
3.72	0.281315406735803\\
3.75	0.284767986499766\\
3.775	0.287681362809307\\
3.8	0.290628037890601\\
3.83	0.294208492177555\\
3.855	0.297229654872719\\
3.88	0.300285245363543\\
3.91	0.303997909747041\\
3.935	0.307130525041215\\
3.96	0.310298745835141\\
3.99	0.314148142842595\\
4.015	0.317396033735455\\
4.045	0.321342156719763\\
4.07	0.324671600211417\\
4.095	0.32803874967065\\
4.125	0.332129671362394\\
4.15	0.335581206357394\\
4.175	0.339071760914883\\
4.205	0.343312525713034\\
4.23	0.34689041348261\\
4.255	0.350508687219353\\
4.285	0.354904542329541\\
4.31	0.358613214648278\\
4.335	0.362363693547376\\
4.365	0.366920094352327\\
4.39	0.370764158094212\\
4.415	0.374651504745458\\
4.445	0.379374120591656\\
4.47	0.383358362668924\\
4.495	0.387387421311852\\
4.525	0.392282141698599\\
4.55	0.396411534332348\\
4.575	0.400587336229311\\
4.605	0.405660277389791\\
4.63	0.409939983704711\\
4.66	0.415139120688698\\
4.685	0.419525264702665\\
4.71	0.423960644641303\\
4.74	0.429348857276454\\
4.765	0.433894477268899\\
4.79	0.438491093821522\\
4.82	0.444075143262269\\
4.845	0.448785945595956\\
4.87	0.453549571344957\\
4.9	0.459336473043973\\
4.925	0.464218378356061\\
4.95	0.469155002335518\\
4.98	0.47515203440512\\
5.005	0.480211184683509\\
5.03	0.485327019555136\\
5.06	0.49154173151017\\
5.085	0.496784497489258\\
5.11	0.502085987038125\\
5.14	0.508526208923197\\
5.165	0.513959197801442\\
5.19	0.5194530247552\\
5.22	0.526126876663319\\
5.245	0.531756940147906\\
5.275	0.538596278447651\\
5.3	0.544365936312634\\
5.325	0.550200178997572\\
5.355	0.557287535513878\\
5.38	0.563266408197994\\
5.405	0.56931219628188\\
5.435	0.57665652087574\\
5.46	0.582852159667638\\
5.485	0.589117130018651\\
5.515	0.596727701192433\\
5.54	0.603147934488745\\
5.565	0.609640004065366\\
5.595	0.617526440417675\\
5.62	0.624179383393132\\
5.645	0.63090675903014\\
5.675	0.639079031151612\\
5.7	0.645973095801507\\
5.725	0.652944284370212\\
5.755	0.661412727193457\\
5.78	0.668556632758011\\
5.805	0.675780451712034\\
5.835	0.684555777328928\\
5.86	0.691958561110103\\
5.89	0.700951284844929\\
5.915	0.708537460386625\\
5.94	0.716208488970245\\
5.97	0.725527064690993\\
5.995	0.733388121603641\\
6.02	0.741337102497315\\
6.05	0.750993323188234\\
6.075	0.759139210337827\\
6.1	0.767376204040185\\
6.13	0.777382291375649\\
6.155	0.785823319172165\\
6.18	0.794358751684204\\
6.21	0.804727371229734\\
6.235	0.8134742244523\\
6.26	0.822318900247047\\
6.29	0.833063177215489\\
6.315	0.842126928305107\\
6.34	0.851292043770011\\
6.37	0.862425579351784\\
6.395	0.871817702187386\\
6.42	0.881314859554296\\
6.45	0.892851747839197\\
6.475	0.902584132012682\\
6.505	0.914406765471842\\
6.53	0.924380199303951\\
6.555	0.93446516491121\\
6.585	0.946716101712322\\
6.61	0.9570508447111\\
6.635	0.967501157676142\\
6.665	0.98019590523176\\
6.69	0.990905039002894\\
6.715	1.00173392695296\\
6.745	1.01488855342639\\
6.77	1.02598563252144\\
6.795	1.03720680123111\\
6.825	1.05083795550971\\
6.85	1.06233702427548\\
6.875	1.07396467470153\\
6.905	1.08808960708677\\
6.93	1.10000521711383\\
6.955	1.1120540630389\\
6.985	1.12669064667458\\
7.01	1.1390378748736\\
7.035	1.15152316117799\\
7.065	1.16668991424552\\
7.09	1.17948438156571\\
7.12	1.19502671429192\\
7.145	1.20813801185707\\
7.17	1.22139590267231\\
7.2	1.23750118138974\\
7.225	1.25108736644192\\
7.25	1.26482544862835\\
7.28	1.2815140412461\\
7.305	1.29559229483457\\
7.33	1.30982794056056\\
7.36	1.32712094891931\\
7.385	1.34170907110682\\
7.41	1.3564602783625\\
7.44	1.37437956438184\\
7.465	1.38949599629659\\
7.49	1.40478141118856\\
7.52	1.42334962394382\\
7.545	1.43901347061414\\
7.57	1.45485241044957\\
7.6	1.47409301420562\\
7.625	1.49032406820445\\
7.65	1.50673654539469\\
7.68	1.52667384863613\\
7.705	1.54349261456327\\
7.735	1.56392345259\\
7.76	1.581158546861\\
7.785	1.59858626669145\\
7.815	1.61975682339029\\
7.84	1.63761591675999\\
7.865	1.65567459609065\\
7.895	1.67761160070464\\
7.92	1.69611723946488\\
7.945	1.71482967468857\\
7.975	1.73756081689083\\
8	1.75673635736723\\
};
\addplot [color=mycolor1, line width=1.4pt, only marks, mark size=0.4pt, mark=*, mark options={solid, fill=mycolor1, mycolor1}, forget plot]
  table[row sep=crcr]{%
1.155	0.12556268994566\\
};
\addplot [color=mycolor3, line width=0.9pt, forget plot]
  table[row sep=crcr]{%
0.005	0.965582581329303\\
0.03	0.818326278313415\\
0.06	0.684460688031333\\
0.085	0.598781461558633\\
0.11	0.530360135731492\\
0.14	0.465259657383791\\
0.165	0.421720303451766\\
0.19	0.385632718155537\\
0.22	0.349976127793189\\
0.245	0.325277056364376\\
0.27	0.304205776160398\\
0.3	0.282783266203527\\
0.325	0.267554638698581\\
0.355	0.251841984789343\\
0.38	0.240524045723158\\
0.405	0.230560465539586\\
0.435	0.220118157844473\\
0.46	0.212492346807002\\
0.485	0.205706523723853\\
0.515	0.198522753167361\\
0.54	0.19323175896303\\
0.565	0.188493810949771\\
0.595	0.183450691846933\\
0.62	0.179721386168037\\
0.645	0.176374134778945\\
0.675	0.172807447797278\\
0.7	0.170171134416928\\
0.725	0.167809379211337\\
0.755	0.165302455530153\\
0.78	0.163460306622273\\
0.805	0.161822142009088\\
0.835	0.160101868430468\\
0.86	0.15885541252102\\
0.885	0.157764938902782\\
0.915	0.156645779330961\\
0.94	0.155858723917831\\
0.97	0.15507503191264\\
0.995	0.154546193231645\\
1.02	0.154122336643403\\
1.05	0.15374286224222\\
1.075	0.15352717088063\\
1.1	0.153397046139891\\
1.13	0.153346901716747\\
1.155	0.153388200054875\\
1.18	0.153500689464753\\
1.21	0.153724454670675\\
1.235	0.153980960342441\\
1.26	0.154297858374241\\
1.29	0.154753919006989\\
1.315	0.155194121581695\\
1.34	0.155686509756298\\
1.37	0.156343250350273\\
1.395	0.156943128419968\\
1.42	0.157588898805006\\
1.45	0.158422081701795\\
1.475	0.159163173035899\\
1.5	0.159945299866474\\
1.53	0.160936223686182\\
1.555	0.161804260414313\\
1.585	0.162895017026846\\
1.61	0.163843752824848\\
1.635	0.164827679229792\\
1.665	0.166053675329918\\
1.69	0.167112182022827\\
1.715	0.168203436587606\\
1.745	0.169555269548493\\
1.77	0.170716382814763\\
1.795	0.171908366054209\\
1.825	0.173378807757287\\
1.85	0.174637038902254\\
1.875	0.175924713929306\\
1.905	0.177508275233884\\
1.93	0.178859481086184\\
1.955	0.180239070407503\\
1.985	0.181931669356999\\
2.01	0.183372799913474\\
2.035	0.184841552929916\\
2.065	0.186640259508985\\
2.09	0.188169163195523\\
2.115	0.189725176298128\\
2.145	0.191628014361183\\
2.17	0.193243286491934\\
2.2	0.195216989065666\\
2.225	0.196891155093106\\
2.25	0.198592019068112\\
2.28	0.200668265625047\\
2.305	0.202427804374581\\
2.33	0.204214017940533\\
2.36	0.206392716427014\\
2.385	0.208237708882683\\
2.41	0.210109489761308\\
2.44	0.212391073848133\\
2.465	0.214322018810009\\
2.49	0.216279985426423\\
2.52	0.218665348658474\\
2.545	0.220683111929097\\
2.57	0.222728236012066\\
2.6	0.225218679099735\\
2.625	0.227324452779959\\
2.65	0.229458021371213\\
2.68	0.23205520992505\\
2.705	0.234250479764952\\
2.73	0.236474064742888\\
2.76	0.239179995370398\\
2.785	0.241466514586513\\
2.815	0.244248522400411\\
2.84	0.246598921384243\\
2.865	0.248978689706051\\
2.895	0.251873486391828\\
2.92	0.254318643468047\\
2.945	0.256793890443779\\
2.975	0.259804242386666\\
3	0.262346537262013\\
3.025	0.264919709059013\\
3.055	0.268048637980242\\
3.08	0.2706906596119\\
3.105	0.273364408704069\\
3.135	0.276615180335525\\
3.16	0.279359718316624\\
3.185	0.282136895575239\\
3.215	0.285513011046204\\
3.24	0.288363049118826\\
3.265	0.291246697918424\\
3.295	0.294751887422836\\
3.32	0.297710598784046\\
3.345	0.300703950810983\\
3.375	0.304342169206298\\
3.4	0.307412913337745\\
3.43	0.311144998607249\\
3.455	0.314294809904984\\
3.48	0.317481130707315\\
3.51	0.32135342115034\\
3.535	0.3246213521029\\
3.56	0.327926977156956\\
3.59	0.331944021652232\\
3.615	0.335333926207669\\
3.64	0.338762767018086\\
3.67	0.342929335501002\\
3.695	0.346445252205777\\
3.72	0.350001405308359\\
3.75	0.354322490427705\\
3.775	0.357968644070841\\
3.8	0.361656392906895\\
3.83	0.36613721256776\\
3.855	0.369918016535818\\
3.88	0.373741834014711\\
3.91	0.378387834718158\\
3.935	0.382307893992205\\
3.96	0.386272445676783\\
3.99	0.391089306584425\\
4.015	0.395153421322757\\
4.045	0.400091166128116\\
4.07	0.404257206849303\\
4.095	0.408470376510162\\
4.125	0.413589106160734\\
4.15	0.417907756925516\\
4.175	0.422275184892209\\
4.205	0.427581228963454\\
4.23	0.432057836283208\\
4.255	0.436584934961904\\
4.285	0.442084876086223\\
4.31	0.446724999261526\\
4.335	0.451417395627863\\
4.365	0.457118076329944\\
4.39	0.461927493294284\\
4.415	0.466791034869924\\
4.445	0.472699564911519\\
4.47	0.477684278494932\\
4.495	0.482725039739063\\
4.525	0.488848803957546\\
4.55	0.494015048557714\\
4.575	0.499239337664187\\
4.605	0.50558600429328\\
4.63	0.510940252945601\\
4.66	0.517444764623541\\
4.685	0.522932148501024\\
4.71	0.528481120342875\\
4.74	0.535222141937508\\
4.765	0.540909015235479\\
4.79	0.546659683360555\\
4.82	0.553645689620336\\
4.845	0.559539204785864\\
4.87	0.565498804487982\\
4.9	0.572738588311337\\
4.925	0.578846165886128\\
4.95	0.58502220329387\\
4.98	0.592524886252033\\
5.005	0.598854223776225\\
5.03	0.605254484880469\\
5.06	0.613029528214293\\
5.085	0.61958860952725\\
5.11	0.626221169606315\\
5.14	0.634278385701133\\
5.165	0.641075490642672\\
5.19	0.647948724090128\\
5.22	0.656298288437991\\
5.245	0.663342002957037\\
5.275	0.671898653951083\\
5.3	0.679117057175355\\
5.325	0.68641628385425\\
5.355	0.695283312577003\\
5.38	0.702763537830423\\
5.405	0.710327506833005\\
5.435	0.719516127760996\\
5.46	0.727267636995012\\
5.485	0.735105917448171\\
5.515	0.744627756631591\\
5.54	0.752660358843097\\
5.565	0.760782870662671\\
5.595	0.770649980103877\\
5.62	0.778973843472769\\
5.645	0.787390869645947\\
5.675	0.797615742229105\\
5.7	0.806241406721312\\
5.725	0.814963606056918\\
5.755	0.825559191045535\\
5.78	0.83449758149366\\
5.805	0.843536001854691\\
5.835	0.854515720973748\\
5.86	0.86377816064868\\
5.89	0.875030011120754\\
5.915	0.8845220167996\\
5.94	0.894120241489327\\
5.97	0.905779993390774\\
5.995	0.915616099295723\\
6.02	0.925562271904632\\
6.05	0.937644701102006\\
6.075	0.947837372528829\\
6.1	0.958144098217002\\
6.13	0.970664517681856\\
6.155	0.981226672910074\\
6.18	0.991907014786238\\
6.21	1.00488129366136\\
6.235	1.01582632004023\\
6.26	1.02689381542874\\
6.29	1.04033839878978\\
6.315	1.05168016940996\\
6.34	1.06314884670992\\
6.37	1.07708077604615\\
6.395	1.08883366701686\\
6.42	1.10071806317221\\
6.45	1.11515499760841\\
6.475	1.12733390595982\\
6.505	1.14212860654091\\
6.53	1.15460932284749\\
6.555	1.16722968507967\\
6.585	1.18256065273901\\
6.61	1.19549375841316\\
6.635	1.20857156971305\\
6.665	1.22445823465056\\
6.69	1.23786012048187\\
6.715	1.25141195463089\\
6.745	1.26787444977026\\
6.77	1.28176209924098\\
6.795	1.2958051292358\\
6.825	1.31286431527067\\
6.85	1.32725532568444\\
6.875	1.34180734509815\\
6.905	1.35948483643673\\
6.93	1.37439744079838\\
6.955	1.38947688590492\\
6.985	1.40779507754001\\
7.01	1.42324816722228\\
7.035	1.4388741399083\\
7.065	1.45785623524882\\
7.09	1.47386938346423\\
7.12	1.49332180626453\\
7.145	1.50973171465407\\
7.17	1.52632520076369\\
7.2	1.54648259699101\\
7.225	1.56348720800344\\
7.25	1.58068204269272\\
7.28	1.60156993406728\\
7.305	1.6191907762398\\
7.33	1.63700872758747\\
7.36	1.65865355571202\\
7.385	1.67691293343433\\
7.41	1.6953765539087\\
7.44	1.717805713055\\
7.465	1.73672673422069\\
7.49	1.75585938863527\\
7.52	1.77910125969621\\
7.545	1.79870786433544\\
7.57	1.81853375879445\\
7.6	1.84261774443602\\
7.625	1.86293473436371\\
7.65	1.88347894623522\\
7.68	1.90843550729792\\
7.705	1.92948857681242\\
7.735	1.95506327104373\\
7.76	1.97663777894933\\
7.785	1.99845354661724\\
7.815	2.02495472411759\\
7.84	2.04731078364516\\
7.865	2.06991682521254\\
7.895	2.09737798115924\\
7.92	2.12054384580851\\
7.945	2.14396872827512\\
7.975	2.17242456173868\\
8	2.19642950038594\\
};
\addplot [color=mycolor1, line width=1.4pt, only marks, mark size=0.4pt, mark=*, mark options={solid, fill=mycolor1, mycolor1}, forget plot]
  table[row sep=crcr]{%
1.13	0.153346901716747\\
};
\addplot [color=mycolor4, line width=0.9pt, forget plot]
  table[row sep=crcr]{%
0.005	0.970872099258847\\
0.03	0.84436814636939\\
0.06	0.726149097399875\\
0.085	0.648419077990237\\
0.11	0.584901552846846\\
0.14	0.52303012939518\\
0.165	0.48073041883824\\
0.19	0.445027928732667\\
0.22	0.409110799326526\\
0.245	0.383820426186479\\
0.27	0.361957401880175\\
0.3	0.339442001133725\\
0.325	0.32325164159219\\
0.355	0.306380206200183\\
0.38	0.294120970048508\\
0.405	0.283254543071193\\
0.435	0.271792314916088\\
0.46	0.263375395623777\\
0.485	0.255854409790341\\
0.515	0.247862884487803\\
0.54	0.241959911168704\\
0.565	0.236663965583826\\
0.595	0.231019708396329\\
0.62	0.226843878542514\\
0.645	0.223097068540184\\
0.675	0.219109556901465\\
0.7	0.216168682186811\\
0.725	0.213541884069224\\
0.755	0.210766035233564\\
0.78	0.20873832327656\\
0.805	0.206947582196063\\
0.835	0.205085231644008\\
0.86	0.203752591675302\\
0.885	0.20260355576056\\
0.915	0.201448549967696\\
0.94	0.200658742312173\\
0.97	0.199902442068615\\
0.995	0.19942067088987\\
1.02	0.199064793906225\\
1.05	0.198793130298012\\
1.075	0.198688089223346\\
1.1	0.198686644555774\\
1.13	0.198813637685341\\
1.155	0.199020654812187\\
1.18	0.199314623488634\\
1.21	0.199776119266107\\
1.235	0.200246709719914\\
1.26	0.200791666930823\\
1.29	0.201539177190986\\
1.315	0.202236555784066\\
1.34	0.202998689254124\\
1.37	0.203995190899594\\
1.395	0.204891188214158\\
1.42	0.205844543462526\\
1.45	0.207061551525301\\
1.475	0.208134450524823\\
1.5	0.209258986919441\\
1.53	0.210674476168237\\
1.555	0.211907466678214\\
1.585	0.213449260301259\\
1.61	0.214784565312426\\
1.635	0.216164621836789\\
1.665	0.217878389142645\\
1.69	0.219353561709496\\
1.715	0.22087061669706\\
1.745	0.222745314183006\\
1.77	0.224351952323476\\
1.795	0.225998279098023\\
1.825	0.228025462757544\\
1.85	0.229757171494328\\
1.875	0.231526918239819\\
1.905	0.233700234448067\\
1.93	0.235552239210275\\
1.955	0.237441073297052\\
1.985	0.239755868933544\\
2.01	0.241724718535085\\
2.035	0.243729552475905\\
2.065	0.246182573620641\\
2.09	0.248265910046099\\
2.115	0.250384688553607\\
2.145	0.252973845733556\\
2.17	0.255170224307426\\
2.2	0.257852278585556\\
2.225	0.260125941541563\\
2.25	0.262434689485689\\
2.28	0.265251499649749\\
2.305	0.267637458058415\\
2.33	0.270058562582784\\
2.36	0.273010359495041\\
2.385	0.275509001521877\\
2.41	0.278043020841298\\
2.44	0.281130689129318\\
2.465	0.283742923429267\\
2.49	0.286390915729788\\
2.52	0.289615914834822\\
2.545	0.292343109607197\\
2.57	0.29510657618078\\
2.6	0.298470877110247\\
2.625	0.301314811249556\\
2.65	0.30419565089948\\
2.68	0.30770168550432\\
2.705	0.31066450942348\\
2.73	0.313664981902643\\
2.76	0.317315602218387\\
2.785	0.320399806463086\\
2.815	0.324151686425611\\
2.84	0.327320948938009\\
2.865	0.330529339122381\\
2.895	0.334431483000966\\
2.92	0.337727032555457\\
2.945	0.341062709463031\\
2.975	0.345118955357877\\
3	0.348544086364364\\
3.025	0.352010430300225\\
3.055	0.356224946304535\\
3.08	0.359783224017916\\
3.105	0.363383882684019\\
3.135	0.367761153618179\\
3.16	0.371456404091417\\
3.185	0.37519528340407\\
3.215	0.379740100873513\\
3.24	0.38357640352088\\
3.265	0.387457660902947\\
3.295	0.392175116159921\\
3.32	0.396156798527242\\
3.345	0.400184838298342\\
3.375	0.405080317486486\\
3.4	0.409211951903104\\
3.43	0.414233098815328\\
3.455	0.418470603904285\\
3.48	0.422757005752241\\
3.51	0.427965927693774\\
3.535	0.432361651556273\\
3.56	0.436807878941188\\
3.59	0.442210745944892\\
3.615	0.446769914011867\\
3.64	0.45138126866068\\
3.67	0.456984543785531\\
3.695	0.461712626362028\\
3.72	0.466494655603757\\
3.75	0.472305097724052\\
3.775	0.477207812645432\\
3.8	0.482166312357634\\
3.83	0.488190980069646\\
3.855	0.493274296255456\\
3.88	0.498415314681008\\
3.91	0.504661571218733\\
3.935	0.509931713040875\\
3.96	0.515261555395891\\
3.99	0.521737074468989\\
4.015	0.527200526866232\\
4.045	0.533838279803104\\
4.07	0.539438533158014\\
4.095	0.545102047441598\\
4.125	0.551982735348573\\
4.15	0.557787849960571\\
4.175	0.563658450265698\\
4.205	0.570790616581771\\
4.23	0.576807808213267\\
4.255	0.582892798560887\\
4.285	0.590285325665049\\
4.31	0.596522095158257\\
4.335	0.602829067098818\\
4.365	0.61049118571582\\
4.39	0.616955327023041\\
4.415	0.623492167839953\\
4.445	0.631433467110158\\
4.47	0.638133075872193\\
4.495	0.644907977327526\\
4.525	0.653138415522818\\
4.55	0.66008189818114\\
4.575	0.667103365793281\\
4.605	0.675633281665995\\
4.63	0.682829365099291\\
4.66	0.691571372594439\\
4.685	0.698946352698278\\
4.71	0.706404094378674\\
4.74	0.715463922634622\\
4.765	0.723106985292293\\
4.79	0.730835783957229\\
4.82	0.740224855324537\\
4.845	0.748145640351917\\
4.87	0.756155247192158\\
4.9	0.765885411725457\\
4.925	0.774093919384788\\
4.95	0.782394449751495\\
4.98	0.792477999472462\\
5.005	0.800984602537482\\
5.03	0.809586548172125\\
5.06	0.820036231963563\\
5.085	0.828851688313587\\
5.11	0.837765930107868\\
5.14	0.848594969245373\\
5.165	0.85773043499371\\
5.19	0.866968256287049\\
5.22	0.878190360621641\\
5.245	0.887657403788438\\
5.275	0.89915795516936\\
5.3	0.908859889014843\\
5.325	0.918670503756016\\
5.355	0.930588410860061\\
5.38	0.94064241711989\\
5.405	0.950809038523604\\
5.435	0.963159409600786\\
5.46	0.973578235388201\\
5.485	0.98411375485242\\
5.515	0.99691225223288\\
5.54	1.00770911188851\\
5.565	1.01862689309467\\
5.595	1.03188975260864\\
5.62	1.04307834409254\\
5.645	1.05439223959204\\
5.675	1.06813629072988\\
5.7	1.0797308126387\\
5.725	1.09145518106796\\
5.755	1.10569786790317\\
5.78	1.11771303712915\\
5.805	1.12986276107278\\
5.835	1.14462216397083\\
5.86	1.15707323402088\\
5.89	1.17219870864921\\
5.915	1.18495859667679\\
5.94	1.19786137667374\\
5.97	1.21353558398058\\
5.995	1.22675838323648\\
6.02	1.24012925762063\\
6.05	1.25637210191452\\
6.075	1.27007460417526\\
6.1	1.28393055313273\\
6.13	1.30076266225012\\
6.155	1.31496226951303\\
6.18	1.32932089014944\\
6.21	1.34676364117914\\
6.235	1.3614783873653\\
6.26	1.37635791565249\\
6.29	1.39443346159663\\
6.315	1.40968203504258\\
6.34	1.42510136856287\\
6.37	1.44383266596644\\
6.395	1.45963443273161\\
6.42	1.47561315424497\\
6.45	1.49502399182395\\
6.475	1.51139901984031\\
6.505	1.53129128722176\\
6.53	1.54807245000698\\
6.555	1.5650415341\\
6.585	1.5856554553136\\
6.61	1.60304540570555\\
6.635	1.62063009351544\\
6.665	1.64199184370744\\
6.69	1.66001266242538\\
6.715	1.67823528125914\\
6.745	1.70037198251992\\
6.77	1.71904654896565\\
6.795	1.7379302336104\\
6.825	1.76086998875669\\
6.85	1.78022200951263\\
6.875	1.79979073104711\\
6.905	1.82356265861528\\
6.93	1.84361669695517\\
6.955	1.86389529259025\\
6.985	1.88852956307731\\
7.01	1.90931106953784\\
7.035	1.93032527351862\\
7.065	1.95585314691962\\
7.09	1.97738849095766\\
7.12	2.00354943434293\\
7.145	2.02561883125252\\
7.17	2.04793533403094\\
7.2	2.07504521110057\\
7.225	2.09791511970998\\
7.25	2.12104108805256\\
7.28	2.14913428332655\\
7.305	2.17283370842145\\
7.33	2.19679846997232\\
7.36	2.22591060780505\\
7.385	2.25046959988747\\
7.41	2.27530353952106\\
7.44	2.30547152832063\\
7.465	2.33092122090249\\
7.49	2.3566558184275\\
7.52	2.38791789644706\\
7.545	2.41429054466725\\
7.57	2.44095841386553\\
7.6	2.47335419661302\\
7.625	2.50068321722407\\
7.65	2.52831814627569\\
7.68	2.5618886756087\\
7.705	2.59020868838046\\
7.735	2.6246114289798\\
7.76	2.65363347667586\\
7.785	2.68298034726407\\
7.815	2.71863046948335\\
7.84	2.74870477834946\\
7.865	2.77911566268384\\
7.895	2.81605829861614\\
7.92	2.84722293820552\\
7.945	2.8787363284466\\
7.975	2.9170182330685\\
8	2.94931264170054\\
};
\addplot [color=mycolor1, line width=1.4pt, only marks, mark size=0.4pt, mark=*, mark options={solid, fill=mycolor1, mycolor1}, forget plot]
  table[row sep=crcr]{%
1.09	0.198675168593949\\
};
\addplot [color=mycolor5, line width=0.9pt, forget plot]
  table[row sep=crcr]{%
0.005	0.977212300208905\\
0.03	0.876116473648082\\
0.06	0.777970765986068\\
0.085	0.711078575945561\\
0.11	0.654762331237564\\
0.14	0.598250476224996\\
0.165	0.558550896556022\\
0.19	0.524296627244214\\
0.22	0.489086707104119\\
0.245	0.463811378605152\\
0.27	0.44162207032704\\
0.3	0.418428945121897\\
0.325	0.401531470453929\\
0.355	0.383725238271462\\
0.38	0.370660577421427\\
0.405	0.358993325863222\\
0.435	0.346601579341817\\
0.46	0.337450497794626\\
0.485	0.329240498847976\\
0.515	0.320488397708855\\
0.54	0.314009865274923\\
0.565	0.30819249298663\\
0.595	0.301993988551332\\
0.62	0.297415011045564\\
0.645	0.293317169666882\\
0.675	0.288975076606471\\
0.7	0.285792262023611\\
0.725	0.282970270869596\\
0.755	0.280019289736618\\
0.78	0.277892717227297\\
0.805	0.27604398575272\\
0.835	0.274163526253946\\
0.86	0.272856718772897\\
0.885	0.271769004109867\\
0.915	0.270732286924046\\
0.94	0.270076618728665\\
0.97	0.269521944329158\\
0.995	0.269240592735811\\
1.02	0.269113323492496\\
1.05	0.269151640447335\\
1.075	0.269333413458741\\
1.1	0.269643660551609\\
1.13	0.270176260684679\\
1.155	0.270746618438265\\
1.18	0.27142613084808\\
1.21	0.272378573668297\\
1.235	0.273281080552667\\
1.26	0.274278013297114\\
1.29	0.275593563436238\\
1.315	0.276785074637812\\
1.34	0.278059692494572\\
1.37	0.27969475877505\\
1.395	0.281142052034368\\
1.42	0.282663705272136\\
1.45	0.284584609965706\\
1.475	0.286261985151827\\
1.5	0.288006942808361\\
1.53	0.290187594017219\\
1.555	0.292075134316837\\
1.585	0.294422338193052\\
1.61	0.296445196310896\\
1.635	0.298527490149893\\
1.665	0.301103081344923\\
1.69	0.303312206019396\\
1.715	0.305577395226815\\
1.745	0.308368397824635\\
1.77	0.310753946402335\\
1.795	0.313193008916274\\
1.825	0.316189605260579\\
1.85	0.318744180234661\\
1.875	0.32135038377521\\
1.905	0.324545316013359\\
1.93	0.327263509407044\\
1.955	0.33003199136236\\
1.985	0.333420099841996\\
2.01	0.336298140743073\\
2.035	0.339225582328245\\
2.065	0.342803446602061\\
2.09	0.34583892632413\\
2.115	0.348923297925906\\
2.145	0.35268895550945\\
2.17	0.355880611780103\\
2.2	0.359774888826654\\
2.225	0.363073697212174\\
2.25	0.366421243675847\\
2.28	0.370502712655474\\
2.305	0.373957708934133\\
2.33	0.377461699760991\\
2.36	0.381731340821399\\
2.385	0.385343592993144\\
2.41	0.389005314915999\\
2.44	0.393464949075241\\
2.465	0.397236196302199\\
2.49	0.401057582257528\\
2.52	0.405709775372223\\
2.545	0.409642353821056\\
2.57	0.413625913739089\\
2.6	0.41847389951593\\
2.625	0.422570682894864\\
2.65	0.426719448204285\\
2.68	0.431767066479078\\
2.705	0.436031418385906\\
2.73	0.440348897565748\\
2.76	0.44560054491112\\
2.785	0.450036280741783\\
2.815	0.455431018624265\\
2.84	0.459987014195093\\
2.865	0.464598370846096\\
2.895	0.470205713715265\\
2.92	0.474940474831603\\
2.945	0.479732086800528\\
2.975	0.485557754583104\\
3	0.490476178016733\\
3.025	0.495453058792852\\
3.055	0.501503219805419\\
3.08	0.506610571126498\\
3.105	0.511778098985557\\
3.135	0.518059354498343\\
3.16	0.523361256477842\\
3.185	0.528725163971934\\
3.215	0.535244536953446\\
3.24	0.540746961254334\\
3.265	0.546313327824487\\
3.295	0.553078255230218\\
3.32	0.55878751689285\\
3.345	0.564562764141958\\
3.375	0.571581092229715\\
3.4	0.577503846551629\\
3.43	0.584701150681505\\
3.455	0.590774718055703\\
3.48	0.596917959540016\\
3.51	0.604382828343009\\
3.535	0.610681888690052\\
3.56	0.617052951796555\\
3.59	0.624794327421153\\
3.615	0.631326450189106\\
3.64	0.637933011672288\\
3.67	0.645960248388202\\
3.695	0.652733347753705\\
3.72	0.659583430409954\\
3.75	0.667906299643418\\
3.775	0.674928639142348\\
3.8	0.68203061682139\\
3.83	0.690659313656618\\
3.855	0.69793951199501\\
3.88	0.705302115758027\\
3.91	0.714247266770981\\
3.935	0.721794304713505\\
3.96	0.729426630040183\\
3.99	0.738699302357892\\
4.015	0.746522530610324\\
4.045	0.756027027409403\\
4.07	0.764045757074586\\
4.095	0.772154905092662\\
4.125	0.782006626888906\\
4.15	0.790318190021379\\
4.175	0.798723375690887\\
4.205	0.808934627061295\\
4.23	0.817549414639613\\
4.255	0.826261154784793\\
4.285	0.836844723591012\\
4.31	0.845773533219371\\
4.335	0.854802754905847\\
4.365	0.865771926177359\\
4.39	0.875025973918723\\
4.415	0.884384026524828\\
4.445	0.895752597255497\\
4.47	0.905343530219954\\
4.495	0.915042198256319\\
4.525	0.926824493094003\\
4.55	0.936764402766184\\
4.575	0.946815919432989\\
4.605	0.959026807248296\\
4.63	0.96932824354409\\
4.66	0.981842705273635\\
4.685	0.992400216340342\\
4.71	1.00307620238944\\
4.74	1.01604563140752\\
4.765	1.02698692815798\\
4.79	1.03805097580918\\
4.82	1.05149179242056\\
4.845	1.06283073206911\\
4.87	1.0742968581216\\
4.9	1.08822609547125\\
4.925	1.09997705176643\\
4.95	1.11185979493307\\
4.98	1.12629511966534\\
5.005	1.13847300035933\\
5.03	1.15078743896944\\
5.06	1.16574717278036\\
5.085	1.17836743784778\\
5.11	1.1911292083011\\
5.14	1.20663235039504\\
5.165	1.21971103101005\\
5.19	1.23293634699415\\
5.22	1.2490025974656\\
5.245	1.26255631575325\\
5.275	1.27902150667106\\
5.3	1.29291177239627\\
5.325	1.30695776197573\\
5.355	1.32402096275601\\
5.38	1.33841571366332\\
5.405	1.3529718410006\\
5.435	1.370654758597\\
5.46	1.38557230905067\\
5.485	1.40065709564101\\
5.515	1.41898223242848\\
5.54	1.43444156756622\\
5.565	1.45007421296854\\
5.595	1.46906489474448\\
5.62	1.48508569411038\\
5.645	1.50128609985855\\
5.675	1.52096650490308\\
5.7	1.53756916697541\\
5.725	1.55435796138939\\
5.755	1.57475315061828\\
5.78	1.59195881825669\\
5.805	1.60935738218235\\
5.835	1.63049333042617\\
5.86	1.64832391726501\\
5.89	1.66998469518437\\
5.915	1.68825803921158\\
5.94	1.70673625703161\\
5.97	1.72918378984734\\
5.995	1.74812085689638\\
6.02	1.76727024400124\\
6.05	1.79053312677229\\
6.075	1.81015804038306\\
6.1	1.83000299101906\\
6.13	1.85411085836734\\
6.155	1.87444861886663\\
6.18	1.89501441371278\\
6.21	1.9199979768651\\
6.235	1.94107449261681\\
6.26	1.96238733038518\\
6.29	1.98827841556049\\
6.315	2.01012053534908\\
6.34	2.03220756553111\\
6.37	2.05903915371959\\
6.395	2.08167470030564\\
6.42	2.10456405710119\\
6.45	2.13237032527842\\
6.475	2.15582813016493\\
6.505	2.18432495805802\\
6.53	2.20836533146354\\
6.555	2.23267527091474\\
6.585	2.26220728418418\\
6.61	2.28712095746523\\
6.635	2.31231398981996\\
6.665	2.34291879554039\\
6.69	2.36873749330075\\
6.715	2.3948456976244\\
6.745	2.42656226407103\\
6.77	2.45331885901247\\
6.795	2.48037547513156\\
6.825	2.51324418040603\\
6.85	2.54097273434787\\
6.875	2.5690122043206\\
6.905	2.60307488671895\\
6.93	2.63181069307447\\
6.955	2.6608687041414\\
6.985	2.6961687143172\\
7.01	2.72594834209628\\
7.035	2.75606187115869\\
7.065	2.7926441261514\\
7.09	2.82350546552539\\
7.12	2.86099615994225\\
7.145	2.89262386439999\\
7.17	2.92460617389048\\
7.2	2.96345861942512\\
7.225	2.99623510330907\\
7.25	3.02937905865911\\
7.28	3.06964266483495\\
7.305	3.10360960747445\\
7.33	3.13795735225337\\
7.36	3.17968331135881\\
7.385	3.21488389574216\\
7.41	3.25047909375028\\
7.44	3.29372044448341\\
7.465	3.3301994109343\\
7.49	3.36708730047001\\
7.52	3.41189899384597\\
7.545	3.44970269560505\\
7.57	3.4879301456536\\
7.6	3.53436911324157\\
7.625	3.57354557401947\\
7.65	3.61316114243708\\
7.68	3.66128636702311\\
7.705	3.70188534059331\\
7.735	3.75120518238313\\
7.76	3.79281192309938\\
7.785	3.83488495500009\\
7.815	3.88599544369549\\
7.84	3.92911275277387\\
7.865	3.97271324419729\\
7.895	4.02567924502256\\
7.92	4.07036183764666\\
7.945	4.11554511158689\\
7.975	4.17043382386159\\
8	4.21673838384699\\
};
\addplot [color=mycolor1, line width=1.4pt, only marks, mark size=0.4pt, mark=*, mark options={solid, fill=mycolor1, mycolor1}, forget plot]
  table[row sep=crcr]{%
1.03	0.26910355910422\\
};
\addplot [color=mycolor6, line width=0.9pt, forget plot]
  table[row sep=crcr]{%
0.005	0.983391670316749\\
0.03	0.907845222731039\\
0.06	0.831248691332097\\
0.085	0.776938399233969\\
0.11	0.729729000395768\\
0.14	0.680854282641563\\
0.165	0.645546565197057\\
0.19	0.614393461955687\\
0.22	0.581673380287177\\
0.245	0.557733215916079\\
0.27	0.536397343242062\\
0.3	0.513775138298066\\
0.325	0.497088436544932\\
0.355	0.479322407715303\\
0.38	0.466174545146756\\
0.405	0.454359040885725\\
0.435	0.441743363464012\\
0.46	0.432392044725438\\
0.485	0.423986087842368\\
0.515	0.415020463299352\\
0.54	0.408391731365805\\
0.565	0.402455715878233\\
0.595	0.396162266348848\\
0.62	0.391546943374917\\
0.645	0.387453523531183\\
0.675	0.38317191188926\\
0.7	0.380085949691638\\
0.725	0.377403015238533\\
0.755	0.374674120910435\\
0.78	0.372777964042798\\
0.805	0.371200087526165\\
0.835	0.369696841577437\\
0.86	0.368746793571849\\
0.885	0.368053377294548\\
0.915	0.367537997490389\\
0.94	0.367355783137467\\
0.97	0.36741462279824\\
0.995	0.367681240508422\\
1.02	0.368134341911133\\
1.05	0.368910644068476\\
1.075	0.369741040828129\\
1.1	0.370729639748103\\
1.13	0.372114439976206\\
1.155	0.373425940816794\\
1.18	0.374874037161207\\
1.21	0.376784111347953\\
1.235	0.378513389266723\\
1.26	0.380362631006151\\
1.29	0.382733931960361\\
1.315	0.384832151570413\\
1.34	0.387037467203052\\
1.37	0.389820451053674\\
1.395	0.392249790552512\\
1.42	0.394776245391548\\
1.45	0.397932489653271\\
1.475	0.400663611013781\\
1.5	0.403484110216946\\
1.53	0.406983830935587\\
1.555	0.409994014003547\\
1.585	0.413716187571605\\
1.61	0.416907806613252\\
1.635	0.420179555506507\\
1.665	0.424209609521081\\
1.69	0.427653238511782\\
1.715	0.431173228915308\\
1.745	0.435496658186524\\
1.77	0.439181362533495\\
1.795	0.44293963833295\\
1.825	0.447545702838274\\
1.85	0.451463474020486\\
1.875	0.45545282854561\\
1.905	0.460333873241663\\
1.93	0.464479109207551\\
1.955	0.468694604702989\\
1.985	0.473845529901894\\
2.01	0.478214629459297\\
2.035	0.482653221143748\\
2.065	0.488071067200213\\
2.09	0.49266211123873\\
2.115	0.497322351279196\\
2.145	0.503005970898987\\
2.17	0.507818470379541\\
2.2	0.51368495941063\\
2.225	0.518650071775333\\
2.25	0.523684767505035\\
2.28	0.529818504856558\\
2.305	0.535006951808657\\
2.33	0.540265654718118\\
2.36	0.546669217578199\\
2.385	0.55208346998092\\
2.41	0.557568933326459\\
2.44	0.564245984374863\\
2.465	0.569889382595614\\
2.49	0.575605199516381\\
2.52	0.582560372960204\\
2.545	0.588437038843438\\
2.57	0.594387560202277\\
2.6	0.60162637039383\\
2.625	0.607741136703629\\
2.65	0.613931405269643\\
2.68	0.621460173561073\\
2.705	0.627818527566061\\
2.73	0.634254225453671\\
2.76	0.642080021880893\\
2.785	0.648688060465886\\
2.815	0.656722531533618\\
2.84	0.663506064510408\\
2.865	0.670370458747813\\
2.895	0.678715462483865\\
2.92	0.685760255924223\\
2.945	0.692888223462292\\
2.975	0.701552640262488\\
3	0.708866274994442\\
3.025	0.716265558951167\\
3.055	0.725258892646354\\
3.08	0.732849463674329\\
3.105	0.740528316982974\\
3.135	0.749860678055897\\
3.16	0.757736782152524\\
3.185	0.765703956511814\\
3.215	0.775386050535945\\
3.24	0.783556777951354\\
3.265	0.791821516693676\\
3.295	0.801864637229496\\
3.32	0.810339566923332\\
3.345	0.818911601244451\\
3.375	0.82932762656951\\
3.4	0.8381168244837\\
3.43	0.848796415968875\\
3.455	0.857807765748331\\
3.48	0.866921785726668\\
3.51	0.877995654055248\\
3.535	0.887339366287269\\
3.56	0.896789253319669\\
3.59	0.908270856444212\\
3.615	0.917958316745189\\
3.64	0.927755612275255\\
3.67	0.939659006609786\\
3.695	0.949702101781417\\
3.72	0.959858850861615\\
3.75	0.972198700578198\\
3.775	0.982609826929547\\
3.8	0.993138587048326\\
3.83	1.00593017531864\\
3.855	1.01672224855855\\
3.88	1.02763609999751\\
3.91	1.04089534201734\\
3.935	1.05208180859332\\
3.96	1.06339436619619\\
3.99	1.07713782386473\\
4.015	1.088732673648\\
4.045	1.10281898056171\\
4.07	1.11470299805719\\
4.095	1.12672077829322\\
4.125	1.14132076144503\\
4.15	1.1536380417712\\
4.175	1.16609387223064\\
4.205	1.18122591599743\\
4.23	1.19399198192418\\
4.255	1.20690157120414\\
4.285	1.22258477359164\\
4.31	1.23581574856958\\
4.335	1.24919541150706\\
4.365	1.26544960549551\\
4.39	1.27916223181601\\
4.415	1.29302890798451\\
4.445	1.30987468416502\\
4.47	1.32408634218981\\
4.495	1.33845761545783\\
4.525	1.35591634593879\\
4.55	1.37064507438292\\
4.575	1.38553919352567\\
4.605	1.40363305709687\\
4.63	1.41889757435691\\
4.66	1.43744138917036\\
4.685	1.45308548325971\\
4.71	1.46890520978755\\
4.74	1.4881234878421\\
4.765	1.5043365621342\\
4.79	1.52073164385703\\
4.82	1.54064886629576\\
4.845	1.55745157961482\\
4.87	1.5744429124098\\
4.9	1.59508447070956\\
4.925	1.61249824927558\\
4.95	1.63010750451232\\
4.98	1.65149973157762\\
5.005	1.66954679528891\\
5.03	1.68779644645364\\
5.06	1.70996664901864\\
5.085	1.72867003883269\\
5.11	1.74758338925442\\
5.14	1.77055988161275\\
5.165	1.78994348804154\\
5.19	1.80954469974103\\
5.22	1.83335683883696\\
5.245	1.85344543158242\\
5.275	1.8778496663813\\
5.3	1.89843777717643\\
5.325	1.91925703593444\\
5.355	1.94454892395968\\
5.38	1.96588590000577\\
5.405	1.98746244646658\\
5.435	2.01367433152515\\
5.46	2.03578746133237\\
5.485	2.05814889240897\\
5.515	2.08531430263156\\
5.54	2.10823187350918\\
5.565	2.13140679582857\\
5.595	2.15956048559262\\
5.62	2.183311819028\\
5.645	2.2073298847232\\
5.675	2.23650787814787\\
5.7	2.26112336651389\\
5.725	2.2860153103523\\
5.755	2.31625494641652\\
5.78	2.34176609104519\\
5.805	2.36756376893723\\
5.835	2.39890374827189\\
5.86	2.42534319893483\\
5.89	2.45746284443172\\
5.915	2.48456006126796\\
5.94	2.51196165711264\\
5.97	2.54525017692058\\
5.995	2.57333351527108\\
6.02	2.60173232725298\\
6.05	2.63623232683077\\
6.075	2.6653377299391\\
6.1	2.69477010480794\\
6.13	2.73052573915296\\
6.155	2.76069045721854\\
6.18	2.79119406309939\\
6.21	2.82825109212148\\
6.235	2.85951372901383\\
6.26	2.89112760264332\\
6.29	2.92953344854451\\
6.315	2.96193401020733\\
6.34	2.99469860587903\\
6.37	3.03450241260887\\
6.395	3.06808235721236\\
6.42	3.10203959748617\\
6.45	3.14329229233824\\
6.475	3.17809458223783\\
6.505	3.22037388670545\\
6.53	3.25604226790318\\
6.555	3.29211142346744\\
6.585	3.33592977869886\\
6.61	3.37289656597936\\
6.635	3.41027872135013\\
6.665	3.45569217000486\\
6.69	3.4940046403605\\
6.715	3.53274760121163\\
6.745	3.57981421618748\\
6.77	3.6195213590435\\
6.795	3.65967466240396\\
6.825	3.70845461949742\\
6.85	3.74960719800609\\
6.875	3.79122217422753\\
6.905	3.8417778273961\\
6.93	3.88442844190665\\
6.955	3.92755827884181\\
6.985	3.97995423814757\\
7.01	4.02415739195589\\
7.035	4.06885720141916\\
7.065	4.12316041375881\\
7.09	4.16897258119109\\
7.12	4.22462711872517\\
7.145	4.27157930050205\\
7.17	4.31905899742463\\
7.2	4.37673929426878\\
7.225	4.42540045728924\\
7.25	4.47460831413231\\
7.28	4.53438802002097\\
7.305	4.58482029217839\\
7.33	4.6358191295503\\
7.36	4.69777455436829\\
7.385	4.75004230727598\\
7.41	4.80289721426742\\
7.44	4.86710742314936\\
7.465	4.92127735234548\\
7.49	4.97605576761703\\
7.52	5.04260267942845\\
7.545	5.0987438874282\\
7.57	5.15551568317025\\
7.6	5.22448417247082\\
7.625	5.28266825476556\\
7.65	5.34150582364058\\
7.68	5.41298382627353\\
7.705	5.47328496037883\\
7.735	5.54654090615117\\
7.76	5.60834192795903\\
7.785	5.67083696561589\\
7.815	5.74675806472699\\
7.84	5.81080742639561\\
7.865	5.87557598904125\\
7.895	5.95425895711626\\
7.92	6.02063824137665\\
7.945	6.08776281899231\\
7.975	6.16930785565121\\
8	6.23810158375973\\
};
\addplot [color=mycolor1, line width=1.4pt, only marks, mark size=0.4pt, mark=*, mark options={solid, fill=mycolor1, mycolor1}, forget plot]
  table[row sep=crcr]{%
0.95	0.367342664722113\\
};
\addplot [color=mycolor7, line width=0.9pt, forget plot]
  table[row sep=crcr]{%
0.005	0.988188115174213\\
0.03	0.933297224004141\\
0.06	0.875576023531237\\
0.085	0.833292776372284\\
0.11	0.795567049673737\\
0.14	0.755514150059352\\
0.165	0.725926714105225\\
0.19	0.69935603804931\\
0.22	0.670976279361158\\
0.245	0.649907649146324\\
0.27	0.630920400573073\\
0.3	0.610583304366803\\
0.325	0.595458843918176\\
0.355	0.579258498623096\\
0.38	0.567219871699123\\
0.405	0.556380485033046\\
0.435	0.544806666872572\\
0.46	0.536246326926567\\
0.485	0.528583343970329\\
0.515	0.520469057920219\\
0.54	0.514531637412156\\
0.565	0.509281598475837\\
0.595	0.503815876262972\\
0.62	0.49990180156505\\
0.645	0.496525586739496\\
0.675	0.493131360703282\\
0.7	0.490810868856768\\
0.725	0.488919629073765\\
0.755	0.487177978508557\\
0.78	0.486137156803564\\
0.805	0.485445361576106\\
0.835	0.485046987277219\\
0.86	0.485052857999419\\
0.885	0.485347649481108\\
0.915	0.48606091144569\\
0.94	0.486938231466949\\
0.97	0.488311178233463\\
0.995	0.489708251087676\\
1.02	0.491323732295695\\
1.05	0.493536678062384\\
1.075	0.495598777545649\\
1.1	0.49785013123429\\
1.13	0.50079080224481\\
1.155	0.503432310335614\\
1.18	0.50624050022641\\
1.21	0.509821985612325\\
1.235	0.5129765203169\\
1.26	0.516280193974404\\
1.29	0.520434958373467\\
1.315	0.524050897455418\\
1.34	0.527802324151152\\
1.37	0.532477819478814\\
1.395	0.536515004696823\\
1.42	0.540677072545136\\
1.45	0.545832480246841\\
1.475	0.550259760484129\\
1.5	0.554803729778826\\
1.53	0.560407514058935\\
1.555	0.565200893125625\\
1.585	0.571098563424848\\
1.61	0.576132747715604\\
1.635	0.581273973716821\\
1.665	0.587582888514423\\
1.69	0.592955129631021\\
1.715	0.598430610674026\\
1.745	0.605136134934962\\
1.77	0.610835545109643\\
1.795	0.616635508884195\\
1.825	0.623727298780909\\
1.85	0.629746341756682\\
1.875	0.635864179248077\\
1.905	0.643335457567791\\
1.93	0.649669392747418\\
1.955	0.656101142307871\\
1.985	0.663948127633501\\
2.01	0.670594571311726\\
2.035	0.677338510691954\\
2.065	0.685559965209026\\
2.09	0.692518542385566\\
2.115	0.699574864518552\\
2.145	0.708171734102766\\
2.17	0.715443800951903\\
2.2	0.724300312278305\\
2.225	0.731789481846982\\
2.25	0.739377903722243\\
2.28	0.74861560416501\\
2.305	0.75642388012056\\
2.33	0.764332853537408\\
2.36	0.773957286257126\\
2.385	0.782089695547001\\
2.41	0.790324605234502\\
2.44	0.800342688502327\\
2.465	0.808805364696391\\
2.49	0.817372668807264\\
2.52	0.827792568288239\\
2.545	0.836592652941687\\
2.57	0.845499791023948\\
2.6	0.856330817488562\\
2.625	0.865476381243447\\
2.65	0.874731700732742\\
2.68	0.885984228505469\\
2.705	0.895484208326466\\
2.73	0.905096906206571\\
2.76	0.916782308636865\\
2.785	0.926646458320827\\
2.815	0.938636554117636\\
2.84	0.948757134289306\\
2.865	0.958995985968639\\
2.895	0.971440279329858\\
2.92	0.981943238265952\\
2.945	0.992568073079052\\
2.975	1.00548040387819\\
3	1.01637752469691\\
3.025	1.0274003511575\\
3.055	1.04079542507624\\
3.08	1.05209920723849\\
3.105	1.06353274611533\\
3.135	1.07742611890579\\
3.16	1.08914976700247\\
3.185	1.10100744174958\\
3.215	1.11541550981121\\
3.24	1.12757292734614\\
3.265	1.13986885930072\\
3.295	1.15480885569977\\
3.32	1.16741464318809\\
3.345	1.18016365089942\\
3.375	1.19565364611571\\
3.4	1.208723102945\\
3.43	1.22460214203451\\
3.455	1.23799961193406\\
3.48	1.25154874153485\\
3.51	1.2680101876852\\
3.535	1.28189874521808\\
3.56	1.2959442631719\\
3.59	1.3130084716955\\
3.615	1.32740532813124\\
3.64	1.34196467291501\\
3.67	1.35965287171101\\
3.695	1.37457596999847\\
3.72	1.38966731611598\\
3.75	1.40800162235289\\
3.775	1.423469651774\\
3.8	1.43911192507223\\
3.83	1.45811536442233\\
3.855	1.47414777731397\\
3.88	1.49036067232169\\
3.91	1.51005720056082\\
3.935	1.52667423096325\\
3.96	1.54347823011641\\
3.99	1.56389275694418\\
4.015	1.5811154410623\\
4.045	1.6020385495578\\
4.07	1.6196902506732\\
4.095	1.63754044906067\\
4.125	1.65922582756837\\
4.15	1.67752056259082\\
4.175	1.69602098398466\\
4.205	1.71849623987368\\
4.23	1.73745731532416\\
4.255	1.75663154207699\\
4.285	1.77992534198334\\
4.31	1.79957695638892\\
4.335	1.81944947145929\\
4.365	1.84359157430788\\
4.39	1.86395884624556\\
4.415	1.88455506171868\\
4.445	1.9095763535736\\
4.47	1.93068535123331\\
4.495	1.95203163835334\\
4.525	1.97796416908062\\
4.55	1.99984194133492\\
4.575	2.02196566204689\\
4.605	2.04884268378381\\
4.63	2.07151729276201\\
4.66	2.09906358491077\\
4.685	2.12230284019108\\
4.71	2.14580339536716\\
4.74	2.17435313089046\\
4.765	2.19843897111584\\
4.79	2.22279566491822\\
4.82	2.25238553173169\\
4.845	2.27734891531789\\
4.87	2.30259306028067\\
4.9	2.33326110618831\\
4.925	2.35913413810649\\
4.95	2.3852982055869\\
4.98	2.41708388539338\\
5.005	2.44389985697641\\
5.03	2.47101751746445\\
5.06	2.50396174201487\\
5.085	2.53175517238358\\
5.11	2.55986133743861\\
5.14	2.59400652462026\\
5.165	2.62281320375785\\
5.19	2.65194406960683\\
5.22	2.68733419736737\\
5.245	2.71719123088408\\
5.275	2.7534636029507\\
5.3	2.78406498514828\\
5.325	2.8150108407815\\
5.355	2.85260605829673\\
5.38	2.88432352433623\\
5.405	2.91639808023094\\
5.435	2.95536459264635\\
5.46	2.98823901928215\\
5.485	3.02148361410129\\
5.515	3.06187164468217\\
5.54	3.09594540472926\\
5.565	3.13040288964842\\
5.595	3.1722644986641\\
5.62	3.20758151431888\\
5.645	3.24329630696416\\
5.675	3.28668545722767\\
5.7	3.32329125537434\\
5.725	3.36030939575857\\
5.755	3.40528202059379\\
5.78	3.44322379015507\\
5.805	3.48159299867973\\
5.835	3.52820707239503\\
5.86	3.56753372378954\\
5.89	3.61531101875877\\
5.915	3.6556190723199\\
5.94	3.69638129910088\\
5.97	3.74590270758327\\
5.995	3.78768225580902\\
6.02	3.82993259854605\\
6.05	3.88126194947158\\
6.075	3.9245668410208\\
6.1	3.96835975351265\\
6.13	4.02156319966137\\
6.155	4.06644924332582\\
6.18	4.1118411612072\\
6.21	4.16698726255265\\
6.235	4.21351229756502\\
6.26	4.26056170939803\\
6.29	4.31772151946315\\
6.315	4.36594548834252\\
6.34	4.41471300927556\\
6.37	4.47396016460593\\
6.395	4.5239451886299\\
6.42	4.57449363670674\\
6.45	4.63590444956147\\
6.475	4.68771490695283\\
6.505	4.75065893776892\\
6.53	4.80376293654425\\
6.555	4.85746554346717\\
6.585	4.9227083454244\\
6.61	4.97775176076117\\
6.635	5.03341565522244\\
6.665	5.10104121702549\\
6.69	5.15809490217486\\
6.715	5.21579173093078\\
6.745	5.28588708849661\\
6.77	5.34502446700013\\
6.795	5.40482847556832\\
6.825	5.47748382135313\\
6.85	5.53878097926727\\
6.875	5.6007691051644\\
6.905	5.6760779014048\\
6.93	5.73961368279932\\
6.955	5.80386565210504\\
6.985	5.88192474809361\\
7.01	5.94778085397969\\
7.035	6.01437928133519\\
7.065	6.09528903488907\\
7.09	6.16355012566038\\
7.12	6.24647980500589\\
7.145	6.31644502109429\\
7.17	6.38719882261936\\
7.2	6.47315682193859\\
7.225	6.54567689547043\\
7.25	6.61901430896905\\
7.28	6.70811105109729\\
7.305	6.78327913210217\\
7.33	6.8592943479216\\
7.36	6.9516442507026\\
7.385	7.02955685890273\\
7.41	7.10834747434391\\
7.44	7.20406909372899\\
7.465	7.28482623922714\\
7.49	7.36649338056126\\
7.52	7.46570955880014\\
7.545	7.54941486705475\\
7.57	7.63406331586519\\
7.6	7.73690133494065\\
7.625	7.82366217617848\\
7.65	7.91140050019393\\
7.68	8.01799224071695\\
7.705	8.10791986393726\\
7.735	8.21717125734739\\
7.76	8.30934265823134\\
7.785	8.40255232995953\\
7.815	8.51579088245232\\
7.84	8.61132598251975\\
7.865	8.70793713048472\\
7.895	8.82530791306886\\
7.92	8.92432908685213\\
7.945	9.02446544875082\\
7.975	9.14611876721904\\
8	9.24875280454203\\
};
\addplot [color=mycolor1, line width=1.4pt, only marks, mark size=0.4pt, mark=*, mark options={solid, fill=mycolor1, mycolor1}, forget plot]
  table[row sep=crcr]{%
0.845	0.48501361022756\\
};
\addplot [color=black!25!mycolor4, line width=0.9pt, forget plot]
  table[row sep=crcr]{%
0.005	0.99143513371636\\
0.03	0.95110262914609\\
0.06	0.907725980713031\\
0.085	0.875308736637338\\
0.11	0.84592249636603\\
0.14	0.814248023358463\\
0.165	0.790542787513805\\
0.19	0.769042957035969\\
0.22	0.745876465182951\\
0.245	0.728560738002519\\
0.27	0.712888792617598\\
0.3	0.696059669191645\\
0.325	0.683540017970218\\
0.355	0.670159853102546\\
0.38	0.660265960072258\\
0.405	0.651421689412732\\
0.435	0.642084897493003\\
0.46	0.635285858242561\\
0.485	0.629311910751724\\
0.515	0.623152342171978\\
0.54	0.618799975758094\\
0.565	0.615107458123739\\
0.595	0.611487975677463\\
0.62	0.609102923176932\\
0.645	0.607254310923859\\
0.675	0.605699269360449\\
0.7	0.604922028636903\\
0.725	0.604587899054242\\
0.755	0.604737705078565\\
0.78	0.605295428164087\\
0.805	0.6062249280102\\
0.835	0.607804809168549\\
0.86	0.60948832803726\\
0.885	0.61148861880727\\
0.915	0.614286751329887\\
0.94	0.616934391209125\\
0.97	0.620472307231683\\
0.995	0.623708021686809\\
1.02	0.627194009924774\\
1.05	0.631694155200536\\
1.075	0.635698100445749\\
1.1	0.639924143700404\\
1.13	0.645278061483597\\
1.155	0.649967132697829\\
1.18	0.65485621463904\\
1.21	0.660978893741726\\
1.235	0.666287945131686\\
1.26	0.671779703842334\\
1.29	0.678604550842571\\
1.315	0.684482594545848\\
1.34	0.690529839231357\\
1.37	0.698004876442926\\
1.395	0.704412198203448\\
1.42	0.710978257998202\\
1.45	0.719063220335928\\
1.475	0.725969170171274\\
1.5	0.733025857766698\\
1.53	0.741689956297811\\
1.555	0.749071241647275\\
1.585	0.758119716661924\\
1.61	0.765817461269434\\
1.635	0.773656818045833\\
1.665	0.783249304047279\\
1.69	0.791396196999084\\
1.715	0.799681347562396\\
1.745	0.809804927408085\\
1.77	0.818391631188754\\
1.795	0.827114462119587\\
1.825	0.837760926501414\\
1.85	0.846781803508532\\
1.875	0.855937713178367\\
1.905	0.867102836629832\\
1.93	0.876555391061705\\
1.955	0.886142769399373\\
1.985	0.897825732137973\\
2.01	0.90771016367575\\
2.035	0.917729975634688\\
2.065	0.929932900084099\\
2.09	0.940251745694663\\
2.115	0.95070719665919\\
2.145	0.963434774796496\\
2.17	0.974192620487778\\
2.2	0.98728480936489\\
2.225	0.998348079928323\\
2.25	1.00955132839467\\
2.28	1.02318110662283\\
2.305	1.03469513669351\\
2.33	1.04635181991168\\
2.36	1.06052945383714\\
2.385	1.0725033060066\\
2.41	1.0846229342047\\
2.44	1.09936041039794\\
2.465	1.11180453920233\\
2.49	1.12439797944494\\
2.52	1.13970886965237\\
2.545	1.15263501734266\\
2.57	1.16571439637866\\
2.6	1.18161374970109\\
2.625	1.19503486324191\\
2.65	1.20861349035108\\
2.68	1.22511774791164\\
2.705	1.23904791355385\\
2.73	1.25314021997903\\
2.76	1.27026714828355\\
2.785	1.28472154077475\\
2.815	1.30228752507482\\
2.84	1.31711167293192\\
2.865	1.33210651699936\\
2.895	1.35032802264253\\
2.92	1.36570437064119\\
2.945	1.38125691283277\\
2.975	1.40015505129262\\
3	1.41610153677793\\
3.025	1.43223002610859\\
3.055	1.45182710600737\\
3.08	1.46836265992443\\
3.105	1.48508633605983\\
3.135	1.50540585179918\\
3.16	1.5225503914489\\
3.185	1.53988947939564\\
3.215	1.56095610722104\\
3.24	1.57873053492256\\
3.265	1.59670624537792\\
3.295	1.61854584593242\\
3.32	1.63697205285351\\
3.345	1.65560658752221\\
3.375	1.67824621410954\\
3.4	1.69734708866218\\
3.43	1.7205530521971\\
3.455	1.74013156189382\\
3.48	1.75993100219817\\
3.51	1.78398538803502\\
3.535	1.80427948739122\\
3.56	1.82480241739971\\
3.59	1.84973556268136\\
3.615	1.87077089200457\\
3.64	1.89204328095533\\
3.67	1.91788678606033\\
3.695	1.93969004610392\\
3.72	1.96173893102972\\
3.75	1.98852568777739\\
3.775	2.01112466422005\\
3.8	2.03397817529082\\
3.83	2.06174239847133\\
3.855	2.08516598910304\\
3.88	2.10885337762618\\
3.91	2.13763063965218\\
3.935	2.16190888411383\\
3.96	2.18646055304452\\
3.99	2.2162878215697\\
4.015	2.24145193352487\\
4.045	2.27202326998574\\
4.07	2.29781514875529\\
4.095	2.32389755040006\\
4.125	2.35558455417358\\
4.15	2.38231773295522\\
4.175	2.40935209179216\\
4.205	2.44219568832819\\
4.23	2.46990471325606\\
4.255	2.49792598108122\\
4.285	2.53196865367438\\
4.31	2.56068938324498\\
4.335	2.5897338373132\\
4.365	2.62501967739179\\
4.39	2.65478932511215\\
4.415	2.68489461145896\\
4.445	2.72146937145731\\
4.47	2.75232655064496\\
4.495	2.78353172955456\\
4.525	2.82144287836429\\
4.55	2.85342764909624\\
4.575	2.88577324272645\\
4.605	2.9250700237345\\
4.63	2.95822394195424\\
4.66	2.99850286120483\\
4.685	3.03248547585366\\
4.71	3.06685164466141\\
4.74	3.10860348842116\\
4.765	3.14382891221907\\
4.79	3.17945203536679\\
4.82	3.222731123838\\
4.845	3.25924519317918\\
4.87	3.29617163080782\\
4.9	3.34103429548799\\
4.925	3.37888454280021\\
4.95	3.41716236946281\\
4.98	3.46366702372526\\
5.005	3.50290273708627\\
5.03	3.54258180210214\\
5.06	3.59078901460088\\
5.085	3.63146129973092\\
5.11	3.67259328992779\\
5.14	3.7225658608846\\
5.165	3.76472770556173\\
5.19	3.80736621046217\\
5.22	3.85916925092188\\
5.245	3.90287559188743\\
5.275	3.95597608977955\\
5.3	4.00077718552924\\
5.325	4.04608497795364\\
5.355	4.10113129106943\\
5.38	4.14757420316143\\
5.405	4.19454249274158\\
5.435	4.25160635898766\\
5.46	4.29975161358019\\
5.485	4.34844161134459\\
5.515	4.40759739865531\\
5.54	4.45750774028389\\
5.565	4.50798290018939\\
5.595	4.56930770082375\\
5.62	4.62104817196443\\
5.645	4.67337427137443\\
5.675	4.73694799969851\\
5.7	4.79058602327683\\
5.725	4.84483124640568\\
5.755	4.91073674027837\\
5.78	4.96634220521444\\
5.805	5.02257722963991\\
5.835	5.09090035552258\\
5.86	5.14854570540652\\
5.89	5.21858238844683\\
5.915	5.27767355365143\\
5.94	5.337433878662\\
5.97	5.41004026651566\\
5.995	5.47129961593379\\
6.02	5.53325274809391\\
6.05	5.6085233975631\\
6.075	5.67203070529069\\
6.1	5.73625731719234\\
6.13	5.81429023547201\\
6.155	5.88012818610547\\
6.18	5.94671189319892\\
6.21	6.02760866213579\\
6.235	6.09586295559505\\
6.26	6.16489042217849\\
6.29	6.24875632701772\\
6.315	6.31951578714932\\
6.34	6.39107683618427\\
6.37	6.47802099896074\\
6.395	6.55137768605506\\
6.42	6.62556541293121\\
6.45	6.71570093065665\\
6.475	6.79175025801177\\
6.505	6.88414757140341\\
6.53	6.96210523622579\\
6.555	7.04094609089986\\
6.585	7.136735025487\\
6.61	7.21755428235276\\
6.635	7.2992891502484\\
6.665	7.39859420954781\\
6.69	7.48238009344847\\
6.715	7.56711518888085\\
6.745	7.67006540523917\\
6.77	7.75692677133041\\
6.795	7.84477217122679\\
6.825	7.95150126881912\\
6.85	8.04155092992178\\
6.875	8.1326207128344\\
6.905	8.24326727650301\\
6.93	8.3366221454896\\
6.955	8.43103453557208\\
6.985	8.54574218568588\\
7.01	8.6425234229852\\
7.035	8.7404009390957\\
7.065	8.85931851266808\\
7.09	8.95965167901389\\
7.12	9.08155271255939\\
7.145	9.18440302741306\\
7.17	9.28841824204324\\
7.2	9.41479271663684\\
7.225	9.52141726594865\\
7.25	9.62924937045353\\
7.28	9.76026109100701\\
7.305	9.87079806104421\\
7.33	9.98258678903783\\
7.36	10.1184055025577\\
7.385	10.2329980917603\\
7.41	10.3488882446\\
7.44	10.4896898536008\\
7.465	10.6084864525748\\
7.49	10.7286280813378\\
7.52	10.8745948648586\\
7.545	10.9977492427718\\
7.57	11.1222978366924\\
7.6	11.2736186791475\\
7.625	11.4012901767529\\
7.65	11.5304068583681\\
7.68	11.6872774865522\\
7.705	11.8196312160702\\
7.735	11.9804345317474\\
7.76	12.1161061717882\\
7.785	12.2533132238851\\
7.815	12.4200128180825\\
7.84	12.5606589845775\\
7.865	12.7028966426977\\
7.895	12.8757078715898\\
7.92	13.0215102337813\\
7.945	13.1689621960825\\
7.975	13.3481082072237\\
8	13.4992550062012\\
};
\addplot [color=mycolor1, line width=1.4pt, only marks, mark size=0.4pt, mark=*, mark options={solid, fill=mycolor1, mycolor1}, forget plot]
  table[row sep=crcr]{%
0.73	0.604572068479323\\
};
\addplot [color=mycolor8, line width=0.9pt, forget plot]
  table[row sep=crcr]{%
0.005	0.99355997261173\\
0.03	0.963035001791439\\
0.06	0.929848153195925\\
0.085	0.904815588005181\\
0.11	0.881967163630303\\
0.14	0.857195556336991\\
0.165	0.838580505695449\\
0.19	0.821663557929446\\
0.22	0.803432018003734\\
0.245	0.789832784183189\\
0.27	0.777574997573161\\
0.3	0.764508215416565\\
0.325	0.75489088750039\\
0.355	0.744763739228993\\
0.38	0.737422461086764\\
0.405	0.731012919686357\\
0.435	0.724471464553997\\
0.46	0.719917602514052\\
0.485	0.716128691103397\\
0.515	0.71253040882027\\
0.54	0.710275227780443\\
0.565	0.70865659661206\\
0.595	0.707506970894844\\
0.62	0.707173076711826\\
0.645	0.707375989807082\\
0.675	0.708290962352853\\
0.7	0.709584445102555\\
0.725	0.711336670085263\\
0.755	0.714015708430949\\
0.78	0.716706056654805\\
0.805	0.719793677018231\\
0.835	0.724000173003728\\
0.86	0.727905586360676\\
0.885	0.732159647216326\\
0.915	0.737706461669511\\
0.94	0.742682973417384\\
0.97	0.749063155401312\\
0.995	0.754708228378213\\
1.02	0.760641555491852\\
1.05	0.768129654589465\\
1.075	0.774666921873494\\
1.1	0.781466293854104\\
1.13	0.789961694139885\\
1.155	0.797313741758421\\
1.18	0.804907182294063\\
1.21	0.814330194126097\\
1.235	0.822435802145267\\
1.26	0.830766489482715\\
1.29	0.841054346630566\\
1.315	0.849865429899482\\
1.34	0.858888866449528\\
1.37	0.869992605066891\\
1.395	0.879471821833831\\
1.42	0.889153616710095\\
1.45	0.901035624334028\\
1.475	0.911154493873324\\
1.5	0.921468613859629\\
1.53	0.934100691970823\\
1.555	0.94483810348428\\
1.585	0.957973641953048\\
1.61	0.969127308451863\\
1.635	0.980468333167837\\
1.665	0.994323560227182\\
1.69	1.00607366015688\\
1.715	1.01800860014459\\
1.745	1.03257381265852\\
1.77	1.04491379093989\\
1.795	1.05743741406885\\
1.825	1.07270799542855\\
1.85	1.08563532426642\\
1.875	1.09874624654359\\
1.905	1.11472197745089\\
1.93	1.12823762217623\\
1.955	1.14193780750881\\
1.985	1.1586223072958\\
2.01	1.17273029303348\\
2.035	1.1870246471445\\
2.065	1.20442491997052\\
2.09	1.21913197947249\\
2.115	1.2340280197773\\
2.145	1.25215408560409\\
2.17	1.26746937282156\\
2.2	1.28610165532988\\
2.225	1.30184155435748\\
2.25	1.31777641154964\\
2.28	1.33715742507066\\
2.305	1.35352586149548\\
2.33	1.37009363019516\\
2.36	1.39024016463062\\
2.385	1.40725176386671\\
2.41	1.42446763196515\\
2.44	1.44539863845059\\
2.465	1.4630697848176\\
2.49	1.48095066356535\\
2.52	1.50268711775358\\
2.545	1.52103584908495\\
2.57	1.53960027501185\\
2.6	1.56216506847666\\
2.625	1.58121099277608\\
2.65	1.60047905048709\\
2.68	1.6238969057323\\
2.705	1.64366113637735\\
2.73	1.66365439893959\\
2.76	1.68795180434193\\
2.785	1.70845691098272\\
2.815	1.73337526658832\\
2.84	1.75440355692033\\
2.865	1.77567352974471\\
2.895	1.80152001178386\\
2.92	1.82333047254511\\
2.945	1.84539069800291\\
2.975	1.87219634161298\\
3	1.89481531029849\\
3.025	1.91769255672061\\
3.055	1.94549005183481\\
3.08	1.96894524299221\\
3.105	1.99266765481269\\
3.135	2.02149134246931\\
3.16	2.04581184720982\\
3.185	2.07040894673246\\
3.215	2.10029482405007\\
3.24	2.12551111648294\\
3.265	2.15101381231425\\
3.295	2.18199954515237\\
3.32	2.20814349497295\\
3.345	2.23458409639881\\
3.375	2.26670903888768\\
3.4	2.29381392948787\\
3.43	2.32674579655174\\
3.455	2.35453138646913\\
3.48	2.38263193804694\\
3.51	2.41677332758117\\
3.535	2.44557929627957\\
3.56	2.47471169351559\\
3.59	2.51010663990232\\
3.615	2.53997019715342\\
3.64	2.57017212110525\\
3.67	2.6068664686482\\
3.695	2.63782634482033\\
3.72	2.66913700874448\\
3.75	2.70717845620535\\
3.775	2.73927494064875\\
3.8	2.7717351295161\\
3.83	2.81117327900536\\
3.855	2.84444826199784\\
3.88	2.87810037552002\\
3.91	2.91898678552875\\
3.935	2.95348380362373\\
3.96	2.98837190263565\\
3.99	3.03076014505144\\
4.015	3.06652442970435\\
4.045	3.10997733309503\\
4.07	3.1466400067737\\
4.095	3.18371853615429\\
4.125	3.22876841150777\\
4.15	3.26677866905479\\
4.175	3.30522022297468\\
4.205	3.35192636525836\\
4.23	3.39133425857381\\
4.255	3.43118947648571\\
4.285	3.47961343701354\\
4.31	3.5204709192817\\
4.335	3.56179236132568\\
4.365	3.61199802226542\\
4.39	3.65435901108751\\
4.415	3.69720122217461\\
4.445	3.74925487470132\\
4.47	3.7931753182714\\
4.495	3.83759489520818\\
4.525	3.89156532131862\\
4.55	3.93710326765587\\
4.575	3.9831589293406\\
4.605	4.03911748735777\\
4.63	4.08633315662176\\
4.66	4.14370133015903\\
4.685	4.19210653114161\\
4.71	4.24106239107958\\
4.74	4.30054520625234\\
4.765	4.35073488898618\\
4.79	4.40149573355407\\
4.82	4.46317192533664\\
4.845	4.51521253033116\\
4.87	4.56784556155229\\
4.9	4.63179678928752\\
4.925	4.68575722331497\\
4.95	4.74033213644698\\
4.98	4.80664308742829\\
5.005	4.86259481100123\\
5.03	4.91918388308477\\
5.06	4.98794238023763\\
5.085	5.04595949853803\\
5.11	5.10463768041577\\
5.14	5.17593479417054\\
5.165	5.23609415150337\\
5.19	5.29693916339679\\
5.22	5.37086932788337\\
5.245	5.43325060575902\\
5.275	5.50904757740387\\
5.3	5.573004170177\\
5.325	5.63768998912735\\
5.355	5.7162873349353\\
5.38	5.78260703000873\\
5.405	5.84968305524656\\
5.435	5.93118485977865\\
5.46	5.99995547892316\\
5.485	6.06951052995867\\
5.515	6.15402470927973\\
5.54	6.22533730628973\\
5.565	6.29746347063008\\
5.595	6.38510191047558\\
5.62	6.45905088774699\\
5.645	6.53384363881297\\
5.675	6.62472233785903\\
5.7	6.70140556727904\\
5.725	6.77896388666527\\
5.755	6.87320310507554\\
5.78	6.95272205336044\\
5.805	7.03314855759105\\
5.835	7.13087297102031\\
5.86	7.2133328296411\\
5.89	7.31352800708631\\
5.915	7.3980727607884\\
5.94	7.48358258065104\\
5.97	7.58748380345942\\
5.995	7.67515580200144\\
6.02	7.76382863302134\\
6.05	7.87157325077948\\
6.075	7.96248837699979\\
6.1	8.054441412766\\
6.13	8.16617180858961\\
6.155	8.260450191477\\
6.18	8.35580491914738\\
6.21	8.47166869177397\\
6.235	8.56943486008178\\
6.26	8.66831721570893\\
6.29	8.78846736817942\\
6.315	8.88985040958898\\
6.34	8.99239093906425\\
6.37	9.1169860742226\\
6.395	9.22211980025402\\
6.42	9.32845382605072\\
6.45	9.45765834908767\\
6.475	9.56668146595436\\
6.505	9.69915342652255\\
6.53	9.81093358818046\\
6.555	9.92398987379587\\
6.585	10.0613624017685\\
6.61	10.1772776164963\\
6.635	10.2945161038447\\
6.665	10.4369702591676\\
6.69	10.5571732825968\\
6.715	10.6787484496936\\
6.745	10.8264719066004\\
6.77	10.9511210739196\\
6.795	11.0771930405774\\
6.825	11.2303803265889\\
6.85	11.3596397541626\\
6.875	11.4903744858319\\
6.905	11.6492272293389\\
6.93	11.7832670235079\\
6.955	11.9188365424138\\
6.985	12.0835637291113\\
7.01	12.2225602025155\\
7.035	12.3631428063335\\
7.065	12.5339610448555\\
7.09	12.6780969404695\\
7.12	12.8532325198369\\
7.145	13.0010112258871\\
7.17	13.1504759490693\\
7.2	13.3320861644476\\
7.225	13.4853279035093\\
7.25	13.6403177623587\\
7.28	13.8286410998469\\
7.305	13.9875470695138\\
7.33	14.148265523791\\
7.36	14.3435491566292\\
7.385	14.5083278824375\\
7.41	14.6749858013837\\
7.44	14.8774859030848\\
7.465	15.0483535026193\\
7.49	15.2211694319772\\
7.52	15.431151500645\\
7.545	15.6083319570623\\
7.57	15.7875323956042\\
7.6	16.0052715897762\\
7.625	16.1889970350803\\
7.65	16.3748167210056\\
7.68	16.6005982074962\\
7.705	16.7911092159049\\
7.735	17.0225905704088\\
7.76	17.2179107375363\\
7.785	17.4154566292929\\
7.815	17.6554851998\\
7.84	17.8580168943757\\
7.865	18.0628560503812\\
7.895	18.311745721809\\
7.92	18.5217537422369\\
7.945	18.7341539298719\\
7.975	18.99222998852\\
8	19.2099887503896\\
};
\addplot [color=mycolor1, line width=1.4pt, only marks, mark size=0.4pt, mark=*, mark options={solid, fill=mycolor1, mycolor1}, forget plot]
  table[row sep=crcr]{%
0.625	0.707171590419843\\
};
\addplot [color=mycolor9, line width=0.9pt, forget plot]
  table[row sep=crcr]{%
0.005	0.994944927918554\\
0.03	0.970923661047974\\
0.06	0.944708754529787\\
0.085	0.92488563283485\\
0.11	0.906774570652525\\
0.14	0.887149147387562\\
0.165	0.872435257452983\\
0.19	0.859116328734616\\
0.22	0.844858572788021\\
0.245	0.83432548095364\\
0.27	0.824943695670296\\
0.3	0.815115101292134\\
0.325	0.808046620588487\\
0.355	0.800829657753971\\
0.38	0.795811258856891\\
0.405	0.791648810975598\\
0.435	0.787723767173696\\
0.46	0.785298190021579\\
0.485	0.783602056637182\\
0.515	0.782481971349988\\
0.54	0.782274478812734\\
0.565	0.782695817022138\\
0.595	0.783993470350333\\
0.62	0.785705431829291\\
0.645	0.787965738135322\\
0.675	0.791371416140704\\
0.7	0.794763568942202\\
0.725	0.798639394590715\\
0.755	0.803904266067633\\
0.78	0.808784146712095\\
0.805	0.814095577147455\\
0.835	0.821019113789338\\
0.86	0.827231509387329\\
0.885	0.833833382042021\\
0.915	0.842253734027874\\
0.94	0.849673351258404\\
0.97	0.859045363352585\\
0.995	0.86723505196531\\
1.02	0.875760834825228\\
1.05	0.886424318721225\\
1.075	0.895662407124938\\
1.1	0.905213094494546\\
1.13	0.917077676638131\\
1.155	0.927294434842606\\
1.18	0.937805060077733\\
1.21	0.950798689151932\\
1.235	0.961938705082073\\
1.26	0.973357818220796\\
1.29	0.987423681486295\\
1.315	0.999443482498763\\
1.34	1.01173093168079\\
1.37	1.02682498703354\\
1.395	1.03969113254593\\
1.42	1.05281628175327\\
1.45	1.06890527147378\\
1.475	1.08259281511955\\
1.5	1.09653310496623\\
1.53	1.11359294913347\\
1.555	1.12808419289357\\
1.585	1.14580166770807\\
1.61	1.16083846431034\\
1.635	1.17612194494087\\
1.665	1.19478692954674\\
1.69	1.21061126222859\\
1.715	1.22668092368683\\
1.745	1.24628818412224\\
1.77	1.26289731333464\\
1.795	1.27975182906684\\
1.825	1.30030156260115\\
1.85	1.31769708427453\\
1.875	1.33533930096689\\
1.905	1.35683651211733\\
1.93	1.3750238715725\\
1.955	1.39346035188503\\
1.985	1.41591433930717\\
2.01	1.43490243145464\\
2.035	1.45414307787471\\
2.065	1.47756701227855\\
2.09	1.49736785277807\\
2.115	1.51742559833029\\
2.145	1.54183617554519\\
2.17	1.5624646306505\\
2.2	1.58756490355904\\
2.225	1.60877234129274\\
2.25	1.63024590646523\\
2.28	1.65636822175739\\
2.305	1.67843422751786\\
2.33	1.70077284679569\\
2.36	1.72794210461447\\
2.385	1.75088821906863\\
2.41	1.77411413866032\\
2.44	1.80235794971406\\
2.465	1.82620793208786\\
2.49	1.850345582863\\
2.52	1.87969413806536\\
2.545	1.90447386592027\\
2.57	1.9295497704334\\
2.6	1.9600357384077\\
2.625	1.98577312940745\\
2.65	2.0118158298655\\
2.68	2.04347427828502\\
2.705	2.07019923081221\\
2.73	2.09723923502549\\
2.76	2.13010757246373\\
2.785	2.15785192248976\\
2.815	2.19157495515358\\
2.84	2.22003960090093\\
2.865	2.24883709296023\\
2.895	2.28383834467763\\
2.92	2.31338042980633\\
2.945	2.34326670029121\\
2.975	2.37958971288364\\
3	2.41024617043402\\
3.025	2.44125874173027\\
3.055	2.47894932053986\\
3.08	2.51075897115037\\
3.105	2.54293725537672\\
3.135	2.58204347742071\\
3.16	2.61504703906659\\
3.185	2.64843235039589\\
3.215	2.68900458244123\\
3.24	2.72324468814943\\
3.265	2.75788026285459\\
3.295	2.79997118845506\\
3.32	2.83549241124076\\
3.345	2.87142343475553\\
3.375	2.91508808944653\\
3.4	2.95193697384843\\
3.43	2.9967167346551\\
3.455	3.0345064282438\\
3.48	3.07273152793939\\
3.51	3.11918334593868\\
3.535	3.15838382319573\\
3.56	3.19803574426616\\
3.59	3.24622120663027\\
3.615	3.28688451253017\\
3.64	3.32801596202834\\
3.67	3.37799919854607\\
3.695	3.4201795153479\\
3.72	3.46284535488089\\
3.75	3.51469310425542\\
3.775	3.55844680839805\\
3.8	3.60270411249307\\
3.83	3.65648579430554\\
3.855	3.70187151795666\\
3.88	3.74777963703488\\
3.91	3.80356742916528\\
3.935	3.85064612627327\\
3.96	3.89826675408046\\
3.99	3.95613567544093\\
4.015	4.00497069217384\\
4.045	4.06431542502895\\
4.07	4.11439593603555\\
4.095	4.16505308552651\\
4.125	4.2266122318071\\
4.15	4.27856159401126\\
4.175	4.33110923535261\\
4.205	4.39496589916714\\
4.23	4.4488542700249\\
4.255	4.50336339148361\\
4.285	4.56960386745457\\
4.31	4.62550409325298\\
4.335	4.68204839970216\\
4.365	4.75076227962283\\
4.39	4.80874998581279\\
4.415	4.86740599038993\\
4.445	4.93868627496468\\
4.47	4.9988399607429\\
4.495	5.05968708076194\\
4.525	5.13363029647415\\
4.55	5.19603143365374\\
4.575	5.25915209082993\\
4.605	5.33585841199947\\
4.63	5.40059154822065\\
4.66	5.47925756976824\\
4.685	5.54564464936015\\
4.71	5.61279751572465\\
4.74	5.69440431609593\\
4.765	5.76327334570276\\
4.79	5.83293696890315\\
4.82	5.91759517188209\\
4.845	5.98903951134135\\
4.87	6.0613083303864\\
4.9	6.14913270399389\\
4.925	6.22324920843374\\
4.95	6.29822119607711\\
4.98	6.38933080025341\\
5.005	6.46621994481184\\
5.03	6.54399673379781\\
5.06	6.63851507392545\\
5.085	6.71828108338971\\
5.11	6.79896809764769\\
5.14	6.89702328405488\\
5.165	6.97977426753095\\
5.19	7.06348085844912\\
5.22	7.16520577208354\\
5.245	7.25105386285054\\
5.275	7.35538135655857\\
5.3	7.44342591521446\\
5.325	7.53248741591503\\
5.355	7.6407201737388\\
5.38	7.7320605970316\\
5.405	7.82445613036208\\
5.435	7.93674072488396\\
5.46	8.03150069589198\\
5.485	8.12735536095657\\
5.515	8.24384380256304\\
5.54	8.34215159165214\\
5.565	8.44159512616461\\
5.595	8.56244506016166\\
5.62	8.66443369131168\\
5.645	8.76760063983868\\
5.675	8.89297555088203\\
5.7	8.99878297419016\\
5.725	9.10581286262039\\
5.755	9.23588228670726\\
5.78	9.34565155726924\\
5.805	9.4566890737275\\
5.835	9.59162881801302\\
5.86	9.70550828138286\\
5.89	9.84390172043666\\
5.915	9.96069583447908\\
5.94	10.0788393189477\\
5.97	10.2224146122747\\
5.995	10.3435817880564\\
6.02	10.4661488039522\\
6.05	10.6150997650492\\
6.075	10.7408035386967\\
6.1	10.867959484989\\
6.13	11.0224870810218\\
6.155	11.1528970236768\\
6.18	11.2848134020864\\
6.21	11.4451260140066\\
6.235	11.5804179512739\\
6.26	11.7172725881783\\
6.29	11.8835862802992\\
6.315	12.0239425196623\\
6.34	12.1659197960674\\
6.37	12.3384585954265\\
6.395	12.4840681619131\\
6.42	12.631359251751\\
6.45	12.810355437588\\
6.475	12.9614143179727\\
6.505	13.1449891287785\\
6.53	13.2999118387517\\
6.555	13.4566232345921\\
6.585	13.6470669568508\\
6.61	13.8077862043475\\
6.635	13.970360794161\\
6.665	14.1679294016255\\
6.69	14.3346611625653\\
6.715	14.5033173734743\\
6.745	14.7082762529923\\
6.77	14.8812444442984\\
6.795	15.0562087334509\\
6.825	15.2688330268101\\
6.85	15.448269794802\\
6.875	15.6297769409581\\
6.905	15.8503518990634\\
6.93	16.0364979181717\\
6.955	16.2247913237739\\
6.985	16.4536126735445\\
7.01	16.6467174558338\\
7.035	16.8420494599156\\
7.065	17.0794237844069\\
7.09	17.2797460001864\\
7.12	17.5231841665519\\
7.145	17.7286233347126\\
7.17	17.9364311446659\\
7.2	18.1889653660141\\
7.225	18.4020801722783\\
7.25	18.6176515715356\\
7.28	18.8796195975172\\
7.305	19.1006950041631\\
7.33	19.3243181786869\\
7.36	19.5960701385096\\
7.385	19.8254015510651\\
7.41	20.0573752468337\\
7.44	20.339274097671\\
7.465	20.5771677431376\\
7.49	20.8178016485578\\
7.52	21.1102236400445\\
7.545	21.3569969586648\\
7.57	21.6066121006611\\
7.6	21.9099472559684\\
7.625	22.1659293070196\\
7.65	22.42485846132\\
7.68	22.7395110754908\\
7.705	23.0050429575944\\
7.735	23.3277185105724\\
7.76	23.6000202297472\\
7.785	23.8754555162115\\
7.815	24.2101645859305\\
7.84	24.4926202610819\\
7.865	24.7783254508063\\
7.895	25.1255132971997\\
7.92	25.4184985834655\\
7.945	25.7148535679524\\
7.975	26.0749817193318\\
8	26.3788859951692\\
};
\addplot [color=mycolor1, line width=1.4pt, only marks, mark size=0.4pt, mark=*, mark options={solid, fill=mycolor1, mycolor1}, forget plot]
  table[row sep=crcr]{%
0.535	0.78226477063024\\
};
\addplot [color=black!50!mycolor2, line width=0.9pt, forget plot]
  table[row sep=crcr]{%
0.005	0.995846249400233\\
0.03	0.976092736561551\\
0.06	0.954520182300844\\
0.085	0.938213760202286\\
0.11	0.923337821957493\\
0.14	0.907267952872717\\
0.165	0.895278537101997\\
0.19	0.884494740617598\\
0.22	0.873061607798609\\
0.245	0.864725362002574\\
0.27	0.857417545738931\\
0.3	0.849939006258308\\
0.325	0.844730961135641\\
0.355	0.839649189390989\\
0.38	0.836343239893616\\
0.405	0.83384413435971\\
0.435	0.831864163828843\\
0.46	0.831027500677612\\
0.485	0.830899822023262\\
0.515	0.831645126796661\\
0.54	0.832985905431917\\
0.565	0.834956242492797\\
0.595	0.838121264615318\\
0.62	0.841402266383925\\
0.645	0.845248094989634\\
0.675	0.850583835929961\\
0.7	0.855611543766713\\
0.725	0.861151186678745\\
0.755	0.868454210165146\\
0.78	0.875070434034673\\
0.805	0.882155334817049\\
0.835	0.891259270093497\\
0.86	0.899334651382368\\
0.885	0.907843361438377\\
0.915	0.918612266930605\\
0.94	0.928041187569231\\
0.97	0.93988917506573\\
0.995	0.950197769241185\\
1.02	0.960894324693448\\
1.05	0.974232815543245\\
1.075	0.985759724254962\\
1.1	0.997654511674233\\
1.13	1.0124062761447\\
1.155	1.02509191236456\\
1.18	1.03812945001454\\
1.21	1.054233044023\\
1.235	1.06803026085836\\
1.26	1.08216690112175\\
1.29	1.09957433468559\\
1.315	1.11444661564381\\
1.34	1.12964884740108\\
1.37	1.14832368463261\\
1.395	1.16424366627089\\
1.42	1.18048672696335\\
1.45	1.20040253020261\\
1.475	1.21735079052314\\
1.5	1.23461752640166\\
1.53	1.2557565768245\\
1.555	1.27372063768858\\
1.585	1.29569441158009\\
1.61	1.31435280624098\\
1.635	1.33332630932105\\
1.665	1.35651034621739\\
1.69	1.37617700211709\\
1.715	1.39615902644154\\
1.745	1.42055429974838\\
1.77	1.44123169949923\\
1.795	1.46222619524509\\
1.825	1.4878393180631\\
1.85	1.50953448289912\\
1.875	1.53154980118495\\
1.905	1.55839249689613\\
1.93	1.58111657279917\\
1.955	1.60416507517463\\
1.985	1.63225373787407\\
2.01	1.65602165207535\\
2.035	1.68011938655073\\
2.065	1.70947471662083\\
2.09	1.7343048930553\\
2.115	1.75947132600835\\
2.145	1.7901180289004\\
2.17	1.81603215333436\\
2.2	1.8475829672485\\
2.225	1.87425648691291\\
2.25	1.90127931412493\\
2.28	1.93417160430288\\
2.305	1.96197254304002\\
2.33	1.99013173981301\\
2.36	2.02439987062744\\
2.385	2.05335782225223\\
2.41	2.08268384100941\\
2.44	2.11836547507003\\
2.465	2.14851274768724\\
2.49	2.17903872649968\\
2.52	2.21617471451275\\
2.545	2.24754624656017\\
2.57	2.27930793198001\\
2.6	2.31794222870514\\
2.625	2.35057552573226\\
2.65	2.3836112154949\\
2.68	2.42379081846924\\
2.705	2.45772590804331\\
2.73	2.49207641117188\\
2.76	2.5338513194174\\
2.785	2.569130721975\\
2.815	2.61203295385888\\
2.84	2.6482625246102\\
2.865	2.68493123842855\\
2.895	2.72951989672083\\
2.92	2.76717115065044\\
2.945	2.8052766494073\\
2.975	2.85160972546378\\
3	2.89073184259325\\
3.025	2.93032408423346\\
3.055	2.97846255710088\\
3.08	3.01910721380344\\
3.105	3.06023865963071\\
3.135	3.11024652291688\\
3.16	3.15246791750936\\
3.185	3.19519356060527\\
3.215	3.24713786020679\\
3.24	3.29099274733284\\
3.265	3.3353701495441\\
3.295	3.38932103138434\\
3.32	3.43486876451948\\
3.345	3.48095810138326\\
3.375	3.53698886844465\\
3.4	3.58429144996942\\
3.43	3.64179633139699\\
3.455	3.69034274111556\\
3.48	3.73946487949586\\
3.51	3.79918069098519\\
3.535	3.8495927473154\\
3.56	3.90060192100812\\
3.59	3.96261079859628\\
3.615	4.01495792478288\\
3.64	4.06792445256361\\
3.67	4.13231197141722\\
3.695	4.18666648296472\\
3.72	4.2416636004263\\
3.75	4.30851886912528\\
3.775	4.36495605442503\\
3.8	4.42205999645066\\
3.83	4.49147575838256\\
3.855	4.55007396520804\\
3.88	4.60936405427915\\
3.91	4.68143679597324\\
3.935	4.74227752372607\\
3.96	4.80383626394168\\
3.99	4.87866633023944\\
4.015	4.94183432786161\\
4.045	5.01862032586244\\
4.07	5.08343922059176\\
4.095	5.14902259009709\\
4.125	5.22874433668645\\
4.15	5.2960411749572\\
4.175	5.36413148055662\\
4.205	5.44690031265028\\
4.23	5.51676910514426\\
4.255	5.58746148797164\\
4.285	5.67339308368905\\
4.31	5.7459315001503\\
4.335	5.81932479862056\\
4.365	5.90853932404685\\
4.39	5.98384881753495\\
4.415	6.06004569338052\\
4.445	6.1526679561712\\
4.47	6.23085389304235\\
4.495	6.3099609630957\\
4.525	6.40612057301575\\
4.55	6.48729236871456\\
4.575	6.56942034253657\\
4.605	6.66925187904631\\
4.63	6.75352313979703\\
4.66	6.85595985845737\\
4.685	6.94243015025441\\
4.71	7.02991882078512\\
4.74	7.13626632865315\\
4.765	7.22603771360842\\
4.79	7.31686623102314\\
4.82	7.42727335008135\\
4.845	7.52047144633027\\
4.87	7.61476690028139\\
4.9	7.72938810343638\\
4.925	7.82614329552338\\
4.95	7.92403759484205\\
4.98	8.04303320858812\\
5.005	8.14348081862277\\
5.03	8.24511086457638\\
5.06	8.36864727969604\\
5.085	8.47292774528367\\
5.11	8.57843561164939\\
5.14	8.70668550215625\\
5.165	8.81494456126242\\
5.19	8.92447768140508\\
5.22	9.05762023199605\\
5.245	9.17000911496802\\
5.275	9.30662280667032\\
5.3	9.4219416183905\\
5.325	9.53861724734433\\
5.355	9.68044139943347\\
5.38	9.80015825001108\\
5.405	9.92128344582277\\
5.435	10.0685159361113\\
5.46	10.1927978259465\\
5.485	10.3185415230752\\
5.515	10.4713876565301\\
5.54	10.6004078519815\\
5.565	10.7309453213845\\
5.595	10.8896181010458\\
5.62	11.0235563630568\\
5.645	11.1590694434662\\
5.675	11.3237898514977\\
5.7	11.4628326727716\\
5.725	11.6035100105951\\
5.755	11.7745072998036\\
5.78	11.9188481508153\\
5.805	12.0648854489645\\
5.835	12.2423974451617\\
5.86	12.3922370292854\\
5.89	12.5743704489384\\
5.915	12.7281107207696\\
5.94	12.8836572388556\\
5.97	13.0727270159698\\
5.995	13.2323218511128\\
6.02	13.3937912758022\\
6.05	13.5900598622407\\
6.075	13.7557307408173\\
6.1	13.9233471009807\\
6.13	14.1270867362938\\
6.155	14.2990633962192\\
6.18	14.4730590718112\\
6.21	14.6845521414751\\
6.235	14.863072880717\\
6.26	15.0436889084389\\
6.29	15.2632283152226\\
6.315	15.4485403051106\\
6.34	15.636026695272\\
6.37	15.8639162441774\\
6.395	16.0562758541648\\
6.42	16.2508919190964\\
6.45	16.487446716343\\
6.475	16.6871198506165\\
6.505	16.9298208626968\\
6.53	17.1346814116283\\
6.555	17.3419438579809\\
6.585	17.593868646355\\
6.61	17.8065140320928\\
6.635	18.0216518234855\\
6.665	18.2831479576869\\
6.69	18.5038714732998\\
6.715	18.7271812625764\\
6.745	18.9986092143358\\
6.77	19.2277150382264\\
6.795	19.4595044864031\\
6.825	19.7412381013487\\
6.85	19.9790416952274\\
6.875	20.2196298739033\\
6.905	20.512056860391\\
6.93	20.7588853815936\\
6.955	21.0086031899253\\
6.985	21.3121256255431\\
7.01	21.5683183543576\\
7.035	21.8275089511034\\
7.065	22.1425438075515\\
7.09	22.4084525899638\\
7.12	22.7316521115998\\
7.145	23.0044515281031\\
7.17	23.2804412345482\\
7.2	23.6158921154616\\
7.225	23.8990311059473\\
7.25	24.1854801066617\\
7.28	24.5336420963049\\
7.305	24.8275085967341\\
7.33	25.1248092542296\\
7.36	25.4861591008683\\
7.385	25.7911553883441\\
7.41	26.0997145676813\\
7.44	26.4747466382111\\
7.465	26.7912898538259\\
7.49	27.1115294512791\\
7.52	27.5007563734167\\
7.545	27.8292790639326\\
7.57	28.1616365546842\\
7.6	28.5655898822257\\
7.625	28.9065405612673\\
7.65	29.2514695668515\\
7.68	29.6707004689421\\
7.705	30.0245441983855\\
7.735	30.4546086886123\\
7.76	30.8175950496811\\
7.785	31.1848140419702\\
7.815	31.6311326864748\\
7.84	32.0078361034659\\
7.865	32.3889303455759\\
7.895	32.8521106398836\\
7.92	33.2430436933168\\
7.945	33.6385315931461\\
7.975	34.1192034109486\\
8	34.5248975600891\\
};
\addplot [color=mycolor1, line width=1.4pt, only marks, mark size=0.4pt, mark=*, mark options={solid, fill=mycolor1, mycolor1}, forget plot]
  table[row sep=crcr]{%
0.475	0.830867556656114\\
};
\addplot [color=black!50!mycolor4, line width=0.9pt, forget plot]
  table[row sep=crcr]{%
0.005	0.996432130227977\\
0.03	0.979455895236085\\
0.06	0.960907314020902\\
0.085	0.946890237774176\\
0.11	0.934116041455791\\
0.14	0.920347042576085\\
0.165	0.910110753044427\\
0.19	0.900947425832824\\
0.22	0.891303788961486\\
0.245	0.884344852954382\\
0.27	0.878323498222011\\
0.3	0.872284139660586\\
0.325	0.868200006904739\\
0.355	0.864390139515078\\
0.38	0.862089954590414\\
0.405	0.860555609581302\\
0.435	0.859689190164708\\
0.46	0.85975133819262\\
0.485	0.860502389656786\\
0.515	0.862283453060009\\
0.54	0.864477762799883\\
0.565	0.867297942625204\\
0.595	0.871484039942424\\
0.62	0.875621723121825\\
0.645	0.880333479110777\\
0.675	0.88672540063003\\
0.7	0.892651244007556\\
0.725	0.899108549852821\\
0.755	0.907542431856876\\
0.78	0.915128772190252\\
0.805	0.923211556592043\\
0.835	0.933552739374712\\
0.86	0.942694755716625\\
0.885	0.952304533870878\\
0.915	0.964442727147802\\
0.94	0.975054721709949\\
0.97	0.988375180771245\\
0.995	0.99995657180615\\
1.02	1.01196900312217\\
1.05	1.0269452719358\\
1.075	1.03988740196823\\
1.1	1.05324449782744\\
1.13	1.06981476570858\\
1.155	1.08407018324878\\
1.18	1.09872800561073\\
1.21	1.11684396895441\\
1.235	1.13237593764127\\
1.26	1.14830080361159\\
1.29	1.16792589644619\\
1.315	1.18470701742628\\
1.34	1.20187420343713\\
1.37	1.22298214280028\\
1.395	1.24099322034137\\
1.42	1.25938589360168\\
1.45	1.28195947828528\\
1.475	1.30118858168889\\
1.5	1.32079691552947\\
1.53	1.34482703038269\\
1.555	1.36526871557855\\
1.585	1.39029870856213\\
1.61	1.4115739429829\\
1.635	1.43322860498458\\
1.665	1.45971584218216\\
1.69	1.48220742789804\\
1.715	1.50508076122294\\
1.745	1.53303413162887\\
1.77	1.55675114080411\\
1.795	1.58085368957735\\
1.825	1.61028773785241\\
1.85	1.63524382721797\\
1.875	1.66059061025516\\
1.905	1.69152511426744\\
1.93	1.71773819467484\\
1.955	1.74434839360402\\
1.985	1.77680801336154\\
2.01	1.8042999755179\\
2.035	1.83219667624058\\
2.065	1.86621066334254\\
2.09	1.89500715125181\\
2.115	1.92421713083154\\
2.145	1.95981908955587\\
2.17	1.9899493147386\\
2.2	2.02666466450541\\
2.225	2.05773055975019\\
2.25	2.08922715050033\\
2.28	2.12759652501919\\
2.305	2.16005313608578\\
2.33	2.19295201703263\\
2.36	2.23301985786376\\
2.385	2.26690531958189\\
2.41	2.30124561760384\\
2.44	2.34306017681287\\
2.465	2.37841577513033\\
2.49	2.41423974832947\\
2.52	2.45785301197461\\
2.545	2.49472312716278\\
2.57	2.53207611415135\\
2.6	2.57754374913474\\
2.625	2.61597581767482\\
2.65	2.65490620486283\\
2.68	2.70228752654158\\
2.705	2.74233201964573\\
2.73	2.78289122459298\\
2.76	2.83224918295712\\
2.785	2.87395960000336\\
2.815	2.92471490167654\\
2.84	2.96760325178094\\
2.865	3.0110361272151\\
2.895	3.0638825508935\\
2.92	3.10853399688008\\
2.945	3.1537489193666\\
2.975	3.20875923915668\\
3	3.25523548644425\\
3.025	3.30229512924271\\
3.055	3.35954583623522\\
3.08	3.40791170177804\\
3.105	3.4568818634816\\
3.135	3.51645321870647\\
3.16	3.56677667843438\\
3.185	3.6177263333542\\
3.215	3.67970243204255\\
3.24	3.73205467751917\\
3.265	3.78505603510975\\
3.295	3.84952488159055\\
3.32	3.90398038535977\\
3.345	3.95910895818732\\
3.375	4.02616255085424\\
3.4	4.08279914003912\\
3.43	4.15168523252306\\
3.455	4.20986824577234\\
3.48	4.26876718357024\\
3.51	4.34040259443029\\
3.535	4.40090586393167\\
3.56	4.46215193871511\\
3.59	4.53663989464606\\
3.615	4.59955068786098\\
3.64	4.66323230966241\\
3.67	4.74068042517292\\
3.695	4.80608970174377\\
3.72	4.87229900594753\\
3.75	4.95281941072247\\
3.775	5.02082193104488\\
3.8	5.08965488915708\\
3.83	5.1733643638826\\
3.855	5.24405880469424\\
3.88	5.31561534211724\\
3.91	5.40263546267122\\
3.935	5.47612453960148\\
3.96	5.55050866034473\\
3.99	5.64096595032344\\
4.015	5.71735654737667\\
4.045	5.81025293280592\\
4.07	5.88870255727449\\
4.095	5.96810586325817\\
4.125	6.06466459832139\\
4.15	6.14620594538233\\
4.175	6.22873759730939\\
4.205	6.32909932289963\\
4.23	6.41385117454438\\
4.255	6.49963140807156\\
4.285	6.60394234749555\\
4.31	6.69202819191342\\
4.335	6.78118199890373\\
4.365	6.88959414991903\\
4.39	6.98114234401166\\
4.415	7.07379963732058\\
4.445	7.18647097289457\\
4.47	7.28161491210175\\
4.495	7.37791069771455\\
4.525	7.49500537566523\\
4.55	7.59388367124832\\
4.575	7.69395822792189\\
4.605	7.8156468096225\\
4.63	7.91840347354432\\
4.66	8.04335273954823\\
4.685	8.14886221844921\\
4.71	8.25564685604991\\
4.74	8.38549300862203\\
4.765	8.49513666340716\\
4.79	8.60610467405024\\
4.82	8.74103664029043\\
4.845	8.85497397434715\\
4.87	8.97028676025428\\
4.9	9.11050075917102\\
4.925	9.22889742719006\\
4.95	9.348722610684\\
4.98	9.49442241741922\\
5.005	9.61745044854282\\
5.03	9.74196209794003\\
5.06	9.89335931848513\\
5.085	10.021197348305\\
5.11	10.1505762127161\\
5.14	10.3078905694095\\
5.165	10.4407240810365\\
5.19	10.575157834335\\
5.22	10.7386174624973\\
5.245	10.8766390370102\\
5.275	11.0444603946958\\
5.3	11.1861642873017\\
5.325	11.3295738638493\\
5.355	11.5039454128625\\
5.38	11.6511791738644\\
5.405	11.8001843197545\\
5.435	11.9813583651627\\
5.46	12.134334968228\\
5.485	12.2891511593054\\
5.515	12.4773896150217\\
5.54	12.6363301412718\\
5.565	12.7971810529348\\
5.595	12.9927557954351\\
5.62	13.1578897319792\\
5.645	13.3250075411989\\
5.675	13.5282007765789\\
5.7	13.6997663262947\\
5.725	13.8733920249864\\
5.755	14.0844966697989\\
5.78	14.2627410727535\\
5.805	14.4431247929102\\
5.835	14.6624448692851\\
5.86	14.847624736168\\
5.89	15.0727752452745\\
5.915	15.2628771326461\\
5.94	15.4552587906567\\
5.97	15.68916410112\\
5.995	15.8866567017846\\
6.02	16.0865165533628\\
6.05	16.3295126072615\\
6.075	16.5346794654048\\
6.1	16.7423043603174\\
6.13	16.9947397938099\\
6.155	17.2078751646886\\
6.18	17.4235627858633\\
6.21	17.6857993980111\\
6.235	17.9072086435606\\
6.26	18.1312679071844\\
6.29	18.4036811461152\\
6.315	18.633681145149\\
6.34	18.8664326155612\\
6.37	19.14941208255\\
6.395	19.3883316560846\\
6.42	19.6301079759271\\
6.45	19.9240579480366\\
6.475	20.1722383002746\\
6.505	20.4739728922533\\
6.53	20.7287246084269\\
6.555	20.9865197839438\\
6.585	21.2999418969043\\
6.61	21.5645595360136\\
6.635	21.8323368937167\\
6.665	22.1578929866808\\
6.69	22.4327533451929\\
6.715	22.7108940496804\\
6.745	23.0490473561608\\
6.77	23.3345413844029\\
6.795	23.6234409152728\\
6.825	23.9746720620601\\
6.85	24.2712053863375\\
6.875	24.5712740661788\\
6.905	24.936081714787\\
6.93	25.24407517849\\
6.955	25.5557387201436\\
6.985	25.9346402316471\\
7.01	26.2545304562099\\
7.035	26.5782305299895\\
7.065	26.9717626541858\\
7.09	27.3040026204701\\
7.12	27.7079150645022\\
7.145	28.0489170317528\\
7.17	28.3939766826213\\
7.2	28.8134714732573\\
7.225	29.1676263737348\\
7.25	29.5259932386097\\
7.28	29.9616628449989\\
7.305	30.3294706729327\\
7.33	30.701650515038\\
7.36	31.1541095604255\\
7.385	31.5360890024589\\
7.41	31.922606489655\\
7.44	32.3924925662985\\
7.465	32.7891816889881\\
7.49	33.1905810889693\\
7.52	33.6785556027123\\
7.545	34.0905125649863\\
7.57	34.5073584657141\\
7.6	35.0141075111484\\
7.625	35.441911302633\\
7.65	35.8747893589548\\
7.68	36.4010246264187\\
7.705	36.8452758325483\\
7.735	37.3853344637737\\
7.76	37.8412532552636\\
7.785	38.3025748501642\\
7.815	38.863381357068\\
7.84	39.3368122449828\\
7.865	39.8158504130546\\
7.895	40.3981900271139\\
7.92	40.8897956449116\\
7.945	41.3872206321887\\
7.975	41.9919078185509\\
8	42.5023754644529\\
};
\addplot [color=mycolor1, line width=1.4pt, only marks, mark size=0.4pt, mark=*, mark options={solid, fill=mycolor1, mycolor1}, forget plot]
  table[row sep=crcr]{%
0.445	0.859629978517138\\
};
\addplot [color=black!50!mycolor6, line width=0.9pt, forget plot]
  table[row sep=crcr]{%
0.005	0.996812577806015\\
0.03	0.981630456151203\\
0.06	0.965012583842477\\
0.085	0.95243549341809\\
0.11	0.940961304671028\\
0.14	0.928584034191075\\
0.165	0.919379972202635\\
0.19	0.911143583105823\\
0.22	0.902486266238931\\
0.245	0.896254277973131\\
0.27	0.890881967967431\\
0.3	0.885528947237165\\
0.325	0.881947655805156\\
0.355	0.878667713883221\\
0.38	0.876754769802773\\
0.405	0.875564066910386\\
0.435	0.875059613231882\\
0.46	0.875386954413968\\
0.485	0.876374788818422\\
0.515	0.878408299511377\\
0.54	0.880791102357556\\
0.565	0.883783742131436\\
0.595	0.888160385705315\\
0.62	0.8924469057939\\
0.645	0.897301626536834\\
0.675	0.903861280628756\\
0.7	0.909926790479369\\
0.725	0.916526294930911\\
0.755	0.925137422196518\\
0.78	0.932879537819142\\
0.805	0.941127635394414\\
0.835	0.951682520659152\\
0.86	0.961017551088977\\
0.885	0.970835770367901\\
0.915	0.983246771398566\\
0.94	0.994106834103057\\
0.97	1.00775205884632\\
0.995	1.0196283017204\\
1.02	1.03195901988524\\
1.05	1.04734995265673\\
1.075	1.06066632246488\\
1.1	1.07442497032103\\
1.13	1.09151470771796\\
1.155	1.10623559763364\\
1.18	1.12138959503353\\
1.21	1.1401428680398\\
1.235	1.15624194526597\\
1.26	1.17276762757468\\
1.29	1.19315942246427\\
1.315	1.21061855403419\\
1.34	1.22850016131597\\
1.37	1.25051458450023\\
1.395	1.2693229516981\\
1.42	1.28855179427771\\
1.45	1.31218113396803\\
1.475	1.33233449938994\\
1.5	1.3529082665584\\
1.53	1.378152214628\\
1.555	1.39965231385156\\
1.585	1.42600966061344\\
1.61	1.44843950504237\\
1.635	1.47129354085965\\
1.665	1.49928001470628\\
1.69	1.52307169741495\\
1.715	1.54729191911398\\
1.745	1.57692419938188\\
1.77	1.60209323249001\\
1.795	1.62769658620546\\
1.825	1.65899686152102\\
1.85	1.68556324890105\\
1.875	1.71257108778668\\
1.905	1.74556673027613\\
1.93	1.77355471342927\\
1.955	1.80199255850367\\
1.985	1.83671585555432\\
2.01	1.86615370491057\\
2.035	1.89605104914323\\
2.065	1.93253898437716\\
2.09	1.963458829513\\
2.115	1.99484897908877\\
2.145	2.03314305573854\\
2.17	2.06558074919536\\
2.2	2.10514290379759\\
2.225	2.13864679785622\\
2.25	2.17264180814067\\
2.28	2.21409001492649\\
2.305	2.24918041468163\\
2.33	2.28477574480751\\
2.36	2.32816304765652\\
2.385	2.36488535063872\\
2.41	2.40212747916015\\
2.44	2.44751107398431\\
2.465	2.48591412443429\\
2.49	2.5248529664297\\
2.52	2.57229416197149\\
2.545	2.61243022428691\\
2.57	2.65311911017279\\
2.6	2.70268330941632\\
2.625	2.74460805816811\\
2.65	2.78710372855497\\
2.68	2.83886042874002\\
2.705	2.88263295367708\\
2.73	2.92699556877476\\
2.76	2.98101837832086\\
2.785	3.02670120303565\\
2.815	3.08232681288258\\
2.84	3.12936103627114\\
2.865	3.17702014941327\\
2.895	3.2350455024778\\
2.92	3.28410341330232\\
2.945	3.33380831140591\\
2.975	3.39431826489385\\
3	3.44547180185943\\
3.025	3.49729554300349\\
3.055	3.56037925298745\\
3.08	3.61370395917513\\
3.105	3.66772322682414\\
3.135	3.73347422889735\\
3.16	3.78904932228795\\
3.185	3.84534449877736\\
3.215	3.91386080008243\\
3.24	3.9717692534078\\
3.265	4.03042450262813\\
3.295	4.10180868476859\\
3.32	4.16213731430794\\
3.345	4.22324067322895\\
3.375	4.2976000056618\\
3.4	4.3604395686728\\
3.43	4.43690903654666\\
3.455	4.50152961101123\\
3.48	4.56697491155534\\
3.51	4.64661152613751\\
3.535	4.71390542231783\\
3.56	4.78205542529304\\
3.59	4.864979694432\\
3.615	4.93504880815142\\
3.64	5.00600676140276\\
3.67	5.09234438508862\\
3.695	5.16529498370466\\
3.72	5.23916854812269\\
3.75	5.32905057788268\\
3.775	5.40499343584842\\
3.8	5.48189482340566\\
3.83	5.57545783122527\\
3.855	5.65450837358002\\
3.88	5.73455449330829\\
3.91	5.8319407500695\\
3.935	5.9142192038754\\
3.96	5.99753181578334\\
3.99	6.09888947927615\\
4.015	6.18452103304041\\
4.045	6.28869813529521\\
4.07	6.37671022261817\\
4.095	6.46582519378076\\
4.125	6.5742376236\\
4.15	6.66582578768306\\
4.175	6.75855979826244\\
4.205	6.87137254107662\\
4.23	6.96667615702234\\
4.255	7.06317031754456\\
4.285	7.1805550232118\\
4.31	7.2797190859665\\
4.335	7.38012018782962\\
4.365	7.5022554049927\\
4.39	7.60543072759976\\
4.415	7.70989144403266\\
4.445	7.83696286487069\\
4.47	7.9443062854476\\
4.495	8.05298538126007\\
4.525	8.18518609687073\\
4.55	8.29686069452911\\
4.575	8.40992324485905\\
4.605	8.5474540115288\\
4.63	8.66362933142879\\
4.66	8.80494507833659\\
4.685	8.92431646634949\\
4.71	9.04516875415798\\
4.74	9.19217148853659\\
4.765	9.31634502662346\\
4.79	9.44205747261601\\
4.82	9.59496995779405\\
4.845	9.72413375739072\\
4.87	9.85489675048462\\
4.9	10.0139504957822\\
4.925	10.1483000475303\\
4.95	10.2843114395316\\
4.98	10.4497470216596\\
5.005	10.5894854668426\\
5.03	10.7309508472263\\
5.06	10.9030182444621\\
5.085	11.0483566572271\\
5.11	11.1954896387777\\
5.14	11.3744485782991\\
5.165	11.525606258951\\
5.19	11.6786287746004\\
5.22	11.8647490934197\\
5.245	12.0219538731855\\
5.275	12.2131595996263\\
5.3	12.3746585043548\\
5.325	12.53814706195\\
5.355	12.7369934945358\\
5.38	12.9049441963068\\
5.405	13.0749623867433\\
5.435	13.2817484313002\\
5.46	13.4564032910947\\
5.485	13.6332064871019\\
5.515	13.8482426151697\\
5.54	14.0298637539663\\
5.565	14.2137172004691\\
5.595	14.4373258801895\\
5.62	14.6261855426793\\
5.645	14.8173647247866\\
5.675	15.0498808681916\\
5.7	15.2462618003006\\
5.725	15.4450528252483\\
5.755	15.6868242524784\\
5.78	15.8910200932094\\
5.805	16.0977200867685\\
5.835	16.349108007863\\
5.86	16.5614236959142\\
5.89	16.8196396836338\\
5.915	17.0377207285752\\
5.94	17.2584729242908\\
5.97	17.5269467128956\\
5.995	17.7536889957878\\
6.02	17.9832065189274\\
6.05	18.262337833755\\
6.075	18.4980787964305\\
6.1	18.7367030803953\\
6.13	19.0269069752089\\
6.155	19.2719969952046\\
6.18	19.5200825582162\\
6.21	19.8217899877948\\
6.235	20.0765928615832\\
6.26	20.3345077954473\\
6.29	20.648166209573\\
6.315	20.9130596537944\\
6.34	21.1811861309306\\
6.37	21.5072600904765\\
6.395	21.7826362618394\\
6.42	22.0613710611388\\
6.45	22.4003428770779\\
6.475	22.6866089115969\\
6.505	23.0347373225216\\
6.53	23.3287343599959\\
6.555	23.6263129322414\\
6.585	23.9881952080629\\
6.61	24.2938046861214\\
6.635	24.6031345891899\\
6.665	24.9793040972015\\
6.69	25.2969762380162\\
6.715	25.6185129850203\\
6.745	26.0095234107924\\
6.77	26.3397255828814\\
6.795	26.6739420306288\\
6.825	27.0803681358503\\
6.85	27.4235854935967\\
6.875	27.7709724893764\\
6.905	28.1934108977199\\
6.93	28.5501470443967\\
6.955	28.9112140966076\\
6.985	29.3502841085426\\
7.01	29.7210617838865\\
7.035	30.0963377573903\\
7.065	30.5526821951095\\
7.09	30.9380439882028\\
7.12	31.4066501225581\\
7.145	31.8023639087043\\
7.17	32.202872997091\\
7.2	32.6898941139866\\
7.225	33.101154720006\\
7.25	33.5173956360398\\
7.28	34.0235422033319\\
7.305	34.4509493021252\\
7.33	34.8835287274307\\
7.36	35.4095381279356\\
7.385	35.8537141047612\\
7.41	36.303261693949\\
7.44	36.8498992261812\\
7.465	37.3114900230329\\
7.49	37.7786592576048\\
7.52	38.3467191732907\\
7.545	38.8263951642208\\
7.57	39.3118642377891\\
7.6	39.9021708173362\\
7.625	40.4006277158682\\
7.65	40.9051004518604\\
7.68	41.518509119323\\
7.705	42.0364689191003\\
7.735	42.6662732210527\\
7.76	43.1980742008067\\
7.785	43.7362861507359\\
7.815	44.3907099308434\\
7.84	44.9432945033183\\
7.865	45.5025361183194\\
7.895	46.1825243445852\\
7.92	46.7566900630543\\
7.945	47.3377681718119\\
7.975	48.0443014315448\\
8	48.6408759055089\\
};
\addplot [color=mycolor1, line width=1.4pt, only marks, mark size=0.4pt, mark=*, mark options={solid, fill=mycolor1, mycolor1}, forget plot]
  table[row sep=crcr]{%
0.435	0.875059613231882\\
};
\addplot [color=black!60!mycolor5, line width=0.9pt, forget plot]
  table[row sep=crcr]{%
0.005	0.997059383032884\\
0.03	0.983027873986184\\
0.06	0.967617164421305\\
0.085	0.955911778173906\\
0.11	0.945195983033454\\
0.14	0.933589899894362\\
0.165	0.924921036007622\\
0.19	0.917129162935545\\
0.22	0.908894157253146\\
0.245	0.902928186651185\\
0.27	0.897749469769383\\
0.3	0.892540141259851\\
0.325	0.889010188724388\\
0.355	0.885716251788451\\
0.38	0.883733397265042\\
0.405	0.882423886466132\\
0.435	0.881717330961814\\
0.46	0.881830569671846\\
0.485	0.882566127391964\\
0.515	0.884250674967864\\
0.54	0.886307413764715\\
0.565	0.888944730709274\\
0.595	0.892859837859969\\
0.62	0.896735184509707\\
0.645	0.901156931175911\\
0.675	0.907171156088137\\
0.7	0.912762889464116\\
0.725	0.918873091965438\\
0.755	0.926879026619838\\
0.78	0.934103734358368\\
0.805	0.941824207941149\\
0.835	0.951734551302004\\
0.86	0.960524623744565\\
0.885	0.96979218622978\\
0.915	0.981536663635001\\
0.94	0.99183792628625\\
0.97	1.00481026803621\\
0.995	1.01612511581142\\
1.02	1.02789494514937\\
1.05	1.04261486789991\\
1.075	1.05537488037348\\
1.1	1.06858069063033\\
1.13	1.08501287503368\\
1.155	1.09919164213855\\
1.18	1.11380966044657\\
1.21	1.13192892382858\\
1.235	1.14750814093018\\
1.26	1.16352240926764\\
1.29	1.18331258779401\\
1.315	1.20028118587347\\
1.34	1.2176827433459\\
1.37	1.23913577470043\\
1.395	1.25748920551432\\
1.42	1.27627541550781\\
1.45	1.2993905831911\\
1.475	1.31913022975264\\
1.5	1.33930422460198\\
1.53	1.3640875337413\\
1.555	1.38522022310047\\
1.585	1.41115710577328\\
1.61	1.43325410607348\\
1.635	1.45579170635045\\
1.665	1.48342071043475\\
1.69	1.50693352966489\\
1.715	1.53089261704113\\
1.745	1.56023549621912\\
1.77	1.58518378501772\\
1.795	1.61058536931014\\
1.825	1.64166909872275\\
1.85	1.66807678186239\\
1.875	1.69494608155545\\
1.905	1.72780260607519\\
1.93	1.75569768710266\\
1.955	1.78406394643004\\
1.985	1.81872997676419\\
2.01	1.84814438045387\\
2.035	1.87804072252927\\
2.065	1.91455757242452\\
2.09	1.94552701979745\\
2.115	1.97699033091737\\
2.145	2.0154037899907\\
2.17	2.04796770162568\\
2.2	2.08771392007243\\
2.225	2.12139878025463\\
2.25	2.15560020434472\\
2.28	2.19733019659587\\
2.305	2.23268424231607\\
2.33	2.26856980455984\\
2.36	2.31234098751171\\
2.385	2.34941328510267\\
2.41	2.38703315799457\\
2.44	2.4329071969916\\
2.465	2.47175034865852\\
2.49	2.5111582381571\\
2.52	2.55920103767755\\
2.545	2.59987117911938\\
2.57	2.64112432701348\\
2.6	2.69140603957168\\
2.625	2.73396285230527\\
2.65	2.7771220528645\\
2.68	2.82971710454675\\
2.705	2.87422384168766\\
2.73	2.9193534714274\\
2.76	2.97434060273423\\
2.785	3.02086412719647\\
2.815	3.07754375685226\\
2.84	3.12549450766271\\
2.865	3.17410534191397\\
2.895	3.23331986886417\\
2.92	3.28340856755245\\
2.945	3.33418095420586\\
2.975	3.39602102333417\\
3	3.44832451060452\\
3.025	3.50133648325249\\
3.055	3.5658973205877\\
3.08	3.62049628250595\\
3.105	3.67582974630943\\
3.135	3.743211259642\\
3.16	3.80019031469953\\
3.185	3.85793113633919\\
3.215	3.92823802692953\\
3.24	3.98768582224032\\
3.265	4.04792392891143\\
3.295	4.12126581394209\\
3.32	4.18327513031059\\
3.345	4.24610461751619\\
3.375	4.32259616301945\\
3.4	4.38726402867047\\
3.43	4.46599017295234\\
3.455	4.53254434068254\\
3.48	4.59997215350846\\
3.51	4.6820533098387\\
3.535	4.75143973810296\\
3.56	4.82173339632739\\
3.59	4.90729853729192\\
3.615	4.97962629843382\\
3.64	5.05289633963798\\
3.67	5.14208006473002\\
3.695	5.21746297087177\\
3.72	5.29382471936683\\
3.75	5.38676743527864\\
3.775	5.46532419144004\\
3.8	5.54489791400134\\
3.83	5.64174602497452\\
3.855	5.723600390589\\
3.88	5.80651146061756\\
3.91	5.9074175689189\\
3.935	5.99269852782318\\
3.96	6.07907759799099\\
3.99	6.18420071463911\\
4.015	6.27304265363055\\
4.045	6.38116063826267\\
4.07	6.47253160301453\\
4.095	6.56507449911246\\
4.125	6.67769297683626\\
4.15	6.77286447720524\\
4.175	6.86925409375996\\
4.205	6.98655042420329\\
4.23	7.08567232008178\\
4.255	7.18606041190794\\
4.285	7.30821923163884\\
4.31	7.41144752075394\\
4.335	7.51599204849042\\
4.365	7.6432055321088\\
4.39	7.75070257084381\\
4.415	7.85956792474559\\
4.445	7.99203605739704\\
4.47	8.10397079120546\\
4.495	8.21732802462252\\
4.525	8.35525888622985\\
4.55	8.47180708992257\\
4.575	8.58983416270385\\
4.605	8.73344422423083\\
4.63	8.85478875236494\\
4.66	9.00243334181291\\
4.685	9.1271852165641\\
4.71	9.25351626425526\\
4.74	9.4072252420435\\
4.765	9.53709879425268\\
4.79	9.66861416870318\\
4.82	9.82862812171283\\
4.845	9.96382655512334\\
4.87	10.1007319582045\\
4.9	10.2673010666383\\
4.925	10.4080356805293\\
4.95	10.5505450027459\\
4.98	10.7239293976911\\
5.005	10.8704198888691\\
5.03	11.0187555150855\\
5.06	11.1992256503426\\
5.085	11.3517004271814\\
5.11	11.5060935544863\\
5.14	11.6939305931183\\
5.165	11.8526271019883\\
5.19	12.013318070092\\
5.22	12.2088142861908\\
5.245	12.3739793507387\\
5.275	12.5749167665562\\
5.3	12.7446771815244\\
5.325	12.9165673551744\\
5.355	13.1256835275622\\
5.38	13.3023512658968\\
5.405	13.4812331980953\\
5.435	13.6988523763818\\
5.46	13.8827012483282\\
5.485	14.0688520649515\\
5.515	14.2953112311656\\
5.54	14.4866257934156\\
5.565	14.680333491853\\
5.595	14.9159828410161\\
5.62	15.1150588021831\\
5.645	15.3166226616126\\
5.675	15.5618261004963\\
5.7	15.7689707421177\\
5.725	15.9787017494899\\
5.755	16.2338374144184\\
5.78	16.4493700280842\\
5.805	16.6675913203964\\
5.835	16.9330521148443\\
5.86	17.1573044559749\\
5.89	17.4300996321992\\
5.915	17.6605459320242\\
5.94	17.893862690927\\
5.97	18.1776810725267\\
5.995	18.4174365212716\\
6.02	18.6601758120437\\
6.05	18.9554528138557\\
6.075	19.2048850477005\\
6.1	19.4574189156942\\
6.13	19.7646068846855\\
6.155	20.0240978309944\\
6.18	20.2868127794286\\
6.21	20.6063816347271\\
6.235	20.8763280522684\\
6.26	21.1496255886918\\
6.29	21.4820634857958\\
6.315	21.7628775252986\\
6.34	22.0471747287027\\
6.37	22.3929887486536\\
6.395	22.6850985344643\\
6.42	22.9808286436755\\
6.45	23.3405455081542\\
6.475	23.6443957417628\\
6.505	24.0139869125244\\
6.53	24.3261755792318\\
6.555	24.642228165433\\
6.585	25.026657641529\\
6.61	25.3513765362214\\
6.635	25.6801112327208\\
6.665	26.079962274491\\
6.69	26.4177038182146\\
6.715	26.7596189096101\\
6.745	27.1754972946284\\
6.77	27.5267729131847\\
6.795	27.8823859100733\\
6.825	28.314920782762\\
6.85	28.680261623752\\
6.875	29.0501099880132\\
6.905	29.4999547410036\\
6.93	29.8799124175307\\
6.955	30.2645543146315\\
6.985	30.7323875029405\\
7.01	31.1275348651164\\
7.035	31.5275499444745\\
7.065	32.0140762338287\\
7.09	32.4250081690109\\
7.12	32.9248087651035\\
7.145	33.3469495237808\\
7.17	33.7742837865613\\
7.2	34.2940285740591\\
7.225	34.7330100774799\\
7.25	35.1773881505789\\
7.28	35.7178569440068\\
7.305	36.1743375039181\\
7.33	36.6364254909971\\
7.36	37.1984280275263\\
7.385	37.6730912096093\\
7.41	38.1535807613234\\
7.44	38.7379578259636\\
7.465	39.2315133993581\\
7.49	39.7311226736813\\
7.52	40.3387472662618\\
7.545	40.8519321882477\\
7.57	41.3714068464072\\
7.6	42.0031853917117\\
7.625	42.5367648288632\\
7.65	43.0768790682021\\
7.68	43.7337526710076\\
7.705	44.2885210581455\\
7.735	44.9632120708113\\
7.76	45.5330244295887\\
7.785	46.1098065679922\\
7.815	46.8112630423441\\
7.84	47.4036744082286\\
7.865	48.0033266043333\\
7.895	48.7325894311958\\
7.92	49.3484784527718\\
7.945	49.9718896990045\\
7.975	50.7300394915482\\
8	51.3703183404126\\
};
\addplot [color=mycolor1, line width=1.4pt, only marks, mark size=0.4pt, mark=*, mark options={solid, fill=mycolor1, mycolor1}, forget plot]
  table[row sep=crcr]{%
0.445	0.881687018461777\\
};
\addplot [color=black!50!mycolor9, line width=0.9pt, forget plot]
  table[row sep=crcr]{%
0.005	0.997219326376158\\
0.03	0.98391990065497\\
0.06	0.969245419156423\\
0.085	0.958041827107726\\
0.11	0.947732032165614\\
0.14	0.936493784075681\\
0.165	0.928037485536111\\
0.19	0.920377432521057\\
0.22	0.91219978111902\\
0.245	0.906202069580544\\
0.27	0.900923129392242\\
0.3	0.895507854370926\\
0.325	0.891738571187541\\
0.355	0.888080593364805\\
0.38	0.88573357303512\\
0.405	0.884007393834799\\
0.435	0.882734961973493\\
0.46	0.882324432356405\\
0.485	0.882491066828259\\
0.515	0.883436159717127\\
0.54	0.884831642300519\\
0.565	0.886768678548807\\
0.595	0.889794326273952\\
0.62	0.892889400253386\\
0.645	0.896496980382918\\
0.675	0.901491468486098\\
0.7	0.906199382720154\\
0.725	0.911396174320586\\
0.755	0.9182686354559\\
0.78	0.92451895389595\\
0.805	0.931239056676209\\
0.835	0.939916069788013\\
0.86	0.94765207193788\\
0.885	0.955842611214909\\
0.915	0.966265544937542\\
0.94	0.975442187125998\\
0.97	0.987038097417593\\
0.995	0.997184334411508\\
1.02	1.00776663397281\\
1.05	1.02103741966912\\
1.075	1.03257041719236\\
1.1	1.04453216002709\\
1.13	1.05944967731101\\
1.155	1.07234863246267\\
1.18	1.0856713310857\\
1.21	1.10221633583508\\
1.235	1.11646757257207\\
1.26	1.13113962052044\\
1.29	1.14930082256623\\
1.315	1.16489705684752\\
1.34	1.18091303874812\\
1.37	1.20068632904099\\
1.395	1.21762606748954\\
1.42	1.23498619383782\\
1.45	1.25637401062665\\
1.475	1.27466105689504\\
1.5	1.29337069975969\\
1.53	1.31638150763619\\
1.555	1.33602455651005\\
1.585	1.36015900486146\\
1.61	1.38074172913076\\
1.635	1.40175403277759\\
1.665	1.42753826427584\\
1.69	1.44950185335553\\
1.715	1.47190094741128\\
1.745	1.49935770652929\\
1.77	1.52272250840948\\
1.795	1.54652998814532\\
1.825	1.57568681080868\\
1.85	1.60047710416266\\
1.875	1.62571844293339\\
1.905	1.65660745425155\\
1.93	1.68285129220745\\
1.955	1.70955569586986\\
1.985	1.7422134487986\\
2.01	1.76994253508488\\
2.035	1.79814282729391\\
2.065	1.83261017652403\\
2.09	1.86185976954567\\
2.115	1.89159230380612\\
2.145	1.92791431025667\\
2.17	1.95872315149595\\
2.2	1.99634852574902\\
2.225	2.02825360856153\\
2.25	2.06066387232282\\
2.28	2.10022930178607\\
2.305	2.13376677626577\\
2.33	2.16782403832753\\
2.36	2.20938561205353\\
2.385	2.24460338314783\\
2.41	2.28035660628952\\
2.44	2.32397447609469\\
2.465	2.36092383470769\\
2.49	2.39842536943319\\
2.52	2.4441637566345\\
2.545	2.48289938997531\\
2.57	2.52220498904557\\
2.6	2.57013220725484\\
2.625	2.61071222217688\\
2.65	2.65188106778667\\
2.68	2.70206956046493\\
2.705	2.74455551877343\\
2.73	2.78765026142002\\
2.76	2.84017665310886\\
2.785	2.88463361799312\\
2.815	2.93881396648632\\
2.84	2.98466558657209\\
2.865	3.03116217946692\\
2.895	3.08781939602483\\
2.92	3.13575991970193\\
2.945	3.18436838700418\\
2.975	3.24359068330516\\
3	3.29369489272476\\
3.025	3.3444911818931\\
3.055	3.40637124481927\\
3.08	3.45871768075941\\
3.105	3.51178152542521\\
3.135	3.57641662282165\\
3.16	3.63108767581388\\
3.185	3.68650268959993\\
3.215	3.75399478638048\\
3.24	3.81107679637513\\
3.265	3.86893057519004\\
3.295	3.93938645953414\\
3.32	3.99896982392082\\
3.345	4.05935405655618\\
3.375	4.13288547601596\\
3.4	4.19506476585349\\
3.43	4.27077808476434\\
3.455	4.33479916014386\\
3.48	4.39967324451569\\
3.51	4.47866232956252\\
3.535	4.54544873365355\\
3.56	4.61312091210209\\
3.59	4.69551158905918\\
3.615	4.76516974835789\\
3.64	4.83574789138568\\
3.67	4.92167153139267\\
3.695	4.99431254464311\\
3.72	5.06790924230821\\
3.75	5.15750294353769\\
3.775	5.23324273592063\\
3.8	5.3099754545783\\
3.83	5.40338223371381\\
3.855	5.4823417192399\\
3.88	5.56233296652813\\
3.91	5.65970196016157\\
3.935	5.74200721237767\\
3.96	5.82538471053192\\
3.99	5.92687138663909\\
4.015	6.01265381776567\\
4.045	6.11706498016675\\
4.07	6.20531706507712\\
4.095	6.29471361373887\\
4.125	6.40351979067093\\
4.15	6.49548344591379\\
4.175	6.58863677363218\\
4.205	6.70201157645378\\
4.23	6.79783351666847\\
4.255	6.89489221331172\\
4.285	7.01301645469609\\
4.31	7.1128494694281\\
4.335	7.21396826784023\\
4.365	7.33703022317653\\
4.39	7.44103339794762\\
4.415	7.54637339797845\\
4.445	7.67456907816591\\
4.47	7.78290802511377\\
4.495	7.8926369273307\\
4.525	8.02617036328258\\
4.55	8.13901746167895\\
4.575	8.25330981099531\\
4.605	8.39239335018356\\
4.63	8.50992799706495\\
4.66	8.65295487970049\\
4.685	8.77382005201345\\
4.71	8.89622892326097\\
4.74	9.04518397115405\\
4.765	9.17105607062328\\
4.79	9.29853339280214\\
4.82	9.45365287084099\\
4.845	9.58473147856793\\
4.87	9.71747931187883\\
4.9	9.87900900519961\\
4.925	10.0155017370766\\
4.95	10.1537302697519\\
4.98	10.3219258418565\\
5.005	10.4640486510488\\
5.03	10.6079765055959\\
5.06	10.7831038708344\\
5.085	10.9310813623971\\
5.11	11.0809359140718\\
5.14	11.2632716249702\\
5.165	11.4173373829588\\
5.19	11.5733550908802\\
5.22	11.7631867408071\\
5.245	11.9235836683596\\
5.275	12.1187415404231\\
5.3	12.283637061439\\
5.325	12.4506177341968\\
5.355	12.6537830695515\\
5.38	12.8254417826362\\
5.405	12.9992688160559\\
5.435	13.2107610003879\\
5.46	13.3894526438537\\
5.485	13.5703990748354\\
5.515	13.7905501611013\\
5.54	13.9765551656525\\
5.565	14.1649048490118\\
5.595	14.3940600394576\\
5.62	14.5876699351839\\
5.645	14.783717955493\\
5.675	15.0222361023824\\
5.7	15.2237539414568\\
5.725	15.4278070407946\\
5.755	15.6760611664084\\
5.78	15.8858019621213\\
5.805	16.0981789843226\\
5.835	16.356556821012\\
5.86	16.5748480036773\\
5.89	16.8404187189553\\
5.915	17.0647849066122\\
5.94	17.2919667969822\\
5.97	17.5683504512639\\
5.995	17.8018490492679\\
6.02	18.0382753500505\\
6.05	18.3259020480348\\
6.075	18.5688962540673\\
6.1	18.8149345602674\\
6.13	19.1142512933006\\
6.155	19.3671185594138\\
6.18	19.6231508881125\\
6.21	19.9346221781365\\
6.235	20.1977547543419\\
6.26	20.4641780946187\\
6.29	20.7882866626766\\
6.315	21.0620921611873\\
6.34	21.339319045026\\
6.37	21.6765665051085\\
6.395	21.9614684860143\\
6.42	22.249927580819\\
6.45	22.6008351600616\\
6.475	22.8972737392277\\
6.505	23.257885552154\\
6.53	23.5625197489984\\
6.555	23.8709522289769\\
6.585	24.2461503031384\\
6.61	24.5631031051566\\
6.635	24.8840046299281\\
6.665	25.274366685998\\
6.69	25.6041258958232\\
6.715	25.9379901303503\\
6.745	26.3441164015523\\
6.77	26.6871888247471\\
6.795	27.0345286600901\\
6.825	27.4570427492198\\
6.85	27.8139549176369\\
6.875	28.1753032027086\\
6.905	28.6148529714635\\
6.93	28.9861518937765\\
6.955	29.3620621946106\\
6.985	29.8193206863432\\
7.01	30.2055746268565\\
7.035	30.596622014485\\
7.065	31.0722884117385\\
7.09	31.4740876985068\\
7.12	31.9628291694168\\
7.145	32.3756701843272\\
7.17	32.7936280475805\\
7.2	33.3020193539809\\
7.225	33.7314542769614\\
7.25	34.1662077190329\\
7.28	34.6950233197862\\
7.305	35.1417060180087\\
7.33	35.59391671582\\
7.36	36.1439610609935\\
7.385	36.6085707162399\\
7.41	37.0789259577069\\
7.44	37.6510346264543\\
7.465	38.1342766954716\\
7.49	38.6234903508552\\
7.52	39.2185312319431\\
7.545	39.7211384426728\\
7.57	40.2299519723006\\
7.6	40.8488264887046\\
7.625	41.3715598737422\\
7.65	41.9007433734314\\
7.68	42.5443877527839\\
7.705	43.0880377214469\\
7.735	43.7492731828691\\
7.76	44.3077775992463\\
7.785	44.8731650497721\\
7.815	45.5608328770152\\
7.84	46.1416574264595\\
7.865	46.7296348995403\\
7.895	47.4447715255267\\
7.92	48.0487911943834\\
7.945	48.660244069797\\
7.975	49.4039258403004\\
8	50.0320493024431\\
};
\addplot [color=mycolor1, line width=1.4pt, only marks, mark size=0.4pt, mark=*, mark options={solid, fill=mycolor1, mycolor1}, forget plot]
  table[row sep=crcr]{%
0.465	0.882311978386633\\
};
\addplot [color=black!75!mycolor4, line width=0.9pt, forget plot]
  table[row sep=crcr]{%
0.005	0.997322858541116\\
0.03	0.984484867929324\\
0.06	0.970244794331434\\
0.085	0.959308375037595\\
0.11	0.94918343133056\\
0.14	0.938063115285451\\
0.165	0.929622234891033\\
0.19	0.921905431824525\\
0.22	0.913568306039056\\
0.245	0.907364587772887\\
0.27	0.901815760431086\\
0.3	0.895995487809219\\
0.325	0.891823589192243\\
0.355	0.887607307019799\\
0.38	0.884734579194285\\
0.405	0.882429599446599\\
0.435	0.880394829838398\\
0.46	0.879294343068193\\
0.485	0.878722856991531\\
0.515	0.878720535735103\\
0.54	0.879276607339482\\
0.565	0.880330143384047\\
0.595	0.882238966497193\\
0.62	0.884357396541437\\
0.645	0.886947602551939\\
0.675	0.890668770973687\\
0.7	0.894272868772949\\
0.725	0.898327868470099\\
0.755	0.903781106301841\\
0.78	0.908808664498382\\
0.805	0.914270270529867\\
0.835	0.921390765118095\\
0.86	0.927791723805513\\
0.885	0.934613277891373\\
0.915	0.943349298608071\\
0.94	0.951083920302269\\
0.97	0.960906495762422\\
0.995	0.96953962440024\\
1.02	0.978577052933941\\
1.05	0.98995248835683\\
1.075	0.999871755478194\\
1.1	1.01018886745602\\
1.13	1.02309233520402\\
1.155	1.03427938695727\\
1.18	1.04585986904343\\
1.21	1.06027438468399\\
1.235	1.0727171476929\\
1.26	1.08555075031945\\
1.29	1.10146633889338\\
1.315	1.1151583731672\\
1.34	1.12924030361918\\
1.37	1.14665332055068\\
1.395	1.16159329166819\\
1.42	1.17692371846029\\
1.45	1.19583629438057\\
1.475	1.21202754157285\\
1.5	1.22861118859886\\
1.53	1.24903077445425\\
1.555	1.26648095656092\\
1.585	1.28794351782963\\
1.61	1.30626576444925\\
1.635	1.32498656971486\\
1.665	1.34797977184594\\
1.69	1.3675828894235\\
1.715	1.38758979809411\\
1.745	1.41213372521012\\
1.77	1.43303568648319\\
1.795	1.45434777747659\\
1.825	1.48046676309199\\
1.85	1.50268902845016\\
1.875	1.52532882314658\\
1.905	1.55305127638671\\
1.93	1.57661865985493\\
1.955	1.60061199740143\\
1.985	1.62997026505693\\
2.01	1.65491082860373\\
2.035	1.68028676822037\\
2.065	1.71131702758939\\
2.09	1.7376619998349\\
2.115	1.76445274648229\\
2.145	1.79719492492282\\
2.17	1.82497864106623\\
2.2	1.85892284558182\\
2.225	1.88771720922582\\
2.25	1.91697706957686\\
2.28	1.95270892482417\\
2.305	1.98300690978799\\
2.33	2.01378335476707\\
2.36	2.05135263959347\\
2.385	2.08319676956935\\
2.41	2.11553326940278\\
2.44	2.15499339794473\\
2.465	2.18842922710327\\
2.49	2.22237228475303\\
2.52	2.26378031506523\\
2.545	2.29885643936869\\
2.57	2.33445560549541\\
2.6	2.37787226217199\\
2.625	2.41464034201483\\
2.65	2.45194824096168\\
2.68	2.49743795063246\\
2.705	2.535952743922\\
2.73	2.57502510990662\\
2.76	2.62265604867284\\
2.785	2.66297545362754\\
2.815	2.71212001104742\\
2.84	2.75371532849748\\
2.865	2.79590043556771\\
2.895	2.8473099399108\\
2.92	2.890814842035\\
2.945	2.9349299944593\\
2.975	2.98868301811434\\
3	3.03416412385951\\
3.025	3.0802769864494\\
3.055	3.13645611830177\\
3.08	3.18398342009235\\
3.105	3.2321650559442\\
3.135	3.29085699661629\\
3.16	3.34050394235938\\
3.185	3.39082889723054\\
3.215	3.45212456344214\\
3.24	3.5039681460601\\
3.265	3.55651454022295\\
3.295	3.62050917805516\\
3.32	3.67463003240024\\
3.345	3.72947965983304\\
3.375	3.79627296675875\\
3.4	3.85275547378705\\
3.43	3.92153307356418\\
3.455	3.97969016417377\\
3.48	4.03862256231953\\
3.51	4.11037753293164\\
3.535	4.17104747753378\\
3.56	4.23252198125714\\
3.59	4.30736665485001\\
3.615	4.37064447391106\\
3.64	4.43475739941199\\
3.67	4.51280908408243\\
3.695	4.57879399041545\\
3.72	4.6456458885518\\
3.75	4.72702703083911\\
3.775	4.79582256750316\\
3.8	4.86551836396701\\
3.83	4.950356720084\\
3.855	5.02207090559889\\
3.88	5.09472004844869\\
3.91	5.1831488643652\\
3.935	5.25789434535151\\
3.96	5.33361095997481\\
3.99	5.42576916050527\\
4.015	5.50366337185415\\
4.045	5.5984691572663\\
4.07	5.67859881056299\\
4.095	5.75976414328671\\
4.125	5.85854718313941\\
4.15	5.94203506754772\\
4.175	6.02659904447689\\
4.205	6.1295145338661\\
4.23	6.2164918064174\\
4.255	6.3045872659161\\
4.285	6.41179685664478\\
4.31	6.50240012096762\\
4.335	6.59416540887281\\
4.365	6.70583744322721\\
4.39	6.80020894726333\\
4.415	6.89578811725123\\
4.445	7.01209787199193\\
4.47	7.1103857142871\\
4.495	7.20992873693073\\
4.525	7.33105867778803\\
4.55	7.43341702265506\\
4.575	7.53708000336788\\
4.605	7.66322005034312\\
4.63	7.76980935209678\\
4.66	7.89950802118969\\
4.685	8.0091025620833\\
4.71	8.12008979865735\\
4.74	8.25513674299109\\
4.765	8.36924792627304\\
4.79	8.48480684250083\\
4.82	8.62541341735422\\
4.845	8.74421978639428\\
4.87	8.86453114727945\\
4.9	9.01091723945028\\
4.925	9.13460453786016\\
4.95	9.25985639212217\\
4.98	9.4122507405418\\
5.005	9.5410121829958\\
5.03	9.67140013727353\\
5.06	9.83004066653279\\
5.085	9.96407722049502\\
5.11	10.0998047245428\\
5.14	10.2649388917739\\
5.165	10.404459570463\\
5.19	10.5457382137186\\
5.22	10.7176233692598\\
5.245	10.8628455362878\\
5.275	11.0395265998227\\
5.3	11.1887991185858\\
5.325	11.3399488045847\\
5.355	11.5238386209585\\
5.38	11.6791992528215\\
5.405	11.8365114845778\\
5.435	12.0278958915857\\
5.46	12.1895860457862\\
5.485	12.3533051882581\\
5.515	12.5524813742943\\
5.54	12.7207520410981\\
5.565	12.8911321527962\\
5.595	13.0984090889263\\
5.62	13.2735212044865\\
5.645	13.4508264066873\\
5.675	13.6665252955858\\
5.7	13.8487501209296\\
5.725	14.0332549812189\\
5.755	14.2577097234507\\
5.78	14.447329238182\\
5.805	14.6393191688578\\
5.835	14.8728768472181\\
5.86	15.0701841584197\\
5.89	15.3102085311254\\
5.915	15.5129772127938\\
5.94	15.7182769796994\\
5.97	15.9680212938298\\
5.995	16.1789988140536\\
6.02	16.3926076890327\\
6.05	16.6524570014706\\
6.075	16.8719685225\\
6.1	17.0942156090018\\
6.13	17.3645701116731\\
6.155	17.5929535717832\\
6.18	17.824180899807\\
6.21	18.1054564976416\\
6.235	18.3430630990463\\
6.26	18.5836261187544\\
6.29	18.8762550297941\\
6.315	19.1234497451789\\
6.34	19.3737178394463\\
6.37	19.678149217781\\
6.395	19.9353113161624\\
6.42	20.1956683325984\\
6.45	20.5123689150789\\
6.475	20.7798925077001\\
6.505	21.1053083696712\\
6.53	21.3801920885192\\
6.555	21.6584866955122\\
6.585	21.9970011336708\\
6.61	22.2829466550615\\
6.635	22.5724377935864\\
6.665	22.9245681964894\\
6.69	23.2220123883317\\
6.715	23.5231421831242\\
6.745	23.8894261266249\\
6.77	24.1988228974661\\
6.795	24.5120507144459\\
6.825	24.8930467318498\\
6.85	25.2148676809354\\
6.875	25.5406707848476\\
6.905	25.9369591662186\\
6.93	26.2716942581067\\
6.955	26.6105684949976\\
6.985	27.0227521166402\\
7.01	27.3709103812988\\
7.035	27.7233708868534\\
7.065	28.1520760721936\\
7.09	28.5141863316111\\
7.12	28.954626019822\\
7.145	29.3266456790107\\
7.17	29.7032573011499\\
7.2	30.1613311479784\\
7.225	30.5482421840303\\
7.25	30.9399258652344\\
7.28	31.4163277291916\\
7.305	31.8187159708424\\
7.33	32.2260645476579\\
7.36	32.7215151988974\\
7.385	33.1399891908807\\
7.41	33.5636184825473\\
7.44	34.0788666266561\\
7.465	34.5140584938338\\
7.49	34.9546081763102\\
7.52	35.4904315166353\\
7.545	35.9429978605314\\
7.57	36.4011323734287\\
7.6	36.9583387132468\\
7.625	37.4289615422319\\
7.65	37.905371029818\\
7.68	38.4847994159614\\
7.705	38.9741871103426\\
7.735	39.569396668824\\
7.76	40.0721103070094\\
7.785	40.5809986204389\\
7.815	41.1999202920035\\
7.84	41.7226567967679\\
7.865	42.2518098949844\\
7.895	42.8953726971353\\
7.92	43.4389163902989\\
7.945	43.9891279288848\\
7.975	44.6582967226775\\
8	45.2234621801596\\
};
\addplot [color=mycolor1, line width=1.4pt, only marks, mark size=0.4pt, mark=*, mark options={solid, fill=mycolor1, mycolor1}, forget plot]
  table[row sep=crcr]{%
0.5	0.878629465830066\\
};
\addplot [color=black!80!mycolor5, line width=0.9pt, forget plot]
  table[row sep=crcr]{%
0.005	0.997389784582587\\
0.03	0.984839228855067\\
0.06	0.970843250004699\\
0.085	0.96002936805559\\
0.11	0.949956103844403\\
0.14	0.938807539459076\\
0.165	0.930270441079108\\
0.19	0.922393507782177\\
0.22	0.913782515884615\\
0.245	0.907284552939983\\
0.27	0.90138321945973\\
0.3	0.895065200948829\\
0.325	0.890417922072026\\
0.355	0.885560484096817\\
0.38	0.882096015145423\\
0.405	0.879148323868547\\
0.435	0.876276560643339\\
0.46	0.874424996946623\\
0.485	0.873054742973068\\
0.515	0.872032201384887\\
0.54	0.871687647819288\\
0.565	0.871795527797381\\
0.595	0.872511132601842\\
0.62	0.873587303094593\\
0.645	0.875092352791944\\
0.675	0.877455511090971\\
0.7	0.879882030910803\\
0.725	0.882718237702335\\
0.755	0.886655140696411\\
0.78	0.890374738833242\\
0.805	0.894488458182755\\
0.835	0.899939247445636\\
0.86	0.904905595564931\\
0.885	0.910253554620068\\
0.915	0.917170123225279\\
0.94	0.923346152024383\\
0.97	0.931247806945423\\
0.995	0.938238155858896\\
1.02	0.945594729213387\\
1.05	0.954902995430293\\
1.075	0.963057902917488\\
1.1	0.971572763875092\\
1.13	0.982263524008764\\
1.155	0.991564961874302\\
1.18	1.00122188785144\\
1.21	1.0132779651017\\
1.235	1.02371345549936\\
1.26	1.03450157638226\\
1.29	1.04791197320164\\
1.315	1.05947399807939\\
1.34	1.07138724259766\\
1.37	1.08614652151876\\
1.395	1.09883204157374\\
1.42	1.11186868574447\\
1.45	1.12797646371576\\
1.475	1.14178652508818\\
1.5	1.15594882269848\\
1.53	1.17340935695133\\
1.555	1.18834876890411\\
1.585	1.20674403936029\\
1.61	1.22246449698065\\
1.635	1.23854156755077\\
1.665	1.25830638360866\\
1.69	1.27517212694742\\
1.715	1.29239844636018\\
1.745	1.31354801865541\\
1.77	1.33157278971605\\
1.795	1.34996305350728\\
1.825	1.37251625572392\\
1.85	1.39171679986745\\
1.875	1.41128866530837\\
1.905	1.43526787411997\\
1.93	1.45566381633265\\
1.955	1.47643778701181\\
1.985	1.50186875399117\\
2.01	1.52348250037937\\
2.035	1.54548183480445\\
2.065	1.57239358567945\\
2.09	1.59525024660686\\
2.115	1.61850088748993\\
2.145	1.64692564501349\\
2.17	1.67105297696763\\
2.2	1.70053825194697\\
2.225	1.72555662577198\\
2.25	1.75098498958242\\
2.28	1.78204474754173\\
2.305	1.80838652225318\\
2.33	1.83514885571248\\
2.36	1.86782368021081\\
2.385	1.89552355608447\\
2.41	1.92365535936208\\
2.44	1.95798890606511\\
2.465	1.98708414245122\\
2.49	2.01662347497347\\
2.52	2.05266247232715\\
2.545	2.08319289176181\\
2.57	2.11418037947056\\
2.6	2.15197464141781\\
2.625	2.18398264347537\\
2.65	2.21646149553896\\
2.68	2.25606394527095\\
2.705	2.28959452934409\\
2.73	2.32361056359406\\
2.76	2.36507726747963\\
2.785	2.40017806383478\\
2.815	2.44296061922059\\
2.84	2.47916995129914\\
2.865	2.51589125986363\\
2.895	2.56063993411303\\
2.92	2.5985059318738\\
2.945	2.63690076831528\\
2.975	2.68368031486253\\
3	2.72325789137425\\
3.025	2.76338204008646\\
3.055	2.81226055857407\\
3.08	2.85360743593061\\
3.105	2.89551950966388\\
3.135	2.94656852015137\\
3.16	2.98974529351429\\
3.185	3.033506799392\\
3.215	3.08680132416544\\
3.24	3.13187153206772\\
3.265	3.1775469438174\\
3.295	3.23316559658274\\
3.32	3.2801957973557\\
3.345	3.32785263361578\\
3.375	3.38587771595033\\
3.4	3.43493757014428\\
3.43	3.49466677997518\\
3.455	3.54516408368455\\
3.48	3.59632644242122\\
3.51	3.6586096203813\\
3.535	3.71126141330094\\
3.56	3.76460241252645\\
3.59	3.82953230435517\\
3.615	3.88441703846834\\
3.64	3.94001619907486\\
3.67	4.00768965538613\\
3.695	4.06488923987729\\
3.72	4.12282957371702\\
3.75	4.19334767966217\\
3.775	4.25294759099485\\
3.8	4.31331571281622\\
3.83	4.38678392486746\\
3.855	4.44887332282261\\
3.88	4.51175956820443\\
3.91	4.58828785814501\\
3.935	4.65295970790899\\
3.96	4.71845825722195\\
3.99	4.79816126345941\\
4.015	4.86551246382274\\
4.045	4.94746708758515\\
4.07	5.01671865866462\\
4.095	5.08685017644344\\
4.125	5.17218399676645\\
4.15	5.24428773261996\\
4.175	5.31730474344373\\
4.205	5.4061457479559\\
4.23	5.48120980279951\\
4.255	5.55722183837814\\
4.285	5.64970330459795\\
4.31	5.72784029382008\\
4.335	5.80696139585619\\
4.365	5.90322207867307\\
4.39	5.98454923741413\\
4.415	6.06689811831667\\
4.445	6.16708244571389\\
4.47	6.25172179435315\\
4.495	6.33742200496759\\
4.525	6.44168028221299\\
4.55	6.52975879901823\\
4.575	6.61893890350352\\
4.605	6.72742752604584\\
4.63	6.8190773271207\\
4.66	6.93056830442333\\
4.685	7.02475276057294\\
4.71	7.12011128735617\\
4.74	7.2361109600177\\
4.765	7.33410182313474\\
4.79	7.4333120376669\\
4.82	7.55399429368308\\
4.845	7.65593843865255\\
4.87	7.75914902688252\\
4.9	7.88469470907467\\
4.925	7.9907448792304\\
4.95	8.09811046255725\\
4.98	8.22870762593335\\
5.005	8.33902265044891\\
5.03	8.4507040058315\\
5.06	8.58654818530901\\
5.085	8.70129320533858\\
5.11	8.81745749422006\\
5.14	8.95875198306814\\
5.165	9.07809868690162\\
5.19	9.19891969325626\\
5.22	9.34587583257427\\
5.245	9.47000270016051\\
5.275	9.62097814443121\\
5.3	9.748498557665\\
5.325	9.87759111472671\\
5.355	10.0346039411879\\
5.38	10.1672218167818\\
5.405	10.3014728989597\\
5.435	10.4647576017465\\
5.46	10.6026709594759\\
5.485	10.7422809867878\\
5.515	10.9120812774781\\
5.54	11.0554959163621\\
5.565	11.2006731783597\\
5.595	11.3772423326438\\
5.62	11.5263721236977\\
5.645	11.6773330756925\\
5.675	11.8609342935598\\
5.7	12.0160014838319\\
5.725	12.1729710545865\\
5.755	12.3638778343768\\
5.78	12.5251133627661\\
5.805	12.6883252740733\\
5.835	12.8868218009634\\
5.86	13.0544656261917\\
5.89	13.2583507001107\\
5.915	13.430544274169\\
5.94	13.6048457153577\\
5.97	13.8168254640135\\
5.995	13.9958536739049\\
6.02	14.1770717641932\\
6.05	14.397461136108\\
6.075	14.5835899055101\\
6.1	14.7719937592194\\
6.13	15.0011199556636\\
6.155	15.1946255499408\\
6.18	15.3904947437509\\
6.21	15.6286976820498\\
6.235	15.8298670993617\\
6.26	16.0334920679569\\
6.29	16.2811248625049\\
6.315	16.4902562397851\\
6.34	16.7019386866497\\
6.37	16.9593681481236\\
6.395	17.176771182491\\
6.42	17.3968245062326\\
6.45	17.6644316598071\\
6.475	17.8904280459658\\
6.505	18.165261030919\\
6.53	18.3973584060463\\
6.555	18.6322822902769\\
6.585	18.9179696886448\\
6.61	19.159231754378\\
6.635	19.4034302056366\\
6.665	19.7003939663021\\
6.69	19.9511769585119\\
6.715	20.2050104574463\\
6.745	20.5136888300917\\
6.77	20.7743627431484\\
6.795	21.0382056894949\\
6.825	21.3590538413269\\
6.85	21.6300029486102\\
6.875	21.9042441877583\\
6.905	22.2377348434713\\
6.93	22.5193582374964\\
6.955	22.8044016068531\\
6.985	23.1510257126662\\
7.01	23.4437378654777\\
7.035	23.7400027614964\\
7.065	24.1002701742189\\
7.09	24.4045015189874\\
7.12	24.7744546306306\\
7.145	25.0868636873306\\
7.17	25.4030612217152\\
7.2	25.787562910384\\
7.225	26.1122554006714\\
7.25	26.440883312866\\
7.28	26.8404978754627\\
7.305	27.1779501813727\\
7.33	27.5194904996021\\
7.36	27.934803915783\\
7.385	28.2855107199799\\
7.41	28.640463986465\\
7.44	29.0720847375753\\
7.465	29.4365597145642\\
7.49	29.8054456845642\\
7.52	30.2540056032738\\
7.545	30.6327821364023\\
7.57	31.0161405035198\\
7.6	31.482295655212\\
7.625	31.8759275804877\\
7.65	32.2743187291074\\
7.68	32.7587503263085\\
7.705	33.1678127040155\\
7.735	33.6652180350239\\
7.76	34.0852338406306\\
7.785	34.5103237559269\\
7.815	35.0272147622306\\
7.84	35.4636818078108\\
7.865	35.9054192009876\\
7.895	36.442549477365\\
7.92	36.8961039138699\\
7.945	37.3551324417364\\
7.975	37.9132844069982\\
8	38.3845867126442\\
};
\addplot [color=mycolor1, line width=1.4pt, only marks, mark size=0.4pt, mark=*, mark options={solid, fill=mycolor1, mycolor1}, forget plot]
  table[row sep=crcr]{%
0.545	0.871673288599548\\
};
\addplot [color=black!80!mycolor9, line width=0.9pt, forget plot]
  table[row sep=crcr]{%
0.005	0.997432976763871\\
0.03	0.985058713025646\\
0.06	0.971189106428493\\
0.085	0.9604117459313\\
0.11	0.950314322654645\\
0.14	0.939058941216902\\
0.165	0.930369738098328\\
0.19	0.922284840923145\\
0.22	0.913352394360669\\
0.245	0.906528065400268\\
0.27	0.900248407288087\\
0.3	0.893409301459442\\
0.325	0.888273025137914\\
0.355	0.882764451852276\\
0.38	0.878704754028244\\
0.405	0.875114847180939\\
0.435	0.871411438652615\\
0.46	0.868816817337329\\
0.485	0.866658710083509\\
0.515	0.86463243916245\\
0.54	0.863403492206052\\
0.565	0.862583918699803\\
0.595	0.862130384319569\\
0.62	0.862185935102065\\
0.645	0.862628648890918\\
0.675	0.863662426254312\\
0.7	0.864936014097158\\
0.725	0.866578575049284\\
0.755	0.869029710326429\\
0.78	0.871466931253676\\
0.805	0.874258260150278\\
0.835	0.878069590112067\\
0.86	0.881626054174195\\
0.885	0.885524546821625\\
0.915	0.890649623071017\\
0.94	0.895289365997157\\
0.97	0.901295477498128\\
0.995	0.906662884702333\\
1.02	0.912357119310586\\
1.05	0.919618525671371\\
1.075	0.926024273092134\\
1.1	0.932750398242081\\
1.13	0.94124229811902\\
1.155	0.948667567293058\\
1.18	0.95640837453124\\
1.21	0.966112155634521\\
1.235	0.974543066586695\\
1.26	0.98328610031343\\
1.29	0.99418862517307\\
1.315	1.00361568245388\\
1.34	1.01335272062141\\
1.37	1.02544574578496\\
1.395	1.03586338718451\\
1.42	1.04659001760177\\
1.45	1.05986971903015\\
1.475	1.07127594628574\\
1.5	1.08299121950126\\
1.53	1.09745780330103\\
1.555	1.10985387890248\\
1.585	1.12513850307531\\
1.61	1.13821740535394\\
1.635	1.15160760311582\\
1.665	1.16808773147389\\
1.69	1.18216533272205\\
1.715	1.19655673247421\\
1.745	1.21424196919133\\
1.77	1.22932717851085\\
1.795	1.24472949729866\\
1.825	1.26363255040245\\
1.85	1.2797368160464\\
1.875	1.2961622697969\\
1.905	1.31629879632197\\
1.93	1.33343598543907\\
1.955	1.35089917625679\\
1.985	1.37228765643687\\
2.01	1.39047395768136\\
2.035	1.40899178288108\\
2.065	1.43165341903731\\
2.09	1.45090726468633\\
2.115	1.47049884441877\\
2.145	1.49445748103206\\
2.17	1.51479948586679\\
2.2	1.53966498494695\\
2.225	1.56076816884725\\
2.25	1.58222108357643\\
2.28	1.60842964791397\\
2.305	1.63066060375706\\
2.33	1.6532492811487\\
2.36	1.68083159421154\\
2.385	1.70421660615781\\
2.41	1.72796797837356\\
2.44	1.75695721239688\\
2.465	1.78152463426503\\
2.49	1.80646770015105\\
2.52	1.83689950470251\\
2.545	1.86267975349095\\
2.57	1.8888455748071\\
2.6	1.92075807568529\\
2.625	1.94778363343499\\
2.65	1.97520533914156\\
2.68	2.00863914666334\\
2.705	2.03694456988589\\
2.73	2.06565736781387\\
2.76	2.1006555936726\\
2.785	2.13027752972952\\
2.815	2.1663776128458\\
2.84	2.1969270070729\\
2.865	2.22790421698615\\
2.895	2.26564721411472\\
2.92	2.29757977123562\\
2.945	2.32995315851295\\
2.975	2.36938904692573\\
3	2.40274716306176\\
3.025	2.43655980814866\\
3.055	2.4777411931874\\
3.08	2.51256946822568\\
3.105	2.54786666834314\\
3.135	2.59084883334716\\
3.16	2.62719411431451\\
3.185	2.66402342840918\\
3.215	2.70886439072346\\
3.24	2.74677582014804\\
3.265	2.78518711878221\\
3.295	2.83194769160863\\
3.32	2.87147676072419\\
3.345	2.91152228075292\\
3.375	2.96026614071375\\
3.4	3.00146674735493\\
3.43	3.05161257748305\\
3.455	3.09399491153351\\
3.48	3.13692342264385\\
3.51	3.18916669367208\\
3.535	3.23331715141407\\
3.56	3.27803247011997\\
3.59	3.33244493169901\\
3.615	3.37842420112518\\
3.64	3.42498783830591\\
3.67	3.48164440016106\\
3.695	3.52951582883703\\
3.72	3.57799197918985\\
3.75	3.63697080426299\\
3.775	3.68680047847344\\
3.8	3.73725610106762\\
3.83	3.79863870404108\\
3.855	3.85049553243775\\
3.88	3.90300043580346\\
3.91	3.96687178705117\\
3.935	4.02082758886091\\
3.96	4.07545452058259\\
3.99	4.14190315561993\\
4.015	4.19803275342064\\
4.045	4.26630653612767\\
4.07	4.32397562882956\\
4.095	4.38235694595662\\
4.125	4.45336582350814\\
4.15	4.51334205940044\\
4.175	4.57405622209578\\
4.205	4.64789889675717\\
4.23	4.71026566581661\\
4.255	4.7733971101648\\
4.285	4.85017630148636\\
4.31	4.91502037988923\\
4.335	4.98065696365716\\
4.365	5.06047954473473\\
4.39	5.12789121013555\\
4.415	5.19612432961509\\
4.445	5.27910146940787\\
4.47	5.34917462128128\\
4.495	5.42009933328797\\
4.525	5.50634664514495\\
4.55	5.57917893038687\\
4.575	5.65289407947871\\
4.605	5.74253177603392\\
4.63	5.81822471988851\\
4.66	5.91026550581247\\
4.685	5.9879861256188\\
4.71	6.06664526811675\\
4.74	6.16229017996681\\
4.765	6.24305195969107\\
4.79	6.32478699773776\\
4.82	6.42416942748237\\
4.845	6.50808500733311\\
4.87	6.59300994265248\\
4.9	6.69626851070677\\
4.925	6.783454940871\\
4.95	6.87168823390541\\
4.98	6.97896697636671\\
5.005	7.06954587515574\\
5.03	7.16121060518538\\
5.06	7.27265916857076\\
5.085	7.36675688689873\\
5.11	7.46198091853283\\
5.14	7.5777547622712\\
5.165	7.675502554647\\
5.19	7.77441871090668\\
5.22	7.89467931778584\\
5.245	7.99621352007717\\
5.275	8.11965568384272\\
5.3	8.22387485720822\\
5.325	8.32933707130365\\
5.355	8.45755277264218\\
5.38	8.56580046323655\\
5.405	8.67533774894139\\
5.435	8.80850576622306\\
5.46	8.92093290132473\\
5.485	9.03469797567297\\
5.515	9.17300396030334\\
5.54	9.2897672658819\\
5.565	9.40791871013749\\
5.595	9.5515554381646\\
5.62	9.67281765134774\\
5.645	9.79552012595209\\
5.675	9.94468775939613\\
5.7	10.0706178490089\\
5.725	10.1980423167522\\
5.755	10.3529486748672\\
5.78	10.483722070432\\
5.805	10.6160460281489\\
5.835	10.776906869021\\
5.86	10.9127056984592\\
5.89	11.0777896679537\\
5.915	11.2171527303647\\
5.94	11.3581660667194\\
5.97	11.5295874327382\\
5.995	11.674299151163\\
6.02	11.8207232686488\\
6.05	11.9987206043601\\
6.075	12.1489823573673\\
6.1	12.3010209970983\\
6.13	12.4858419483304\\
6.155	12.6418627698465\\
6.18	12.7997274155767\\
6.21	12.9916290359012\\
6.235	13.1536258980327\\
6.26	13.3175360622008\\
6.29	13.516785160672\\
6.315	13.6849832671358\\
6.34	13.8551667877422\\
6.37	14.0620402895763\\
6.395	14.236673380129\\
6.42	14.4133667286595\\
6.45	14.6281520494756\\
6.475	14.8094627157283\\
6.505	15.029859772551\\
6.53	15.2159067506637\\
6.555	15.4041467637209\\
6.585	15.6329655301567\\
6.61	15.8261204485739\\
6.635	16.0215510988113\\
6.665	16.2591091477487\\
6.69	16.4596400450842\\
6.715	16.6625324947065\\
6.745	16.9091593942267\\
6.77	17.1173444308309\\
6.795	17.3279800794698\\
6.825	17.5840178371515\\
6.85	17.8001456700187\\
6.875	18.0188165339436\\
6.905	18.284620057656\\
6.93	18.5089902292826\\
6.955	18.7359993346579\\
6.985	19.0119369104462\\
7.01	19.2448602520177\\
7.035	19.4805220428587\\
7.065	19.7669758305731\\
7.09	20.0087748802514\\
7.12	20.3026877014548\\
7.145	20.5507821886043\\
7.17	20.801791625202\\
7.2	21.1068984546648\\
7.225	21.364440695989\\
7.25	21.6250077887217\\
7.28	21.9417306722873\\
7.305	22.2090768624441\\
7.33	22.4795619372808\\
7.36	22.8083387860909\\
7.385	23.0858585070564\\
7.41	23.366635430422\\
7.44	23.7079206085275\\
7.465	23.9959973255006\\
7.49	24.2874540056568\\
7.52	24.6417189393356\\
7.545	24.94075051484\\
7.57	25.2432894218771\\
7.6	25.6110232314016\\
7.625	25.921422458349\\
7.65	26.2354611633383\\
7.68	26.6171713181002\\
7.705	26.9393664724773\\
7.735	27.3309895961797\\
7.76	27.6615512040538\\
7.785	27.9959866181255\\
7.815	28.4024859545619\\
7.84	28.7456029222718\\
7.865	29.0927395786399\\
7.895	29.5146753459284\\
7.92	29.8708204592042\\
7.945	30.2311366068783\\
7.975	30.6690900032178\\
8	31.0387537508139\\
};
\addplot [color=mycolor1, line width=1.4pt, only marks, mark size=0.4pt, mark=*, mark options={solid, fill=mycolor1, mycolor1}, forget plot]
  table[row sep=crcr]{%
0.605	0.862105739970475\\
};
\addplot [color=black!90!mycolor9, line width=0.9pt, forget plot]
  table[row sep=crcr]{%
0.005	0.997460795963732\\
0.03	0.985192374205362\\
0.06	0.971378102889025\\
0.085	0.960588678719563\\
0.11	0.950427615759807\\
0.14	0.939029390806281\\
0.165	0.930166998967769\\
0.19	0.921860736381515\\
0.22	0.912600347669706\\
0.245	0.90545181455714\\
0.27	0.898802509214617\\
0.3	0.891460878791883\\
0.325	0.885857512521619\\
0.355	0.879731397239968\\
0.38	0.875110241514545\\
0.405	0.870916875961441\\
0.435	0.866434582841816\\
0.46	0.863145823554311\\
0.485	0.860253064034538\\
0.515	0.857292282818926\\
0.54	0.855240891913777\\
0.565	0.853559509762834\\
0.595	0.852020260547621\\
0.62	0.851128409134919\\
0.645	0.850585213313513\\
0.675	0.850385368510386\\
0.7	0.850589048606791\\
0.725	0.851123794247377\\
0.755	0.852195707608458\\
0.78	0.853442176046907\\
0.805	0.855005214923116\\
0.835	0.857293147472\\
0.86	0.859538965139959\\
0.885	0.86208943614462\\
0.915	0.865547565327128\\
0.94	0.868757074438401\\
0.97	0.872997529770665\\
0.995	0.876852355412535\\
1.02	0.880996483819377\\
1.05	0.886348114109104\\
1.075	0.891120891178209\\
1.1	0.896176204144513\\
1.13	0.902612991545979\\
1.155	0.908283711358913\\
1.18	0.914231652913328\\
1.21	0.921733164319825\\
1.235	0.928286262705942\\
1.26	0.935112543897903\\
1.29	0.943663248024007\\
1.315	0.95108707558696\\
1.34	0.958781177804837\\
1.37	0.968369889532731\\
1.395	0.976656279751191\\
1.42	0.985211049068042\\
1.45	0.995830474807459\\
1.475	1.00497439005025\\
1.5	1.01438570426509\\
1.53	1.02603206823902\\
1.555	1.03603130839243\\
1.585	1.0483831808407\\
1.61	1.05897053859287\\
1.635	1.06982550043983\\
1.665	1.08320508292984\\
1.69	1.09464983663332\\
1.715	1.10636331472232\\
1.745	1.12077490189388\\
1.77	1.13308139706262\\
1.795	1.14565841876282\\
1.825	1.161108924831\\
1.85	1.17428364316663\\
1.875	1.18773133211018\\
1.905	1.20423013735677\\
1.93	1.21828157696776\\
1.955	1.2326090402549\\
1.985	1.25016786438677\\
2.01	1.265106441116\\
2.035	1.2803246765543\\
2.065	1.29895747477465\\
2.09	1.31479544143117\\
2.115	1.33091726179049\\
2.145	1.35064014021494\\
2.17	1.36739152107764\\
2.2	1.3878742571652\\
2.225	1.40526264059643\\
2.25	1.42294317907215\\
2.28	1.44454786924771\\
2.305	1.46287720941623\\
2.33	1.48150432709203\\
2.36	1.50425256030403\\
2.385	1.52354141401553\\
2.41	1.54313417966717\\
2.44	1.56704951121361\\
2.465	1.58731806630817\\
2.49	1.60789717366924\\
2.52	1.63300510185715\\
2.545	1.65427516040923\\
2.57	1.67586291372599\\
2.6	1.70219086549117\\
2.625	1.72448583422332\\
2.65	1.74710614074008\\
2.68	1.77468346532108\\
2.705	1.79802835247486\\
2.73	1.82170672120779\\
2.76	1.85056469107634\\
2.785	1.87498610914766\\
2.815	1.90474366598444\\
2.84	1.92992147377168\\
2.865	1.95544764760146\\
2.895	1.98654300711728\\
2.92	2.01284591711446\\
2.945	2.03950669692857\\
2.975	2.07197635450664\\
3	2.09943534433109\\
3.025	2.12726222333356\\
3.055	2.16114465442712\\
3.08	2.18979235929554\\
3.105	2.21881849558253\\
3.135	2.25415418368617\\
3.16	2.28402492112209\\
3.185	2.3142851636274\\
3.215	2.35111663303385\\
3.24	2.38224643147676\\
3.265	2.41377734976811\\
3.295	2.45214920285104\\
3.32	2.48457583401145\\
3.345	2.51741575176901\\
3.375	2.5573747106502\\
3.4	2.59113772543564\\
3.43	2.63221604541256\\
3.455	2.66692170989392\\
3.48	2.70206247780922\\
3.51	2.74481168366792\\
3.535	2.78092461985449\\
3.56	2.81748637181228\\
3.59	2.86195912116577\\
3.615	2.8995238662347\\
3.64	2.93755174067086\\
3.67	2.98380300009168\\
3.695	3.02286603236209\\
3.72	3.06240712463049\\
3.75	3.11049423028056\\
3.775	3.15110402083298\\
3.8	3.19220743583972\\
3.83	3.2421901579663\\
3.855	3.28439722549121\\
3.88	3.32711413358595\\
3.91	3.37905474488994\\
3.935	3.42291171379328\\
3.96	3.46729540983415\\
3.99	3.52125875774066\\
4.015	3.56682041946764\\
4.045	3.62221338968594\\
4.07	3.66897996798024\\
4.095	3.71630334557527\\
4.125	3.77383453489748\\
4.15	3.82240336204484\\
4.175	3.87154777802932\\
4.205	3.93128934169511\\
4.23	3.98172136219418\\
4.255	4.03274850974665\\
4.285	4.09477547808492\\
4.31	4.14713405798229\\
4.335	4.20010807549357\\
4.365	4.26449844479694\\
4.39	4.31884944859983\\
4.415	4.37383699765214\\
4.445	4.4406718251894\\
4.47	4.49708369622808\\
4.495	4.55415404308987\\
4.525	4.6235175463421\\
4.55	4.68206139070941\\
4.575	4.74128649136529\\
4.605	4.81326615157514\\
4.63	4.87401582566419\\
4.66	4.94784658734592\\
4.685	5.01015708464529\\
4.71	5.07318928640224\\
4.74	5.14979153354353\\
4.765	5.21443898591324\\
4.79	5.27983334988008\\
4.82	5.35930388268723\\
4.845	5.42637000074264\\
4.87	5.49420919702677\\
4.9	5.57664850029632\\
4.925	5.64621810055533\\
4.95	5.71658793776213\\
4.98	5.80210030610849\\
5.005	5.8742614167545\\
5.03	5.94725094906594\\
5.06	6.03594461713582\\
5.085	6.11078858798789\\
5.11	6.18649022646663\\
5.14	6.27847750429757\\
5.165	6.35609912109469\\
5.19	6.43460874935307\\
5.22	6.53000616298632\\
5.245	6.61050376615379\\
5.275	6.70831550303959\\
5.3	6.79084929810766\\
5.325	6.87432490606347\\
5.355	6.97575338647639\\
5.38	7.06133752387829\\
5.405	7.14789696110829\\
5.435	7.25307078126635\\
5.46	7.34181377897465\\
5.485	7.43156680897607\\
5.515	7.5406193481608\\
5.54	7.63263375746589\\
5.565	7.72569422073471\\
5.595	7.83876380896877\\
5.62	7.93416635605229\\
5.645	8.03065231327235\\
5.675	8.14788240536862\\
5.7	8.24679413813964\\
5.725	8.34682801874519\\
5.755	8.46836737483872\\
5.78	8.57091381530384\\
5.805	8.67462257160751\\
5.835	8.80062544442531\\
5.86	8.90693674672764\\
5.89	9.0361007130425\\
5.915	9.14507834287602\\
5.94	9.25528945720515\\
5.97	9.389190241502\\
5.995	9.50216333177343\\
6.02	9.61641417433911\\
6.05	9.75522177339774\\
6.075	9.87233375631128\\
6.1	9.99076938559761\\
6.13	10.1346600414267\\
6.155	10.256059616377\\
6.18	10.3788304158901\\
6.21	10.5279868377333\\
6.235	10.6538281586347\\
6.26	10.7810900263089\\
6.29	10.9357016208571\\
6.315	11.0661444903373\\
6.34	11.1980590349647\\
6.37	11.3583221449084\\
6.395	11.4935322156057\\
6.42	11.6302669603865\\
6.45	11.7963851117385\\
6.475	11.9365340949465\\
6.505	12.10679950363\\
6.53	12.2504468469185\\
6.555	12.3957127288582\\
6.585	12.5721936367428\\
6.61	12.7210840003353\\
6.635	12.8716512662234\\
6.665	13.05457180083\\
6.69	13.2088942722693\\
6.715	13.3649541386325\\
6.745	13.5545466031614\\
6.77	13.7144971658397\\
6.795	13.8762478211901\\
6.825	14.0727529869282\\
6.85	14.2385347667378\\
6.875	14.4061816209877\\
6.905	14.6098490302867\\
6.93	14.7816725512304\\
6.955	14.9554284941079\\
6.985	15.1665167743203\\
7.01	15.3446002232741\\
7.035	15.5246858921738\\
7.065	15.7434630809453\\
7.09	15.9280325821224\\
7.12	16.1522564294854\\
7.145	16.341420521793\\
7.17	16.5327104210218\\
7.2	16.7650977742297\\
7.225	16.9611482991635\\
7.25	17.1594014584425\\
7.28	17.4002473493994\\
7.305	17.6034332002174\\
7.33	17.8089013028552\\
7.36	18.058511477994\\
7.385	18.2690905853888\\
7.41	18.4820344522353\\
7.44	18.7407257565385\\
7.465	18.9589654126413\\
7.49	19.1796553295358\\
7.52	19.4477561043751\\
7.545	19.6739332983165\\
7.57	19.9026493554105\\
7.6	20.1804998506717\\
7.625	20.4149016162091\\
7.65	20.6519340594608\\
7.68	20.939886860521\\
7.705	21.1828106361229\\
7.735	21.477919920958\\
7.76	21.7268807012791\\
7.785	21.9786347032701\\
7.815	22.2844704935013\\
7.84	22.5424798511198\\
7.865	22.8033834600026\\
7.895	23.1203337284166\\
7.92	23.3877189505647\\
7.945	23.6581031113284\\
7.975	23.9865698771823\\
8	24.2636700990563\\
};
\addplot [color=mycolor1, line width=1.4pt, only marks, mark size=0.4pt, mark=*, mark options={solid, fill=mycolor1, mycolor1}, forget plot]
  table[row sep=crcr]{%
0.67	0.850384806860622\\
};
\addplot [color=black, line width=0.9pt, forget plot]
  table[row sep=crcr]{%
0.005	0.99747866917676\\
0.03	0.985271857259535\\
0.06	0.971471522136607\\
0.085	0.960644831918275\\
0.11	0.950403009165189\\
0.14	0.938851611712845\\
0.165	0.929815583444764\\
0.19	0.921294625986456\\
0.22	0.911723327124952\\
0.245	0.904272018549605\\
0.27	0.897280827763574\\
0.3	0.889478045488961\\
0.325	0.88344854721743\\
0.355	0.876761536938085\\
0.38	0.871632172705005\\
0.405	0.866893918267437\\
0.435	0.861709833645405\\
0.46	0.857796688725455\\
0.485	0.854243895530285\\
0.515	0.850444396702932\\
0.54	0.847655463435543\\
0.565	0.845201673771011\\
0.595	0.84268979192997\\
0.62	0.840949500614533\\
0.645	0.839523565051046\\
0.675	0.83821934714184\\
0.7	0.837465281488113\\
0.725	0.83700835409711\\
0.755	0.836845587440901\\
0.78	0.837026001784662\\
0.805	0.837489258660076\\
0.835	0.83841297222746\\
0.86	0.83948488874656\\
0.885	0.840827772463013\\
0.915	0.842792305964601\\
0.94	0.844720038199347\\
0.97	0.847377754048995\\
0.995	0.849876385015484\\
1.02	0.852630379298577\\
1.05	0.856268923081687\\
1.075	0.859576570281873\\
1.1	0.863132487409479\\
1.13	0.867724588489495\\
1.155	0.871820062277475\\
1.18	0.876158042959349\\
1.21	0.881681545382851\\
1.235	0.886547715546441\\
1.26	0.891651788035106\\
1.29	0.898088989156911\\
1.315	0.903712271796128\\
1.34	0.90956987067824\\
1.37	0.91690698265319\\
1.395	0.923276921449582\\
1.42	0.92987849420112\\
1.45	0.93810520279437\\
1.475	0.945214125287308\\
1.5	0.952552808756734\\
1.53	0.961661909935226\\
1.555	0.96950464392472\\
1.585	0.979217681682817\\
1.61	0.987563096251358\\
1.635	0.996136734800498\\
1.665	1.00672624751341\\
1.69	1.01580177923413\\
1.715	1.02510548528163\\
1.745	1.0365712750766\\
1.77	1.0463774137277\\
1.795	1.05641225722153\\
1.825	1.06875635525771\\
1.85	1.07929540154707\\
1.875	1.09006422347179\\
1.905	1.10329073656246\\
1.93	1.11456668199432\\
1.955	1.12607398081551\\
1.985	1.14018896335015\\
2.01	1.15220738993072\\
2.035	1.16445922684223\\
2.065	1.17947057398519\\
2.09	1.19223857067011\\
2.115	1.20524249071202\\
2.145	1.22115984960946\\
2.17	1.23468594293735\\
2.2	1.25123265585208\\
2.225	1.26528560598894\\
2.25	1.27957970285539\\
2.28	1.29705248153985\\
2.305	1.31188102222321\\
2.33	1.32695436308519\\
2.36	1.34536723514068\\
2.385	1.36098350124975\\
2.41	1.37684862236831\\
2.44	1.39621715704447\\
2.465	1.41263455794349\\
2.49	1.42930526150379\\
2.52	1.44964653682338\\
2.545	1.46687973131993\\
2.57	1.48437106248181\\
2.6	1.50570364030693\\
2.625	1.5237685179853\\
2.65	1.54209674802666\\
2.68	1.56444065618504\\
2.705	1.58335432437137\\
2.73	1.60253693930243\\
2.76	1.62591365998892\\
2.785	1.64569443576829\\
2.815	1.66979442919138\\
2.84	1.69018259366724\\
2.865	1.71085000045005\\
2.895	1.73602234143064\\
2.92	1.75731125923315\\
2.945	1.77888602747081\\
2.975	1.80515609147558\\
3	1.82736732526765\\
3.025	1.84987139691502\\
3.055	1.87726601607251\\
3.08	1.90042234514947\\
3.105	1.92387888164535\\
3.135	1.95242635566033\\
3.16	1.9765517861647\\
3.185	2.00098517962936\\
3.215	2.03071529050555\\
3.24	2.05583506872401\\
3.265	2.08127095618584\\
3.295	2.11221498648118\\
3.32	2.13835561505921\\
3.345	2.16482089552945\\
3.375	2.19701164999255\\
3.4	2.22420090687494\\
3.43	2.25726861631973\\
3.455	2.28519559150162\\
3.48	2.31346255148711\\
3.51	2.34783586541833\\
3.535	2.37686129732528\\
3.56	2.40623635615792\\
3.59	2.44195230162117\\
3.615	2.4721075093679\\
3.64	2.50262241023401\\
3.67	2.53971964278841\\
3.695	2.57103731253305\\
3.72	2.60272517491951\\
3.75	2.64124401424001\\
3.775	2.6737582296205\\
3.8	2.70665358046895\\
3.83	2.746636048671\\
3.855	2.78038232351053\\
3.88	2.81452113066423\\
3.91	2.85601099332755\\
3.935	2.89102630661275\\
3.96	2.9264460144345\\
3.99	2.96948882432484\\
4.015	3.00581165738544\\
4.045	3.0499494935034\\
4.07	3.08719436807271\\
4.095	3.12486474394755\\
4.125	3.17063658059425\\
4.15	3.20925742966382\\
4.175	3.24831695471093\\
4.205	3.2957733847342\\
4.23	3.33581292846446\\
4.255	3.37630483615681\\
4.285	3.42549842362874\\
4.31	3.46700104069505\\
4.335	3.50897023764255\\
4.365	3.55995557400938\\
4.39	3.60296734916664\\
4.415	3.64646046337787\\
4.445	3.69929422630706\\
4.47	3.74386300020089\\
4.495	3.7889284316214\\
4.525	3.84366944658272\\
4.55	3.88984486786461\\
4.575	3.93653284119845\\
4.605	3.99324214581682\\
4.63	4.04107572556986\\
4.66	4.0991748175641\\
4.685	4.14817925667005\\
4.71	4.19772442520267\\
4.74	4.25790001267071\\
4.765	4.30865391786011\\
4.79	4.35996608390123\\
4.82	4.42228549049418\\
4.845	4.4748456662792\\
4.87	4.52798227864756\\
4.9	4.59251530224755\\
4.925	4.64694063707064\\
4.95	4.70196124883398\\
4.98	4.76878023961209\\
5.005	4.82513177177296\\
5.03	4.88209810699538\\
5.06	4.95127804857983\\
5.085	5.00961903478285\\
5.11	5.06859505789984\\
5.14	5.14021365173889\\
5.165	5.20060963831792\\
5.19	5.26166162621728\\
5.22	5.33579937919739\\
5.245	5.39831827610646\\
5.275	5.47423615262384\\
5.3	5.53825520843187\\
5.325	5.60296736603826\\
5.355	5.68154687390622\\
5.38	5.74780901322588\\
5.405	5.81478730202816\\
5.435	5.89611694831318\\
5.46	5.96469682201928\\
5.485	6.0340167153446\\
5.515	6.11818816548164\\
5.54	6.18916308590914\\
5.565	6.26090274569049\\
5.595	6.34801092769372\\
5.62	6.4214609565961\\
5.645	6.49570132231582\\
5.675	6.58584453485606\\
5.7	6.66185257456565\\
5.725	6.73867745546061\\
5.755	6.83195747988238\\
5.78	6.91060936782257\\
5.805	6.99010553845959\\
5.835	7.08662775487509\\
5.86	7.16801236150102\\
5.89	7.26682664124072\\
5.915	7.35014316840219\\
5.94	7.43435249868319\\
5.97	7.53659529725502\\
5.995	7.62280167423188\\
6.02	7.70993097278495\\
6.05	7.81571795027227\\
6.075	7.90491171188271\\
6.1	7.99505957187989\\
6.13	8.10451045263188\\
6.155	8.19679255936578\\
6.18	8.29006103533224\\
6.21	8.40329974494754\\
6.235	8.49877469850889\\
6.26	8.59526942311201\\
6.29	8.71242423032597\\
6.315	8.81120019327938\\
6.34	8.91103049830543\\
6.37	9.03223416186206\\
6.395	9.13442308151599\\
6.42	9.23770212309797\\
6.45	9.36309204384411\\
6.475	9.46880978050989\\
6.505	9.59715994385811\\
6.53	9.7053730471282\\
6.555	9.81473950691299\\
6.585	9.94751875921447\\
6.61	10.0594654391424\\
6.635	10.1726047113747\\
6.665	10.3099637248731\\
6.69	10.4257710290058\\
6.715	10.5428115285389\\
6.745	10.6849062261599\\
6.77	10.8047056282623\\
6.795	10.9257802427308\\
6.825	11.0727719769132\\
6.85	11.1966995277513\\
6.875	11.321945769987\\
6.905	11.4740015077543\\
6.93	11.6021979911676\\
6.955	11.731758156927\\
6.985	11.8890506714192\\
7.01	12.0216617657872\\
7.035	12.15568309703\\
7.065	12.3183911657191\\
7.09	12.4555676111312\\
7.12	12.6221057964516\\
7.145	12.7625110747214\\
7.17	12.9044088459139\\
7.2	13.07667846282\\
7.225	13.2219154287564\\
7.25	13.3686959140807\\
7.28	13.546892913053\\
7.305	13.6971267839194\\
7.33	13.8489569705078\\
7.36	14.0332840987505\\
7.385	14.188685821604\\
7.41	14.3457384870249\\
7.44	14.5364055211086\\
7.465	14.6971519689954\\
7.49	14.8596058794609\\
7.52	15.0568298660315\\
7.545	15.2231040409301\\
7.57	15.3911441564313\\
7.6	15.5951496612536\\
7.625	15.7671409040495\\
7.65	15.9409585905847\\
7.68	16.1519779562294\\
7.705	16.3298821639395\\
7.735	16.5458625043572\\
7.76	16.7279490254333\\
7.785	16.9119688760238\\
7.815	17.1353735371811\\
7.84	17.3237190961946\\
7.865	17.5140642810501\\
7.895	17.7451478523446\\
7.92	17.939967101464\\
7.945	18.1368545659649\\
7.975	18.3758804431218\\
8	18.5773954586942\\
};
\addplot [color=mycolor1, line width=1.4pt, only marks, mark size=0.4pt, mark=*, mark options={solid, fill=mycolor1, mycolor1}, forget plot]
  table[row sep=crcr]{%
0.75	0.836843809772079\\
};
\node[right, align=left, inner sep=0]
at (rel axis cs:-0.15,1.15) {$(a)$};
\end{axis}
\end{tikzpicture}%
  \end{minipage}
  \hfill
  \begin{minipage}{0.48\textwidth}
    \centering
    % This file was created by matlab2tikz.
%
\begin{tikzpicture}

\begin{axis}[%
name=ax_main,width=11.252in,
height=7.131in,
at={(1.887in,0.962in)},
scale only axis,
clip=true,
separate axis lines,
every outer x axis line/.append style={black},
every x tick label/.append style={font=\color{black}},
every x tick/.append style={black},
xmode=log,
xmin=0.5,
xmax=30,
xminorticks=true,
xlabel={$\lambda$},
every outer y axis line/.append style={black},
every y tick label/.append style={font=\color{black}},
every y tick/.append style={black},
ymin=0,
ymax=8,
ylabel={$t_c$},
axis background/.style={fill=white},
clip=true,
clip mode=individual,
axis lines=box,
xtick pos=bottom,
ytick pos=left,
tick align=inside,
major tick length=2pt,
width=0.8\linewidth,
height=0.8\linewidth,
every axis/.append style={font=\fontsize{9}{9}\selectfont},xlabel style={font=\fontsize{9}{9}\selectfont},title style={font=\fontsize{9}{9}\selectfont},
legend style={font=\fontsize{8}{9}\selectfont,inner sep=1pt,row sep=1pt,column sep=2pt,nodes={scale=1.00},draw=none},legend image post style={xscale=0.35},
legend columns=1,
scaled x ticks=false,
tick label style={/pgf/number format/fixed,/pgf/number format/precision=2},
every x tick label/.append style={font=\fontsize{9}{9}\selectfont\color{black}},
every y tick label/.append style={font=\fontsize{9}{9}\selectfont\color{black}},
,,
colormap={sigmamap}{rgb(0.0000)=(0.0200,0.1880,0.3800); rgb(0.1000)=(0.0743,0.3456,0.5616); rgb(0.2000)=(0.1918,0.4980,0.7060); rgb(0.3000)=(0.3890,0.6708,0.8226); rgb(0.4000)=(0.7552,0.8682,0.9228); rgb(0.5000)=(0.9690,0.9689,0.9689); rgb(0.6000)=(0.9425,0.8044,0.7614); rgb(0.7000)=(0.8767,0.4910,0.4020); rgb(0.8000)=(0.7869,0.2866,0.2470); rgb(0.9000)=(0.6186,0.1239,0.1568); rgb(1.0000)=(0.4040,0.0000,0.1220)}
]
\addplot [color=black, line width=1.0pt, mark size=1.5pt, mark=*, mark options={solid, fill=black, black}, forget plot]
  table[row sep=crcr]{%
0.5	7.07717180368207\\
0.620235917330657	6.71101473416963\\
0.769385186294003	6.19875905971079\\
0.954400653603358	5.5221767475147\\
1.18390712977731	4.7028236385544\\
1.46860344934348	3.81614126277316\\
1.82176121519704	2.96231615336066\\
2.25984347693029	2.22870480860547\\
2.80327218387512	1.67712469715036\\
3.4773801889866	1.31614219420198\\
4.31359218284711	1.10484080806223\\
5.35088960903706	0.995091690978841\\
6.63762785039236	0.953131764000139\\
8.23379039737524	0.959050136183402\\
10.2137850804488	1.00174839730939\\
12.6699127175806	1.07531561691457\\
15.716669873776	1.17688257535553\\
19.4960863130891	1.30550680676947\\
24.1843459575129	1.46148748713865\\
30	1.64605071648603\\
};
\node[right, align=left, inner sep=0]
at (rel axis cs:-0.15,1.15) {$(b)$};
\end{axis}

\begin{axis}[%
name=ax_inset,
width=0.304\linewidth,
height=0.28\linewidth,
at={(ax_main.north east)}, anchor=north east, xshift=-10pt, yshift=-10pt,
scale only axis,
separate axis lines,
every outer x axis line/.append style={black},
every x tick label/.append style={font=\color{black}},
every x tick/.append style={black},
xmode=log,
xmin=0.5,
xmax=30,
xminorticks=true,
xlabel={$\lambda$},
every outer y axis line/.append style={black},
every y tick label/.append style={font=\color{black}},
every y tick/.append style={black},
ymode=log,
ymin=0.108667651227639,
ymax=0.882311978386633,
yminorticks=true,
ylabel={$\mathscr{G}^{\mathrm{mean}}_{\min}$},
axis background/.style={fill=white},
clip=true,
clip mode=individual,
axis lines=box,
xtick pos=bottom,
ytick pos=left,
tick align=inside,
major tick length=2pt,
every axis/.append style={font=\fontsize{6}{7}\selectfont},xlabel style={font=\fontsize{6}{7}\selectfont},title style={font=\fontsize{6}{7}\selectfont},
legend style={font=\fontsize{6}{7}\selectfont,inner sep=1pt,row sep=1pt,column sep=2pt,nodes={scale=1.00},draw=none},legend image post style={xscale=0.35},
legend columns=1,
scaled x ticks=false,
tick label style={/pgf/number format/fixed,/pgf/number format/precision=2},
every x tick label/.append style={font=\fontsize{6}{7}\selectfont\color{black}},
every y tick label/.append style={font=\fontsize{6}{7}\selectfont\color{black}},
,,
colormap={sigmamap}{rgb(0.0000)=(0.0200,0.1880,0.3800); rgb(0.1000)=(0.0743,0.3456,0.5616); rgb(0.2000)=(0.1918,0.4980,0.7060); rgb(0.3000)=(0.3890,0.6708,0.8226); rgb(0.4000)=(0.7552,0.8682,0.9228); rgb(0.5000)=(0.9690,0.9689,0.9689); rgb(0.6000)=(0.9425,0.8044,0.7614); rgb(0.7000)=(0.8767,0.4910,0.4020); rgb(0.8000)=(0.7869,0.2866,0.2470); rgb(0.9000)=(0.6186,0.1239,0.1568); rgb(1.0000)=(0.4040,0.0000,0.1220)}
]
\addplot [color=black, line width=0.8pt, forget plot]
  table[row sep=crcr]{%
0.5	0.108667651227639\\
0.620235917330657	0.12556268994566\\
0.769385186294003	0.153346901716747\\
0.954400653603358	0.198675168593949\\
1.18390712977731	0.26910355910422\\
1.46860344934348	0.367342664722113\\
1.82176121519704	0.48501361022756\\
2.25984347693029	0.604572068479323\\
2.80327218387512	0.707171590419843\\
3.4773801889866	0.78226477063024\\
4.31359218284711	0.830867556656114\\
5.35088960903706	0.859629978517138\\
6.63762785039236	0.875059613231882\\
8.23379039737524	0.881687018461777\\
10.2137850804488	0.882311978386633\\
12.6699127175806	0.878629465830066\\
15.716669873776	0.871673288599548\\
19.4960863130891	0.862105739970475\\
24.1843459575129	0.850384806860622\\
30	0.836843809772079\\
};
\end{axis}
\end{tikzpicture}%
  \end{minipage}
  \caption{Three-dimensional mean energy growth in the Boussinesq limit for the \(n=1\) initial
  spectrum. (a)~Evolution of \(\Gm(t)\) for 20 logarithmically spaced integral correlation
  lengths \(0.5\leq\lambda\leq20\). The curves progress from light gray to black as \(\lambda\)
  increases. The horizontal dashed line marks \(\Gm=1\), and the small red markers identify
  \(\Gm_{\min}\) for each history. (b)~Critical time \(t_c\) as a function of \(\lambda\);
  the inset shows the corresponding minimum \(\Gm_{\min}\).}
  \label{fig:BQ_P3_Gmean_analysis}
\end{figure}

The light-to-dark family in \cref{fig:BQ_P3_Gmean_analysis}(a) shows that changing
\(\lambda\) modifies both the depth of the initial decay and the time required for the mean energy
to recover. The red markers trace the non-monotonic variation of \(\Gm_{\min}\) across
the family. The physical mechanisms responsible for this behaviour can be
understood directly from the structure of~\eqref{eq:bq_gmean_isotropic}. For the \(n=1\) spectrum,
increasing the integral correlation length is equivalent to increasing the Gaussian reference scale
in that equation. The exponential kernel acts as a low-pass filter on the wavenumber contributions
to \(\Gm\). As \(\lambda\) increases from zero, the strongly damped
large-\(k\) contributions are progressively filtered out, causing \(\Gm_{\min}\) to increase until
a maximum is reached. Beyond this optimal \(\lambda\), amplifying contributions at lower
wavenumbers are also suppressed, leading to a subsequent decrease in \(\Gm_{\min}\). This
non-monotonic behaviour highlights the sensitivity of the mean energy growth to the horizontal
correlation length. Correspondingly, \cref{fig:BQ_P3_Gmean_analysis}(b) shows that \(t_c\) reaches
a minimum at an intermediate correlation length. Perturbations recover most rapidly when the initial
spectrum reduces the strongly damped small-scale content without concentrating all of its energy
at the largest scales. The inset confirms that the same intermediate range of \(\lambda\) also
produces the shallowest transient energy loss.

\clearpage
\section{Variable-density results}
\label{sec:variable_density}
\citet{khorasani-szekely-ascher-2015}

\subsection{Linearized operator at zero horizontal wavenumber}

For completeness, we record the form taken by the variable-density linearized
operator for a perturbation that is uniform in the horizontal plane. Setting
\(k=0\) gives \(\nabla_{xy}^{2}=0\) and \(\nabla^{2}=\partial_{zz}\). It is
convenient to introduce the corresponding one-dimensional divergence operators
\begin{equation}
  \mathscr{L}_{\rho}^{(0)}(\star)
  =-\Sch^{-1}\partial_{zz}\!\left(\frac{\star}{\rho_0}\right),
  \qquad
  \dot{\mathscr{L}}_{\rho}^{(0)}(\star)
  =-\Sch^{-1}\partial_{zz}\!\left[
    \frac{\partial}{\partial t}\left(\frac{\star}{\rho_0}\right)
  \right].
  \label{eq:vd_k0_divergence_operators}
\end{equation}
The operator blocks defined in \cref{eq:A_B_operators} then reduce to
\begin{subequations}
\label{eq:vd_k0_operators}
\begin{align}
  \widehat{\mathscr{B}}_{ww}^{(0)}
  &= -\left(\rho_0'\partial_z+\rho_0\partial_{zz}\right), \\
  \widehat{\mathscr{B}}_{w\rho}^{(0)}
  &= \left(\rho_0'+\rho_0\partial_z\right)
     \mathscr{L}_{\rho}^{(0)}, \\
  \widehat{\mathscr{A}}_{ww}^{(0)}
  &= -\mu_0\partial_{zzzz}-2\mu_0'\partial_{zzz}
     -\mu_0''\partial_{zz}, \\
  \widehat{\mathscr{A}}_{w_0w}^{(0)}
  &= \rho_0w_0\partial_{zzz}+(\rho_0w_0)'\partial_{zz}, \\
  \widehat{\mathscr{A}}_{w\rho}^{(0)}
  &= \left(\mu_0\partial_{zzz}+2\mu_0'\partial_{zz}
     +\mu_0''\partial_z\right)\mathscr{L}_{\rho}^{(0)}, \\
  \widehat{\mathscr{A}}_{w_0\rho}^{(0)}
  &= -\left((\rho_0w_0)'\partial_z+\rho_0w_0\partial_{zz}\right)
     \mathscr{L}_{\rho}^{(0)}
     -\left(\rho_0'+\rho_0\partial_z\right)
     \dot{\mathscr{L}}_{\rho}^{(0)}, \\
  \widehat{\mathscr{A}}_{\rho w}^{(0)}
  &= -\rho_0', \\
  \widehat{\mathscr{A}}_{\rho\rho}^{(0)}
  &= -\left(w_0\partial_z+w_0'\right)
     -\rho_0\mathscr{L}_{\rho}^{(0)}.
\end{align}
\end{subequations}
The two divergence operators can now be expanded explicitly as
\begin{equation}
  \mathscr{L}_{\rho}^{(0)}(\rho)
  =-\Sch^{-1}\partial_{zz}\!\left(\frac{\rho}{\rho_0}\right),
  \qquad
  \dot{\mathscr{L}}_{\rho}^{(0)}(\rho)
  =\Sch^{-1}\partial_{zz}\!\left(
    \frac{\rho\,\partial_t\rho_0}{\rho_0^2}
  \right),
  \label{eq:vd_k0_expanded_divergence_operators}
\end{equation}
where the overdot differentiates the time-dependent operator while holding its
argument fixed. Hence
\begin{equation}
  \partial_t\!\left[\mathscr{L}_{\rho}^{(0)}(\rho)\right]
  =\mathscr{L}_{\rho}^{(0)}(\partial_t\rho)
   +\dot{\mathscr{L}}_{\rho}^{(0)}(\rho).
  \label{eq:vd_k0_operator_product_rule}
\end{equation}

At zero horizontal wavenumber, the divergence constraint itself becomes
\begin{equation}
  \partial_z w
  =-\Sch^{-1}\partial_{zz}\!\left(\frac{\rho}{\rho_0}\right).
  \label{eq:vd_k0_divergence_constraint}
\end{equation}
Using this constraint and its time derivative in the vertical-velocity equation
makes all of its terms cancel. For example, its remaining spatial terms are
\begin{equation}
  \partial_z\!\left(\mu_0\partial_{zz}w\right)
  -\rho_0w_0\partial_{zz}w
  -\left(\mu_0\partial_{zz}+\mu_0'\partial_z
  -\rho_0w_0\partial_z\right)
  \mathscr{L}_{\rho}^{(0)}(\rho)=0,
\end{equation}
because \(\partial_z w=\mathscr{L}_{\rho}^{(0)}(\rho)\). Thus, at \(k=0\),
the pressure-eliminated vertical-momentum equation supplies no additional
dynamics.

Integrating \eqref{eq:vd_k0_divergence_constraint} once and applying the
far-field conditions gives
\begin{subequations}
  \label{eq:vd_k0_linearized_system}
  \begin{align}
    w
    &=-\Sch^{-1}\partial_z\!\left(\frac{\rho}{\rho_0}\right),
    \label{eq:vd_k0_vertical_velocity}\\
    \partial_t\rho+\partial_z(w_0\rho)
    &=\Sch^{-1}\partial_z\!\left[
      \rho_0\partial_z\!\left(\frac{\rho}{\rho_0}\right)
    \right].
    \label{eq:vd_k0_density}
  \end{align}
\end{subequations}
The \(k=0\) problem therefore reduces to a single scalar equation for \(\rho\),
after which \(w\) follows directly from \eqref{eq:vd_k0_vertical_velocity}.

To examine its modal behaviour, freeze the base state at a time \(t_0\) and set
\(\rho(z,t)=\phi(z)e^{\sigma t}\). The eigenvalue problem is simply
\begin{equation}
  \sigma\phi
  =-\partial_z(w_0\phi)
  +\Sch^{-1}\partial_z\!\left[
    \rho_0\partial_z\!\left(\frac{\phi}{\rho_0}\right)
  \right].
  \label{eq:vd_k0_eigenvalue_problem}
\end{equation}
This is a conservative one-dimensional advection--diffusion operator and hence
has no eigenvalues with positive real part. In the uniform far field,
\(w_0\to0\) and \(\rho_0\) is constant, so a vertical Fourier mode
\(\phi\sim e^{\mathrm{i}mz}\) has the growth rate
\begin{equation}
  \sigma(m)=-\frac{m^2}{\Sch}\leq0.
  \label{eq:vd_k0_growth_rate}
\end{equation}
Thus, on an unbounded domain, the spectral bound is \(\sigma_{\max}=0\), reached
in the limit \(m\to0\); every finite vertical wavenumber decays. On a finite
domain the largest eigenvalue is strictly negative and depends on the vertical
boundary conditions.

Expanding the diffusive flux in \eqref{eq:vd_k0_density} writes the same
equation in conservative advection--diffusion form,
\begin{equation}
  \partial_t\rho
  =
  \Sch^{-1}\partial_{zz}\rho
  -
  \partial_z\!\left[
    \left(w_0+\Sch^{-1}\left(\ln\rho_0\right)'\right)\rho
  \right],
  \label{eq:vd_k0_conservative_form}
\end{equation}
or, equivalently, as an advection--reaction--diffusion equation
\(\partial_t\rho=a\,\partial_{zz}\rho-c\,\partial_z\rho+r\,\rho\) with
\begin{subequations}
  \label{eq:vd_k0_ard_coefficients}
  \begin{align}
    a &= \Sch^{-1},
    \label{eq:vd_k0_ard_a}\\
    c &= w_0+\Sch^{-1}\frac{\rho_0'}{\rho_0},
    \label{eq:vd_k0_ard_c}\\
    r &= -\left[w_0'+\Sch^{-1}\left(\ln\rho_0\right)''\right]=-c',
    \label{eq:vd_k0_ard_r}
  \end{align}
\end{subequations}
where \(\left(\ln\rho_0\right)''=\rho_0''/\rho_0-\left(\rho_0'/\rho_0\right)^2\).
The reaction coefficient is minus the derivative of the advection coefficient,
which is the conservative structure already used above. The base state
\eqref{eq:vd_base_state_w0} gives \(w_0=\rho_0'/\rho_0=\left(\ln\rho_0\right)'\),
so that
\begin{equation}
  \partial_t\rho
  =
  \Sch^{-1}\partial_{zz}\rho
  -
  \left(1+\Sch^{-1}\right)\partial_z\!\left(w_0\rho\right),
  \qquad
  c=\left(1+\Sch^{-1}\right)w_0,
  \qquad
  r=-\left(1+\Sch^{-1}\right)w_0' ,
  \label{eq:vd_k0_ard_base_state}
\end{equation}
and, with \(w_0'=-2zw_0/\delta^{2}-w_0^{2}\),
\(r=\left(1+\Sch^{-1}\right)w_0\left(2z/\delta^{2}+w_0\right)\).

%k-star and sigma-star grids at delta0 = 2: rows sweep Sc and columns sweep At,
%So a row shares one ordinate scale set by its leftmost panel. Panels are
%Generated by codes/4_VD/P2_sigma_kstar.m.
\newcommand{\VDStarGrid}[2]{%
\begin{figure}[p]
  \centering
  \begin{minipage}[t]{0.24\textwidth}\centering
    \tikzfig{#1_VD_Sc_0p25_d0_2_At_0p05}
  \end{minipage}\hfill
  \begin{minipage}[t]{0.24\textwidth}\centering
    \tikzfig{#1_VD_Sc_0p25_d0_2_At_0p2}
  \end{minipage}\hfill
  \begin{minipage}[t]{0.24\textwidth}\centering
    \tikzfig{#1_VD_Sc_0p25_d0_2_At_0p5}
  \end{minipage}\hfill
  \begin{minipage}[t]{0.24\textwidth}\centering
    \tikzfig{#1_VD_Sc_0p25_d0_2_At_0p7}
  \end{minipage}
  \\[2pt]
  \begin{minipage}[t]{0.24\textwidth}\centering
    \tikzfig{#1_VD_Sc_1_d0_2_At_0p05}
  \end{minipage}\hfill
  \begin{minipage}[t]{0.24\textwidth}\centering
    \tikzfig{#1_VD_Sc_1_d0_2_At_0p2}
  \end{minipage}\hfill
  \begin{minipage}[t]{0.24\textwidth}\centering
    \tikzfig{#1_VD_Sc_1_d0_2_At_0p5}
  \end{minipage}\hfill
  \begin{minipage}[t]{0.24\textwidth}\centering
    \tikzfig{#1_VD_Sc_1_d0_2_At_0p7}
  \end{minipage}
  \\[2pt]
  \begin{minipage}[t]{0.24\textwidth}\centering
    \tikzfig{#1_VD_Sc_10_d0_2_At_0p05}
  \end{minipage}\hfill
  \begin{minipage}[t]{0.24\textwidth}\centering
    \tikzfig{#1_VD_Sc_10_d0_2_At_0p2}
  \end{minipage}\hfill
  \begin{minipage}[t]{0.24\textwidth}\centering
    \tikzfig{#1_VD_Sc_10_d0_2_At_0p5}
  \end{minipage}\hfill
  \begin{minipage}[t]{0.24\textwidth}\centering
    \tikzfig{#1_VD_Sc_10_d0_2_At_0p7}
  \end{minipage}
  \caption{#2 in the variable-density regime at \(\delta_0=2\), as defined
  in~\eqref{eq:maximum_growth_definitions}. Rows correspond to \(\Sch=0.25\),
  \(1\), and \(10\); columns correspond to \(\Atw=0.05\), \(0.2\), \(0.5\), and
  \(0.7\), as labelled above each panel. Each panel compares the three viscosity
  contrasts \(\Amu=-0.9\), \(0\), and \(0.9\). The ordinate range and the time
  window are shared along each row, since both change strongly with \(\Sch\) and
  weakly with \(\Atw\).}
  \label{fig:VD_#1_grid}
\end{figure}%
}

\VDStarGrid{sigmastar}{Maximum mean growth rate \(\overline{\sigma}_\ast(t)\)}


% Three-dimensional growth-rate grids at delta0 = 2, one figure per viscosity
% contrast: rows sweep Sc and columns sweep At. Panels are generated by
% codes/4_VD/P3_growth_3D.m.
\newcommand{\VDThreeDGrid}[2]{%
\begin{figure}[p]
  \centering
  \begin{minipage}[t]{0.24\textwidth}\centering
    \tikzfig{sigma3D_VD_Sc_0p25_d0_2_At_0p05_Amu_#1}
  \end{minipage}\hfill
  \begin{minipage}[t]{0.24\textwidth}\centering
    \tikzfig{sigma3D_VD_Sc_0p25_d0_2_At_0p2_Amu_#1}
  \end{minipage}\hfill
  \begin{minipage}[t]{0.24\textwidth}\centering
    \tikzfig{sigma3D_VD_Sc_0p25_d0_2_At_0p5_Amu_#1}
  \end{minipage}\hfill
  \begin{minipage}[t]{0.24\textwidth}\centering
    \tikzfig{sigma3D_VD_Sc_0p25_d0_2_At_0p7_Amu_#1}
  \end{minipage}
  \\[2pt]
  \begin{minipage}[t]{0.24\textwidth}\centering
    \tikzfig{sigma3D_VD_Sc_1_d0_2_At_0p05_Amu_#1}
  \end{minipage}\hfill
  \begin{minipage}[t]{0.24\textwidth}\centering
    \tikzfig{sigma3D_VD_Sc_1_d0_2_At_0p2_Amu_#1}
  \end{minipage}\hfill
  \begin{minipage}[t]{0.24\textwidth}\centering
    \tikzfig{sigma3D_VD_Sc_1_d0_2_At_0p5_Amu_#1}
  \end{minipage}\hfill
  \begin{minipage}[t]{0.24\textwidth}\centering
    \tikzfig{sigma3D_VD_Sc_1_d0_2_At_0p7_Amu_#1}
  \end{minipage}
  \\[2pt]
  \begin{minipage}[t]{0.24\textwidth}\centering
    \tikzfig{sigma3D_VD_Sc_10_d0_2_At_0p05_Amu_#1}
  \end{minipage}\hfill
  \begin{minipage}[t]{0.24\textwidth}\centering
    \tikzfig{sigma3D_VD_Sc_10_d0_2_At_0p2_Amu_#1}
  \end{minipage}\hfill
  \begin{minipage}[t]{0.24\textwidth}\centering
    \tikzfig{sigma3D_VD_Sc_10_d0_2_At_0p5_Amu_#1}
  \end{minipage}\hfill
  \begin{minipage}[t]{0.24\textwidth}\centering
    \tikzfig{sigma3D_VD_Sc_10_d0_2_At_0p7_Amu_#1}
  \end{minipage}
  \\[2pt]
  \begin{tikzpicture}
\begin{axis}[
hide axis, scale only axis,
width=0.75\textwidth, height=0pt,
colormap={sigmamap}{rgb(0.0000)=(0.8800,0.8800,0.8800); rgb(0.1000)=(0.7903,0.7903,0.7903); rgb(0.2000)=(0.7040,0.7040,0.7040); rgb(0.3000)=(0.6143,0.6143,0.6143); rgb(0.4000)=(0.5280,0.5280,0.5280); rgb(0.5000)=(0.4383,0.4383,0.4383); rgb(0.6000)=(0.3520,0.3520,0.3520); rgb(0.7000)=(0.2623,0.2623,0.2623); rgb(0.8000)=(0.1760,0.1760,0.1760); rgb(0.9000)=(0.0863,0.0863,0.0863); rgb(1.0000)=(0.0000,0.0000,0.0000)},
point meta min=-0.30103, point meta max=1.30103,
colorbar horizontal,
colorbar style={
  width=0.75\textwidth, height=7pt,
  at={(0,0)}, anchor=north west,
  xlabel={$\lambda$},
  xlabel style={font=\footnotesize},
  tick label style={font=\footnotesize},
  xtick pos=bottom,
  xtick={-0.30103,0,0.30103,0.69897,1,1.30103},
  xticklabels={0.5,1,2,5,10,20},
},
]
\addplot[draw=none] coordinates {(0,0)};
\end{axis}
\end{tikzpicture}

  \caption{Three-dimensional mean growth rate \(\overline{\sigma}_{3D}(t)\) in
  the variable-density regime at \(\Amu=#2\) and \(\delta_0=2\). Rows
  correspond to \(\Sch=0.25\), \(1\), and \(10\); columns correspond to
  \(\Atw=0.05\), \(0.2\), \(0.5\), and \(0.7\), as labelled above each panel.
  Each panel contains 20 logarithmically spaced integral correlation lengths
  \(0.5\leq\lambda\leq20\), progressing from light gray to black as \(\lambda\)
  increases, and the (\dashlegend) line marks neutral growth. The ordinate range
  is shared along each row.}
  \label{fig:VD_sigma3D_grid_Amu#1}
\end{figure}%
}

\VDThreeDGrid{m0p9}{-0.9}
\VDThreeDGrid{0}{0}
\VDThreeDGrid{0p9}{0.9}



\clearpage
\begin{appendices}
\crefname{appendix}{Appendix}{Appendices}
\Crefname{appendix}{Appendix}{Appendices}
\crefalias{section}{appendix}
\section{Continuous linearized operators}
\label{app:operators}

This appendix collects the continuous differential operators used in the linearized
variable-density system~\eqref{eq:vd_lin_continuous}. These operators are defined
before the vertical discretization described in \eqref{eq:semidiscrete_system}; after
discretization they produce the matrices \(\mathsfbi{A}_{k}(t)\) and \(\mathsfbi{B}_{k}(t)\) that
define the discrete propagator in \eqref{eq:discrete_propagator}.

Primes denote derivatives with respect to \(z\), overdots denote derivatives with respect to time,
and
\begin{equation}
  \nabla^2=\partial_{zz}+\nabla_{xy}^2,
  \qquad
  \nabla_{xy}^2=\partial_{xx}+\partial_{yy}.
  \label{eq:appendix_laplacian_definitions}
\end{equation}
The base-state viscosity is
\begin{equation}
  \mu_0(z,t)
  =
  1+\Amu\,\erf\!\left(\frac{z}{\delta(t)}\right),
  \label{eq:appendix_base_viscosity}
\end{equation}
and viscosity perturbations are related to density perturbations by
\begin{equation}
  \mu' = C_{\mu\rho}\rho',
  \qquad
  C_{\mu\rho}=\frac{\Amu}{\Atw}.
  \label{eq:appendix_viscosity_density_coupling}
\end{equation}
The linearized divergence constraint entering the operators is
\begin{equation}
  \bnabla\bcdot\boldsymbol{u}
  =
  \mathscr{L}_{\rho}\,\rho,
  \qquad
  \mathscr{L}_{\rho}(\star)
  =
  -\Sch^{-1}\nabla^2\!\left(\frac{\star}{\rho_0}\right),
  \label{eq:linearized_divergence}
\end{equation}
where \(\star\) denotes the argument on which the operator acts. We also define
\begin{equation}
  \dot{\mathscr{L}}_{\rho}(\star)
  =
  -\Sch^{-1}
  \nabla^2
  \left[
    \frac{\partial}{\partial t}
    \left(
      \frac{\star}{\rho_0}
    \right)
  \right].
  \label{eq:ldot_rho}
\end{equation}
With this notation, the operator blocks in \eqref{eq:vd_lin_continuous} are
\begin{subequations}
\label{eq:A_B_operators}
\begin{align}
  \mathscr{B}_{ww}
  &=
  -\Big(\rho_0'\,\partial_z+\rho_0\,\nabla^2\Big), \\
  \mathscr{B}_{w\rho}
  &=
  \Big(\rho_0'+\rho_0\,\partial_z\Big)\mathscr{L}_{\rho}, \\
  \mathscr{A}_{ww}
  &=
  -\mu_0\,\nabla^4
  -2\,\mu_0'\,\partial_z\nabla^2
  +\mu_0''\,\nabla^2
  -2\,\mu_0''\,\partial_{zz}, \\
  \mathscr{A}_{w_0 w}
  &=
  \Big(\rho_0 w_0'+\rho_0 w_0\,\partial_z\Big)\nabla^2
  +w_0\rho_0'\,\partial_{zz}, \\
  \mathscr{A}_{w\rho}
  &=
  \Big(\mu_0\,\partial_z\nabla^2+2\,\mu_0'\,\nabla^2+\mu_0''\,\partial_z\Big)
  \mathscr{L}_{\rho}
  -2\,\partial_z\!\left(w_0'C_{\mu\rho}\,\nabla_{xy}^2\star\right)
  +\nabla_{xy}^2, \\
  \mathscr{A}_{w_0\rho}
  &=
  -\Big((\rho_0 w_0)'\,\partial_z+\rho_0 w_0\,\partial_{zz}\Big)
  \mathscr{L}_{\rho}
  +\left(\dot{w}_0+w_0 w_0'\right)\nabla_{xy}^2
  -\Big(\rho_0'+\rho_0\,\partial_z\Big)\dot{\mathscr{L}}_{\rho}, \\
  \mathscr{A}_{\rho w}
  &=
  -\,\partial_z\rho_0, \\
  \mathscr{A}_{\rho\rho}
  &=
  -\left(w_0\,\partial_z+w_0'\right)-\rho_0\,\mathscr{L}_{\rho}.
\end{align}
\end{subequations}

\clearpage
\section{Trace-estimation strategies and computational cost}
\label{sec:appendix_main}
In this appendix we discuss the practical computation of the mean energy growth introduced in
\cref{subsec:mean_energy_growth} and in particular we focus on the verification of the methods against test cases for which solutions are available. The formula for the per-wavenumber mean energy growth is
\begin{equation}
  g(k,t)
  =
  \frac{
    \tr\left\{\mathsfbi{T}_{t,k}\mathsfbi{C}_{00}^{z}\right\}
  }{
    \tr\left\{\mathsfbi{C}_{00}^{z}\mathsfbi{W}_{0,k}\right\}
  },
  \qquad
  \mathsfbi{T}_{t,k}
  =
  \boldsymbol{\Phi}_{t,k}^{*}\mathsfbi{W}_{t,k}\boldsymbol{\Phi}_{t,k}.
  \label{eq:appendix_modal_growth}
\end{equation}
To make the notation less cumbersome, we refer to this quantity with the wavenumber \(k\) held
fixed and the time set to the final time \(T\), so that \(\boldsymbol{\Phi}_{T,k}\equiv\boldsymbol{\Phi}\),
\(\mathsfbi{W}_{T,k}\equiv\mathsfbi{W}\), \(\mathsfbi{C}_{00}^{z}\equiv\mathsfbi{C}_{00}\), and
\begin{equation}
  g(T)
  =
  \frac{
    \tr\left\{\mathsfbi{T}\mathsfbi{C}_{00}\right\}
  }{
    \tr\left\{\mathsfbi{C}_{00}\mathsfbi{W}\right\}
  },
  \qquad
  \mathsfbi{T}
  =
  \boldsymbol{\Phi}^{*}\mathsfbi{W}\boldsymbol{\Phi}.
  \label{eq:appendix_growth_functional}
\end{equation}
We verify the SSTE algorithm on canonical autonomous test cases and on an analytical non-autonomous problem. As a measure of error we define
\begin{equation}
  \varepsilon_{\mathrm{rel}}(T)
  =
  \frac{
    \left|g(T)-g_{\mathrm{ref}}(T)\right|
  }{
    \left|g_{\mathrm{ref}}(T)\right|
  },
  \label{eq:verification_relative_error}
\end{equation}
where \(g_{\mathrm{ref}}\) denotes the reference value. The reference stays exact in every case considered, although its origin differs. For autonomous problems the linearized operators are constant in time, and the reference is therefore given by the trace evaluated from the exponential propagator. For the fully non-autonomous case we derive a closed-form solution, and the reference is the resulting analytical expression.

In this appendix we analyze and compare two methodologies, the deterministic covariance truncation,
which we abbreviate DCT and describe in \cref{sec:appendix_deterministic_truncation}, and the SSTE
already described in \cref{subsec:randomized_trace_estimation}. In particular, we evaluate the
computational cost of each method and discuss when DCT or SSTE must be preferred.

We measure cost in stepper calls, since the time integration dominates every strategy considered
here. The symbol \(\mathcal{C}_{X}(T)\) counts the forward integrations that strategy \(X\) needs to
evaluate \(g(T)\) at a single target time, whereas the symbol \(\mathcal{W}_{X}\) counts the
cumulative stepping work that the same strategy needs over a history of \(N_t\) target intervals.

\subsection{Deterministic covariance truncation}
\label{sec:appendix_deterministic_truncation}
\label{sec:appendix_necessity_sste}

The idea of deterministic truncation is to exploit the spectral structure of the initial covariance.
The singular value decomposition of such a matrix is given by
\begin{equation}
  \mathsfbi{C}_{00}
  =
  \mathsfbi{U}\boldsymbol{\Sigma}\mathsfbi{U}^{*},
  \qquad
  \boldsymbol{\Sigma}
  =
  \operatorname{diag}
  \left(
    \sigma_1,\dots,\sigma_N
  \right),
  \qquad
  \sigma_1\geq\dots\geq\sigma_N>0,
  \label{eq:C00_eigendecomposition_sste_necessity}
\end{equation}
and, since \(\mathsfbi{C}_{00}\in\mathbb{C}^{N\times N}\) is Hermitian positive definite, this
decomposition is also its eigendecomposition, which delivers the square-root factor
\(\mathsfbi{L}\) in the form
\begin{equation}
  \mathsfbi{L}
  =
  \mathsfbi{U}\boldsymbol{\Sigma}^{1/2},
  \qquad
  \boldsymbol{\Sigma}^{1/2}
  =
  \operatorname{diag}
  \left(
    \sqrt{\sigma_1},\dots,\sqrt{\sigma_N}
  \right),
  \label{eq:C00_square_root_factor_sste_necessity}
\end{equation}
so that each eigenpair carries a single column of the factor into the dynamics.

DCT replaces \(\mathsfbi{C}_{00}\) by the rank-\(r\) truncation \(\mathsfbi{C}_{00}^{(r)}\) that
retains the leading \(r\) eigenpairs and propagates every retained column up to the target time
\(T\), as reported in \cref{alg:deterministic_truncation}. The cost, in terms of time-stepper
calls, is
\begin{equation}
  \mathcal{C}_{\mathrm{DCT}}(T)=r,
  \label{eq:deterministic_fixed_time_cost}
\end{equation}
and, because we need to repeat this for all the \(N_t\) target times, the total cost is
\begin{equation}
  \mathcal{W}_{\mathrm{DCT}}=N_t\,r.
  \label{eq:deterministic_history_cost}
\end{equation}
Whenever the covariance has rank one, or an exact rank that a modest \(r\) already captures, DCT
returns \(g(T)\) exactly, and we prefer it to every alternative; whenever the spectrum extends
beyond the retained rank, DCT introduces the error
\begin{equation}
  \varepsilon_r
  =
  \frac{1}{\tr\left\{\mathsfbi{W}\mathsfbi{C}_{00}\right\}}
  \tr
  \left\{
    \mathsfbi{T}
    \left(
      \mathsfbi{C}_{00}
      -
      \mathsfbi{C}_{00}^{(r)}
    \right)
  \right\}
  =
  \frac{1}{\tr\left\{\mathsfbi{W}\mathsfbi{C}_{00}\right\}}
  \sum_{j=r+1}^{N}
  \sigma_j\,
  \mathbf{u}_j^{*}\mathsfbi{T}\mathbf{u}_j.
  \label{eq:trace_error_1}
\end{equation}

\begin{algorithm}[H]
  \caption{Deterministic truncation (DCT) estimate of \(g(T)\)}
  \label{alg:deterministic_truncation}

  \begin{algorithmic}[1]
    \Require Initial covariance \(\mathsfbi{C}_{00}\), truncation rank \(r\),
             target time \(T\), operators \(\mathsfbi{A}(t)\) and \(\mathsfbi{B}(t)\),
             energy weight \(\mathsfbi{W}\)
    \Ensure Truncated growth functional \(g^{(r)}(T)\)

    \State Compute \((\sigma_j,\mathbf{u}_j)\) of \(\mathsfbi{C}_{00}\)
      and order them so that \(\sigma_1\geq\dots\geq\sigma_N\)

    \State \(s\gets0\)
      \Comment{Trace accumulator}

    \For{\(j=1,\dots,r\)}
      \State \(\hat{\mathbf{q}}_j(0)\gets\sqrt{\sigma_j}\,\mathbf{u}_j\)
        \Comment{Retained column of \(\mathsfbi{L}\)}

      \State Solve
        \(\mathsfbi{B}(t)\,\partial_t\hat{\mathbf{q}}_j(t)
        =
        \mathsfbi{A}(t)\hat{\mathbf{q}}_j(t)\),
        \(0\leq t\leq T\)
        \Comment{Forward solve}

      \State \(s\gets s+\hat{\mathbf{q}}_j^{*}(T)\mathsfbi{W}\hat{\mathbf{q}}_j(T)\)
        \Comment{Energy of the propagated column}
    \EndFor

    \State \(g^{(r)}(T)\gets s/\tr\left\{\mathsfbi{C}_{00}\mathsfbi{W}\right\}\)
      \Comment{Normalize by the initial energy}

    \State \Return \(g^{(r)}(T)\)
  \end{algorithmic}
\end{algorithm}

\Cref{eq:trace_error_1} shows that two mechanisms govern the relative error. The first one is that
choosing \(r\) such that \(\sigma_r\simeq0\) makes the error small, while the second contribution
comes from the amplification of the operator. In fact, \eqref{eq:trace_error_1} can be rearranged as
\begin{equation}
  \varepsilon_r
  =
  \frac{1}{\tr\left\{\mathsfbi{W}\mathsfbi{C}_{00}\right\}}
  \sum_{j=r+1}^{N}
  \sigma_j
  \left\|\boldsymbol{\Phi}\mathbf{u}_j\right\|_{\mathsfbi{W}}^{2},
  \label{eq:trace_error_phi_tail}
\end{equation}
which requires \(\left\|\boldsymbol{\Phi}\mathbf{u}_j\right\|_{\mathsfbi{W}}^{2}\) to be small as
well.

This suggests that small neglected eigenvalues guarantee little on their own, since a discarded
eigendirection that aligns with a strongly amplified direction of the dynamics still contributes
substantially to the error. Consider in fact the simple case where
\begin{equation}
  \boldsymbol{\Phi}
  =
  \begin{pmatrix}
    1 & 0 \\
    0 & \sqrt{\kappa/\epsilon}
  \end{pmatrix},
  \qquad
  \mathsfbi{C}_{00}
  =
  \begin{pmatrix}
    1 & 0 \\
    0 & \epsilon
  \end{pmatrix},
  \qquad
  \mathsfbi{C}_{00}^{(r)}
  =
  \begin{pmatrix}
    1 & 0 \\
    0 & 0
  \end{pmatrix},
  \qquad
  \epsilon \ll 1 ,
  \label{eq:low_rank_failure_example}
\end{equation}
with \(\kappa\) the gain ratio and \(\epsilon\) the small covariance eigenvalue. Being
\(g^{(r)}\) the growth functional computed with \(\mathsfbi{C}_{00}^{(r)}\), the relative error is
\begin{equation}
  \frac{
    g(T)-g^{(r)}(T)
  }{
    g(T)
  }
  =
  \frac{\kappa}{\kappa+1}\sim O(1).
  \label{eq:low_rank_failure_relative_error}
\end{equation}
The discarded covariance eigenvalue \(\epsilon\) stays small while its contribution to the
propagated energy does not, precisely because the dynamics amplify the corresponding
eigendirection.

We therefore argue that a low-rank truncation of the initial covariance fails exactly when the
low-energy directions of \(\mathsfbi{C}_{00}\) intersect the high-gain directions of the evolution
operator, and that DCT delivers its substantial savings only when the covariance spectrum decays
rapidly and the dynamics stay benign, namely for orthogonal operators. In this sense one does not
know \emph{a priori} the value of \(r\) that must be chosen, and reaching a reasonable estimate for
covariances with a slowly decaying spectrum may require \(r\sim N\), for which
\begin{equation}
  \mathcal{W}_{\mathrm{DCT}}=N\,N_t.
  \label{eq:deterministic_full_rank_cost}
\end{equation}

Summarizing, DCT carries two advantages, namely that it requires no adjoint problem, and that a
low-rank or rapidly decaying spectrum with \(\mathrm{rank}(\mathsfbi{C}_{00})\ll N\) keeps
\(r\sim O(1)\). It carries one disadvantage, namely that a correlation with a broad or full
spectrum, \(\mathrm{rank}(\mathsfbi{C}_{00})=N\), leads to large errors unless we choose the
directions of large gain carefully, and those directions are in principle not known \emph{a priori}
for non-orthogonal operators.
\subsection{Stochastic stepping trace estimation}
\label{sec:appendix_sste_method}

SSTE has been described in \cref{subsec:randomized_trace_estimation} and is not repeated here. Let
\(N_p\) be the number of probes used for the estimation, and let the cost account for the adjoint
stepper as well, so that for a fixed time
\begin{equation}
  \mathcal{C}_{\mathrm{SSTE}}(T)=6N_p .
  \label{eq:sste_fixed_time_cost}
\end{equation}
The real advantage at a single target time is then guaranteed whenever
\(\mathcal{C}_{\mathrm{SSTE}}\ll\mathcal{C}_{\mathrm{DCT}}\), that is, whenever \(N_p\ll r/6\). We
expect this in problems whose dynamics are non-orthogonal, because the sketch captures those
directions through the action on random probes, and this is especially suited for fluid-dynamics
applications, where the operators are known to be highly non-normal.

One major bottleneck nevertheless remains for an evaluation across \(N_t\) intervals, where the
cost grows as \(O(N_t^{2})\), because the backward loop must be performed several times over the
same instants in time, as \cref{fig:sste_time_history_cost} reports. The cost in this case is given by
\begin{equation}
  \sum_{j=1}^{N_t} j
  =
  \frac{N_t(N_t+1)}{2}
  \sim
  \frac{1}{2}N_t^{2},
  \label{eq:sste_backward_sum}
\end{equation}
which brings the total stepping work to
\begin{equation}
  \mathcal{W}_{\mathrm{SSTE}}\sim3N_pN_t^{2}.
  \label{eq:sste_multiple_time_cost}
\end{equation}

\begin{figure}[t]
  \centering
  \begin{tikzpicture}[x=1.15cm, y=1cm, >=Stealth, line width=0.7pt]
    % one row per target time: the forward sweep is stored once, while the
    % adjoint loop restarts from every target time and repeats the same intervals
    \foreach \row/\nback/\col in {0/1/blue, 1/2/green!55!black, 2/3/magenta}{%
      \begin{scope}[yshift={-\row*2.0cm}]
        \draw[black] (-0.3,0) -- (7.3,0);
        \foreach \i in {0,...,7}{\fill (\i,0) circle (1.5pt);}
        \foreach \i in {0,...,6}{\draw[orange,->] (\i,0) to[bend left=55] (\i+1,0);}
        \node[anchor=south, font=\footnotesize] at (3.5,0.62)
          {forward: solved once, \(\hat{\mathbf{q}}(t)\) stored};
        \node[anchor=west, font=\footnotesize] at (7.6,0) {\(T_{\the\numexpr\row+1\relax}\)};
      \end{scope}
    }
    \foreach \j in {1}{\draw[blue,->] (\j,-0.1) to[bend left=55] (\j-1,-0.1);}
    \begin{scope}[yshift=-2.0cm]
      \foreach \j in {1,2}{\draw[green!55!black,->] (\j,-0.1) to[bend left=55] (\j-1,-0.1);}
    \end{scope}
    \begin{scope}[yshift=-4.0cm]
      \foreach \j in {1,2,3}{\draw[magenta,->] (\j,-0.1) to[bend left=55] (\j-1,-0.1);}
    \end{scope}
  \end{tikzpicture}
  \caption{Stepping work of SSTE over a history of target times. The forward solve is performed
  once and stored, whereas the adjoint loop restarts from each target time \(T_j\) and traverses
  the same intervals again, so that the total work grows as \(O(N_t^{2})\).}
  \label{fig:sste_time_history_cost}
\end{figure}

The method must therefore always be preferred whenever
\(\mathcal{W}_{\mathrm{SSTE}}\ll\mathcal{W}_{\mathrm{DCT}}\), that is, whenever
\(3N_pN_t\ll r\). This can be useful when \(r\sim N\) and \(N_p\sim O(1)\), meaning for
problems whose dynamics are low rank and whose initial disturbances carry a slowly decaying
correlation spectrum. The metric that suggests which method to use over the other is therefore
\begin{equation}
  3N_pN_t\ll r ,
  \label{eq:sste_vs_dct_criterion}
\end{equation}
which means that, with \(N_p\sim O(1)\) and \(N_t\sim O(10)\), we must prefer SSTE only when
\(r>O(100)\). This holds particularly for large two- and three-dimensional problems, where the
initial correlation carries a slowly decaying spectrum.

What we show below is that \(N_p\sim O(1)\) for every case considered here, so that a small
\(N_t\) makes SSTE preferable. The real advantage of SSTE is that it needs no \emph{a priori}
knowledge of \(\mathsfbi{C}_{00}\), which makes it flexible over a large class of problems;
for \(N_t\gg1\), however, the cost becomes very large and parallelization in time becomes
necessary. Adjoint methods are known to carry this bottleneck and, as the structure of the problem
shows, they are very amenable to parallelization in time, as discussed by
\citet{skene-eggl-schmid-2021}; we do not discuss such strategies here.
\subsection{Verification}
\label{sec:appendix_verification}
The verification cases are organized as follows. \Cref{sec:appendix_verification_gl} uses an autonomous complex Ginzburg--Landau problem, for which the exact matrix-exponential reference is available, and the same is done in \cref{sec:appendix_verification_poiseuille} for plane Poiseuille flow, providing a fluid-dynamical test of the same trace construction. \Cref{sec:appendix_verification_nonautonomous} uses a fully non-autonomous scalar problem with a closed-form trace reference as a final test for a fully non-autonomous problem.

\subsubsection{Complex Ginzburg--Landau}
\label{sec:appendix_verification_gl}
\begin{figure}
    \centering
    \begin{minipage}[t]{0.49\textwidth}
        \centering
        \vspace{0pt}
        % This file was created by matlab2tikz.
%
\definecolor{mycolor1}{rgb}{0.20000,0.40000,0.60000}%
%
\begin{tikzpicture}

\begin{axis}[%
width=11.252in,
height=7.131in,
at={(1.887in,0.962in)},
scale only axis,
clip=true,
separate axis lines,
every outer x axis line/.append style={black},
every x tick label/.append style={font=\color{black}},
every x tick/.append style={black},
xmin=0,
xmax=80,
xlabel={stepper calls},
every outer y axis line/.append style={black},
every y tick label/.append style={font=\color{black}},
every y tick/.append style={black},
ymode=log,
ymin=0.0001,
ymax=1,
yminorticks=true,
ylabel={$\epsilon_{\mathrm{rel}}(T)$},
axis background/.style={fill=white},
clip=true,
clip mode=individual,
axis lines=box,
xtick pos=bottom,
ytick pos=left,
tick align=inside,
major tick length=2pt,
width=0.7\linewidth,
height=0.65\linewidth,
every axis/.append style={font=\fontsize{9}{9}\selectfont},xlabel style={font=\fontsize{9}{9}\selectfont},title style={font=\fontsize{9}{9}\selectfont},
legend style={font=\fontsize{8}{9}\selectfont,inner sep=1pt,row sep=1pt,column sep=2pt,nodes={scale=1.00},draw=none},legend image post style={xscale=0.35},
legend columns=1,
scaled x ticks=false,
tick label style={/pgf/number format/fixed,/pgf/number format/precision=2},
every x tick label/.append style={font=\fontsize{9}{9}\selectfont\color{black}},
every y tick label/.append style={font=\fontsize{9}{9}\selectfont\color{black}},
,,
colormap={sigmamap}{rgb(0.0000)=(0.0200,0.1880,0.3800); rgb(0.1000)=(0.0743,0.3456,0.5616); rgb(0.2000)=(0.1918,0.4980,0.7060); rgb(0.3000)=(0.3890,0.6708,0.8226); rgb(0.4000)=(0.7552,0.8682,0.9228); rgb(0.5000)=(0.9690,0.9689,0.9689); rgb(0.6000)=(0.9425,0.8044,0.7614); rgb(0.7000)=(0.8767,0.4910,0.4020); rgb(0.8000)=(0.7869,0.2866,0.2470); rgb(0.9000)=(0.6186,0.1239,0.1568); rgb(1.0000)=(0.4040,0.0000,0.1220)}
]
\addplot [color=black, line width=1.3pt, forget plot]
  table[row sep=crcr]{%
1	0.985165464502917\\
2	0.971048359830997\\
3	0.951205179050231\\
4	0.919681047700775\\
5	0.899249412820878\\
6	0.880805994430103\\
7	0.868459742494046\\
8	0.865199492680536\\
9	0.861114197062287\\
10	0.832525113113597\\
11	0.822398689108786\\
12	0.788670148690936\\
13	0.775021625627415\\
14	0.764549352578276\\
15	0.747746973613479\\
16	0.745712139733251\\
17	0.724965063859842\\
18	0.713974519830427\\
19	0.701444244155202\\
20	0.688482774619424\\
21	0.663182110301451\\
22	0.65435357265626\\
23	0.64407505382054\\
24	0.624129625964055\\
25	0.618414678614345\\
26	0.614229876610716\\
27	0.600565092768903\\
28	0.597031889384245\\
29	0.591240547176463\\
30	0.579989136718529\\
31	0.574332747631998\\
32	0.558906071472685\\
33	0.551424866729417\\
34	0.539182540890269\\
35	0.533524060767838\\
36	0.52717094141623\\
37	0.523571656178599\\
38	0.516042372207813\\
39	0.512332268723539\\
40	0.509670785497178\\
41	0.500782949538627\\
42	0.496546792309842\\
43	0.494342949858862\\
44	0.474330601756383\\
45	0.45263687517021\\
46	0.436463402022355\\
47	0.41765607247181\\
48	0.397465607198713\\
49	0.385282284592009\\
50	0.337283356987566\\
51	0.328513798162926\\
52	0.309670329668666\\
53	0.301926775381416\\
54	0.299026111763803\\
55	0.288612391923131\\
56	0.264457019190453\\
57	0.254980536533236\\
58	0.250205782373461\\
59	0.245812576451499\\
60	0.237374855379344\\
61	0.230549355537883\\
62	0.224861889966195\\
63	0.212799059810108\\
64	0.209293996703853\\
65	0.205196913687323\\
66	0.203571124582919\\
67	0.188035812467885\\
68	0.177630310730338\\
69	0.16372654312346\\
70	0.138302991694355\\
71	0.122002256923464\\
72	0.119097970006177\\
73	0.107519625210511\\
74	0.105067727551673\\
75	0.0992037278314537\\
76	0.0754210347718126\\
77	0.0602465034532608\\
78	0.0366705573794107\\
79	0.019419822005287\\
80	0.000197351436241734\\
};
\addplot [color=black, dashed, line width=1.1pt, mark size=2.0pt, mark=o, mark options={solid, black}, forget plot]
  table[row sep=crcr]{%
6	0.335159987761167\\
18	0.118531569790602\\
30	0.00739365023044105\\
42	0.00417927746508909\\
54	0.00280275867247948\\
66	0.001325154662104\\
78	0.000274288373235347\\
};
\addplot [color=white!35!black, line width=1.3pt, forget plot]
  table[row sep=crcr]{%
1	0.978599504948092\\
2	0.958939501992847\\
3	0.939693066967005\\
4	0.89261026284617\\
5	0.867840917394025\\
6	0.861406077656702\\
7	0.852698239057536\\
8	0.85017245615521\\
9	0.847032144185355\\
10	0.80729432471463\\
11	0.793416282120727\\
12	0.765630153360495\\
13	0.74988185038214\\
14	0.740470036675323\\
15	0.716426184465566\\
16	0.715341748661657\\
17	0.691584384991789\\
18	0.683455781101978\\
19	0.670675562334103\\
20	0.654454550999417\\
21	0.637112149433065\\
22	0.633257778866529\\
23	0.620317618750451\\
24	0.588720020973351\\
25	0.585945706960771\\
26	0.582267838430227\\
27	0.56742924037083\\
28	0.565364720145365\\
29	0.562807616623031\\
30	0.555634500442104\\
31	0.553314809191816\\
32	0.54238243410962\\
33	0.539039247828093\\
34	0.532643096219939\\
35	0.530584870874834\\
36	0.522539364472817\\
37	0.521445234589325\\
38	0.513154688684804\\
39	0.508839959353012\\
40	0.505939514085901\\
41	0.500655814429845\\
42	0.499712265513766\\
43	0.498404391560835\\
44	0.478878943258887\\
45	0.445626597424573\\
46	0.428464859611121\\
47	0.398258275036155\\
48	0.376128564976229\\
49	0.362789243735769\\
50	0.307295386037411\\
51	0.293601467871381\\
52	0.263957921658544\\
53	0.251759619862065\\
54	0.249676687234887\\
55	0.242489853789401\\
56	0.219495168292962\\
57	0.207859611903422\\
58	0.205754176866974\\
59	0.204741582868334\\
60	0.199458164115086\\
61	0.193520310988052\\
62	0.190931637714247\\
63	0.173087985360088\\
64	0.169340146062514\\
65	0.167072941945809\\
66	0.166542467767846\\
67	0.151614680982998\\
68	0.145368741746296\\
69	0.139618579434393\\
70	0.115866788291259\\
71	0.104232963892324\\
72	0.101843105302551\\
73	0.0955708834689444\\
74	0.0938174024887371\\
75	0.0906079543094753\\
76	0.0657956360552881\\
77	0.0418853536153708\\
78	0.0228826914806503\\
79	0.00138771895001788\\
80	0.0156460030015282\\
};
\addplot [color=white!35!black, dashed, line width=1.1pt, mark size=2.0pt, mark=o, mark options={solid, white!35!black}, forget plot]
  table[row sep=crcr]{%
6	0.178634994647454\\
18	0.0272494655305371\\
30	0.0157452937165598\\
42	0.0156391308471297\\
54	0.0156432097817025\\
66	0.0156460997728516\\
78	0.0156460065903751\\
};
\addplot [color=white!60!black, line width=1.3pt, forget plot]
  table[row sep=crcr]{%
1	0.983939814725101\\
2	0.958621292050389\\
3	0.941410908329715\\
4	0.890755899497388\\
5	0.861457015736498\\
6	0.859985130675222\\
7	0.845697130140114\\
8	0.84384524991711\\
9	0.840712634135543\\
10	0.806066660950854\\
11	0.785220616474763\\
12	0.776763649448303\\
13	0.761103340890232\\
14	0.75261633371115\\
15	0.740564558251317\\
16	0.739651588083636\\
17	0.686174570596409\\
18	0.676538220042228\\
19	0.65747615287434\\
20	0.651761925296852\\
21	0.640299018596555\\
22	0.636199216480547\\
23	0.610879445340382\\
24	0.573000938865489\\
25	0.570269997442903\\
26	0.568735320710756\\
27	0.550363133482551\\
28	0.546783209951033\\
29	0.545147664528803\\
30	0.530625903954825\\
31	0.528389559501782\\
32	0.508318851954079\\
33	0.503980652374716\\
34	0.499613192675354\\
35	0.49763637312336\\
36	0.490959101835251\\
37	0.490643493508232\\
38	0.472132561173985\\
39	0.46529930178891\\
40	0.459836536188394\\
41	0.454070526874192\\
42	0.453084808068927\\
43	0.450921515878575\\
44	0.446915469987828\\
45	0.423482412433708\\
46	0.395423667568545\\
47	0.367204866683023\\
48	0.356292784685205\\
49	0.341663964243568\\
50	0.306891857435353\\
51	0.285139234618715\\
52	0.246301072832208\\
53	0.231967547990899\\
54	0.230547619059275\\
55	0.221083463022862\\
56	0.192621792750176\\
57	0.187439295663882\\
58	0.184303648641886\\
59	0.182138524821364\\
60	0.179570703186825\\
61	0.171154224360516\\
62	0.167433487565166\\
63	0.145003386402656\\
64	0.139611153942317\\
65	0.139222997248754\\
66	0.138620251364076\\
67	0.130288551990365\\
68	0.113200339183079\\
69	0.112642085997531\\
70	0.101776616986649\\
71	0.0890035156036651\\
72	0.0854020580954648\\
73	0.077391766903902\\
74	0.0762840444580965\\
75	0.0753165799713176\\
76	0.0680914430609462\\
77	0.0523770125267109\\
78	0.0446863165468323\\
79	0.0113900519737571\\
80	0.00109306870352377\\
};
\addplot [color=white!60!black, dashed, line width=1.1pt, mark size=2.0pt, mark=o, mark options={solid, white!60!black}, forget plot]
  table[row sep=crcr]{%
6	0.0432217565843151\\
18	0.00104058419757173\\
30	0.00109300796967562\\
42	0.00109306851668768\\
54	0.00109306870455549\\
66	0.00109306870351742\\
78	0.00109306870352434\\
};
\addplot [color=mycolor1, line width=1.3pt, forget plot]
  table[row sep=crcr]{%
1	0.989811891709964\\
2	0.960431859012181\\
3	0.947709764979468\\
4	0.902376300114507\\
5	0.881075920619318\\
6	0.876674455805922\\
7	0.856885350128844\\
8	0.855090827494591\\
9	0.851761640525362\\
10	0.819320189184422\\
11	0.797292548273364\\
12	0.78806122225905\\
13	0.77276006590294\\
14	0.764282501804924\\
15	0.76071644434043\\
16	0.759703269063204\\
17	0.691619520587041\\
18	0.67967319361489\\
19	0.666621616328036\\
20	0.665504318632122\\
21	0.658670748743334\\
22	0.653705159150912\\
23	0.621556267060521\\
24	0.579768859923312\\
25	0.576997067249284\\
26	0.575868845616789\\
27	0.562810814189649\\
28	0.55797512494194\\
29	0.555658984500403\\
30	0.533488707883595\\
31	0.530480999056352\\
32	0.498123892171704\\
33	0.491116268851211\\
34	0.487948530398078\\
35	0.485573568159683\\
36	0.479551659861754\\
37	0.479359417296565\\
38	0.452077992888222\\
39	0.440933079350967\\
40	0.436152475342342\\
41	0.428890724024518\\
42	0.426217957695805\\
43	0.42258233338666\\
44	0.422254731169116\\
45	0.409952512737682\\
46	0.381810352992008\\
47	0.358308180091778\\
48	0.352175878101747\\
49	0.333306977167566\\
50	0.30424016970697\\
51	0.284178935698881\\
52	0.245756387656101\\
53	0.232000781397229\\
54	0.231310588412865\\
55	0.220004381709227\\
56	0.191907614754701\\
57	0.190463142463489\\
58	0.184969423564124\\
59	0.180402954768267\\
60	0.174144315703504\\
61	0.165465646833326\\
62	0.158712453070042\\
63	0.139617470060995\\
64	0.135879800133814\\
65	0.135207205933744\\
66	0.134716855584887\\
67	0.130151348694189\\
68	0.103306056164728\\
69	0.103176492744065\\
70	0.0941986632569092\\
71	0.0758320657443923\\
72	0.0710547597258276\\
73	0.0563570401697191\\
74	0.0558645644393192\\
75	0.0547663759366945\\
76	0.0531627847504146\\
77	0.0449587446280415\\
78	0.036338931906693\\
79	0.00714351915577783\\
80	0.002055929390514\\
};
\addplot [color=mycolor1, dashed, line width=1.1pt, mark size=2.0pt, mark=o, mark options={solid, mycolor1}, forget plot]
  table[row sep=crcr]{%
6	0.0221710184949555\\
18	0.00205572676526485\\
30	0.00205592952111303\\
42	0.00205592939052286\\
54	0.0020559293905168\\
66	0.00205592939051692\\
78	0.0020559293905112\\
};
\addplot [color=red!60!teal, line width=1.3pt, forget plot]
  table[row sep=crcr]{%
1	0.992325618839863\\
2	0.961016066817814\\
3	0.951208342921825\\
4	0.909270854860793\\
5	0.89254895682782\\
6	0.884815513639991\\
7	0.862290154659604\\
8	0.860415877005582\\
9	0.857853229850865\\
10	0.822332385875221\\
11	0.801433628859122\\
12	0.787904543949321\\
13	0.774317514752259\\
14	0.765141215814968\\
15	0.762530829118749\\
16	0.761325433853939\\
17	0.694682953325543\\
18	0.682341731535508\\
19	0.674346656399483\\
20	0.672786566648059\\
21	0.668200119625061\\
22	0.663293053353033\\
23	0.629081819149951\\
24	0.586888228823027\\
25	0.584317228262044\\
26	0.581314029745595\\
27	0.571490377868498\\
28	0.566770176814566\\
29	0.563935424311444\\
30	0.537455260735946\\
31	0.53422339882455\\
32	0.497334856023746\\
33	0.488809171678629\\
34	0.486072051456463\\
35	0.483504725005103\\
36	0.478304543183782\\
37	0.478240020525966\\
38	0.446736796169955\\
39	0.433120419580013\\
40	0.429548615138443\\
41	0.422298926512988\\
42	0.417698915367519\\
43	0.413275061006734\\
44	0.412531808256351\\
45	0.40388650613073\\
46	0.379494267012634\\
47	0.358844828047536\\
48	0.354638192262945\\
49	0.333453830121188\\
50	0.307032096596556\\
51	0.289435507933016\\
52	0.251135032099835\\
53	0.238193191630314\\
54	0.237783085631403\\
55	0.226677193556239\\
56	0.199360448424778\\
57	0.198959877325978\\
58	0.192612715272565\\
59	0.186689944541466\\
60	0.177117654749482\\
61	0.168016765002079\\
62	0.158442443672403\\
63	0.142320577140866\\
64	0.140036078567272\\
65	0.138761648579695\\
66	0.138314633879418\\
67	0.132805255890634\\
68	0.105232681267841\\
69	0.105081998102503\\
70	0.0959096699196229\\
71	0.0759144093102118\\
72	0.070894907510888\\
73	0.0519731430204992\\
74	0.0517514147346468\\
75	0.0501827934196301\\
76	0.049479744716153\\
77	0.0439304260189536\\
78	0.0333957003281456\\
79	0.00955117552443943\\
80	0.000288500251500087\\
};
\addplot [color=red!60!teal, dashed, line width=1.1pt, mark size=2.0pt, mark=o, mark options={solid, red!60!teal}, forget plot]
  table[row sep=crcr]{%
6	0.000156112389094138\\
18	0.000288490466867784\\
30	0.000288500251553511\\
42	0.00028850025150066\\
54	0.000288500251500851\\
66	0.000288500251497798\\
78	0.000288500251498943\\
};
\node[right, align=left, inner sep=0]
at (rel axis cs:-0.15,1.1) {$(a)$};
\end{axis}
\end{tikzpicture}%
    \end{minipage}
    \hfill
    \begin{minipage}[t]{0.49\textwidth}
        \centering
        \vspace{0pt}
        % This file was created by matlab2tikz.
%
\definecolor{mycolor1}{rgb}{0.20000,0.40000,0.60000}%
%
\begin{tikzpicture}

\begin{axis}[%
width=11.252in,
height=7.131in,
at={(1.887in,0.962in)},
scale only axis,
clip=true,
separate axis lines,
every outer x axis line/.append style={black},
every x tick label/.append style={font=\color{black}},
every x tick/.append style={black},
xmin=0,
xmax=80,
xlabel={stepper calls},
every outer y axis line/.append style={black},
every y tick label/.append style={font=\color{black}},
every y tick/.append style={black},
ymode=log,
ymin=1e-05,
ymax=1,
yminorticks=true,
axis background/.style={fill=white},
legend style={legend cell align=left, align=left},
clip=true,
clip mode=individual,
axis lines=box,
xtick pos=bottom,
ytick pos=left,
tick align=inside,
major tick length=2pt,
width=0.7\linewidth,
height=0.65\linewidth,
every axis/.append style={font=\fontsize{9}{9}\selectfont},xlabel style={font=\fontsize{9}{9}\selectfont},title style={font=\fontsize{9}{9}\selectfont},
legend style={font=\fontsize{8}{9}\selectfont,inner sep=1pt,row sep=1pt,column sep=2pt,nodes={scale=1.00},draw=none},legend image post style={xscale=0.35},
legend columns=2,
scaled x ticks=false,
tick label style={/pgf/number format/fixed,/pgf/number format/precision=2},
every x tick label/.append style={font=\fontsize{9}{9}\selectfont\color{black}},
every y tick label/.append style={font=\fontsize{9}{9}\selectfont\color{black}},
,,
colormap={sigmamap}{rgb(0.0000)=(0.0200,0.1880,0.3800); rgb(0.1000)=(0.0743,0.3456,0.5616); rgb(0.2000)=(0.1918,0.4980,0.7060); rgb(0.3000)=(0.3890,0.6708,0.8226); rgb(0.4000)=(0.7552,0.8682,0.9228); rgb(0.5000)=(0.9690,0.9689,0.9689); rgb(0.6000)=(0.9425,0.8044,0.7614); rgb(0.7000)=(0.8767,0.4910,0.4020); rgb(0.8000)=(0.7869,0.2866,0.2470); rgb(0.9000)=(0.6186,0.1239,0.1568); rgb(1.0000)=(0.4040,0.0000,0.1220)}
]
\addplot [color=black, line width=1.3pt]
  table[row sep=crcr]{%
1	0.832142165809318\\
2	0.698807622707636\\
3	0.546807535646839\\
4	0.353256308468127\\
5	0.2533545113672\\
6	0.181829801602638\\
7	0.143961931458514\\
8	0.136069474952372\\
9	0.128276084439533\\
10	0.0853515771514262\\
11	0.0733972032655607\\
12	0.0421167159491625\\
13	0.032179511831854\\
14	0.026197264594499\\
15	0.0186703519100572\\
16	0.0179558484860287\\
17	0.0122477491436808\\
18	0.00987930822292705\\
19	0.00776495043096266\\
20	0.00605286229772802\\
21	0.00343739612353565\\
22	0.00272330657881337\\
23	0.00207295030764483\\
24	0.00108590094293048\\
25	0.000864740726484522\\
26	0.000738121574372372\\
27	0.000414905706821819\\
28	0.00034958297313094\\
29	0.000265902270606906\\
30	0.00013885897151188\\
31	8.89550707338264e-05\\
32	1.73790854390435e-05\\
33	5.76632853457698e-05\\
34	0.000109156501001198\\
35	0.00012774614898567\\
36	0.000144046922404943\\
37	0.00015125888994871\\
38	0.000163039715967061\\
39	0.000167572474343433\\
40	0.000170111260942139\\
41	0.000176730335302363\\
42	0.000179193232154384\\
43	0.000180193469332686\\
44	0.000187283454047588\\
45	0.000193282521925197\\
46	0.000196773393056463\\
47	0.000199941632828961\\
48	0.00020259610362408\\
49	0.000203846130160141\\
50	0.000207689293072777\\
51	0.000208237218434015\\
52	0.000209155919455236\\
53	0.000209450502219311\\
54	0.000209536602094388\\
55	0.000209777778221061\\
56	0.000210214244888818\\
57	0.000210347836618805\\
58	0.000210400349410914\\
59	0.000210438042734826\\
60	0.000210494518788903\\
61	0.00021053015716279\\
62	0.000210553322293998\\
63	0.000210591647673277\\
64	0.000210600334145073\\
65	0.000210608254088011\\
66	0.000210610705413184\\
67	0.000210628975313441\\
68	0.000210638519721798\\
69	0.00021064846641863\\
70	0.000210662651566689\\
71	0.000210669744848924\\
72	0.000210670730477753\\
73	0.00021067379487986\\
74	0.00021067430095941\\
75	0.000210675244846075\\
76	0.000210678230149818\\
77	0.000210679715527429\\
78	0.000210681515133548\\
79	0.000210682541956614\\
80	0.000210683452492522\\
};
\addlegendentry{$T=2$}

\addplot [color=black, dashed, line width=1.1pt, mark size=2.0pt, mark=o, mark options={solid, black}, forget plot]
  table[row sep=crcr]{%
6	0.38634485377839\\
18	0.0742417278942887\\
30	0.00635197919233409\\
42	0.00924881520992697\\
54	0.00128367921167339\\
66	0.000206791538716321\\
78	0.000207192745182013\\
};
\addplot [color=white!35!black, line width=1.3pt]
  table[row sep=crcr]{%
1	0.794201241174241\\
2	0.636392197071579\\
3	0.511097089381333\\
4	0.265419020726859\\
5	0.162490441026057\\
6	0.141282246946287\\
7	0.118583735195212\\
8	0.113387263024999\\
9	0.108295982662846\\
10	0.0575898627114288\\
11	0.0436663052599052\\
12	0.0217657425047832\\
13	0.0120211522516976\\
14	0.0074518699460532\\
15	0.00170193007565686\\
16	0.00202554643688287\\
17	0.00758054743888711\\
18	0.00906925911452958\\
19	0.0109020267657692\\
20	0.0127229895482339\\
21	0.0142466131427801\\
22	0.014511566017313\\
23	0.0152074097413436\\
24	0.0165363359023962\\
25	0.0166275793436142\\
26	0.0167221533865549\\
27	0.017020440194329\\
28	0.017052879094671\\
29	0.0170842802529692\\
30	0.0171531143988702\\
31	0.0171705074859301\\
32	0.0172345496823139\\
33	0.0172498491050995\\
34	0.0172727133096399\\
35	0.0172784599704133\\
36	0.0172960038585485\\
37	0.0172978670507015\\
38	0.01730889148876\\
39	0.0173133715233898\\
40	0.0173157228794942\\
41	0.0173190670670061\\
42	0.0173195332858691\\
43	0.017320037760821\\
44	0.0173259167076527\\
45	0.0173337316030692\\
46	0.0173368796634658\\
47	0.0173412042420974\\
48	0.0173436768691506\\
49	0.0173448400265504\\
50	0.0173486162131316\\
51	0.0173493433608628\\
52	0.0173505716335023\\
53	0.017350966015995\\
54	0.0173510185610251\\
55	0.0173511600156901\\
56	0.0173515131306214\\
57	0.0173516525332058\\
58	0.0173516722123954\\
59	0.0173516795960179\\
60	0.0173517096502393\\
61	0.0173517359992392\\
62	0.017351744959992\\
63	0.0173517931405994\\
64	0.0173518010342846\\
65	0.0173518047589665\\
66	0.0173518054387212\\
67	0.0173518203585172\\
68	0.0173518252274636\\
69	0.0173518287235146\\
70	0.0173518399862865\\
71	0.0173518442887167\\
72	0.0173518449779994\\
73	0.0173518463888206\\
74	0.0173518466964082\\
75	0.0173518471354517\\
76	0.0173518497824039\\
77	0.0173518517715081\\
78	0.0173518530042554\\
79	0.0173518540916225\\
80	0.0173518547635486\\
};
\addlegendentry{$T=5$}

\addplot [color=white!35!black, dashed, line width=1.1pt, mark size=2.0pt, mark=o, mark options={solid, white!35!black}, forget plot]
  table[row sep=crcr]{%
6	0.114627266417663\\
18	0.0245110260665455\\
30	0.0173550749753947\\
42	0.0173553372422539\\
54	0.0173517356824942\\
66	0.017351859167501\\
78	0.0173518546728774\\
};
\addplot [color=white!60!black, line width=1.3pt]
  table[row sep=crcr]{%
1	0.84876248212394\\
2	0.649751643309823\\
3	0.540037042350577\\
4	0.281205903590061\\
5	0.161982206992004\\
6	0.157231809946868\\
7	0.120760737936852\\
8	0.117029821354168\\
9	0.112056444044837\\
10	0.0687652892723018\\
11	0.048285017579275\\
12	0.0417577423341335\\
13	0.0322687478482075\\
14	0.0282339730743595\\
15	0.0237409519465848\\
16	0.0234741596608912\\
17	0.0112296011081531\\
18	0.00950138804760071\\
19	0.00682450635838535\\
20	0.00619634569975223\\
21	0.00521017233758741\\
22	0.00493419836620897\\
23	0.00360091672537679\\
24	0.00204089843628622\\
25	0.00195294587640533\\
26	0.00191430181873946\\
27	0.00155264879828137\\
28	0.00149756666011107\\
29	0.00147789912083002\\
30	0.00134143911241692\\
31	0.00132501903343994\\
32	0.00120988481005397\\
33	0.00119044401344739\\
34	0.00117515581497276\\
35	0.00116975101375087\\
36	0.0011554929174147\\
37	0.00115496662479683\\
38	0.00113086248234519\\
39	0.00112391470680959\\
40	0.00111957803700095\\
41	0.00111600433861349\\
42	0.00111552739337057\\
43	0.00111471028631988\\
44	0.00111352913783512\\
45	0.0011081362684778\\
46	0.00110309615758645\\
47	0.00109914002544802\\
48	0.00109794609075457\\
49	0.00109669697086278\\
50	0.00109437995187906\\
51	0.00109324886331059\\
52	0.00109167301886039\\
53	0.00109121922217887\\
54	0.00109118414593429\\
55	0.00109100173467356\\
56	0.00109057373932463\\
57	0.00109051293803442\\
58	0.00109048423798858\\
59	0.00109046877807932\\
60	0.00109045447447922\\
61	0.00109041790193447\\
62	0.00109040528986774\\
63	0.0010903459823417\\
64	0.00109033486099181\\
65	0.00109033423654462\\
66	0.00109033348021297\\
67	0.00109032532582996\\
68	0.00109031228143468\\
69	0.0010903119490669\\
70	0.00109030690376152\\
71	0.00109030227805846\\
72	0.00109030126088857\\
73	0.0010902994965216\\
74	0.00109029930624357\\
75	0.00109029917664437\\
76	0.00109029842187397\\
77	0.00109029714172199\\
78	0.00109029665316491\\
79	0.00109029500376933\\
80	0.00109029460601764\\
};
\addlegendentry{$T=10$}

\addplot [color=white!60!black, dashed, line width=1.1pt, mark size=2.0pt, mark=o, mark options={solid, white!60!black}, forget plot]
  table[row sep=crcr]{%
6	0.0296986393320929\\
18	0.00102941666107282\\
30	0.00109027663329274\\
42	0.00109029472584195\\
54	0.00109029460615259\\
66	0.00109029460601256\\
78	0.00109029460601087\\
};
\addplot [color=mycolor1, line width=1.3pt]
  table[row sep=crcr]{%
1	0.894315116322509\\
2	0.639924565807262\\
3	0.550585211968434\\
4	0.29541887879764\\
5	0.19993958879838\\
6	0.184291417981773\\
7	0.128648029366811\\
8	0.12466547313007\\
9	0.118843197055609\\
10	0.0741895421110601\\
11	0.0503504136921993\\
12	0.042501822485729\\
13	0.0322887958668297\\
14	0.0278491723748038\\
15	0.0263846834294881\\
16	0.0260585377750902\\
17	0.00888618726500854\\
18	0.00652609261204634\\
19	0.00450711030892488\\
20	0.00437181197927593\\
21	0.00372419760204234\\
22	0.0033559953999986\\
23	0.00149116875713173\\
24	0.000404631592314746\\
25	0.000502966377477281\\
26	0.000534261061203403\\
27	0.000817412827319104\\
28	0.000899373545622831\\
29	0.000930054060066246\\
30	0.00115954606743714\\
31	0.00118387273760401\\
32	0.00138833905161327\\
33	0.00142293184311428\\
34	0.0014351466626162\\
35	0.00144229952283777\\
36	0.00145646421016995\\
37	0.00145681734349956\\
38	0.00149595012924684\\
39	0.00150843272503628\\
40	0.00151261330844982\\
41	0.00151757116846259\\
42	0.00151899574832936\\
43	0.00152050844936937\\
44	0.00152061485037717\\
45	0.00152373363012612\\
46	0.0015293021478584\\
47	0.00153293167946044\\
48	0.00153367078561144\\
49	0.00153544559513626\\
50	0.00153757916197315\\
51	0.00153872824909911\\
52	0.00154044556888507\\
53	0.0015409253005438\\
54	0.00154094408191749\\
55	0.00154118412931707\\
56	0.00154164954964419\\
57	0.00154166821745001\\
58	0.00154172360755486\\
59	0.0015417595257933\\
60	0.00154179792928781\\
61	0.00154183947136366\\
62	0.00154186468740169\\
63	0.00154192030448154\\
64	0.00154192879628973\\
65	0.00154192998822334\\
66	0.00154193066601282\\
67	0.00154193558818563\\
68	0.00154195816204324\\
69	0.0015419582470162\\
70	0.0015419628392149\\
71	0.00154197016611782\\
72	0.00154197165242503\\
73	0.00154197521858553\\
74	0.00154197531177196\\
75	0.00154197547382402\\
76	0.00154197565835642\\
77	0.00154197639456516\\
78	0.00154197699776023\\
79	0.0015419785909014\\
80	0.00154197898234938\\
};
\addlegendentry{$T=15$}

\addplot [color=mycolor1, dashed, line width=1.1pt, mark size=2.0pt, mark=o, mark options={solid, mycolor1}, forget plot]
  table[row sep=crcr]{%
6	0.00203990605396475\\
18	0.00154172974518009\\
30	0.0015419790075907\\
42	0.00154197898235307\\
54	0.00154197898235694\\
66	0.00154197898235731\\
78	0.00154197898235675\\
};
\addplot [color=red!60!teal, line width=1.3pt]
  table[row sep=crcr]{%
1	0.916655916931551\\
2	0.632839416655546\\
3	0.560734559075273\\
4	0.313608263459002\\
5	0.235135400577056\\
6	0.206351388400978\\
7	0.140042582106633\\
8	0.135687874596275\\
9	0.130995901556129\\
10	0.0798097946470546\\
11	0.0561312542461006\\
12	0.0440889906468997\\
13	0.0345946063616061\\
14	0.0295636019494487\\
15	0.0284412871383848\\
16	0.0280350603307864\\
17	0.0104376190014379\\
18	0.00788512184256926\\
19	0.0065903174958666\\
20	0.00639253620729769\\
21	0.00593748810345597\\
22	0.00555655425078409\\
23	0.00347899564883621\\
24	0.00147495861576603\\
25	0.00137946800118453\\
26	0.00129225679938062\\
27	0.00106924562969753\\
28	0.000985488853024056\\
29	0.00094617686728798\\
30	0.000659211687506794\\
31	0.000631845660007868\\
32	0.000387808663954991\\
33	0.00034374747666017\\
34	0.000332697949419378\\
35	0.000324602961382195\\
36	0.000311797257064023\\
37	0.000311673173719314\\
38	0.000264364495329873\\
39	0.00024839828505949\\
40	0.000245128238914654\\
41	0.000239946394987256\\
42	0.000237379566910402\\
43	0.000235452543395795\\
44	0.000235199818794693\\
45	0.000232905289889007\\
46	0.000227852328983014\\
47	0.000224513740367466\\
48	0.000223982941141446\\
49	0.000221896853141786\\
50	0.000219866449830847\\
51	0.00021881124715571\\
52	0.000217019068638015\\
53	0.00021654654148235\\
54	0.000216534858195951\\
55	0.000216288001097712\\
56	0.000215814272028608\\
57	0.000215808852322244\\
58	0.000215741855002887\\
59	0.000215693082893171\\
60	0.000215631590872829\\
61	0.000215585983924185\\
62	0.000215548556690648\\
63	0.000215499396187992\\
64	0.000215493962406096\\
65	0.000215491597973077\\
66	0.000215490951095293\\
67	0.000215484732640243\\
68	0.000215460459439456\\
69	0.000215460355978811\\
70	0.00021545544417535\\
71	0.00021544709332156\\
72	0.000215445458395092\\
73	0.000215440651942681\\
74	0.000215440608018712\\
75	0.0002154403656884\\
76	0.000215440280989658\\
77	0.000215439759644969\\
78	0.000215438987861761\\
79	0.000215437625663258\\
80	0.000215437213033184\\
};
\addlegendentry{$T=20$}

\addplot [color=red!60!teal, dashed, line width=1.1pt, mark size=2.0pt, mark=o, mark options={solid, red!60!teal}, forget plot]
  table[row sep=crcr]{%
6	0.00173146060343864\\
18	0.000215427460460972\\
30	0.000215437213018904\\
42	0.000215437213023171\\
54	0.000215437213029573\\
66	0.000215437213026782\\
78	0.000215437213026126\\
};
\node[right, align=left, inner sep=0]
at (rel axis cs:-0.15,1.1) {$(b)$};
\end{axis}
\end{tikzpicture}%
    \end{minipage}
    \caption{Ginzburg--Landau relative error versus stepper calls for \((a)\) slow and \((b)\) fast covariance-spectrum decay. (\solidlegend) deterministic truncation; (\dashlegend\circlegend) SSTE.}
    \label{fig:GL_eig_vs_sste}
\end{figure}
We compare the cost of SSTE against the deterministic truncation strategy on the linear complex Ginzburg--Landau system used in the numerical experiments. In continuous form, the governing equation is
\begin{equation}
  \frac{\partial q}{\partial t}
  =
  -\nu\frac{\partial q}{\partial x}
  +
  \gamma\frac{\partial^2 q}{\partial x^2}
  +
  \mu(x)q,
  \label{eq:GL_continuous_equation_sste_necessity}
\end{equation}
where \(q(x,t)\) is a complex amplitude, \(\nu\) is the convection coefficient, \(\gamma\) is the complex diffusion coefficient, and \(\mu(x)\) is a spatially varying growth-rate profile. In the present computations, the latter is taken in quadratic form,
\begin{equation}
  \mu(x)
  =
  \mu_0+\mu_2x^2,
  \label{eq:GL_mu_profile_sste_necessity}
\end{equation}
with parameter values
\begin{equation}
  \nu = 2+0.4\mathrm{i},
  \qquad
  \gamma = 1-\mathrm{i},
  \qquad
  \mu_2 = -0.01,
  \qquad
  \mu_0 = 0.39.
  \label{eq:GL_parameters_sste_necessity}
\end{equation}
We use \(n_x=80\), \(\Delta t=0.08\), and \(2\leq T\leq20\).
The initial covariance is generated through
\begin{equation}
  \mathsfbi{C}_{00}^{z}
  =
  \mathsfbi{Q}
  \operatorname{diag}
  \left(
    \lambda_1,\dots,\lambda_n
  \right)
  \mathsfbi{Q}^{*},
  \label{eq:GL_covariance_construction_sste_necessity}
\end{equation}
where \(\mathsfbi{Q}\) is a fixed orthonormal basis and the eigenvalues \(\lambda_j\) are chosen to impose a prescribed spectral decay. We consider two representative cases. The first is a slow algebraic decay,
\begin{equation}
  \lambda_j
  \propto
  j^{-1/10},
  \label{eq:GL_slow_covariance_decay}
\end{equation}
and the second is a fast exponential decay,
\begin{equation}
  \lambda_j
  \propto
  \exp
  \left[
    -0.25(j-1)
  \right].
  \label{eq:GL_fast_covariance_decay}
\end{equation}
Both spectra are chosen to have the same total variance.
\subsubsection{Plane Poiseuille flow}
\label{sec:appendix_verification_poiseuille}
\begin{figure}[t]
    \centering
    \begin{minipage}[t]{0.49\textwidth}
        \centering
        \vspace{0pt}
        % This file was created by matlab2tikz.
%
\definecolor{mycolor1}{rgb}{0.20000,0.40000,0.60000}%
%
\begin{tikzpicture}

\begin{axis}[%
width=11.252in,
height=7.131in,
at={(1.887in,0.962in)},
scale only axis,
clip=true,
separate axis lines,
every outer x axis line/.append style={black},
every x tick label/.append style={font=\color{black}},
every x tick/.append style={black},
xmin=0,
xmax=257,
xlabel={stepper calls},
every outer y axis line/.append style={black},
every y tick label/.append style={font=\color{black}},
every y tick/.append style={black},
ymode=log,
ymin=1e-08,
ymax=1,
yminorticks=true,
ylabel={$\epsilon_{\mathrm{rel}}(T)$},
axis background/.style={fill=white},
clip=true,
clip mode=individual,
axis lines=box,
xtick pos=bottom,
ytick pos=left,
tick align=inside,
major tick length=2pt,
width=0.7\linewidth,
height=0.65\linewidth,
every axis/.append style={font=\fontsize{9}{9}\selectfont},xlabel style={font=\fontsize{9}{9}\selectfont},title style={font=\fontsize{9}{9}\selectfont},
legend style={font=\fontsize{8}{9}\selectfont,inner sep=1pt,row sep=1pt,column sep=2pt,nodes={scale=1.00},draw=none},legend image post style={xscale=0.35},
legend columns=1,
scaled x ticks=false,
tick label style={/pgf/number format/fixed,/pgf/number format/precision=2},
every x tick label/.append style={font=\fontsize{9}{9}\selectfont\color{black}},
every y tick label/.append style={font=\fontsize{9}{9}\selectfont\color{black}},
,,
colormap={sigmamap}{rgb(0.0000)=(0.0200,0.1880,0.3800); rgb(0.1000)=(0.0743,0.3456,0.5616); rgb(0.2000)=(0.1918,0.4980,0.7060); rgb(0.3000)=(0.3890,0.6708,0.8226); rgb(0.4000)=(0.7552,0.8682,0.9228); rgb(0.5000)=(0.9690,0.9689,0.9689); rgb(0.6000)=(0.9425,0.8044,0.7614); rgb(0.7000)=(0.8767,0.4910,0.4020); rgb(0.8000)=(0.7869,0.2866,0.2470); rgb(0.9000)=(0.6186,0.1239,0.1568); rgb(1.0000)=(0.4040,0.0000,0.1220)}
]
\addplot [color=black, line width=1.3pt, forget plot]
  table[row sep=crcr]{%
1	0.982966330608578\\
2	0.97604883720077\\
3	0.972269958891538\\
4	0.970503289638935\\
5	0.958721405225083\\
6	0.957618750114593\\
7	0.949409965780001\\
8	0.946631301384151\\
9	0.943171012824689\\
10	0.941395877849633\\
11	0.940290265059246\\
12	0.938680845066753\\
13	0.933745419964527\\
14	0.930398404968719\\
15	0.922469639890038\\
16	0.920422703961159\\
17	0.899687438112519\\
18	0.890297123781761\\
19	0.86969413750324\\
20	0.865349827896549\\
21	0.863113422223451\\
22	0.856767046091549\\
23	0.84879289854553\\
24	0.840831855227552\\
25	0.837942514843372\\
26	0.837238991301756\\
27	0.835937564003596\\
28	0.834883569615981\\
29	0.821679151704008\\
30	0.821641156782489\\
31	0.819587506270877\\
32	0.815963187418029\\
33	0.800080935248923\\
34	0.792732055253705\\
35	0.791347539085762\\
36	0.789765437092575\\
37	0.789171658957207\\
38	0.785660559843496\\
39	0.785257112321728\\
40	0.785097583433237\\
41	0.778586531726282\\
42	0.768184102914871\\
43	0.766969040934379\\
44	0.758926927712864\\
45	0.756396495819457\\
46	0.755374614863844\\
47	0.74950069971379\\
48	0.748573883440318\\
49	0.747114265567734\\
50	0.746920553620366\\
51	0.74381281171101\\
52	0.743470394194\\
53	0.734622775759791\\
54	0.725486251213745\\
55	0.712441438373767\\
56	0.712120976527085\\
57	0.695902684640726\\
58	0.69471188185694\\
59	0.694228151426545\\
60	0.693348470801481\\
61	0.691123211794459\\
62	0.683595567934903\\
63	0.670517708929653\\
64	0.669396278255485\\
65	0.659972844110037\\
66	0.659284011573087\\
67	0.654108239315489\\
68	0.648872022132782\\
69	0.647004651894603\\
70	0.646898315315815\\
71	0.645320035064132\\
72	0.641495156999121\\
73	0.638336522912256\\
74	0.636245270124167\\
75	0.63561283511298\\
76	0.630013905355875\\
77	0.629300933483904\\
78	0.629273608383272\\
79	0.626588102812313\\
80	0.6249785610807\\
81	0.618051779287091\\
82	0.61760161399641\\
83	0.610116289153122\\
84	0.598733061118022\\
85	0.596754701744156\\
86	0.578336557744394\\
87	0.577520095902458\\
88	0.57668039214578\\
89	0.569192866936287\\
90	0.568684616196599\\
91	0.567781390196765\\
92	0.545174934342963\\
93	0.537206580369471\\
94	0.534168567143395\\
95	0.531660070357362\\
96	0.515279920603679\\
97	0.515160316290744\\
98	0.51217502354736\\
99	0.510102845790237\\
100	0.508441823338196\\
101	0.507293324179209\\
102	0.50625248442349\\
103	0.504533478349842\\
104	0.50418367932901\\
105	0.50101186912038\\
106	0.500475996711498\\
107	0.495556169726679\\
108	0.491526093163103\\
109	0.474939220277814\\
110	0.471976517209076\\
111	0.471397729050771\\
112	0.464951351514588\\
113	0.462909553414059\\
114	0.462755087835895\\
115	0.459685754382438\\
116	0.456650219837631\\
117	0.4533773073742\\
118	0.45239030088586\\
119	0.448771128351402\\
120	0.441373111119473\\
121	0.437878899357908\\
122	0.436567986771897\\
123	0.42766600490534\\
124	0.426962658684734\\
125	0.424411307266562\\
126	0.423093294104886\\
127	0.421487056457706\\
128	0.40845747323444\\
129	0.403086637641859\\
130	0.402033847441137\\
131	0.39963805823747\\
132	0.395906034136358\\
133	0.393585528707658\\
134	0.391587290882353\\
135	0.390649900206814\\
136	0.386536948297603\\
137	0.38607559989746\\
138	0.380611582691387\\
139	0.378319261067495\\
140	0.377583277377092\\
141	0.371470401535299\\
142	0.356947519756303\\
143	0.355342294240175\\
144	0.352648624936454\\
145	0.352217842979882\\
146	0.351991646559007\\
147	0.350450262027681\\
148	0.348728664558269\\
149	0.34495729136732\\
150	0.340417054042536\\
151	0.314960478632387\\
152	0.31224391750994\\
153	0.30858129443841\\
154	0.307276851922746\\
155	0.304599755220858\\
156	0.304382343507257\\
157	0.300892308138577\\
158	0.298336980594795\\
159	0.297890704814015\\
160	0.297098357212825\\
161	0.296242533695882\\
162	0.294960908507088\\
163	0.289463715265222\\
164	0.284838137674088\\
165	0.277855806654834\\
166	0.274504727859216\\
167	0.272738580449843\\
168	0.271597079745049\\
169	0.250008657723613\\
170	0.236175429837658\\
171	0.229519436480829\\
172	0.223306566476213\\
173	0.217704853889386\\
174	0.217075226705583\\
175	0.216849872756759\\
176	0.212100689986535\\
177	0.197922224306309\\
178	0.195473806039263\\
179	0.192569664433095\\
180	0.191886950398255\\
181	0.176321829335418\\
182	0.176204493054979\\
183	0.172463392607111\\
184	0.170621903239322\\
185	0.165314988403719\\
186	0.162816088242698\\
187	0.162526474650984\\
188	0.161678528039627\\
189	0.159804474035487\\
190	0.158932133711831\\
191	0.157800302896849\\
192	0.157306968318194\\
193	0.156737123015743\\
194	0.153987930466226\\
195	0.152326731562892\\
196	0.146403680036494\\
197	0.143332220703391\\
198	0.142404957423671\\
199	0.138318017143347\\
200	0.136175805378474\\
201	0.134661682318101\\
202	0.13436358968688\\
203	0.129508102469386\\
204	0.128403683267132\\
205	0.12543655257506\\
206	0.124809848549587\\
207	0.124055903093489\\
208	0.113575556045366\\
209	0.110608590113858\\
210	0.109552532273075\\
211	0.107969507101389\\
212	0.101027338240515\\
213	0.097803717398423\\
214	0.0955306143288746\\
215	0.0932514775609851\\
216	0.0930263447886719\\
217	0.0929092471308039\\
218	0.0899489553254883\\
219	0.0886650201623336\\
220	0.0877413646642759\\
221	0.0864065450407232\\
222	0.0838469939164371\\
223	0.0830610312376052\\
224	0.0820146517143478\\
225	0.0813842215555414\\
226	0.077576766810229\\
227	0.0745877437127478\\
228	0.0738979308037943\\
229	0.0716342426149204\\
230	0.0709763409928529\\
231	0.0699835156519321\\
232	0.0699233608336728\\
233	0.0693814538777797\\
234	0.0663670778628009\\
235	0.0618109747664514\\
236	0.0603545842356326\\
237	0.0603415309260477\\
238	0.0584113841832503\\
239	0.0556902643722308\\
240	0.0517211154263843\\
241	0.0511581131244168\\
242	0.0430304290108372\\
243	0.0421833016960384\\
244	0.0395771986799776\\
245	0.0388023074808236\\
246	0.0376868444208258\\
247	0.0364514261558841\\
248	0.0353058581982174\\
249	0.0321373391485221\\
250	0.0275465832028203\\
251	0.0192319951327931\\
252	0.0177656087315\\
253	0.0160570296940422\\
254	0.011589220689924\\
255	0.00827908511913625\\
256	0.00148125784234795\\
257	0.000799887133449348\\
};
\addplot [color=black, dashed, line width=1.1pt, mark size=2.0pt, mark=o, mark options={solid, black}, forget plot]
  table[row sep=crcr]{%
6	0.540776137824498\\
18	0.00134857593831537\\
30	0.000106169089263514\\
42	0.00103472753047721\\
54	0.000818869801991856\\
78	0.000795027686100323\\
102	0.00080652247982178\\
138	0.000797136215325989\\
186	0.000799978080205416\\
246	0.000799887317381151\\
};
\addplot [color=white!35!black, line width=1.3pt, forget plot]
  table[row sep=crcr]{%
1	0.983721731500279\\
2	0.976660572325747\\
3	0.972839765959342\\
4	0.971301944669699\\
5	0.959799388973549\\
6	0.958845785394696\\
7	0.950915762436377\\
8	0.948041225904038\\
9	0.944346971973855\\
10	0.942540613520157\\
11	0.94138602620319\\
12	0.939761896169868\\
13	0.934356346896206\\
14	0.931111268656859\\
15	0.923302424930562\\
16	0.921478664417873\\
17	0.899741624448408\\
18	0.890432187571784\\
19	0.869300494849979\\
20	0.865058581388375\\
21	0.862798872541602\\
22	0.856474746027135\\
23	0.848536004128235\\
24	0.840371331165609\\
25	0.837276405728466\\
26	0.836594436298206\\
27	0.835141284275067\\
28	0.834298718469335\\
29	0.821337155438688\\
30	0.821295692243148\\
31	0.819295008323054\\
32	0.815413158213949\\
33	0.79952914798548\\
34	0.792159017364284\\
35	0.790726732613222\\
36	0.789309976737867\\
37	0.78871132463182\\
38	0.785295804424289\\
39	0.784868817221122\\
40	0.784665853107149\\
41	0.778426668236252\\
42	0.767800574043802\\
43	0.766544033207126\\
44	0.758861657858685\\
45	0.756266684319795\\
46	0.755297466880924\\
47	0.749736860798518\\
48	0.748832270444804\\
49	0.747316222540201\\
50	0.747124160165675\\
51	0.743815661846753\\
52	0.74352402584819\\
53	0.734804186102495\\
54	0.726012073172219\\
55	0.713160171571788\\
56	0.712819168031805\\
57	0.696035404032751\\
58	0.694899194569886\\
59	0.694452784105282\\
60	0.693545123847677\\
61	0.691468075089052\\
62	0.683859289397706\\
63	0.670697255063324\\
64	0.669618074520227\\
65	0.659894134355704\\
66	0.659257138344594\\
67	0.654192946539639\\
68	0.649121825806224\\
69	0.647377004423068\\
70	0.647257749951999\\
71	0.645637324810665\\
72	0.641644297563263\\
73	0.638644926606291\\
74	0.636659888389678\\
75	0.636016696380125\\
76	0.630442991047505\\
77	0.629691127532882\\
78	0.629663571456078\\
79	0.626962015004404\\
80	0.625303929218444\\
81	0.618518590554813\\
82	0.618136407226562\\
83	0.610763401580046\\
84	0.599703852657755\\
85	0.597767926376595\\
86	0.579307343374107\\
87	0.578556908511256\\
88	0.577775110223124\\
89	0.570280511409948\\
90	0.569645275139811\\
91	0.568672815433737\\
92	0.545940826157171\\
93	0.537489471573664\\
94	0.534622988925008\\
95	0.53196759861026\\
96	0.515165461570201\\
97	0.515059546573259\\
98	0.512015654184668\\
99	0.509984678533576\\
100	0.508322690590231\\
101	0.507173934668139\\
102	0.505891996014098\\
103	0.504161690246267\\
104	0.503836246353716\\
105	0.500755667095347\\
106	0.500148189553613\\
107	0.495232399677385\\
108	0.490992829688539\\
109	0.47444739476465\\
110	0.471364808374564\\
111	0.470852558846438\\
112	0.46402871467548\\
113	0.462129610430649\\
114	0.461934175426679\\
115	0.458858161668526\\
116	0.455926172006135\\
117	0.45239406251629\\
118	0.451405565888917\\
119	0.447739464588823\\
120	0.440197321956762\\
121	0.436709575562376\\
122	0.435369708156115\\
123	0.426486648442205\\
124	0.425768117268333\\
125	0.423295217565363\\
126	0.422069531227627\\
127	0.42042549594331\\
128	0.407182141326848\\
129	0.401976405552608\\
130	0.400983187764168\\
131	0.398802595546109\\
132	0.395141527383881\\
133	0.392762339182165\\
134	0.390518532446232\\
135	0.389605097450214\\
136	0.385188083104472\\
137	0.384738975160471\\
138	0.379129008690447\\
139	0.376921521850993\\
140	0.376061413808015\\
141	0.36998203054877\\
142	0.355684850507635\\
143	0.354064865128139\\
144	0.351380734335976\\
145	0.351012165164004\\
146	0.35079893919004\\
147	0.349273321309767\\
148	0.347519892386844\\
149	0.343947425478915\\
150	0.33925139899578\\
151	0.312776128560529\\
152	0.310061507704277\\
153	0.306345270919523\\
154	0.305119137125808\\
155	0.302424694587555\\
156	0.302228278258648\\
157	0.298976198549163\\
158	0.296447995791092\\
159	0.295975825222196\\
160	0.295207626119482\\
161	0.294301927680087\\
162	0.293193725519402\\
163	0.287807795861337\\
164	0.283371540934813\\
165	0.276144097437816\\
166	0.272913117956184\\
167	0.270969306791585\\
168	0.269810561212348\\
169	0.24831889256419\\
170	0.234361131190009\\
171	0.227591528571758\\
172	0.22169980762679\\
173	0.216046146159817\\
174	0.21553423309152\\
175	0.215277881870524\\
176	0.211039071269715\\
177	0.197209807991937\\
178	0.194842712537043\\
179	0.191786914851731\\
180	0.191000872839168\\
181	0.175555094863494\\
182	0.17541158351594\\
183	0.171679137315959\\
184	0.169957135630688\\
185	0.164699124574965\\
186	0.162299552708448\\
187	0.162032621884314\\
188	0.161045389626439\\
189	0.159113085063453\\
190	0.158212445369588\\
191	0.157121746924538\\
192	0.156609347346084\\
193	0.156058045476637\\
194	0.153399999616301\\
195	0.151714045363992\\
196	0.145689946537044\\
197	0.142333898656792\\
198	0.141511082028109\\
199	0.137248939827936\\
200	0.135154630010122\\
201	0.133792968433812\\
202	0.133480603711473\\
203	0.128739464224185\\
204	0.12758439541454\\
205	0.124594275787022\\
206	0.123812912840757\\
207	0.12294785729522\\
208	0.11282879460601\\
209	0.110043883983169\\
210	0.109072275811829\\
211	0.107451659079384\\
212	0.100326949657664\\
213	0.0971007378046024\\
214	0.0949502053810914\\
215	0.0925531958522484\\
216	0.0923145226876021\\
217	0.0921888158672689\\
218	0.0890373140921958\\
219	0.0877737359853838\\
220	0.0867448467395237\\
221	0.0855261552854014\\
222	0.0830294432371949\\
223	0.0821235807980739\\
224	0.0811546294500055\\
225	0.0805359064013733\\
226	0.0769383803051045\\
227	0.0740384553315421\\
228	0.0732644994468261\\
229	0.0708651061704157\\
230	0.0702382799068362\\
231	0.0693089736284029\\
232	0.0692472339575829\\
233	0.068615548551236\\
234	0.0658212553165535\\
235	0.0615423709681331\\
236	0.0601802785121967\\
237	0.0601565782594766\\
238	0.0582951298019207\\
239	0.0553340172415406\\
240	0.0512981201613456\\
241	0.050804381136601\\
242	0.0427730740333626\\
243	0.0418490418795195\\
244	0.0391742378748484\\
245	0.0384458175943264\\
246	0.0373723632556348\\
247	0.0361811635809749\\
248	0.0349470450914869\\
249	0.0316221908211703\\
250	0.0270292946524804\\
251	0.0184742093757056\\
252	0.0167652606864055\\
253	0.0151053150621221\\
254	0.0108364747434359\\
255	0.00745626623942512\\
256	0.000740077970959662\\
257	9.5638055352644e-06\\
};
\addplot [color=white!35!black, dashed, line width=1.1pt, mark size=2.0pt, mark=o, mark options={solid, white!35!black}, forget plot]
  table[row sep=crcr]{%
6	0.0425726702224608\\
18	0.000771103631284935\\
30	5.76953839186773e-06\\
42	0.000202269655849491\\
54	4.23633235146686e-05\\
78	1.04201734344765e-05\\
102	9.67580865857114e-06\\
138	9.56323720505195e-06\\
186	9.56385438447046e-06\\
246	9.56380553744547e-06\\
};
\addplot [color=white!60!black, line width=1.3pt, forget plot]
  table[row sep=crcr]{%
1	0.984438196041516\\
2	0.977250889648189\\
3	0.973538493665872\\
4	0.972148413672703\\
5	0.961423255261589\\
6	0.960612714763109\\
7	0.953107091090466\\
8	0.950137039255104\\
9	0.946412008453305\\
10	0.944533892675864\\
11	0.943321070645482\\
12	0.941737052353571\\
13	0.935695871106321\\
14	0.932460740045847\\
15	0.924767228785686\\
16	0.92313703400658\\
17	0.90062716364848\\
18	0.891501706078173\\
19	0.869638793498805\\
20	0.8654779956666\\
21	0.86325844878803\\
22	0.857115601289254\\
23	0.84932205742142\\
24	0.841035860613004\\
25	0.837665480488096\\
26	0.83695350630032\\
27	0.835391985009747\\
28	0.834743232211089\\
29	0.822335686089452\\
30	0.822301185753557\\
31	0.820254611640779\\
32	0.816190336478491\\
33	0.800397102494706\\
34	0.793158896356638\\
35	0.791676728111207\\
36	0.790352938215864\\
37	0.789762860283539\\
38	0.786480983312684\\
39	0.786049064116082\\
40	0.785784931736868\\
41	0.779768185955226\\
42	0.768917697806536\\
43	0.767619499987759\\
44	0.760549230896024\\
45	0.757906960883045\\
46	0.756961630810407\\
47	0.751725870175965\\
48	0.750876000381561\\
49	0.74938229640779\\
50	0.749167114182767\\
51	0.745710591379976\\
52	0.745448345492567\\
53	0.737000677068402\\
54	0.728474359858584\\
55	0.715816734997077\\
56	0.715488433371106\\
57	0.698245694043465\\
58	0.697147951172875\\
59	0.696815944938359\\
60	0.695851276611026\\
61	0.693901924020054\\
62	0.686342534755478\\
63	0.673477326207909\\
64	0.672413996437189\\
65	0.662493941968813\\
66	0.661943474750917\\
67	0.656974357751622\\
68	0.652060813851287\\
69	0.650363623729753\\
70	0.650230757616197\\
71	0.648570424500669\\
72	0.644464350084388\\
73	0.641634077075451\\
74	0.639716488977034\\
75	0.639098013133221\\
76	0.633564825552144\\
77	0.632806637693438\\
78	0.632776550301435\\
79	0.630219953325068\\
80	0.6285708404064\\
81	0.621986233148742\\
82	0.621668453414749\\
83	0.614192848184915\\
84	0.603543736133059\\
85	0.601669208031684\\
86	0.582932989772586\\
87	0.582286403322812\\
88	0.581512908065036\\
89	0.573952430093363\\
90	0.573139861399523\\
91	0.57209301972925\\
92	0.549103064089356\\
93	0.54033116477838\\
94	0.537577968940599\\
95	0.534915347241\\
96	0.517616901939132\\
97	0.517510809643035\\
98	0.514522911031465\\
99	0.512563753520298\\
100	0.51095934474651\\
101	0.509812722419265\\
102	0.508199748534845\\
103	0.506433159136401\\
104	0.506115569764798\\
105	0.503014449648495\\
106	0.502329747679228\\
107	0.497637009916537\\
108	0.493165216820319\\
109	0.476377079194557\\
110	0.473122492268175\\
111	0.472685148938307\\
112	0.465691811942206\\
113	0.463867352465157\\
114	0.463609502935609\\
115	0.460618049109904\\
116	0.45777893861093\\
117	0.45404316895493\\
118	0.45300442944444\\
119	0.449400474349613\\
120	0.441498328011525\\
121	0.437897702841981\\
122	0.436533108068428\\
123	0.427400020490414\\
124	0.426577349100371\\
125	0.424238488406351\\
126	0.42304493184229\\
127	0.42129500581268\\
128	0.407733754124686\\
129	0.402654198527134\\
130	0.401633639134831\\
131	0.399591756914907\\
132	0.395978191742509\\
133	0.39349350552974\\
134	0.391049197418753\\
135	0.390162295741211\\
136	0.385401729792171\\
137	0.384903159989808\\
138	0.379140019772928\\
139	0.377020153547278\\
140	0.376069270322217\\
141	0.370188593224544\\
142	0.355990483916006\\
143	0.354338771145094\\
144	0.35152189727626\\
145	0.351174905869468\\
146	0.350971969632866\\
147	0.34943728555485\\
148	0.347614029527348\\
149	0.344316925450285\\
150	0.339589430361029\\
151	0.31200244779414\\
152	0.309349513623591\\
153	0.305599118283937\\
154	0.304490752736916\\
155	0.301817861873773\\
156	0.301613212636632\\
157	0.298830274524587\\
158	0.296209518035964\\
159	0.295711681745862\\
160	0.294908794975138\\
161	0.293987690196305\\
162	0.29302567247614\\
163	0.287781366912243\\
164	0.283483767268958\\
165	0.276078772779\\
166	0.273024521850267\\
167	0.270860292779637\\
168	0.26963627986392\\
169	0.24878026495233\\
170	0.234663904031939\\
171	0.227625376320666\\
172	0.222048459659075\\
173	0.216426018521787\\
174	0.215956378299498\\
175	0.215649679586962\\
176	0.211832137828534\\
177	0.198379916166883\\
178	0.196041223667668\\
179	0.192831703145837\\
180	0.191922116616984\\
181	0.176283371394621\\
182	0.176123960411809\\
183	0.17240105608116\\
184	0.170792278019413\\
185	0.165513729176464\\
186	0.163255797306021\\
187	0.163019596586466\\
188	0.161906219225444\\
189	0.159869397011314\\
190	0.158983100137573\\
191	0.157926242709283\\
192	0.157351962738762\\
193	0.156762067413474\\
194	0.154242565931781\\
195	0.152413185728589\\
196	0.146447092608558\\
197	0.142796630104112\\
198	0.142091969748519\\
199	0.137818322954891\\
200	0.135667532177481\\
201	0.134455206141531\\
202	0.134118431059526\\
203	0.129495223557602\\
204	0.128341356051089\\
205	0.125389329104555\\
206	0.124446816815473\\
207	0.123515257270778\\
208	0.113668762799612\\
209	0.111215702833093\\
210	0.110275861221697\\
211	0.10859704427882\\
212	0.101191039165533\\
213	0.0979610317117518\\
214	0.0958271902795158\\
215	0.0934967273819717\\
216	0.0932187566098985\\
217	0.0930894329309905\\
218	0.0897343556271128\\
219	0.0885532491878702\\
220	0.0873821287136454\\
221	0.0862683575198378\\
222	0.0837680736032251\\
223	0.0827345169872991\\
224	0.0818125387952693\\
225	0.0812705140289651\\
226	0.0779686836325485\\
227	0.0750624213772202\\
228	0.0742069811382484\\
229	0.0715578663909939\\
230	0.0710022322983575\\
231	0.0701770196729075\\
232	0.0701255086579608\\
233	0.0693857139429908\\
234	0.0666911203793207\\
235	0.0627521722941033\\
236	0.0614352276051345\\
237	0.0613958389125091\\
238	0.0595915114672801\\
239	0.0562807018656572\\
240	0.052245665098149\\
241	0.0518373984863463\\
242	0.0439786944431006\\
243	0.0429853017281701\\
244	0.0402388125541034\\
245	0.039614055313959\\
246	0.0385837413106389\\
247	0.0373770466608731\\
248	0.0360533616938746\\
249	0.0325465516189133\\
250	0.0279032988140638\\
251	0.0189491124805734\\
252	0.0170987124454109\\
253	0.015374038065023\\
254	0.011135449353681\\
255	0.00764530793416698\\
256	0.000876771473112565\\
257	1.81205973380823e-07\\
};
\addplot [color=white!60!black, dashed, line width=1.1pt, mark size=2.0pt, mark=o, mark options={solid, white!60!black}, forget plot]
  table[row sep=crcr]{%
6	0.0882449204147932\\
18	0.000158546263894275\\
30	2.08026673424925e-05\\
42	5.99655533441666e-06\\
54	1.18271275381314e-06\\
78	2.51119372926873e-07\\
102	1.81122582715968e-07\\
138	1.81205705172834e-07\\
186	1.81205973621476e-07\\
246	1.8120597277919e-07\\
};
\addplot [color=mycolor1, line width=1.3pt, forget plot]
  table[row sep=crcr]{%
1	0.984979046555728\\
2	0.977762344004689\\
3	0.97422401229873\\
4	0.972884323185155\\
5	0.96292101771676\\
6	0.96221221502663\\
7	0.95507705901735\\
8	0.95204147921312\\
9	0.948398477070442\\
10	0.946491555842803\\
11	0.945238564051925\\
12	0.943691161494871\\
13	0.937169556810045\\
14	0.933869692225292\\
15	0.926183218351588\\
16	0.924651961028707\\
17	0.901834364874588\\
18	0.892870704900758\\
19	0.870339965364721\\
20	0.866226465999293\\
21	0.864120420075511\\
22	0.858142792076508\\
23	0.850569104298476\\
24	0.842224706742624\\
25	0.838619107958216\\
26	0.837853191711839\\
27	0.836238969994582\\
28	0.835667104456301\\
29	0.82372614104189\\
30	0.82369652070334\\
31	0.821573882421239\\
32	0.817409367389452\\
33	0.801808125178607\\
34	0.794683038345624\\
35	0.793181382670936\\
36	0.791818234478965\\
37	0.791243373231866\\
38	0.788054360831366\\
39	0.787633855577587\\
40	0.787314953936164\\
41	0.7813942862848\\
42	0.770451948979776\\
43	0.769145671081232\\
44	0.762547957028142\\
45	0.759848507579816\\
46	0.758896491276212\\
47	0.75373735398808\\
48	0.752926393580488\\
49	0.751474420510401\\
50	0.751224684931217\\
51	0.74768682190135\\
52	0.747432045065007\\
53	0.739215153579631\\
54	0.730741856624636\\
55	0.718251744097016\\
56	0.717945160103779\\
57	0.700475592877424\\
58	0.699400408800778\\
59	0.699148654996058\\
60	0.698141474989377\\
61	0.696235321372734\\
62	0.688724467762319\\
63	0.676295194163577\\
64	0.675226318290551\\
65	0.665131571982671\\
66	0.664644140202097\\
67	0.659722433671249\\
68	0.654818228942684\\
69	0.653153817850148\\
70	0.653007641941007\\
71	0.651302051787327\\
72	0.647171360157201\\
73	0.644428665952818\\
74	0.642533930379262\\
75	0.641944709879106\\
76	0.636358546322258\\
77	0.635620410323294\\
78	0.635589002573583\\
79	0.633230428423561\\
80	0.631566352721277\\
81	0.625105356934239\\
82	0.624819433926981\\
83	0.617193423369116\\
84	0.606761561087893\\
85	0.604942077515237\\
86	0.585910737120139\\
87	0.585354755763201\\
88	0.584558113129068\\
89	0.57690954610101\\
90	0.575962073984441\\
91	0.574876078457187\\
92	0.551679474045353\\
93	0.54277775514811\\
94	0.540085586282079\\
95	0.537518566703271\\
96	0.519938485078341\\
97	0.519820133341328\\
98	0.516958080447775\\
99	0.515059595902115\\
100	0.513472502000987\\
101	0.512356171007906\\
102	0.510461430657877\\
103	0.508644578583082\\
104	0.508304201583455\\
105	0.505101532669923\\
106	0.5043806645722\\
107	0.499942508568545\\
108	0.495333663902341\\
109	0.478302757850775\\
110	0.474921376942366\\
111	0.474528430990486\\
112	0.467586367565165\\
113	0.465793435489065\\
114	0.465467953046149\\
115	0.462568208964251\\
116	0.459797928202949\\
117	0.455964773777906\\
118	0.454869797625158\\
119	0.45130971303161\\
120	0.443041974238248\\
121	0.439311926695577\\
122	0.437954179702021\\
123	0.428430028009257\\
124	0.427524541704305\\
125	0.425293750091169\\
126	0.424101137941698\\
127	0.42227032662227\\
128	0.408469047332927\\
129	0.403448330541981\\
130	0.402371704887521\\
131	0.400358585849785\\
132	0.396735912596379\\
133	0.394188652681824\\
134	0.391613935666266\\
135	0.390719507989045\\
136	0.385689956717372\\
137	0.385120312857692\\
138	0.379219669216968\\
139	0.377155557595791\\
140	0.376167328700289\\
141	0.370478188334312\\
142	0.356225269152565\\
143	0.354572441223781\\
144	0.351623697580591\\
145	0.351273355101006\\
146	0.351088957635978\\
147	0.349536220169029\\
148	0.347665807773466\\
149	0.344560084985144\\
150	0.339891398283475\\
151	0.311607784935495\\
152	0.308997170972734\\
153	0.305273260268609\\
154	0.304262235409794\\
155	0.301624640061336\\
156	0.301381844461452\\
157	0.298981738642644\\
158	0.29621966436292\\
159	0.295712339725817\\
160	0.294858147616419\\
161	0.293945455038766\\
162	0.293051717217513\\
163	0.287940587214346\\
164	0.283691325724683\\
165	0.276229816787661\\
166	0.273334831207778\\
167	0.27099021824102\\
168	0.26967483017751\\
169	0.249434647747071\\
170	0.235232915505818\\
171	0.227899162115936\\
172	0.22248472836704\\
173	0.216975515205276\\
174	0.216485224849204\\
175	0.216134259034742\\
176	0.212515031032913\\
177	0.199393337861495\\
178	0.197057030744913\\
179	0.193719996199678\\
180	0.192734203007479\\
181	0.176873074202265\\
182	0.176712071163416\\
183	0.172985884993098\\
184	0.171470484867237\\
185	0.166145831300897\\
186	0.163990022236533\\
187	0.163778009236257\\
188	0.162592436394901\\
189	0.160474154423473\\
190	0.159594731680062\\
191	0.158544459561454\\
192	0.157923933091928\\
193	0.157302353713511\\
194	0.154880705850032\\
195	0.152914961917974\\
196	0.147103595341219\\
197	0.143313380398314\\
198	0.142676392494879\\
199	0.138499157541125\\
200	0.136270746320468\\
201	0.135152063583766\\
202	0.134805381075761\\
203	0.130241567905919\\
204	0.129113765969502\\
205	0.126217814402612\\
206	0.125187584642449\\
207	0.124248050002468\\
208	0.114534936621194\\
209	0.112309681228619\\
210	0.111365525733231\\
211	0.109637075264217\\
212	0.10199628904251\\
213	0.0987476835255234\\
214	0.0965924402735988\\
215	0.0944102614600353\\
216	0.0941069510326786\\
217	0.0939857699234567\\
218	0.0904304668020894\\
219	0.0893056510396403\\
220	0.0880288447750878\\
221	0.0869605466384646\\
222	0.0844626015664507\\
223	0.0833520920569109\\
224	0.0824528169168948\\
225	0.0819796059547138\\
226	0.0789008265040805\\
227	0.0759542219551092\\
228	0.0750366789352135\\
229	0.072192194602008\\
230	0.0716666527877808\\
231	0.0709148409376931\\
232	0.0708748927584018\\
233	0.0700820929893892\\
234	0.0673355996351619\\
235	0.0635833424187202\\
236	0.0622695817566305\\
237	0.0622174633754054\\
238	0.0604557905050335\\
239	0.0568972805276117\\
240	0.0529054405543381\\
241	0.052560270358218\\
242	0.0449722813065116\\
243	0.0439648012284682\\
244	0.0411645281301937\\
245	0.0406104148561503\\
246	0.0395931373367396\\
247	0.0383370285193731\\
248	0.036962773764747\\
249	0.0333541819944153\\
250	0.0286548000269635\\
251	0.0193314836680381\\
252	0.0174665855347717\\
253	0.0156442370669313\\
254	0.0113335626555241\\
255	0.00774594840363584\\
256	0.00098353506740542\\
257	1.33425588507537e-07\\
};
\addplot [color=mycolor1, dashed, line width=1.1pt, mark size=2.0pt, mark=o, mark options={solid, mycolor1}, forget plot]
  table[row sep=crcr]{%
6	0.534141369727649\\
18	0.000603297263871021\\
30	0.000158252570681296\\
42	2.57632504846352e-05\\
54	1.3491696228713e-06\\
78	1.26922344317802e-07\\
102	1.33403191147048e-07\\
138	1.33425587454495e-07\\
186	1.33425587454495e-07\\
246	1.33425587875712e-07\\
};
\addplot [color=red!60!teal, line width=1.3pt, forget plot]
  table[row sep=crcr]{%
1	0.985433437119212\\
2	0.978226156497124\\
3	0.974858303052713\\
4	0.973550546771614\\
5	0.964247406184525\\
6	0.963623432606279\\
7	0.956808094356886\\
8	0.953726234515284\\
9	0.950205084093044\\
10	0.948293377255706\\
11	0.947013372594169\\
12	0.945489499806989\\
13	0.938598123614141\\
14	0.935219561788285\\
15	0.927518026658643\\
16	0.926056410768445\\
17	0.903151132485688\\
18	0.894314395687705\\
19	0.871164705531516\\
20	0.867099550786276\\
21	0.865136717211747\\
22	0.859308292797039\\
23	0.851950399134391\\
24	0.843556525761834\\
25	0.839736938986317\\
26	0.838916327227756\\
27	0.837263020267531\\
28	0.836719412542575\\
29	0.825128983856816\\
30	0.825098591648919\\
31	0.822902187627803\\
32	0.818656387389979\\
33	0.803279876749618\\
34	0.796232376526153\\
35	0.794721826619264\\
36	0.793273621329155\\
37	0.792717892623931\\
38	0.789602869956297\\
39	0.789198755343686\\
40	0.78882751280897\\
41	0.782955436766487\\
42	0.772001252617009\\
43	0.770700950316691\\
44	0.764459455809137\\
45	0.761695619345811\\
46	0.76073080580687\\
47	0.755523726348044\\
48	0.754735933327556\\
49	0.75331906234388\\
50	0.753030866432908\\
51	0.749439059179947\\
52	0.749185950289736\\
53	0.741147387109837\\
54	0.732645470562572\\
55	0.72030362247721\\
56	0.720021706908331\\
57	0.702391721462863\\
58	0.701335934842103\\
59	0.701131743001314\\
60	0.700093890471457\\
61	0.698189507550848\\
62	0.690700367240771\\
63	0.678684390844238\\
64	0.677606667363141\\
65	0.66736162077697\\
66	0.666913865770498\\
67	0.661992867867015\\
68	0.657014928278315\\
69	0.655392820471167\\
70	0.655235180726789\\
71	0.653477271851709\\
72	0.649339932332809\\
73	0.646644367839009\\
74	0.644739551402431\\
75	0.6441779734909\\
76	0.638493352307471\\
77	0.637793507593545\\
78	0.637761723658338\\
79	0.63559798209538\\
80	0.63389248406268\\
81	0.627501826019807\\
82	0.627233463760944\\
83	0.619476303550847\\
84	0.609147286040542\\
85	0.607370153335081\\
86	0.588100278978302\\
87	0.587621956281122\\
88	0.586788539095416\\
89	0.579047294613075\\
90	0.578003526702037\\
91	0.576898142952121\\
92	0.553535517920229\\
93	0.544575541131496\\
94	0.541939496342758\\
95	0.539501412645691\\
96	0.521788435831489\\
97	0.521651868948416\\
98	0.518936252825635\\
99	0.517082194697159\\
100	0.515475569016346\\
101	0.514399056453652\\
102	0.512267157650943\\
103	0.510398899650529\\
104	0.51002026397758\\
105	0.506699003989169\\
106	0.505967133711821\\
107	0.501754696811466\\
108	0.497064799783851\\
109	0.479874078940073\\
110	0.476411983562801\\
111	0.476047926363015\\
112	0.469207975370984\\
113	0.467440193595513\\
114	0.467046881035241\\
115	0.464230090089265\\
116	0.461519757857268\\
117	0.457631501579576\\
118	0.456475829616928\\
119	0.452936695411869\\
120	0.444333557304978\\
121	0.440504118887919\\
122	0.439168219348771\\
123	0.429207646760505\\
124	0.428249039844844\\
125	0.42610705725604\\
126	0.424918929002331\\
127	0.423027762785162\\
128	0.409034872396845\\
129	0.404068787771389\\
130	0.402932959953179\\
131	0.400901932676183\\
132	0.397235282075766\\
133	0.394661769879551\\
134	0.391980003206608\\
135	0.391063523409152\\
136	0.385804242550817\\
137	0.385154315679774\\
138	0.379104587079604\\
139	0.377082838167835\\
140	0.376073571436321\\
141	0.370549011900347\\
142	0.356202119590363\\
143	0.354559247275338\\
144	0.351504259506116\\
145	0.351145583225759\\
146	0.350983927853728\\
147	0.349404611196558\\
148	0.347507407207439\\
149	0.344524741916775\\
150	0.339952470091902\\
151	0.311203188593182\\
152	0.308607707059191\\
153	0.304924699519752\\
154	0.303994093796389\\
155	0.301386599565027\\
156	0.301094025301825\\
157	0.298986256517665\\
158	0.29607842834256\\
159	0.295564851025265\\
160	0.294665512796261\\
161	0.293771072889431\\
162	0.292914439035968\\
163	0.287918909096498\\
164	0.283688506155204\\
165	0.27621304802726\\
166	0.273459028709306\\
167	0.270958678588185\\
168	0.269545712351546\\
169	0.249798842261022\\
170	0.235564833646911\\
171	0.227963230580533\\
172	0.222649540042387\\
173	0.217297247411937\\
174	0.216771390406564\\
175	0.216380358704446\\
176	0.212869231091148\\
177	0.200060267050675\\
178	0.197726999266613\\
179	0.194285003709634\\
180	0.193249786046618\\
181	0.177195854596434\\
182	0.177039310260157\\
183	0.173315955183564\\
184	0.171885681664132\\
185	0.166514285965975\\
186	0.164439634284395\\
187	0.164247406683519\\
188	0.163011281683419\\
189	0.160838696539045\\
190	0.159955783174774\\
191	0.158902714066542\\
192	0.158249135208618\\
193	0.157610099675219\\
194	0.155255248386711\\
195	0.153185885662299\\
196	0.147545624913399\\
197	0.143702460210388\\
198	0.143107044731602\\
199	0.139035476561537\\
200	0.136741788204889\\
201	0.135692940679979\\
202	0.135347911567647\\
203	0.130813656049871\\
204	0.129712698552638\\
205	0.126864838581168\\
206	0.125787810906675\\
207	0.124867500266527\\
208	0.115235152681392\\
209	0.113152857852916\\
210	0.112198966711064\\
211	0.110432483055322\\
212	0.102592050027853\\
213	0.0993206979993164\\
214	0.0971492692560892\\
215	0.0951114990600909\\
216	0.0947928153077289\\
217	0.0946842148808844\\
218	0.0909308577034882\\
219	0.0898379553398797\\
220	0.0884799328125141\\
221	0.0874242009969079\\
222	0.0849354947973515\\
223	0.083776234564642\\
224	0.0828960494705763\\
225	0.0824802128228838\\
226	0.0795775145620272\\
227	0.0765971141195632\\
228	0.07563403267502\\
229	0.0726457804379045\\
230	0.0721284567956434\\
231	0.0714205397119096\\
232	0.0713897745286732\\
233	0.0705771041384048\\
234	0.0677154252126164\\
235	0.0640714638959657\\
236	0.0627516735698862\\
237	0.0626862300935229\\
238	0.0609622469260856\\
239	0.0572368720306183\\
240	0.0532932426352169\\
241	0.0529931673690846\\
242	0.0457152683014643\\
243	0.0447195165207952\\
244	0.0418837260880643\\
245	0.0413776602428935\\
246	0.0403496399225189\\
247	0.0390411974980232\\
248	0.0376308716299813\\
249	0.0339716810339769\\
250	0.0292202093054891\\
251	0.0195664188710238\\
252	0.0177335883602452\\
253	0.0158138534254714\\
254	0.0114353126106502\\
255	0.00776224987044037\\
256	0.00106397003185377\\
257	1.27647786424573e-07\\
};
\addplot [color=red!60!teal, dashed, line width=1.1pt, mark size=2.0pt, mark=o, mark options={solid, red!60!teal}, forget plot]
  table[row sep=crcr]{%
6	0.363577994174019\\
18	0.00147537380809706\\
30	2.23616822497886e-05\\
42	3.58373598322532e-06\\
54	8.34578580063978e-08\\
78	1.27557398345942e-07\\
102	1.27647705941126e-07\\
138	1.27647786623298e-07\\
186	1.27647786225849e-07\\
246	1.27647786225849e-07\\
};
\node[right, align=left, inner sep=0]
at (rel axis cs:-0.15,1.1) {$(a)$};
\end{axis}
\end{tikzpicture}%
    \end{minipage}
    \hfill
    \begin{minipage}[t]{0.49\textwidth}
        \centering
        \vspace{0pt}
        % This file was created by matlab2tikz.
%
\definecolor{mycolor1}{rgb}{0.20000,0.40000,0.60000}%
%
\begin{tikzpicture}

\begin{axis}[%
width=11.252in,
height=7.131in,
at={(1.887in,0.962in)},
scale only axis,
clip=true,
separate axis lines,
every outer x axis line/.append style={black},
every x tick label/.append style={font=\color{black}},
every x tick/.append style={black},
xmin=0,
xmax=257,
xlabel={stepper calls},
every outer y axis line/.append style={black},
every y tick label/.append style={font=\color{black}},
every y tick/.append style={black},
ymode=log,
ymin=1e-09,
ymax=1,
yminorticks=true,
axis background/.style={fill=white},
legend style={legend cell align=left, align=left},
clip=true,
clip mode=individual,
axis lines=box,
xtick pos=bottom,
ytick pos=left,
tick align=inside,
major tick length=2pt,
width=0.7\linewidth,
height=0.65\linewidth,
every axis/.append style={font=\fontsize{9}{9}\selectfont},xlabel style={font=\fontsize{9}{9}\selectfont},title style={font=\fontsize{9}{9}\selectfont},
legend style={font=\fontsize{8}{9}\selectfont,inner sep=1pt,row sep=1pt,column sep=2pt,nodes={scale=1.00},draw=none},legend image post style={xscale=0.35},
legend columns=2,
scaled x ticks=false,
tick label style={/pgf/number format/fixed,/pgf/number format/precision=2},
every x tick label/.append style={font=\fontsize{9}{9}\selectfont\color{black}},
every y tick label/.append style={font=\fontsize{9}{9}\selectfont\color{black}},
,,
colormap={sigmamap}{rgb(0.0000)=(0.0200,0.1880,0.3800); rgb(0.1000)=(0.0743,0.3456,0.5616); rgb(0.2000)=(0.1918,0.4980,0.7060); rgb(0.3000)=(0.3890,0.6708,0.8226); rgb(0.4000)=(0.7552,0.8682,0.9228); rgb(0.5000)=(0.9690,0.9689,0.9689); rgb(0.6000)=(0.9425,0.8044,0.7614); rgb(0.7000)=(0.8767,0.4910,0.4020); rgb(0.8000)=(0.7869,0.2866,0.2470); rgb(0.9000)=(0.6186,0.1239,0.1568); rgb(1.0000)=(0.4040,0.0000,0.1220)}
]
\addplot [color=black, line width=1.3pt]
  table[row sep=crcr]{%
1	0.549669373328673\\
2	0.397017758326252\\
3	0.329385940378878\\
4	0.304042611548967\\
5	0.169444010329133\\
6	0.159452986665028\\
7	0.100626867648435\\
8	0.0849104695143835\\
9	0.0694873979297496\\
10	0.0632602086455798\\
11	0.0602107033749527\\
12	0.0567233098120105\\
13	0.0483275747489066\\
14	0.0438603616932953\\
15	0.0355617152947995\\
16	0.0338823885271802\\
17	0.0205533187125209\\
18	0.0158253009137569\\
19	0.00770254065852706\\
20	0.00636178928833088\\
21	0.00582162810062492\\
22	0.00462227982143794\\
23	0.0034434240593133\\
24	0.00252292974444555\\
25	0.00226168408256191\\
26	0.00221194947129454\\
27	0.0021400268633986\\
28	0.00209449770537478\\
29	0.00164871731744055\\
30	0.00164771495230215\\
31	0.00160538211639674\\
32	0.00154701309238332\\
33	0.00134719680001395\\
34	0.00127497590904898\\
35	0.00126434855840259\\
36	0.00125486412740125\\
37	0.00125208430335636\\
38	0.00123924856778656\\
39	0.00123809692125212\\
40	0.0012377413739315\\
41	0.00122641194599775\\
42	0.00121228121879711\\
43	0.00121099274233937\\
44	0.00120433582689565\\
45	0.00120270089422069\\
46	0.0012021855630201\\
47	0.00119987363566855\\
48	0.00119958893998503\\
49	0.0011992390365001\\
50	0.00119920279807839\\
51	0.00119874912374614\\
52	0.00119871011832831\\
53	0.00119792370841877\\
54	0.00119729006958437\\
55	0.00119658420419771\\
56	0.0011965706750862\\
57	0.001196036488324\\
58	0.00119600588916133\\
59	0.00119599619206095\\
60	0.00119598243518513\\
61	0.00119595528832575\\
62	0.00119588365243689\\
63	0.0011957865724953\\
64	0.00119578007904377\\
65	0.00119573751801717\\
66	0.00119573509137369\\
67	0.00119572086982291\\
68	0.00119570964810775\\
69	0.00119570652683007\\
70	0.00119570638820676\\
71	0.00119570478355878\\
72	0.00119570175073589\\
73	0.00119569979750223\\
74	0.00119569878899582\\
75	0.00119569855114903\\
76	0.00119569690909236\\
77	0.0011956967460316\\
78	0.00119569674115822\\
79	0.00119569636767616\\
80	0.00119569619312675\\
81	0.00119569560737483\\
82	0.0011956955776915\\
83	0.00119569519282965\\
84	0.00119569473647112\\
85	0.00119569467462868\\
86	0.00119569422571662\\
87	0.00119569421020061\\
88	0.00119569419775862\\
89	0.00119569411125754\\
90	0.00119569410667959\\
91	0.00119569410033646\\
92	0.0011956939765608\\
93	0.00119569394254596\\
94	0.0011956939324354\\
95	0.00119569392592682\\
96	0.00119569389279282\\
97	0.00119569389260429\\
98	0.00119569388893417\\
99	0.00119569388694803\\
100	0.00119569388570688\\
101	0.00119569388503786\\
102	0.00119569388456522\\
103	0.00119569388395673\\
104	0.00119569388386028\\
105	0.00119569388317799\\
106	0.00119569388308807\\
107	0.00119569388244497\\
108	0.00119569388203437\\
109	0.0011956938807168\\
110	0.00119569388053328\\
111	0.00119569388050541\\
112	0.00119569388026288\\
113	0.00119569388020301\\
114	0.00119569388019953\\
115	0.00119569388014489\\
116	0.00119569388010265\\
117	0.00119569388006716\\
118	0.00119569388005889\\
119	0.00119569388003516\\
120	0.00119569387999728\\
121	0.00119569387998335\\
122	0.00119569387997921\\
123	0.00119569387995766\\
124	0.00119569387995635\\
125	0.00119569387995265\\
126	0.00119569387995113\\
127	0.0011956938799496\\
128	0.00119569387994046\\
129	0.00119569387993763\\
130	0.00119569387993719\\
131	0.00119569387993632\\
132	0.00119569387993545\\
133	0.00119569387993502\\
134	0.0011956938799348\\
135	0.00119569387993458\\
136	0.00119569387993414\\
137	0.00119569387993414\\
138	0.00119569387993393\\
139	0.00119569387993393\\
140	0.00119569387993393\\
141	0.00119569387993371\\
142	0.00119569387993349\\
143	0.00119569387993349\\
144	0.00119569387993349\\
145	0.00119569387993349\\
146	0.00119569387993349\\
147	0.00119569387993349\\
148	0.00119569387993349\\
149	0.00119569387993349\\
150	0.00119569387993349\\
151	0.00119569387993349\\
152	0.00119569387993349\\
153	0.00119569387993349\\
154	0.00119569387993349\\
155	0.00119569387993349\\
156	0.00119569387993349\\
157	0.00119569387993349\\
158	0.00119569387993349\\
159	0.00119569387993349\\
160	0.00119569387993349\\
161	0.00119569387993349\\
162	0.00119569387993349\\
163	0.00119569387993349\\
164	0.00119569387993349\\
165	0.00119569387993349\\
166	0.00119569387993349\\
167	0.00119569387993349\\
168	0.00119569387993349\\
169	0.00119569387993349\\
170	0.00119569387993349\\
171	0.00119569387993349\\
172	0.00119569387993349\\
173	0.00119569387993349\\
174	0.00119569387993349\\
175	0.00119569387993349\\
176	0.00119569387993349\\
177	0.00119569387993349\\
178	0.00119569387993349\\
179	0.00119569387993349\\
180	0.00119569387993349\\
181	0.00119569387993349\\
182	0.00119569387993349\\
183	0.00119569387993349\\
184	0.00119569387993349\\
185	0.00119569387993349\\
186	0.00119569387993349\\
187	0.00119569387993349\\
188	0.00119569387993349\\
189	0.00119569387993349\\
190	0.00119569387993349\\
191	0.00119569387993349\\
192	0.00119569387993349\\
193	0.00119569387993349\\
194	0.00119569387993349\\
195	0.00119569387993349\\
196	0.00119569387993349\\
197	0.00119569387993349\\
198	0.00119569387993349\\
199	0.00119569387993349\\
200	0.00119569387993349\\
201	0.00119569387993349\\
202	0.00119569387993349\\
203	0.00119569387993349\\
204	0.00119569387993349\\
205	0.00119569387993349\\
206	0.00119569387993349\\
207	0.00119569387993349\\
208	0.00119569387993349\\
209	0.00119569387993349\\
210	0.00119569387993349\\
211	0.00119569387993349\\
212	0.00119569387993349\\
213	0.00119569387993349\\
214	0.00119569387993349\\
215	0.00119569387993349\\
216	0.00119569387993349\\
217	0.00119569387993349\\
218	0.00119569387993349\\
219	0.00119569387993349\\
220	0.00119569387993349\\
221	0.00119569387993349\\
222	0.00119569387993349\\
223	0.00119569387993349\\
224	0.00119569387993349\\
225	0.00119569387993349\\
226	0.00119569387993349\\
227	0.00119569387993349\\
228	0.00119569387993349\\
229	0.00119569387993349\\
230	0.00119569387993349\\
231	0.00119569387993349\\
232	0.00119569387993349\\
233	0.00119569387993349\\
234	0.00119569387993349\\
235	0.00119569387993349\\
236	0.00119569387993349\\
237	0.00119569387993349\\
238	0.00119569387993349\\
239	0.00119569387993349\\
240	0.00119569387993349\\
241	0.00119569387993349\\
242	0.00119569387993349\\
243	0.00119569387993349\\
244	0.00119569387993349\\
245	0.00119569387993349\\
246	0.00119569387993349\\
247	0.00119569387993349\\
248	0.00119569387993349\\
249	0.00119569387993349\\
250	0.00119569387993349\\
251	0.00119569387993349\\
252	0.00119569387993349\\
253	0.00119569387993349\\
254	0.00119569387993349\\
255	0.00119569387993349\\
256	0.00119569387993349\\
257	0.00119569387993349\\
};
\addlegendentry{$T=2$}

\addplot [color=black, dashed, line width=1.1pt, mark size=2.0pt, mark=o, mark options={solid, black}, forget plot]
  table[row sep=crcr]{%
6	0.500186391700627\\
18	0.000953091266489265\\
30	0.000937895531436231\\
42	0.00114613720263973\\
54	0.00120561859610211\\
78	0.00119455605801567\\
102	0.00119526527591131\\
138	0.00119569325590729\\
186	0.00119569419389845\\
246	0.00119569387985054\\
};
\addplot [color=white!35!black, line width=1.3pt]
  table[row sep=crcr]{%
1	0.559033137196902\\
2	0.399370562646894\\
3	0.329302895120021\\
4	0.306698715275885\\
5	0.172052357335026\\
6	0.16319890522921\\
7	0.104969779021928\\
8	0.0883103836569877\\
9	0.0714386473743471\\
10	0.064945741622158\\
11	0.0616826625826992\\
12	0.0580766532695019\\
13	0.0486545368239233\\
14	0.0442166251248846\\
15	0.0358420473896074\\
16	0.0343089376356351\\
17	0.0199915060764878\\
18	0.0151886799857187\\
19	0.006652133182784\\
20	0.00531071631804672\\
21	0.00475147435290793\\
22	0.00352687356420163\\
23	0.0023243251692648\\
24	0.00135701802760233\\
25	0.0010702867649284\\
26	0.00102088761475115\\
27	0.000938600650181147\\
28	0.000901307452132813\\
29	0.000452940546045711\\
30	0.000451819721460005\\
31	0.000409562223026839\\
32	0.000345504828198503\\
33	0.000140740903948513\\
34	6.65259626557667e-05\\
35	5.52609723132258e-05\\
36	4.65584268039886e-05\\
37	4.36867067747377e-05\\
38	3.0892628351158e-05\\
39	2.96437461122934e-05\\
40	2.91802439066835e-05\\
41	1.80562903197952e-05\\
42	3.26595915914533e-06\\
43	1.90065585801823e-06\\
44	4.61522082985248e-06\\
45	6.3331791828639e-06\\
46	6.83399937297337e-06\\
47	9.07655446816112e-06\\
48	9.36127163991361e-06\\
49	9.73366034250314e-06\\
50	9.77047575220323e-06\\
51	1.02653611725188e-05\\
52	1.02994007819449e-05\\
53	1.10935563276389e-05\\
54	1.1718338308664e-05\\
55	1.24309056312515e-05\\
56	1.24456567959361e-05\\
57	1.30120941042536e-05\\
58	1.30420100392807e-05\\
59	1.30511795746009e-05\\
60	1.30657238665149e-05\\
61	1.30916871846514e-05\\
62	1.31658799261434e-05\\
63	1.32659928977106e-05\\
64	1.32723957246663e-05\\
65	1.33173964641523e-05\\
66	1.33196958058559e-05\\
67	1.33339537340158e-05\\
68	1.33450894980249e-05\\
69	1.33480778203941e-05\\
70	1.3348237115529e-05\\
71	1.3349925219261e-05\\
72	1.33531694088253e-05\\
73	1.33550698724182e-05\\
74	1.33560507515186e-05\\
75	1.33562986059396e-05\\
76	1.33579735548693e-05\\
77	1.33581497486299e-05\\
78	1.33581547843016e-05\\
79	1.33585397589709e-05\\
80	1.33587240048085e-05\\
81	1.33593119382798e-05\\
82	1.33593377601764e-05\\
83	1.33597261906152e-05\\
84	1.33601805009195e-05\\
85	1.33602425083795e-05\\
86	1.33607035448835e-05\\
87	1.33607181576611e-05\\
88	1.33607300272703e-05\\
89	1.33608187441448e-05\\
90	1.33608246069358e-05\\
91	1.3360831604489e-05\\
92	1.33609591350746e-05\\
93	1.33609961007948e-05\\
94	1.3361005875671e-05\\
95	1.3361012935255e-05\\
96	1.33610477606609e-05\\
97	1.33610479318107e-05\\
98	1.33610517662945e-05\\
99	1.33610537608238e-05\\
100	1.3361055033209e-05\\
101	1.33610557188127e-05\\
102	1.33610563152629e-05\\
103	1.3361056942854e-05\\
104	1.33610570348957e-05\\
105	1.33610577139699e-05\\
106	1.33610578183173e-05\\
107	1.33610584767982e-05\\
108	1.3361058919427e-05\\
109	1.33610602661487e-05\\
110	1.33610604617843e-05\\
111	1.33610604871492e-05\\
112	1.33610607502151e-05\\
113	1.33610608072231e-05\\
114	1.33610608117436e-05\\
115	1.33610608678727e-05\\
116	1.33610609095614e-05\\
117	1.33610609487388e-05\\
118	1.33610609572774e-05\\
119	1.33610609820144e-05\\
120	1.33610610216941e-05\\
121	1.33610610360089e-05\\
122	1.33610610402782e-05\\
123	1.33610610623782e-05\\
124	1.33610610637595e-05\\
125	1.33610610675266e-05\\
126	1.33610610689078e-05\\
127	1.33610610704146e-05\\
128	1.33610610798323e-05\\
129	1.33610610827203e-05\\
130	1.3361061083097e-05\\
131	1.33610610838504e-05\\
132	1.3361061084855e-05\\
133	1.33610610853573e-05\\
134	1.3361061085734e-05\\
135	1.33610610858595e-05\\
136	1.33610610862362e-05\\
137	1.33610610862362e-05\\
138	1.3361061086613e-05\\
139	1.33610610867385e-05\\
140	1.33610610867385e-05\\
141	1.33610610868641e-05\\
142	1.33610610871152e-05\\
143	1.33610610871152e-05\\
144	1.33610610871152e-05\\
145	1.33610610871152e-05\\
146	1.33610610871152e-05\\
147	1.33610610871152e-05\\
148	1.33610610871152e-05\\
149	1.33610610871152e-05\\
150	1.33610610871152e-05\\
151	1.33610610871152e-05\\
152	1.33610610871152e-05\\
153	1.33610610871152e-05\\
154	1.33610610871152e-05\\
155	1.33610610871152e-05\\
156	1.33610610871152e-05\\
157	1.33610610871152e-05\\
158	1.33610610871152e-05\\
159	1.33610610871152e-05\\
160	1.33610610871152e-05\\
161	1.33610610871152e-05\\
162	1.33610610871152e-05\\
163	1.33610610871152e-05\\
164	1.33610610871152e-05\\
165	1.33610610871152e-05\\
166	1.33610610871152e-05\\
167	1.33610610871152e-05\\
168	1.33610610871152e-05\\
169	1.33610610871152e-05\\
170	1.33610610871152e-05\\
171	1.33610610871152e-05\\
172	1.33610610871152e-05\\
173	1.33610610871152e-05\\
174	1.33610610871152e-05\\
175	1.33610610871152e-05\\
176	1.33610610871152e-05\\
177	1.33610610871152e-05\\
178	1.33610610871152e-05\\
179	1.33610610871152e-05\\
180	1.33610610871152e-05\\
181	1.33610610871152e-05\\
182	1.33610610871152e-05\\
183	1.33610610871152e-05\\
184	1.33610610871152e-05\\
185	1.33610610871152e-05\\
186	1.33610610871152e-05\\
187	1.33610610871152e-05\\
188	1.33610610871152e-05\\
189	1.33610610871152e-05\\
190	1.33610610871152e-05\\
191	1.33610610871152e-05\\
192	1.33610610871152e-05\\
193	1.33610610871152e-05\\
194	1.33610610871152e-05\\
195	1.33610610871152e-05\\
196	1.33610610871152e-05\\
197	1.33610610871152e-05\\
198	1.33610610871152e-05\\
199	1.33610610871152e-05\\
200	1.33610610871152e-05\\
201	1.33610610871152e-05\\
202	1.33610610871152e-05\\
203	1.33610610871152e-05\\
204	1.33610610871152e-05\\
205	1.33610610871152e-05\\
206	1.33610610871152e-05\\
207	1.33610610871152e-05\\
208	1.33610610871152e-05\\
209	1.33610610871152e-05\\
210	1.33610610871152e-05\\
211	1.33610610871152e-05\\
212	1.33610610871152e-05\\
213	1.33610610871152e-05\\
214	1.33610610871152e-05\\
215	1.33610610871152e-05\\
216	1.33610610871152e-05\\
217	1.33610610871152e-05\\
218	1.33610610871152e-05\\
219	1.33610610871152e-05\\
220	1.33610610871152e-05\\
221	1.33610610871152e-05\\
222	1.33610610871152e-05\\
223	1.33610610871152e-05\\
224	1.33610610871152e-05\\
225	1.33610610871152e-05\\
226	1.33610610871152e-05\\
227	1.33610610871152e-05\\
228	1.33610610871152e-05\\
229	1.33610610871152e-05\\
230	1.33610610871152e-05\\
231	1.33610610871152e-05\\
232	1.33610610871152e-05\\
233	1.33610610871152e-05\\
234	1.33610610871152e-05\\
235	1.33610610871152e-05\\
236	1.33610610871152e-05\\
237	1.33610610871152e-05\\
238	1.33610610871152e-05\\
239	1.33610610871152e-05\\
240	1.33610610871152e-05\\
241	1.33610610871152e-05\\
242	1.33610610871152e-05\\
243	1.33610610871152e-05\\
244	1.33610610871152e-05\\
245	1.33610610871152e-05\\
246	1.33610610871152e-05\\
247	1.33610610871152e-05\\
248	1.33610610871152e-05\\
249	1.33610610871152e-05\\
250	1.33610610871152e-05\\
251	1.33610610871152e-05\\
252	1.33610610871152e-05\\
253	1.33610610871152e-05\\
254	1.33610610871152e-05\\
255	1.33610610871152e-05\\
256	1.33610610871152e-05\\
257	1.33610610871152e-05\\
};
\addlegendentry{$T=5$}

\addplot [color=white!35!black, dashed, line width=1.1pt, mark size=2.0pt, mark=o, mark options={solid, white!35!black}, forget plot]
  table[row sep=crcr]{%
6	0.173444432589479\\
18	0.000322957642254964\\
30	0.000101687665819906\\
42	3.26495648006814e-05\\
54	1.62064647237235e-05\\
78	1.34198389411002e-05\\
102	1.33590527996498e-05\\
138	1.33610355388322e-05\\
186	1.33610611655955e-05\\
246	1.33610610878686e-05\\
};
\addplot [color=white!60!black, line width=1.3pt]
  table[row sep=crcr]{%
1	0.564607223850406\\
2	0.396758973927432\\
3	0.326445192001137\\
4	0.305342094296728\\
5	0.175675676464535\\
6	0.167903493405876\\
7	0.110982018363616\\
8	0.0932041795791994\\
9	0.0756335859496363\\
10	0.0686612058289143\\
11	0.0651210578167606\\
12	0.0614886904025271\\
13	0.0506130701614662\\
14	0.0460435679071081\\
15	0.037521905756448\\
16	0.0361065407031926\\
17	0.0207935070318107\\
18	0.0159310964064198\\
19	0.00680931860605786\\
20	0.00545037290468793\\
21	0.00488304372519022\\
22	0.00365450952483655\\
23	0.00243521282985079\\
24	0.00142129122639922\\
25	0.00109879304218316\\
26	0.00104552799072529\\
27	0.00095420262223343\\
28	0.000924545529589524\\
29	0.00048125790828469\\
30	0.000480294696618071\\
31	0.000435649338525187\\
32	0.000366380621243458\\
33	0.000156105502138586\\
34	8.08270566108681e-05\\
35	6.87871673626111e-05\\
36	6.03888192089379e-05\\
37	5.74653367814186e-05\\
38	4.47684286374736e-05\\
39	4.34636625254548e-05\\
40	4.28406773344513e-05\\
41	3.17612705488361e-05\\
42	1.61629785633702e-05\\
43	1.47061215135406e-05\\
44	8.51261199210358e-06\\
45	6.705935367087e-06\\
46	6.20142788109269e-06\\
47	4.02058535815122e-06\\
48	3.74431287273559e-06\\
49	3.36537183448011e-06\\
50	3.32277108203417e-06\\
51	2.7887766861761e-06\\
52	2.75716296231618e-06\\
53	1.96254667866241e-06\\
54	1.33676872918941e-06\\
55	6.11941895989095e-07\\
56	5.9727412945619e-07\\
57	3.75051588558086e-09\\
58	3.36021656519237e-08\\
59	4.06455718167724e-08\\
60	5.6610641453684e-08\\
61	8.17774088124832e-08\\
62	1.57907492776768e-07\\
63	2.58974101885202e-07\\
64	2.65489923131242e-07\\
65	3.12904836386055e-07\\
66	3.14957046412481e-07\\
67	3.29406420299031e-07\\
68	3.4055024922441e-07\\
69	3.4355238615584e-07\\
70	3.43735687345281e-07\\
71	3.45522129333168e-07\\
72	3.48967645564236e-07\\
73	3.50819817227821e-07\\
74	3.51798462951844e-07\\
75	3.52044614345287e-07\\
76	3.53761955151238e-07\\
77	3.53945461897642e-07\\
78	3.53951140478496e-07\\
79	3.54327414265969e-07\\
80	3.54516676793354e-07\\
81	3.55105940892562e-07\\
82	3.5512811593124e-07\\
83	3.55534876327613e-07\\
84	3.559866825326e-07\\
85	3.56048693766554e-07\\
86	3.56531969905423e-07\\
87	3.56544973569101e-07\\
88	3.56557102546048e-07\\
89	3.56649536396507e-07\\
90	3.56657281894477e-07\\
91	3.56665061992985e-07\\
92	3.5679827257809e-07\\
93	3.56837899495915e-07\\
94	3.56847596203823e-07\\
95	3.56854907334817e-07\\
96	3.56891937910082e-07\\
97	3.56892115005312e-07\\
98	3.56896002543236e-07\\
99	3.56897989655884e-07\\
100	3.56899258342296e-07\\
101	3.56899965235019e-07\\
102	3.56900740398696e-07\\
103	3.56901402273514e-07\\
104	3.5690149509969e-07\\
105	3.5690220106229e-07\\
106	3.56902322536224e-07\\
107	3.56902971761382e-07\\
108	3.56903453936628e-07\\
109	3.56904865303756e-07\\
110	3.56905078487719e-07\\
111	3.56905100810647e-07\\
112	3.56905379289175e-07\\
113	3.56905435840593e-07\\
114	3.56905442165422e-07\\
115	3.56905498530816e-07\\
116	3.56905540200282e-07\\
117	3.56905582985894e-07\\
118	3.56905592287114e-07\\
119	3.56905617400408e-07\\
120	3.5690566018602e-07\\
121	3.56905675440021e-07\\
122	3.56905679904606e-07\\
123	3.56905703343681e-07\\
124	3.569057050179e-07\\
125	3.56905708738389e-07\\
126	3.56905710226584e-07\\
127	3.56905711900803e-07\\
128	3.56905721946121e-07\\
129	3.56905724922511e-07\\
130	3.5690572529456e-07\\
131	3.56905726038658e-07\\
132	3.5690572696878e-07\\
133	3.56905727526853e-07\\
134	3.56905727898902e-07\\
135	3.56905728084926e-07\\
136	3.56905728642999e-07\\
137	3.56905728642999e-07\\
138	3.56905729015048e-07\\
139	3.56905729201073e-07\\
140	3.56905729201073e-07\\
141	3.56905729387097e-07\\
142	3.56905729759146e-07\\
143	3.56905729759146e-07\\
144	3.56905729759146e-07\\
145	3.56905729759146e-07\\
146	3.56905729759146e-07\\
147	3.56905729759146e-07\\
148	3.56905729759146e-07\\
149	3.56905729759146e-07\\
150	3.56905729759146e-07\\
151	3.56905729759146e-07\\
152	3.56905729759146e-07\\
153	3.56905729759146e-07\\
154	3.56905729759146e-07\\
155	3.56905729759146e-07\\
156	3.56905729759146e-07\\
157	3.56905729759146e-07\\
158	3.56905729759146e-07\\
159	3.56905729759146e-07\\
160	3.56905729759146e-07\\
161	3.56905729759146e-07\\
162	3.56905729759146e-07\\
163	3.56905729759146e-07\\
164	3.56905729759146e-07\\
165	3.56905729759146e-07\\
166	3.56905729759146e-07\\
167	3.56905729759146e-07\\
168	3.56905729759146e-07\\
169	3.56905729759146e-07\\
170	3.56905729759146e-07\\
171	3.56905729759146e-07\\
172	3.56905729759146e-07\\
173	3.56905729759146e-07\\
174	3.56905729759146e-07\\
175	3.56905729759146e-07\\
176	3.56905729759146e-07\\
177	3.56905729759146e-07\\
178	3.56905729759146e-07\\
179	3.56905729759146e-07\\
180	3.56905729759146e-07\\
181	3.56905729759146e-07\\
182	3.56905729759146e-07\\
183	3.56905729759146e-07\\
184	3.56905729759146e-07\\
185	3.56905729759146e-07\\
186	3.56905729759146e-07\\
187	3.56905729759146e-07\\
188	3.56905729759146e-07\\
189	3.56905729759146e-07\\
190	3.56905729759146e-07\\
191	3.56905729759146e-07\\
192	3.56905729759146e-07\\
193	3.56905729759146e-07\\
194	3.56905729759146e-07\\
195	3.56905729759146e-07\\
196	3.56905729759146e-07\\
197	3.56905729759146e-07\\
198	3.56905729759146e-07\\
199	3.56905729759146e-07\\
200	3.56905729759146e-07\\
201	3.56905729759146e-07\\
202	3.56905729759146e-07\\
203	3.56905729759146e-07\\
204	3.56905729759146e-07\\
205	3.56905729759146e-07\\
206	3.56905729759146e-07\\
207	3.56905729759146e-07\\
208	3.56905729759146e-07\\
209	3.56905729759146e-07\\
210	3.56905729759146e-07\\
211	3.56905729759146e-07\\
212	3.56905729759146e-07\\
213	3.56905729759146e-07\\
214	3.56905729759146e-07\\
215	3.56905729759146e-07\\
216	3.56905729759146e-07\\
217	3.56905729759146e-07\\
218	3.56905729759146e-07\\
219	3.56905729759146e-07\\
220	3.56905729759146e-07\\
221	3.56905729759146e-07\\
222	3.56905729759146e-07\\
223	3.56905729759146e-07\\
224	3.56905729759146e-07\\
225	3.56905729759146e-07\\
226	3.56905729759146e-07\\
227	3.56905729759146e-07\\
228	3.56905729759146e-07\\
229	3.56905729759146e-07\\
230	3.56905729759146e-07\\
231	3.56905729759146e-07\\
232	3.56905729759146e-07\\
233	3.56905729759146e-07\\
234	3.56905729759146e-07\\
235	3.56905729759146e-07\\
236	3.56905729759146e-07\\
237	3.56905729759146e-07\\
238	3.56905729759146e-07\\
239	3.56905729759146e-07\\
240	3.56905729759146e-07\\
241	3.56905729759146e-07\\
242	3.56905729759146e-07\\
243	3.56905729759146e-07\\
244	3.56905729759146e-07\\
245	3.56905729759146e-07\\
246	3.56905729759146e-07\\
247	3.56905729759146e-07\\
248	3.56905729759146e-07\\
249	3.56905729759146e-07\\
250	3.56905729759146e-07\\
251	3.56905729759146e-07\\
252	3.56905729759146e-07\\
253	3.56905729759146e-07\\
254	3.56905729759146e-07\\
255	3.56905729759146e-07\\
256	3.56905729759146e-07\\
257	3.56905729759146e-07\\
};
\addlegendentry{$T=10$}

\addplot [color=white!60!black, dashed, line width=1.1pt, mark size=2.0pt, mark=o, mark options={solid, white!60!black}, forget plot]
  table[row sep=crcr]{%
6	0.195015805283674\\
18	0.000603102194024149\\
30	9.17721279357794e-05\\
42	8.00760381008745e-06\\
54	9.51141711308382e-08\\
78	3.62838592575086e-07\\
102	3.5667661094585e-07\\
138	3.56905727898902e-07\\
186	3.56905730503243e-07\\
246	3.56905728642999e-07\\
};
\addplot [color=mycolor1, line width=1.3pt]
  table[row sep=crcr]{%
1	0.566656855065321\\
2	0.392875725684473\\
3	0.323772573360371\\
4	0.302801358166162\\
5	0.178595981845843\\
6	0.171587776627804\\
7	0.115791401523224\\
8	0.0970557094110618\\
9	0.0793371199946198\\
10	0.0720374269239434\\
11	0.068266173476437\\
12	0.0646073113866519\\
13	0.0525013358978267\\
14	0.0476953080114036\\
15	0.0389164095610798\\
16	0.0375455584116823\\
17	0.0215399832699583\\
18	0.0166151049257236\\
19	0.00692206348940702\\
20	0.00553674357961893\\
21	0.00498166842603597\\
22	0.00374896230271129\\
23	0.00252717681739777\\
24	0.0014743491939089\\
25	0.00111860408408585\\
26	0.00105951972411547\\
27	0.000962173331887127\\
28	0.000935217272346847\\
29	0.000495319031786381\\
30	0.000494466321345801\\
31	0.000446720213088636\\
32	0.000373533603640214\\
33	0.000159348556174509\\
34	8.29398214375151e-05\\
35	7.03619109225198e-05\\
36	6.14446590014436e-05\\
37	5.85079068285983e-05\\
38	4.57862100247197e-05\\
39	4.447638060846e-05\\
40	4.37008017074584e-05\\
41	3.24589285546669e-05\\
42	1.6238922261738e-05\\
43	1.47273643345182e-05\\
44	8.76789612878521e-06\\
45	6.86466479015716e-06\\
46	6.34077289656015e-06\\
47	4.12495145987292e-06\\
48	3.85312104754706e-06\\
49	3.47330018353092e-06\\
50	3.4223196416228e-06\\
51	2.85874500519026e-06\\
52	2.82707558876772e-06\\
53	2.03010681674796e-06\\
54	1.3888612900811e-06\\
55	6.51362104441099e-07\\
56	6.3723823857289e-07\\
57	9.35149208097676e-09\\
58	2.07968674990126e-08\\
59	2.6304002017418e-08\\
60	4.34915151341377e-08\\
61	6.88666338902863e-08\\
62	1.46862583115573e-07\\
63	2.47544103202229e-07\\
64	2.54297800845049e-07\\
65	3.04049677345911e-07\\
66	3.05923451899042e-07\\
67	3.20680472215724e-07\\
68	3.32149361028681e-07\\
69	3.35185165238863e-07\\
70	3.35393106301297e-07\\
71	3.37285369297977e-07\\
72	3.40859442088952e-07\\
73	3.42710173612424e-07\\
74	3.43707257990091e-07\\
75	3.43949065854969e-07\\
76	3.45736820472359e-07\\
77	3.45921035283644e-07\\
78	3.45927147681899e-07\\
79	3.4628508272022e-07\\
80	3.46482007633499e-07\\
81	3.47078208696351e-07\\
82	3.47098781888717e-07\\
83	3.4752664288393e-07\\
84	3.47983009323369e-07\\
85	3.48045073364963e-07\\
86	3.48551242790067e-07\\
87	3.4856277244648e-07\\
88	3.48575653269912e-07\\
89	3.48672074959069e-07\\
90	3.48681387586927e-07\\
91	3.48689709795425e-07\\
92	3.48828301862211e-07\\
93	3.48869767027683e-07\\
94	3.48879543977565e-07\\
95	3.48886811976983e-07\\
96	3.48925617111071e-07\\
97	3.48925820734577e-07\\
98	3.48929660372203e-07\\
99	3.48931645911217e-07\\
100	3.48932940003223e-07\\
101	3.4893364957994e-07\\
102	3.48934588463919e-07\\
103	3.48935290318722e-07\\
104	3.4893539288588e-07\\
105	3.48936144765254e-07\\
106	3.48936276709257e-07\\
107	3.48936909738308e-07\\
108	3.48937422238359e-07\\
109	3.48938898466807e-07\\
110	3.48939126934731e-07\\
111	3.48939147582456e-07\\
112	3.48939432621791e-07\\
113	3.4893949003254e-07\\
114	3.48939498090189e-07\\
115	3.48939554493732e-07\\
116	3.48939596460654e-07\\
117	3.4893964178493e-07\\
118	3.48939651856991e-07\\
119	3.4893967737288e-07\\
120	3.48939723704361e-07\\
121	3.48939739987527e-07\\
122	3.48939744519955e-07\\
123	3.48939769700108e-07\\
124	3.48939771546653e-07\\
125	3.48939775071874e-07\\
126	3.48939776582683e-07\\
127	3.48939778429228e-07\\
128	3.48939789004892e-07\\
129	3.4893979202651e-07\\
130	3.48939792530114e-07\\
131	3.48939793201584e-07\\
132	3.4893979420879e-07\\
133	3.48939794712394e-07\\
134	3.48939795215997e-07\\
135	3.48939795383864e-07\\
136	3.48939795887467e-07\\
137	3.48939795887467e-07\\
138	3.48939796223203e-07\\
139	3.4893979639107e-07\\
140	3.4893979639107e-07\\
141	3.48939796558938e-07\\
142	3.48939796894673e-07\\
143	3.48939796894673e-07\\
144	3.48939796894673e-07\\
145	3.48939796894673e-07\\
146	3.48939796894673e-07\\
147	3.48939796894673e-07\\
148	3.48939796894673e-07\\
149	3.48939796894673e-07\\
150	3.48939796894673e-07\\
151	3.48939796894673e-07\\
152	3.48939796894673e-07\\
153	3.48939796894673e-07\\
154	3.48939796894673e-07\\
155	3.48939796894673e-07\\
156	3.48939796894673e-07\\
157	3.48939796894673e-07\\
158	3.48939796894673e-07\\
159	3.48939796894673e-07\\
160	3.48939796894673e-07\\
161	3.48939796894673e-07\\
162	3.48939796894673e-07\\
163	3.48939796894673e-07\\
164	3.48939796894673e-07\\
165	3.48939796894673e-07\\
166	3.48939796894673e-07\\
167	3.48939796894673e-07\\
168	3.48939796894673e-07\\
169	3.48939796894673e-07\\
170	3.48939796894673e-07\\
171	3.48939796894673e-07\\
172	3.48939796894673e-07\\
173	3.48939796894673e-07\\
174	3.48939796894673e-07\\
175	3.48939796894673e-07\\
176	3.48939796894673e-07\\
177	3.48939796894673e-07\\
178	3.48939796894673e-07\\
179	3.48939796894673e-07\\
180	3.48939796894673e-07\\
181	3.48939796894673e-07\\
182	3.48939796894673e-07\\
183	3.48939796894673e-07\\
184	3.48939796894673e-07\\
185	3.48939796894673e-07\\
186	3.48939796894673e-07\\
187	3.48939796894673e-07\\
188	3.48939796894673e-07\\
189	3.48939796894673e-07\\
190	3.48939796894673e-07\\
191	3.48939796894673e-07\\
192	3.48939796894673e-07\\
193	3.48939796894673e-07\\
194	3.48939796894673e-07\\
195	3.48939796894673e-07\\
196	3.48939796894673e-07\\
197	3.48939796894673e-07\\
198	3.48939796894673e-07\\
199	3.48939796894673e-07\\
200	3.48939796894673e-07\\
201	3.48939796894673e-07\\
202	3.48939796894673e-07\\
203	3.48939796894673e-07\\
204	3.48939796894673e-07\\
205	3.48939796894673e-07\\
206	3.48939796894673e-07\\
207	3.48939796894673e-07\\
208	3.48939796894673e-07\\
209	3.48939796894673e-07\\
210	3.48939796894673e-07\\
211	3.48939796894673e-07\\
212	3.48939796894673e-07\\
213	3.48939796894673e-07\\
214	3.48939796894673e-07\\
215	3.48939796894673e-07\\
216	3.48939796894673e-07\\
217	3.48939796894673e-07\\
218	3.48939796894673e-07\\
219	3.48939796894673e-07\\
220	3.48939796894673e-07\\
221	3.48939796894673e-07\\
222	3.48939796894673e-07\\
223	3.48939796894673e-07\\
224	3.48939796894673e-07\\
225	3.48939796894673e-07\\
226	3.48939796894673e-07\\
227	3.48939796894673e-07\\
228	3.48939796894673e-07\\
229	3.48939796894673e-07\\
230	3.48939796894673e-07\\
231	3.48939796894673e-07\\
232	3.48939796894673e-07\\
233	3.48939796894673e-07\\
234	3.48939796894673e-07\\
235	3.48939796894673e-07\\
236	3.48939796894673e-07\\
237	3.48939796894673e-07\\
238	3.48939796894673e-07\\
239	3.48939796894673e-07\\
240	3.48939796894673e-07\\
241	3.48939796894673e-07\\
242	3.48939796894673e-07\\
243	3.48939796894673e-07\\
244	3.48939796894673e-07\\
245	3.48939796894673e-07\\
246	3.48939796894673e-07\\
247	3.48939796894673e-07\\
248	3.48939796894673e-07\\
249	3.48939796894673e-07\\
250	3.48939796894673e-07\\
251	3.48939796894673e-07\\
252	3.48939796894673e-07\\
253	3.48939796894673e-07\\
254	3.48939796894673e-07\\
255	3.48939796894673e-07\\
256	3.48939796894673e-07\\
257	3.48939796894673e-07\\
};
\addlegendentry{$T=15$}

\addplot [color=mycolor1, dashed, line width=1.1pt, mark size=2.0pt, mark=o, mark options={solid, mycolor1}, forget plot]
  table[row sep=crcr]{%
6	0.816868180679755\\
18	0.000246780445434153\\
30	7.22780370615801e-05\\
42	2.28708801699798e-07\\
54	7.29981302074725e-07\\
78	3.48475742115126e-07\\
102	3.48938569614007e-07\\
138	3.48939797734012e-07\\
186	3.48939797566144e-07\\
246	3.48939797734012e-07\\
};
\addplot [color=red!60!teal, line width=1.3pt]
  table[row sep=crcr]{%
1	0.567622633937709\\
2	0.389053383852605\\
3	0.321379058864112\\
4	0.300316175834723\\
5	0.180989383558685\\
6	0.174641644194835\\
7	0.119806200761475\\
8	0.100235228438087\\
9	0.0826144240859506\\
10	0.0750849512158574\\
11	0.0711210713964446\\
12	0.0674137274725377\\
13	0.0542517024716162\\
14	0.049188870992894\\
15	0.0401385995807562\\
16	0.0387922841897729\\
17	0.0222609311499474\\
18	0.0172654941549523\\
19	0.00701838794087909\\
20	0.0056097896523553\\
21	0.00507751122920044\\
22	0.00384084280965212\\
23	0.00261957051818231\\
24	0.00152989766845017\\
25	0.00114214985799079\\
26	0.00107701695694854\\
27	0.000974432491025389\\
28	0.000948067991377848\\
29	0.000508745144464383\\
30	0.000507844931534173\\
31	0.000457011942330344\\
32	0.000380240768015184\\
33	0.000163041080796977\\
34	8.52805210804461e-05\\
35	7.22625095956086e-05\\
36	6.2515093026299e-05\\
37	5.95940461329791e-05\\
38	4.68084354259516e-05\\
39	4.55132877191635e-05\\
40	4.45843247689138e-05\\
41	3.31125370677244e-05\\
42	1.64057710327918e-05\\
43	1.48576497539357e-05\\
44	9.05703472463093e-06\\
45	7.05210037598223e-06\\
46	6.5058243837547e-06\\
47	4.20478891782225e-06\\
48	3.93309372534775e-06\\
49	3.55174526691527e-06\\
50	3.49121351060755e-06\\
51	2.90251236554197e-06\\
52	2.87014115404481e-06\\
53	2.06793938504974e-06\\
54	1.40593602753245e-06\\
55	6.56133632421202e-07\\
56	6.42770920735969e-07\\
57	9.19148907733855e-09\\
58	3.96513874145718e-08\\
59	4.42471699258408e-08\\
60	6.24698826862706e-08\\
61	8.85539850274795e-08\\
62	1.68571700956209e-07\\
63	2.68717910376782e-07\\
64	2.75724281197083e-07\\
65	3.27675940855192e-07\\
66	3.29446927821112e-07\\
67	3.44628179249239e-07\\
68	3.56605889577629e-07\\
69	3.59650027244835e-07\\
70	3.59880755931486e-07\\
71	3.61887419756041e-07\\
72	3.65570686960882e-07\\
73	3.67442175664221e-07\\
74	3.68473529888335e-07\\
75	3.68710653128182e-07\\
76	3.70582486681397e-07\\
77	3.70762192254e-07\\
78	3.70768556657279e-07\\
79	3.71106412566622e-07\\
80	3.71314071187735e-07\\
81	3.71920821994808e-07\\
82	3.71940689650113e-07\\
83	3.72388485036143e-07\\
84	3.72853409324095e-07\\
85	3.72915780378445e-07\\
86	3.73443103589051e-07\\
87	3.73453309417857e-07\\
88	3.73467174142956e-07\\
89	3.7356758412684e-07\\
90	3.73578139694698e-07\\
91	3.73586855302681e-07\\
92	3.73730472563973e-07\\
93	3.73773415173991e-07\\
94	3.73783264907817e-07\\
95	3.73790367250716e-07\\
96	3.73830595390258e-07\\
97	3.73830837241494e-07\\
98	3.73834585772679e-07\\
99	3.73836580882403e-07\\
100	3.73837928665234e-07\\
101	3.73838632706567e-07\\
102	3.73839719570376e-07\\
103	3.73840462073227e-07\\
104	3.7384057941345e-07\\
105	3.73841381564247e-07\\
106	3.73841519439008e-07\\
107	3.73842137756791e-07\\
108	3.73842674262363e-07\\
109	3.73844207507934e-07\\
110	3.73844448218362e-07\\
111	3.73844467938038e-07\\
112	3.73844756888336e-07\\
113	3.73844815069529e-07\\
114	3.73844825173826e-07\\
115	3.73844881562322e-07\\
116	3.73844923772207e-07\\
117	3.73844971034241e-07\\
118	3.73844981953401e-07\\
119	3.73845008029006e-07\\
120	3.73845057572655e-07\\
121	3.73845074684771e-07\\
122	3.73845079410974e-07\\
123	3.73845106627387e-07\\
124	3.73845108583057e-07\\
125	3.73845112168453e-07\\
126	3.73845113635206e-07\\
127	3.73845115590876e-07\\
128	3.73845126510036e-07\\
129	3.73845129606514e-07\\
130	3.73845130095432e-07\\
131	3.73845130910294e-07\\
132	3.73845132051102e-07\\
133	3.73845132702992e-07\\
134	3.7384513319191e-07\\
135	3.73845133354882e-07\\
136	3.738451338438e-07\\
137	3.738451338438e-07\\
138	3.73845134169745e-07\\
139	3.73845134332717e-07\\
140	3.73845134332717e-07\\
141	3.7384513449569e-07\\
142	3.73845134821635e-07\\
143	3.73845134821635e-07\\
144	3.73845134821635e-07\\
145	3.73845134821635e-07\\
146	3.73845134821635e-07\\
147	3.73845134821635e-07\\
148	3.73845134821635e-07\\
149	3.73845134821635e-07\\
150	3.73845134821635e-07\\
151	3.73845134821635e-07\\
152	3.73845134821635e-07\\
153	3.73845134821635e-07\\
154	3.73845134821635e-07\\
155	3.73845134821635e-07\\
156	3.73845134821635e-07\\
157	3.73845134821635e-07\\
158	3.73845134821635e-07\\
159	3.73845134821635e-07\\
160	3.73845134821635e-07\\
161	3.73845134821635e-07\\
162	3.73845134821635e-07\\
163	3.73845134821635e-07\\
164	3.73845134821635e-07\\
165	3.73845134821635e-07\\
166	3.73845134821635e-07\\
167	3.73845134821635e-07\\
168	3.73845134821635e-07\\
169	3.73845134821635e-07\\
170	3.73845134821635e-07\\
171	3.73845134821635e-07\\
172	3.73845134821635e-07\\
173	3.73845134821635e-07\\
174	3.73845134821635e-07\\
175	3.73845134821635e-07\\
176	3.73845134821635e-07\\
177	3.73845134821635e-07\\
178	3.73845134821635e-07\\
179	3.73845134821635e-07\\
180	3.73845134821635e-07\\
181	3.73845134821635e-07\\
182	3.73845134821635e-07\\
183	3.73845134821635e-07\\
184	3.73845134821635e-07\\
185	3.73845134821635e-07\\
186	3.73845134821635e-07\\
187	3.73845134821635e-07\\
188	3.73845134821635e-07\\
189	3.73845134821635e-07\\
190	3.73845134821635e-07\\
191	3.73845134821635e-07\\
192	3.73845134821635e-07\\
193	3.73845134821635e-07\\
194	3.73845134821635e-07\\
195	3.73845134821635e-07\\
196	3.73845134821635e-07\\
197	3.73845134821635e-07\\
198	3.73845134821635e-07\\
199	3.73845134821635e-07\\
200	3.73845134821635e-07\\
201	3.73845134821635e-07\\
202	3.73845134821635e-07\\
203	3.73845134821635e-07\\
204	3.73845134821635e-07\\
205	3.73845134821635e-07\\
206	3.73845134821635e-07\\
207	3.73845134821635e-07\\
208	3.73845134821635e-07\\
209	3.73845134821635e-07\\
210	3.73845134821635e-07\\
211	3.73845134821635e-07\\
212	3.73845134821635e-07\\
213	3.73845134821635e-07\\
214	3.73845134821635e-07\\
215	3.73845134821635e-07\\
216	3.73845134821635e-07\\
217	3.73845134821635e-07\\
218	3.73845134821635e-07\\
219	3.73845134821635e-07\\
220	3.73845134821635e-07\\
221	3.73845134821635e-07\\
222	3.73845134821635e-07\\
223	3.73845134821635e-07\\
224	3.73845134821635e-07\\
225	3.73845134821635e-07\\
226	3.73845134821635e-07\\
227	3.73845134821635e-07\\
228	3.73845134821635e-07\\
229	3.73845134821635e-07\\
230	3.73845134821635e-07\\
231	3.73845134821635e-07\\
232	3.73845134821635e-07\\
233	3.73845134821635e-07\\
234	3.73845134821635e-07\\
235	3.73845134821635e-07\\
236	3.73845134821635e-07\\
237	3.73845134821635e-07\\
238	3.73845134821635e-07\\
239	3.73845134821635e-07\\
240	3.73845134821635e-07\\
241	3.73845134821635e-07\\
242	3.73845134821635e-07\\
243	3.73845134821635e-07\\
244	3.73845134821635e-07\\
245	3.73845134821635e-07\\
246	3.73845134821635e-07\\
247	3.73845134821635e-07\\
248	3.73845134821635e-07\\
249	3.73845134821635e-07\\
250	3.73845134821635e-07\\
251	3.73845134821635e-07\\
252	3.73845134821635e-07\\
253	3.73845134821635e-07\\
254	3.73845134821635e-07\\
255	3.73845134821635e-07\\
256	3.73845134821635e-07\\
257	3.73845134821635e-07\\
};
\addlegendentry{$T=20$}

\addplot [color=red!60!teal, dashed, line width=1.1pt, mark size=2.0pt, mark=o, mark options={solid, red!60!teal}, forget plot]
  table[row sep=crcr]{%
6	0.173764830829518\\
18	0.000712138773012073\\
30	1.84486100969018e-05\\
42	1.5095497681292e-06\\
54	2.72572906789643e-07\\
78	3.73748170728945e-07\\
102	3.73845113146288e-07\\
138	3.7384513579947e-07\\
186	3.73845136614333e-07\\
246	3.73845136940278e-07\\
};
\node[right, align=left, inner sep=0]
at (rel axis cs:-0.15,1.1) {$(b)$};
\end{axis}
\end{tikzpicture}%
    \end{minipage}
    \caption{Plane-Poiseuille relative error versus stepper calls for \((a)\) slow and \((b)\) fast covariance-spectrum decay. (\solidlegend) deterministic truncation; (\dashlegend\circlegend) SSTE.}
    \label{fig:Poiseuille_eig_vs_sste}
\end{figure}
To demonstrate that the proposed framework applies to fluid-dynamical systems of practical relevance, we next consider plane Poiseuille flow. Plane Poiseuille flow is the steady laminar flow between two plates separated by a distance \(2h\) in the wall-normal direction. The details can be found in \citet{schmid-henningson-2001} and are not repeated here.
The final operator is given by
\begin{equation}
  \frac{\partial}{\partial t}
  \begin{bmatrix}
    \hat{v}\\
    \hat{\eta}
  \end{bmatrix}
  =
  -i
  \begin{bmatrix}
    \mathcal{L}_{OS} & 0\\
    \mathcal{L}_{C} & \mathcal{L}_{SQ}
  \end{bmatrix}
  \begin{bmatrix}
    \hat{v}\\
    \hat{\eta}
  \end{bmatrix},
\end{equation}
where \(\hat{\eta}\) denotes the wall-normal vorticity, \(\hat{v}\) the vertical velocity, and
\begin{equation}
  k^2 = k_x^2 + k_z^2,
\end{equation}
where \(k_x\) and \(k_z\) are the horizontal wavenumbers. The corresponding operators are given by
\begin{align}
  \mathcal{L}_{OS}
  &=
  -\left(
    \frac{\partial^2}{\partial y^2}-k^2
  \right)^{-1}
  \left[
    \frac{1}{i\Rey}
    \left(
      \frac{\partial^2}{\partial y^2}-k^2
    \right)^2
    - k_x U
    \left(
      \frac{\partial^2}{\partial y^2}-k^2
    \right)
    + k_x U''
  \right], \\
  \mathcal{L}_{C}
  &= k_z U'', \\
  \mathcal{L}_{SQ}
  &= k_x U - \frac{1}{i\Rey}
  \left(
    \frac{\partial^2}{\partial y^2}-k^2
  \right).
\end{align}
The disturbance energy in the velocity-vorticity variables is given by
\begin{equation}
  e
  =
  \frac{1}{4\pi^2}
  \int_{-\infty}^{\infty}
  \int_{-\infty}^{\infty}
  \frac{1}{2k^2}
  \int_{-1}^{1}
  \left(
    \left|\frac{\partial \hat{v}}{\partial y}\right|^2
    +
    k^2 |\hat{v}|^2
    +
    |\hat{\eta}|^2
  \right)
  dy\, dk_x\, dk_z.
\end{equation}
The Orr--Sommerfeld/Squire system is discretized with \(N=256\) wall-normal modes, the Reynolds number is set to \(R=1000\), the streamwise wavenumber is fixed to \(k_x=0\), and the spanwise wavenumber is set to \(k_z=2\). Chebyshev collocation differentiation matrices are used in the wall-normal direction, yielding the generalized semi-discrete system \(\mathsfbi{B}\dot{\mathbf{q}}=\mathsfbi{A}\mathbf{q}\). Since the coefficients are time independent, the reference propagator is computed as \(\boldsymbol{\Phi}(T)=\exp(-\ii\mathsfbi{B}^{-1}\mathsfbi{A}T)\).
The initial covariance is chosen in the wall-normal direction from a Gaussian kernel of the form
\begin{equation}
  C_{00}^{z}(y,y')
  \propto
  \exp\!\left(
    -\frac{(y-y')^2}{\lambda^2}
  \right),
\end{equation}
with correlation length \(\lambda=1\). In the present example, only the Orr--Sommerfeld block of the covariance is populated, while the Squire block is initially set to zero. This choice provides a simple but meaningful test in which the initial perturbation statistics are smooth in the wall-normal direction and the subsequent evolution is entirely determined by the linearized Poiseuille dynamics. Results are reported in \cref{fig:Poiseuille_eig_vs_sste}.
\subsubsection{Non-autonomous scalar problem}
\label{sec:appendix_verification_nonautonomous}
\begin{figure}[t]
    \centering
    \begin{minipage}[t]{0.49\textwidth}
        \centering
        \vspace{0pt}
        % This file was created by matlab2tikz.
%
\begin{tikzpicture}

\begin{axis}[%
width=11.252in,
height=7.131in,
at={(1.887in,0.962in)},
scale only axis,
clip=true,
separate axis lines,
every outer x axis line/.append style={black},
every x tick label/.append style={font=\color{black}},
every x tick/.append style={black},
xmin=0,
xmax=2,
xlabel={$t$},
every outer y axis line/.append style={black},
every y tick label/.append style={font=\color{black}},
every y tick/.append style={black},
ymin=0.3,
ymax=1,
ylabel={$g(t)$},
axis background/.style={fill=white},
legend style={legend cell align=left, align=left},
clip=true,
clip mode=individual,
axis lines=box,
xtick pos=bottom,
ytick pos=left,
tick align=inside,
major tick length=2pt,
width=0.7\linewidth,
height=0.65\linewidth,
every axis/.append style={font=\fontsize{9}{9}\selectfont},xlabel style={font=\fontsize{9}{9}\selectfont},title style={font=\fontsize{9}{9}\selectfont},
legend style={font=\fontsize{8}{9}\selectfont,inner sep=1pt,row sep=1pt,column sep=2pt,nodes={scale=1.00},draw=none},legend image post style={xscale=0.35},
legend columns=1,
scaled x ticks=false,
tick label style={/pgf/number format/fixed,/pgf/number format/precision=2},
every x tick label/.append style={font=\fontsize{9}{9}\selectfont\color{black}},
every y tick label/.append style={font=\fontsize{9}{9}\selectfont\color{black}},
,,
colormap={sigmamap}{rgb(0.0000)=(0.0200,0.1880,0.3800); rgb(0.1000)=(0.0743,0.3456,0.5616); rgb(0.2000)=(0.1918,0.4980,0.7060); rgb(0.3000)=(0.3890,0.6708,0.8226); rgb(0.4000)=(0.7552,0.8682,0.9228); rgb(0.5000)=(0.9690,0.9689,0.9689); rgb(0.6000)=(0.9425,0.8044,0.7614); rgb(0.7000)=(0.8767,0.4910,0.4020); rgb(0.8000)=(0.7869,0.2866,0.2470); rgb(0.9000)=(0.6186,0.1239,0.1568); rgb(1.0000)=(0.4040,0.0000,0.1220)}
]
\addplot [color=black, line width=1.5pt]
  table[row sep=crcr]{%
0	1\\
0.0408163265306122	0.66330675395804\\
0.0816326530612245	0.536824941129494\\
0.122448979591837	0.467466882696835\\
0.163265306122449	0.423266896705372\\
0.204081632653061	0.392783951145584\\
0.244897959183673	0.370777690358497\\
0.285714285714286	0.354468604631329\\
0.326530612244898	0.342230432206107\\
0.36734693877551	0.333042432038574\\
0.408163265306122	0.326229415862261\\
0.448979591836735	0.321326430782801\\
0.489795918367347	0.318003124620606\\
0.530612244897959	0.316019002335715\\
0.571428571428571	0.31519570292885\\
0.612244897959184	0.315399149096094\\
0.653061224489796	0.316527675802755\\
0.693877551020408	0.318503915093072\\
0.73469387755102	0.3212691167896\\
0.775510204081633	0.324779093288403\\
0.816326530612245	0.32900127412074\\
0.857142857142857	0.333912535699938\\
0.897959183673469	0.339497583456231\\
0.938775510204082	0.345747734884182\\
0.979591836734694	0.352659998591456\\
1.02040816326531	0.360236375476831\\
1.06122448979592	0.368483329251084\\
1.10204081632653	0.377411388089313\\
1.14285714285714	0.387034849440937\\
1.18367346938776	0.3973715673234\\
1.22448979591837	0.408442806704931\\
1.26530612244898	0.420273153452165\\
1.30612244897959	0.432890471193891\\
1.3469387755102	0.446325898617112\\
1.38775510204082	0.460613882363757\\
1.42857142857143	0.475792241975417\\
1.46938775510204	0.491902264338924\\
1.51020408163265	0.508988825889449\\
1.55102040816327	0.527100541482599\\
1.59183673469388	0.546289939391666\\
1.63265306122449	0.566613662349894\\
1.6734693877551	0.588132694962395\\
1.71428571428571	0.61091261817534\\
1.75510204081633	0.635023891824554\\
1.79591836734694	0.660542166602299\\
1.83673469387755	0.687548627088745\\
1.87755102040816	0.716130367800471\\
1.91836734693878	0.746380804518932\\
1.95918367346939	0.778400123482621\\
2	0.812295771362957\\
};
\addlegendentry{true}

\addplot [color=black, dashed, line width=1.1pt, mark size=2.0pt, mark=o, mark options={solid, black}]
  table[row sep=crcr]{%
0	1\\
0.0408163265306122	0.669131978952054\\
0.0816326530612245	0.530815625761453\\
0.122448979591837	0.461572038731577\\
0.163265306122449	0.428126499381252\\
0.204081632653061	0.396595186200975\\
0.244897959183673	0.370207055435072\\
0.285714285714286	0.34837607826492\\
0.326530612244898	0.340799996480598\\
0.36734693877551	0.335145490191288\\
0.408163265306122	0.325876802705544\\
0.448979591836735	0.321768597745276\\
0.489795918367347	0.319912049594763\\
0.530612244897959	0.31398495334479\\
0.571428571428571	0.313997204515448\\
0.612244897959184	0.314478299077611\\
0.653061224489796	0.3167772597657\\
0.693877551020408	0.318928667484235\\
0.73469387755102	0.320297180575997\\
0.775510204081633	0.323813698807683\\
0.816326530612245	0.328247661412773\\
0.857142857142857	0.333389122926578\\
0.897959183673469	0.33744213721697\\
0.938775510204082	0.345392063179751\\
0.979591836734694	0.351249172134897\\
1.02040816326531	0.35971284408624\\
1.06122448979592	0.366584347482057\\
1.10204081632653	0.376741786601441\\
1.14285714285714	0.386389034390955\\
1.18367346938776	0.395513055004814\\
1.22448979591837	0.406947617761363\\
1.26530612244898	0.419615540307281\\
1.30612244897959	0.432545766317335\\
1.3469387755102	0.445343882754692\\
1.38775510204082	0.459584101510697\\
1.42857142857143	0.475003561481353\\
1.46938775510204	0.491028098241514\\
1.51020408163265	0.507798149804996\\
1.55102040816327	0.525774415952527\\
1.59183673469388	0.544826733015256\\
1.63265306122449	0.565393610281706\\
1.6734693877551	0.586699025730275\\
1.71428571428571	0.609343836169864\\
1.75510204081633	0.632867063009338\\
1.79591836734694	0.659313548617411\\
1.83673469387755	0.685568801617742\\
1.87755102040816	0.714051110686585\\
1.91836734693878	0.744499584323107\\
1.95918367346939	0.776204500056826\\
2	0.810132941789997\\
};
\addlegendentry{SST}

\node[right, align=left, inner sep=0]
at (rel axis cs:-0.15,1.1) {$(a)$};
\end{axis}
\end{tikzpicture}%
    \end{minipage}
    \hfill
    \begin{minipage}[t]{0.49\textwidth}
        \centering
        \vspace{0pt}
        % This file was created by matlab2tikz.
%
\begin{tikzpicture}

\begin{axis}[%
width=11.252in,
height=7.131in,
at={(1.887in,0.962in)},
scale only axis,
clip=true,
separate axis lines,
every outer x axis line/.append style={black},
every x tick label/.append style={font=\color{black}},
every x tick/.append style={black},
xmin=0,
xmax=2,
xlabel={$t$},
every outer y axis line/.append style={black},
every y tick label/.append style={font=\color{black}},
every y tick/.append style={black},
ymode=log,
ymin=0.000788505972857121,
ymax=0.0171877742818613,
yminorticks=true,
ylabel={$\varepsilon_{\mathrm{rel}}$},
axis background/.style={fill=white},
clip=true,
clip mode=individual,
axis lines=box,
xtick pos=bottom,
ytick pos=left,
tick align=inside,
major tick length=2pt,
width=0.7\linewidth,
height=0.65\linewidth,
every axis/.append style={font=\fontsize{9}{9}\selectfont},xlabel style={font=\fontsize{9}{9}\selectfont},title style={font=\fontsize{9}{9}\selectfont},
legend style={font=\fontsize{8}{9}\selectfont,inner sep=1pt,row sep=1pt,column sep=2pt,nodes={scale=1.00},draw=none},legend image post style={xscale=0.35},
legend columns=1,
scaled x ticks=false,
tick label style={/pgf/number format/fixed,/pgf/number format/precision=2},
every x tick label/.append style={font=\fontsize{9}{9}\selectfont\color{black}},
every y tick label/.append style={font=\fontsize{9}{9}\selectfont\color{black}},
,,
colormap={sigmamap}{rgb(0.0000)=(0.0200,0.1880,0.3800); rgb(0.1000)=(0.0743,0.3456,0.5616); rgb(0.2000)=(0.1918,0.4980,0.7060); rgb(0.3000)=(0.3890,0.6708,0.8226); rgb(0.4000)=(0.7552,0.8682,0.9228); rgb(0.5000)=(0.9690,0.9689,0.9689); rgb(0.6000)=(0.9425,0.8044,0.7614); rgb(0.7000)=(0.8767,0.4910,0.4020); rgb(0.8000)=(0.7869,0.2866,0.2470); rgb(0.9000)=(0.6186,0.1239,0.1568); rgb(1.0000)=(0.4040,0.0000,0.1220)}
]
\addplot [color=black, line width=1.1pt, mark size=2.0pt, mark=o, mark options={solid, black}, forget plot]
  table[row sep=crcr]{%
0.0408163265306122	0.00878209811562205\\
0.0816326530612245	0.0111941806492774\\
0.122448979591837	0.0126101851991099\\
0.163265306122449	0.0114811782204232\\
0.204081632653061	0.00970313334920131\\
0.244897959183673	0.00153902173259956\\
0.285714285714286	0.0171877742818613\\
0.326530612244898	0.00417974437950421\\
0.36734693877551	0.00631468530853708\\
0.408163265306122	0.00108087480641458\\
0.448979591836735	0.00137606782422971\\
0.489795918367347	0.00600284974065746\\
0.530612244897959	0.00643647684440219\\
0.571428571428571	0.00380239451954645\\
0.612244897959184	0.00291963380726255\\
0.653061224489796	0.000788505972857121\\
0.693877551020408	0.00133358609120693\\
0.73469387755102	0.00302530234874586\\
0.775510204081633	0.00297246497902733\\
0.816326530612245	0.00229060726278548\\
0.857142857142857	0.00156751459559027\\
0.897959183673469	0.00605437664190624\\
0.938775510204082	0.00102870291992012\\
0.979591836734694	0.00400052873076924\\
1.02040816326531	0.00145329962832689\\
1.06122448979592	0.00515350795621087\\
1.10204081632653	0.00177419523894462\\
1.14285714285714	0.00166862247912742\\
1.18367346938776	0.00467701383645614\\
1.22448979591837	0.00366070578064619\\
1.26530612244898	0.00156472793820272\\
1.30612244897959	0.0007962865886255\\
1.3469387755102	0.00220022155439051\\
1.38775510204082	0.00223567046606397\\
1.42857142857143	0.00165761528769134\\
1.46938775510204	0.00177711338366069\\
1.51020408163265	0.00233929710023111\\
1.55102040816327	0.00251588724675035\\
1.59183673469388	0.00267844283941867\\
1.63265306122449	0.00215323446866404\\
1.6734693877551	0.00243766286825963\\
1.71428571428571	0.0025679319084316\\
1.75510204081633	0.00339645302008812\\
1.79591836734694	0.00186001446540253\\
1.83673469387755	0.00287954246870691\\
1.87755102040816	0.00290346172621076\\
1.91836734693878	0.00252045629313518\\
1.95918367346939	0.00282068740684575\\
2	0.00266261335982454\\
};
\node[right, align=left, inner sep=0]
at (rel axis cs:-0.15,1.1) {$(b)$};
\end{axis}
\end{tikzpicture}%
    \end{minipage}
    \caption{Non-autonomous verification of the SSTE estimator. \((a)\) Normalized mean energy growth \(g(T)\) compared with the closed-form prediction. \((b)\) Relative error of the SSTE estimate as a function of final time.}
    \label{fig:nonautonomous_hutchpp}
\end{figure}
We finally validate the SSTE construction against a fully non-autonomous linear problem for which \(g(T)\) is available in closed form. The model problem is
\begin{equation}
  b(t)\frac{\partial q}{\partial t}
  =
  a(t)\frac{\partial^2 q}{\partial z^2}
  -
  c(t)\frac{\partial q}{\partial z}
  +
  r(t)q .
\end{equation}
The initial covariance is chosen as
\begin{equation}
  C_{00}^{z}(z,z')
  =
  \exp\left(
    -\frac{(z-z')^2}{\lambda^2}
  \right).
\end{equation}
The normalized mean energy growth can be shown to be
\begin{equation}
  g(T)
  =
  \frac{\tr\left\{\mathsfbi{C}(T)\right\}}{\tr\left\{\mathsfbi{C}(0)\right\}}
  =
  \frac{
    \exp\left(2S(T)\right)
  }{
    \sqrt{1+8\tau(T)/\lambda^2}
  },
  \qquad
  \tau(T)=\int_0^T \frac{a(t)}{b(t)}\,dt,
  \qquad
  S(T)=\int_0^T \frac{r(t)}{b(t)}\,dt .
\end{equation}
This provides a direct reference for the matrix-free SSTE estimate. In the numerical test shown below, the interval is \(z\in[-64,64)\), discretized with \(N=256\) periodic grid points. The first and second spatial derivatives are approximated with eighth-order periodic finite-difference matrices, giving \(\mathsfbi{B}(t)=b(t)\mathsfbi{I}\) and \(\mathsfbi{A}(t)=a(t)\mathsfbi{D}_2-c(t)\mathsfbi{D}_1+r(t)\mathsfbi{I}\). The final time range is \(0\le T\le 2\), with time-step target \(\Delta t=0.01\), and the initial covariance length is \(\lambda=0.5\). The scalar coefficients are
\begin{equation}
  a(t)=0.2+0.8e^{-t/5},\qquad
  c(t)=0.25+0.15\cos(1.3t),
\end{equation}
\begin{equation}
  r(t)=0.5\sqrt{t}+0.04\sin(0.9t),\qquad
  b(t)=1.0+0.2\sin(0.7t).
\end{equation}
The SSTE estimate uses \(50\) probes. The method recovers the true solution, verifying the implementation of the method.


\subsubsection{Diagonal and cross-diagonal estimation}
\label{sec:appendix_verification_diagonal}

The estimator described so far returns a single scalar per target time, the trace of the
covariance-weighted operator. In several cases the quantity of interest is not the total energy but
its distribution: the pointwise variance profile across the inhomogeneous direction, or the
correlation between two state components at the same location. Both are diagonal entries of the
propagated covariance
\begin{equation}
  \mathsfbi{C}(T)
  =
  \boldsymbol{\Phi}(T)\,\mathsfbi{C}_{00}\,\boldsymbol{\Phi}^{*}(T)
  =
  \boldsymbol{\Phi}(T)\mathsfbi{L}
  \left[\boldsymbol{\Phi}(T)\mathsfbi{L}\right]^{*},
  \label{eq:diag_propagated_covariance}
\end{equation}
where \(\mathsfbi{C}_{00}=\mathsfbi{L}\mathsfbi{L}^{*}\) as in
\cref{subsec:randomized_trace_estimation}. The same matrix-free machinery applies: the loop-adjoint
procedure of \cref{alg:loop-adjoint} supplies the action
\begin{equation}
  \mathsfbi{D}(\mathbf{v})
  =
  \boldsymbol{\Phi}(T)\mathsfbi{L}
  \left(
    \mathsfbi{L}^{*}\boldsymbol{\Phi}^{*}(T)\mathsfbi{W}\mathbf{v}
  \right)
  =
  \mathsfbi{C}(T)\mathsfbi{W}\mathbf{v},
  \label{eq:diag_kernel}
\end{equation}
one adjoint sweep followed by one forward sweep per probe, without ever forming
\(\boldsymbol{\Phi}(T)\). This plays the role that \(\mathsfbi{K}\) plays in
\cref{alg:stochastic_stepping_trace_estimator}.

The starting point is the stochastic diagonal estimator of
\citet{bekas-kokiopoulou-saad-2007}. For probes \(\mathbf{g}_i\sim\mathcal{N}(\mathbf{0},\mathsfbi{I})\),
the identity \(\mathbb{E}\left[\mathbf{g}\mathbf{g}^{*}\right]=\mathsfbi{I}\) gives
\begin{equation}
  \mathrm{diag}\left\{\mathsfbi{C}\right\}
  =
  \mathbb{E}\left[
    \overline{\mathbf{g}}\odot\mathsfbi{D}(\mathbf{g})
  \right]
  \approx
  \frac{1}{N_p}\sum_{i=1}^{N_p}
  \overline{\mathbf{g}}_i\odot\mathsfbi{D}(\mathbf{g}_i),
  \label{eq:diag_hutchinson}
\end{equation}
with \(\odot\) the elementwise product and the overbar the complex conjugate. This is the diagonal
analogue of the Hutchinson estimator and converges at the same \(O(N_p^{-1/2})\) rate.

The deflation of \cref{alg:stochastic_stepping_trace_estimator} carries over unchanged, and with the
same three-stage split of the probe budget. Writing \(s=N_p/3\), a sketch
\(\mathsfbi{S}=[\mathbf{g}_1\ \cdots\ \mathbf{g}_s]\) is propagated to give
\(\mathsfbi{Y}=\mathsfbi{D}(\mathsfbi{S})\), whose thin QR factorization
\(\mathsfbi{Y}=\mathsfbi{Q}\mathsfbi{R}\) spans the dominant range of \(\mathsfbi{C}\). The diagonal
then splits into a deflated contribution evaluated exactly on that subspace and a stochastic
remainder confined to its orthogonal complement,
\begin{equation}
  \mathrm{diag}\left\{\mathsfbi{C}\right\}
  \approx
  \underbrace{
    \mathrm{diag}\left\{\mathsfbi{Q}\mathsfbi{B}\mathsfbi{Q}^{*}\right\}
  }_{\mathbf{d}_d}
  +
  \underbrace{
    \frac{1}{s}\sum_{i=1}^{s}
    \overline{\mathbf{r}}_i\odot\mathsfbi{D}(\mathbf{r}_i)
  }_{\mathbf{d}_r},
  \qquad
  \mathsfbi{B}=\mathsfbi{Q}^{*}\mathsfbi{D}(\mathsfbi{Q}),
  \qquad
  \mathbf{r}_i=\left(\mathsfbi{I}-\mathsfbi{Q}\mathsfbi{Q}^{*}\right)\mathbf{g}_i,
  \label{eq:diag_hutchpp}
\end{equation}
in direct correspondence with the \(t_d\) and \(t_r\) contributions of the trace estimator. Because
the probes \(\mathbf{r}_i\) carry no component along the leading subspace, the variance of
\(\mathbf{d}_r\) is governed by the trailing spectrum of \(\mathsfbi{C}\), which is precisely the
regime in which \eqref{eq:diag_hutchinson} is slowest.

The complete procedure is summarized in
\cref{alg:stochastic_stepping_diagonal_estimator}.

\begin{algorithm}[H]
  \caption{Stochastic stepping diagonal estimator for
           \(\mathrm{diag}\left\{\mathsfbi{C}(T)\right\}\)}
  \label{alg:stochastic_stepping_diagonal_estimator}

  \begin{algorithmic}[1]
    \Require Number of probes \(N_p\), loop-direct-adjoint action \(\LDA(\cdot)\),
             factor \(\mathsfbi{L}\) s.t. \(\mathsfbi{C}_{00}=\mathsfbi{L}\mathsfbi{L}^{*}\)
    \Ensure Diagonal estimate \(\widehat{\mathbf{d}}_{\mathrm{H++}}\)

    \State Define kernel
      \(\mathsfbi{D}(\mathbf{v})=\boldsymbol{\Phi}\mathsfbi{L}
      \left(\mathsfbi{L}^{*}\LDA(\mathbf{v})\right)\)
      \Comment{Action \(\mathsfbi{C}(T)\mathsfbi{W}\mathbf{v}\), one adjoint and one forward sweep}

    \State Set \(s\gets N_p/3\)
      \Comment{Number of probes per stage}

    \For{\(i=1,\ldots,s\)}
      \State Draw \(\mathbf{g}_i\sim\mathcal{N}(\mathbf{0},\mathsfbi{I})\)
        \Comment{Gaussian probe}
      \State \(\mathbf{y}_i\gets\mathsfbi{D}(\mathbf{g}_i)\)
        \Comment{Sketch action}
    \EndFor

    \State Form \(\mathsfbi{Y}\gets[\,\mathbf{y}_1\ \cdots\ \mathbf{y}_s\,]\)
      \Comment{Range sketch}

    \State Compute \(\mathsfbi{Y}=\mathsfbi{Q}\mathsfbi{R}\)
      \Comment{Thin QR factorization}

    \For{\(j=1,\ldots,s\)}
      \State \(\mathsfbi{B}(:,j)\gets\mathsfbi{Q}^{*}\mathsfbi{D}(\mathbf{q}_j)\)
        \Comment{Compressed operator on the captured subspace}
    \EndFor

    \State \(\mathbf{d}_d\gets
      \mathrm{diag}\left\{\mathsfbi{Q}\mathsfbi{B}\mathsfbi{Q}^{*}\right\}\)
      \Comment{Deflated contribution}

    \State \(\mathbf{d}_r\gets\mathbf{0}\)
    \For{\(i=1,\ldots,s\)}
      \State Draw \(\mathbf{g}_i\sim\mathcal{N}(\mathbf{0},\mathsfbi{I})\)
        \Comment{Residual probe}
      \State \(\mathbf{r}_i\gets(\mathsfbi{I}-\mathsfbi{Q}\mathsfbi{Q}^{*})\mathbf{g}_i\)
        \Comment{Orthogonal projection}
      \State \(\mathbf{d}_r\gets\mathbf{d}_r+
        s^{-1}\,\overline{\mathbf{r}}_i\odot\mathsfbi{D}(\mathbf{r}_i)\)
        \Comment{Stochastic residual contribution}
    \EndFor

    \State \(\widehat{\mathbf{d}}_{\mathrm{H++}}\gets\mathbf{d}_d+\mathbf{d}_r\)
      \Comment{Hutch++ diagonal estimate}

    \State \Return \(\widehat{\mathbf{d}}_{\mathrm{H++}}\)
  \end{algorithmic}
\end{algorithm}

The algorithm returns the whole diagonal of \(\mathsfbi{C}(T)\), from which the individual blocks
are read off afterwards: for a state vector partitioned into wall-normal velocity and wall-normal
vorticity, the variances \(C_{ww}(z)\) and \(C_{\eta\eta}(z)\) are the two halves of
\(\widehat{\mathbf{d}}_{\mathrm{H++}}\). The same sketch and residual probes also give the
cross-covariance between two blocks at coincident points, obtained by pairing the entries of one
block against those of the other in the last step instead of with themselves; for plane Poiseuille
flow this is the velocity--vorticity correlation \(C_{w\eta}(z)\), and it requires no additional
sweeps.

The cost follows the same accounting as \eqref{eq:sste_fixed_time_cost}: every one of the \(N_p\)
probes requires a full forward and adjoint sweep to the target time, so a diagonal estimate costs
the same as a trace estimate at that time. The difference is that the result is now required at
every grid point rather than as a single number, and the \(O(N_t^{2})\) growth of
\eqref{eq:sste_multiple_time_cost} applies unchanged. For this reason the verification below fixes a
small number of target times, \(N_t=5\), and overlays the resulting profiles rather than resolving a
\((z,T)\) map.

\Cref{fig:diag_verification} collects the four cases. Panel \((a)\) is the complex
Ginzburg--Landau problem with the Gaussian initial covariance of
\cref{sec:appendix_verification_gl}, using \(N_p=30\) probes; the estimate tracks the exact profile
obtained from the matrix exponential at every sampled time, the two being indistinguishable at plot
accuracy. Panels \((b)\)--\((d)\) show plane Poiseuille flow: the wall-normal velocity variance
\(\mathsfbi{C}_{ww}(z)\), the wall-normal vorticity variance \(\mathsfbi{C}_{\eta\eta}(z)\), and the
velocity--vorticity correlation \(\mathsfbi{C}_{w\eta}(z)\) at coincident wall-normal locations. The
errors in \((b)\)--\((d)\) sit three to four orders of magnitude below \((a)\), because the leading
spectrum of the propagated covariance is captured almost entirely by the deflation subspace, leaving
very little for the stochastic residual to account for.

\begin{figure}
    \centering
    \begin{minipage}[t]{0.49\textwidth}
        \centering
        \vspace{0pt}
        \tikzfig{GL_diag_overlay_hutchpp}
    \end{minipage}\hfill
    \begin{minipage}[t]{0.49\textwidth}
        \centering
        \vspace{0pt}
        \tikzfig{Poiseuille_diag_OS_overlay}
    \end{minipage}

    \vspace{1em}

    \begin{minipage}[t]{0.49\textwidth}
        \centering
        \vspace{0pt}
        \tikzfig{Poiseuille_diag_SQ_overlay}
    \end{minipage}\hfill
    \begin{minipage}[t]{0.49\textwidth}
        \centering
        \vspace{0pt}
        \tikzfig{Poiseuille_diag_cross_overlay}
    \end{minipage}
    \caption{Diagonal and cross-diagonal estimation at \(N_t=5\) fixed final times, exact profiles
    from the matrix exponential (solid) against the Hutch++ estimate (markers), one colour per time.
    \((a)\) complex Ginzburg--Landau, \(C(z)\); \((b)\) plane Poiseuille wall-normal velocity
    variance \(C_{ww}(z)\); \((c)\) wall-normal vorticity variance \(C_{\eta\eta}(z)\);
    \((d)\) velocity--vorticity correlation \(C_{w\eta}(z)\). The relative error integrated over the
    domain, \(\|\widehat{\mathbf{d}}-\mathbf{d}\|_{2}/\|\mathbf{d}\|_{2}\), averaged over the five
    times, is \(5.8\times10^{-4}\), \(6.2\times10^{-8}\), \(2.4\times10^{-7}\) and
    \(5.1\times10^{-8}\) respectively. The legend of \((a)\) applies to all four panels.}
    \label{fig:diag_verification}
\end{figure}

\clearpage
\end{appendices}




\clearpage
\bibliographystyle{plainnat}
\bibliography{references,manual_references}


\end{document}
